% !TeX program = xelatex
% Compile twice: xelatex -interaction=nonstopmode main.tex
\documentclass[UTF8,10.5pt,a4paper,oneside,fontset=fandol,scheme=plain]{ctexbook}
% ==================== Art-Book Editorial Preamble for AI Series Volume 1 ====================
\usepackage[top=24mm,bottom=24mm,left=23mm,right=23mm,headheight=22pt,headsep=7mm]{geometry}
\usepackage{amsmath,amssymb,mathtools,amsthm}
\usepackage{booktabs,tabularx,array,multirow,longtable}
\usepackage{graphicx}
\usepackage{xcolor}
\usepackage{tikz}
\usepackage{pgfplots}
\pgfplotsset{compat=1.18}
\usetikzlibrary{arrows.meta,positioning,shapes.geometric,calc,decorations.pathreplacing,backgrounds}
\usepackage[most]{tcolorbox}
\usepackage{fancyhdr}
\usepackage{listings}
\usepackage{enumitem}
\usepackage{hyperref}

% ==================== Curated Gallery Color Palette ====================
\definecolor{ink}{HTML}{1B2430}        % Deep charcoal-slate ink
\definecolor{teal}{HTML}{0B5555}       % Refined verdigris teal
\definecolor{tealDark}{HTML}{083B3B}   % Deep gallery spine teal
\definecolor{tealLight}{HTML}{F2F8F8}  % Delicate ice-teal parchment
\definecolor{tealLine}{HTML}{A3C6C5}   % Soft architectural rule
\definecolor{tealMist}{HTML}{DCEBEB}   % Watermark chapter numeral tint
\definecolor{navy}{HTML}{1D3B5E}       % Muted Prussian indigo
\definecolor{navyLight}{HTML}{F2F5FA}  % Soft indigo tint
\definecolor{ochre}{HTML}{9E4A22}      % Burnished terracotta-ochre
\definecolor{gold}{HTML}{B97A24}       % Fine burnished amber-gold accent
\definecolor{ochreLight}{HTML}{FCF5F0} % Warm parchment tint
\definecolor{plum}{HTML}{5B3A6B}       % Muted editorial plum
\definecolor{plumLight}{HTML}{F8F4FB}  % Soft plum tint
\definecolor{slate}{HTML}{526071}      % Muted architectural gray
\definecolor{codebg}{HTML}{F7F9FB}     % Clean light code canvas
\definecolor{codeborder}{HTML}{D4DCE5}

\hypersetup{
  unicode=true,
  colorlinks=true,
  urlcolor=teal,
  linkcolor=teal,
  citecolor=teal,
  pdftitle={天书中的智能 · 第1卷：Python/NumPy编程、AI线性代数与多元微积分},
  pdfauthor={天书中的智能：人工智能的七卷手记}
}

% ==================== Editorial Chapter Openers & Typography ====================
\makeatletter
\newcommand{\chapWatermarkNum}{%
  \@ifundefined{inappendix}{\ifnum\value{chapter}<10 0\fi\thechapter}{\thechapter}%
}
\makeatother

\ctexset{
  contentsname={目\quad 录},
  chapter={
    name={第,章},
    number=\arabic{chapter},
    format=\LARGE\sffamily\bfseries\color{tealDark},
    nameformat=\small\sffamily\bfseries\color{gold},
    aftername=\hfill{\fontsize{38}{38}\selectfont\sffamily\bfseries\color{tealMist}\chapWatermarkNum}\par\vspace{-4mm},
    aftertitle=\par\vspace{2.8mm}{\noindent\color{gold}\rule{26mm}{1.5pt}}{\color{tealLine}\rule{\dimexpr\linewidth-26mm\relax}{0.55pt}}\par,
    beforeskip=6pt,
    afterskip=10pt
  },
  section={
    format=\Large\sffamily\bfseries\color{teal},
    beforeskip=14pt plus 3pt minus 2pt,
    afterskip=6pt plus 1pt
  },
  subsection={
    format=\large\sffamily\bfseries\color{ink},
    beforeskip=10pt plus 2pt minus 1pt,
    afterskip=4pt plus 1pt
  },
  subsubsection={
    format=\normalsize\sffamily\bfseries\color{navy},
    beforeskip=7pt,
    afterskip=2pt
  }
}
\renewcommand{\lstlistingname}{代码}
\renewcommand{\figurename}{图}
\renewcommand{\tablename}{表}

% Artisanal Chapter Epigraph macro (quote + attribution)
\newcommand{\chapEpigraph}[2]{%
  \vspace{-2mm}%
  \begin{tcolorbox}[
    enhanced,
    colback=ochreLight!45!white,
    colframe=ochreLight!45!white,
    borderline west={1.8pt}{0pt}{gold},
    boxrule=0pt,
    arc=0.8mm,
    left=4mm,right=3.5mm,top=2mm,bottom=2mm,
    before skip=0pt,after skip=8pt
  ]
    {\kaishu\normalsize\color{ink!85}“#1”}\par\vspace{1mm}%
    \hfill{\footnotesize\sffamily\color{slate}——\,#2}%
  \end{tcolorbox}%
}

% Exercise tier header macro
\newcommand{\exTier}[2]{%
  \par\vspace{3.8mm}\noindent
  \colorbox{tealLight}{\strut\;\sffamily\bfseries\footnotesize\color{tealDark}#1\;}\;%
  {\small\sffamily\bfseries\color{slate}#2}\par\nopagebreak\vspace{1.2mm}%
}
\newcommand{\exHint}[1]{\par\smallskip\noindent{\footnotesize\sffamily\color{navy}\textbf{【提示】}\,#1}}

\setlength{\parindent}{0pt}
\setlength{\parskip}{5pt plus 1.5pt minus 0.5pt}
\linespread{1.26}
\setlist[itemize]{leftmargin=1.8em,topsep=3pt,itemsep=2.5pt,parsep=1pt}
\setlist[enumerate]{leftmargin=2.0em,topsep=3pt,itemsep=2.5pt,parsep=1pt}
\setlength{\tabcolsep}{5.5pt}
\renewcommand{\arraystretch}{1.22}

% ==================== Running Headers & Footers ====================
\pagestyle{fancy}
\fancyhf{}
\fancyhead[L]{\small\sffamily\color{slate}天书中的智能 \;\textcolor{gold}{$\cdot$}\; 第 1 卷}
\fancyhead[R]{\small\sffamily\color{slate}\nouppercase{\leftmark}}
\fancyfoot[C]{\small\sffamily\color{slate}\textcolor{gold}{$\diamond$}\quad\textcolor{tealDark}{\thepage}\quad\textcolor{gold}{$\diamond$}}
\renewcommand{\headrulewidth}{0.4pt}
\renewcommand{\headrule}{\hbox to\headwidth{\color{tealLine}\leaders\hrule height \headrulewidth\hfill}}
\renewcommand{\chaptermark}[1]{\markboth{第 \thechapter\ 章\quad #1}{}}

\fancypagestyle{plain}{
  \fancyhf{}
  \fancyfoot[C]{\small\sffamily\color{slate}\textcolor{gold}{$\diamond$}\quad\textcolor{tealDark}{\thepage}\quad\textcolor{gold}{$\diamond$}}
  \renewcommand{\headrulewidth}{0pt}
}

% ==================== Theorem & Pedagogical Box Environments ====================
\theoremstyle{definition}
\newtheorem{theorem}{定理}[chapter]
\newtheorem{lemma}[theorem]{引理}
\newtheorem{proposition}[theorem]{命题}
\newtheorem{corollary}[theorem]{推论}
\newtheorem{definition}[theorem]{定义}
\newtheorem{example}[theorem]{例}
\newtheorem{remark}[theorem]{注记}
\renewcommand{\proofname}{\textbf{\sffamily\color{teal}证明}}

\newtcolorbox{thmBox}[1][]{
  enhanced,breakable,
  colback=tealLight,
  colframe=tealLine,
  borderline west={2.4pt}{0pt}{teal},
  boxrule=0.45pt,
  arc=1.5mm,
  left=4mm,right=3.8mm,top=2.8mm,bottom=2.8mm,
  before skip=6.5pt,after skip=6.5pt,
  #1
}

\newtcolorbox{defBox}[1][]{
  enhanced,breakable,
  colback=navyLight,
  colframe=navy!30,
  borderline west={2.4pt}{0pt}{navy},
  boxrule=0.45pt,
  arc=1.5mm,
  left=4mm,right=3.8mm,top=2.8mm,bottom=2.8mm,
  before skip=6.5pt,after skip=6.5pt,
  #1
}

\newtcolorbox{insightBox}[1]{
  enhanced,breakable,
  colback=white,
  colframe=navy!35,
  borderline west={2.4pt}{0pt}{gold},
  coltitle=navy,
  fonttitle=\sffamily\bfseries\small,
  attach title to upper={\quad},
  title={【直觉与下游坐标】#1},
  boxrule=0.45pt,
  arc=1.2mm,
  left=4mm,right=3.8mm,top=2.6mm,bottom=2.6mm,
  before skip=6.5pt,after skip=6.5pt
}

\newtcolorbox{recallBox}[1]{
  enhanced,breakable,
  colback=codebg,
  colframe=codeborder,
  borderline west={2.2pt}{0pt}{slate},
  coltitle=slate,
  fonttitle=\sffamily\bfseries\small,
  attach title to upper={\quad},
  title={【前置速忆】#1},
  boxrule=0.4pt,
  arc=1.2mm,
  left=3.8mm,right=3.5mm,top=2.4mm,bottom=2.4mm,
  before skip=5.5pt,after skip=5.5pt
}

\newtcolorbox{failBox}[1]{
  enhanced,breakable,
  colback=ochreLight,
  colframe=ochre!35,
  borderline west={2.4pt}{0pt}{ochre},
  coltitle=ochre,
  fonttitle=\sffamily\bfseries\small,
  attach title to upper={\quad},
  title={【反例、数值陷阱与排错】#1},
  boxrule=0.45pt,
  arc=1.2mm,
  left=4mm,right=3.8mm,top=2.6mm,bottom=2.6mm,
  before skip=6.5pt,after skip=6.5pt
}

\newtcolorbox{extBox}[1]{
  enhanced,breakable,
  colback=plumLight,
  colframe=plum!30,
  borderline west={2.4pt}{0pt}{plum},
  coltitle=plum,
  fonttitle=\sffamily\bfseries\small,
  attach title to upper={\quad},
  title={【选读拓展 · 初读可跳过】#1},
  boxrule=0.45pt,
  arc=1.2mm,
  left=4mm,right=3.8mm,top=2.6mm,bottom=2.6mm,
  before skip=6.5pt,after skip=6.5pt
}

\newtcolorbox{chapIntro}[1]{
  enhanced,
  colback=tealLight,
  colframe=tealLine,
  borderline west={2.8pt}{0pt}{tealDark},
  boxrule=0.45pt,
  arc=1.5mm,
  left=4.2mm,right=4mm,top=3mm,bottom=3mm,
  before skip=2pt,after skip=10pt,
  title={\sffamily\bfseries\small\color{tealDark}#1},
  colbacktitle=tealLight,
  attach boxed title to top left={yshift=-1mm,xshift=2mm},
  boxed title style={empty,size=minimal}
}

% ==================== Code Listing Style ====================
\lstdefinestyle{pycode}{
  language=Python,
  basicstyle=\footnotesize\ttfamily\color{ink},
  keywordstyle=\bfseries\color{teal},
  commentstyle=\itshape\color{slate},
  stringstyle=\color{ochre},
  backgroundcolor=\color{codebg},
  frame=single,
  rulecolor=\color{codeborder},
  framesep=5pt,
  xleftmargin=4pt,
  xrightmargin=4pt,
  numbers=left,
  numberstyle=\tiny\ttfamily\color{slate!80},
  numbersep=6pt,
  columns=fullflexible,
  keepspaces=true,
  showstringspaces=false,
  breaklines=true,
  aboveskip=6.5pt,
  belowskip=6.5pt
}
\lstset{style=pycode}

% ==================== Math Shortcuts ====================
\newcommand{\R}{\mathbb{R}}
\newcommand{\C}{\mathbb{C}}
\newcommand{\N}{\mathbb{N}}
\newcommand{\E}{\mathbb{E}}
\newcommand{\col}{\operatorname{col}}
\newcommand{\row}{\operatorname{row}}
\newcommand{\nul}{\operatorname{null}}
\newcommand{\rank}{\operatorname{rank}}
\newcommand{\diag}{\operatorname{diag}}
\newcommand{\tr}{\operatorname{tr}}
\newcommand{\spanop}{\operatorname{span}}
\newcommand{\argmin}{\operatorname*{arg\,min}}
\newcommand{\argmax}{\operatorname*{arg\,max}}
\newcommand{\cond}{\kappa}
\newcommand{\norm}[1]{\left\|#1\right\|}
\newcommand{\abs}[1]{\left|#1\right|}
\newcommand{\inner}[2]{\left\langle #1,\, #2 \right\rangle}
\newcommand{\had}{\odot}

\makeatletter
\setlength{\@fpsep}{16pt}
\makeatother


\begin{document}
\color{ink}

% ==================== Cover Page (Gallery Art-Book Edition) ====================
\frontmatter
\thispagestyle{empty}
\begin{tikzpicture}[remember picture, overlay]
  % Architectural deep teal spine + burnished gold hairline + vertical monograph label
  \fill[tealDark] (current page.north west) rectangle ([xshift=11mm]current page.south west);
  \fill[gold] ([xshift=11mm]current page.north west) rectangle ([xshift=12.2mm]current page.south west);
  \node[rotate=90, anchor=center, font=\sffamily\bfseries\scriptsize, color=tealMist, opacity=0.85]
    at ([xshift=5.5mm]current page.west) {INTELLIGENCE FROM THE BOOK \quad $\cdot$ \quad VOLUME I};

  % Bespoke fine-art geometric SVD / orthogonal projection / vector-constellation composition (top-right)
  \begin{scope}[shift={([xshift=-95mm, yshift=-20mm]current page.north east)}, scale=0.48, transform shape]
    % Delicate sheared lattice backdrop (linear transformation motif)
    \begin{scope}[opacity=0.11]
      \foreach \k in {-2,-1,0,1,2} {
        \draw[teal, line width=0.45pt] ({-2.4 + 0.35*\k}, {-1.15 + 0.55*\k}) -- ({2.4 + 0.35*\k}, {-1.15 + 0.55*\k});
        \draw[navy, line width=0.45pt] ({-1.9 + 0.95*\k}, -1.35) -- ({-0.9 + 0.95*\k}, 1.45);
      }
    \end{scope}

    % Concentric rotated SVD singular ellipses
    \begin{scope}[rotate=24]
      \fill[teal, opacity=0.05] (0,0) ellipse (2.65cm and 1.18cm);
      \draw[teal, opacity=0.30, line width=1.05pt] (0,0) ellipse (2.65cm and 1.18cm);
      \draw[navy, opacity=0.22, line width=0.80pt] (0,0) ellipse (1.90cm and 0.85cm);
      \draw[teal, opacity=0.18, line width=0.60pt] (0,0) ellipse (1.12cm and 0.50cm);
      % Orthogonal principal singular axes (sigma_1 u_1, sigma_2 u_2)
      \draw[gold, opacity=0.55, line width=1.15pt, -{Stealth[length=2.2mm]}] (0,0) -- (2.60, 0);
      \draw[ochre, opacity=0.48, line width=1.05pt, -{Stealth[length=2.2mm]}] (0,0) -- (0, 1.15);
      % Subtle orthogonal projection triangle inside the singular frame
      \draw[navy, opacity=0.32, line width=0.8pt, dashed] (1.65, 0.92) -- (1.65, 0);
      \draw[navy, opacity=0.30, line width=0.65pt] (1.65, 0.18) -- (1.47, 0.18) -- (1.47, 0);
      \draw[tealDark, opacity=0.35, line width=0.85pt, -{Stealth[length=2.0mm]}] (0,0) -- (1.65, 0.92);
      % Vector constellation nodes along the singular shell
      \foreach \ang in {18, 62, 134, 198, 255, 318} {
        \draw[tealLine, opacity=0.35, thin] (0,0) -- ({\ang}:2.65cm and 1.18cm);
        \fill[gold, opacity=0.55] ({\ang}:2.65cm and 1.18cm) circle (1.5pt);
      }
      \fill[tealDark, opacity=0.50] (0,0) circle (1.8pt);
      \fill[ochre, opacity=0.55] (1.65, 0.92) circle (1.6pt);
    \end{scope}
  \end{scope}
\end{tikzpicture}
\vspace*{1mm}
\noindent{\small\sffamily\bfseries\color{teal}天书中的智能 \;\textcolor{gold}{$\cdot$}\; 人工智能的七卷手记 \hfill 第 1 卷 \quad 2026 年 11 月版}
\vspace{1.5mm}
\noindent{\color{gold}\rule{30mm}{1.8pt}}{\color{tealLine}\rule{\dimexpr\linewidth-30mm\relax}{0.7pt}}
\vspace{6.5mm}

{\fontsize{24}{32}\selectfont\sffamily\bfseries\color{ink}Python/NumPy 编程、\\AI 线性代数与多元微积分}\par
\vspace{2.5mm}
{\large\sffamily\color{teal}将基础微积分与 C 语言功底，顺滑升维为人工智能的计算与几何语言}\par
\vspace{4.5mm}
{\small\color{slate}融合北美四大（CMU / MIT / Stanford / UC Berkeley）与清北交南本科 AI 培养方案重构}\par

\vspace{5.5mm}

\begin{tcolorbox}[
  enhanced,
  colback=tealLight,
  colframe=tealLine,
  arc=2mm,
  boxrule=0.6pt,
  left=4.5mm,right=4.5mm,top=3mm,bottom=3mm
]
\textbf{\sffamily\color{tealDark}本卷定位与四大编写红线}\par\vspace{1.2mm}
\small
本卷面向\textbf{具备基础微积分与 C 语言编程基础}的学习者（线性代数核心前置直觉已在本卷内自洽补齐），不堆砌枯燥的纯技巧题，也绝不在任何 AI 核心数学上跳步或炫技：
\begin{itemize}
  \item \textbf{零先修倒挂与完全自洽}：内置“零障碍线性代数前置精华”（张成、基、线性变换、秩与零空间、行列式体积直觉、特征值几何），先打通 NumPy 向量化与正交投影（第 2 章即做出首个数据拟合模型），再证对称阵谱定理并顺推半正定矩阵，彻底消灭超前引用。
  \item \textbf{AI 视角重塑线代与微积分}：聚焦传统教材讲得浅但 AI 天天用的核心——矩阵乘法四视角与外积/张量、高维近正交、超定 vs 欠定、PSD 三大 AI 化身、SVD/PCA/伪逆、Jacobian 链式法则与矩阵微分迹技巧。
  \item \textbf{主干扎实、选读不卡人}：每节先给 2D/3D 几何画面与手算小例，再给严格推导，明确标注“后续在哪一卷、哪个模型里用到”；高阶数值拓展统一归入可跳过的【选读拓展】框。
  \item \textbf{四层阶梯习题 + 实测闭环}：全卷共 48 道精选习题（概念反例 / 小例手算 / 实用证明 / 可自检编码）附完整精解，所有实验数据与曲线均由本机 Python 实测生成。
\end{itemize}
\end{tcolorbox}

\vspace{4.5mm}

\noindent\textbf{\sffamily\color{tealDark}第 1 卷章节导览}
\vspace{1.5mm}

{\footnotesize\renewcommand{\arraystretch}{1.15}
\begin{tabularx}{\linewidth}{@{}p{15mm}p{47mm}X@{}}
\toprule
\textbf{章节} & \textbf{标题} & \textbf{核心能力、下游用途与代码实战} \\
\midrule
序\quad 章 & 七卷本全景与八校覆盖对照 & 北美四大 + 清北交南课程共识、七卷体系坐标、本卷先修依赖 DAG \\
第 1 章 & Python/NumPy 与矩阵乘法四视角 & C $\to$ Python 过渡、内存布局、广播防坑、外积秩 1、批处理张量与 \texttt{einsum} \\
第 2 章 & 范数、高维几何与线性最小二乘 & 前置线代直觉、$\ell_1/\ell_2/\ell_\infty$、高维近正交、四大子空间、正交投影与多项式/岭回归 \\
第 3 章 & 特征分解、谱定理与半正定矩阵 & 特征值几何与特征多项式、矩阵幂与 Jordan 直觉、实对称阵谱定理、PSD、条件数与特殊阵 \\
第 4 章 & 奇异值分解、伪逆与 PCA 降维 & SVD 几何、Eckart-Young 短证、欠定最小范数解、PCA 双视角等价与图像压缩 \\
第 5 章 & AI 多元微积分与矩阵微分 & 梯度与最速下降、Jacobian 链式法则、Hessian、拉格朗日乘子、迹技巧与 Hadamard 求导 \\
第 6 章 & 二次型梯度下降与有限差分校验 & $0<\eta<2/\lambda_{\max}$、最优步长与病态锯齿、特征标准化实测、SGD 直觉与中心差分校验 \\
附\quad 录 & 全卷 48 道分层习题完整精解 & 第 1--6 章 A/B/C/D 四层习题分步推导、反例验算与代码自检要点 \\
\bottomrule
\end{tabularx}}

\vfill
\noindent{\color{tealLine}\rule{\linewidth}{0.5pt}}
{\footnotesize\color{slate}配套实验脚本 \texttt{code/vol1\_experiments.py} 与实测数据随本卷提供。下一卷：第 2 卷《AI 概率论、数理统计、信息论与首个回归/分类闭环》。}

% ==================== Table of Contents ====================
\clearpage
\tableofcontents

% ==================== Main Matter ====================
\mainmatter
\chapter*{序章：八校培养共识、七卷全景与本卷使用指南}
\addcontentsline{toc}{chapter}{序章：八校培养共识、七卷全景与本卷使用指南}
\markboth{序章\quad 八校培养共识与七卷全景}{}

\begin{chapIntro}{本章定位}
本序章说明整套七卷本本科人工智能自学讲义的课程设计依据（综合北美四大 CMU、MIT、Stanford、UC Berkeley 与国内清北交南最新培养方案）、第 1 卷在整套体系中的精确位置与八校考点覆盖对照，以及阅读与习题自检守则。
\end{chapIntro}

\section*{0.1\quad 从八校培养方案到七卷本阶梯体系}
\addcontentsline{toc}{section}{0.1\quad 从八校培养方案到七卷本阶梯体系}

对照北美四大（CMU BSAI、MIT Course 6-4、Stanford CS AI Track、UC Berkeley EECS/CS）与国内四校（清华智班、北大通班、上交国智班、南京大学 2025 版 AI 培养方案），顶尖本科 AI 教育已形成高度一致的能力结构：
\begin{enumerate}
  \item \textbf{AI 定制化的数理与编程底座}（MIT 6.100+18.C06、Stanford Math 51+CS109、清华《AI 应用数学》）：不按纯数学系那样深挖多项式环或曲面积分，而是把\textbf{Python/NumPy 数值编程、AI 线性代数、多元矩阵微积分、概率统计与信息论}打通为可直接建模的底层语言。
  \item \textbf{经典 AI 决策与现代统计机器学习并重}（CMU 07-280/07-380、Stanford CS221+CS229、Berkeley CS188+189+EECS127、南大最优化与 ML）：先用搜索、概率图与 MDP 建立智能体决策直觉，再系统攻克凸优化与经典机器学习的泛化理论。
  \item \textbf{深度学习前置、408 升级为 AI 系统与科研复现}（交大《深度学习及应用》、Stanford CS224n / CS231n / CS234 / CS336、CMU 15-213、MIT SERC）：尽早吃透自动微分与 Transformer，将计算机组成原理/OS/网络知识升级为 GPU 显存与算术强度直觉，最终通过顶会复现与消融实验完成本科科研训练。
\end{enumerate}

据此，本系列划分为职责纯粹、零跨卷越界的\textbf{七卷本体系}：
\begin{center}
\small
\begin{tabularx}{\linewidth}{@{}p{13mm}p{52mm}X@{}}
\toprule
\textbf{卷次} & \textbf{卷名} & \textbf{核心主题与对应课程} \\
\midrule
\textbf{第 1 卷} & \textbf{Python/NumPy 编程、AI 线性代数与多元微积分} & \textbf{（本卷）}向量化与广播、矩阵乘法四视角、正交投影与最小二乘、谱定理与 PSD、SVD/PCA、矩阵微分与二次型梯度下降 \\
第 2 卷 & AI 概率论、数理统计、信息论与首个回归/分类闭环 & 条件独立与 Bayes、多元高斯、浓度不等式、熵/交叉熵/KL 散度、MLE/MAP、纯 NumPy 线性/逻辑/Softmax 回归全实验 \\
第 3 卷 & 经典人工智能：搜索、博弈、概率图与序列决策 & A* 搜索、Minimax/$\alpha$-$\beta$/MCTS、CSP、贝叶斯网络与 HMM 推断、MDP 与 Bellman 方程（衔接 408 算法） \\
第 4 卷 & 凸优化与经典机器学习 & $L$-光滑/强凸收敛率、拉格朗日对偶与 KKT、偏差-方差与泛化、SVM/核方法、树与 GBDT、聚类与 EM 算法 \\
第 5 卷 & 深度学习、自动微分与现代架构 & 手写 Micro-Autograd、残差/归一化/AdamW/排错、CNN 与 ViT、Transformer 与迷你 GPT、VAE 与扩散模型 \\
第 6 卷 & 强化学习、大模型后训练与 AI 系统 & DQN/PPO、Scaling Laws/SFT/LoRA/RLHF/DPO/GRPO、GPU HBM/SRAM 与 Roofline、FlashAttention、$16\Phi$ 显存与多卡并行 \\
第 7 卷 & AI 科研实践、严谨评测、安全伦理与顶会复现 & 论文三遍精读、多随机种子与消融实验、误差分析、分布偏移与安全伦理边界、完整顶会复现与技术报告 \\
\bottomrule
\end{tabularx}
\end{center}

\section*{0.2\quad 第 1 卷八校核心考点覆盖对照表}
\addcontentsline{toc}{section}{0.2\quad 第 1 卷八校核心考点覆盖对照表}

为确保本卷\textbf{不偷工减料、不留基础漏洞，同时不堆砌冷门炫技内容}，下表列出八校对“AI 数值编程 + 线代 + 微积分”模块要求的全部核心考点及其在本卷的对应章节：

\begin{center}
\footnotesize
\begin{tabularx}{\linewidth}{@{}p{34mm}p{34mm}p{16mm}X@{}}
\toprule
\textbf{八校核心考点} & \textbf{主要对标课程} & \textbf{本卷位置} & \textbf{下游直接用途} \\
\midrule
NumPy 内存/广播/向量化与 \texttt{einsum} & MIT 6.100 / Stanford CS231n & 第 1 章 & 第 2--6 卷全部代码实现与高维张量排错 \\
矩阵乘法四视角（列/行/外积/批处理） & MIT 18.06 / 18.C06 & 第 1 章 & Attention 加权、SVD/LoRA 秩 1 分解、反传外积梯度 \\
向量空间、秩、行列式与高维几何 & MIT 18.06 / 清华 AI 数学 & 第 2 章 & 线性变换直觉、$L_1/L_2$ 正则化、Embedding 检索与维度灾难 \\
四大子空间、超定最小二乘与岭回归 & MIT 18.06 / Stanford Math 51 & 第 2 章 & 线性回归解析解、过拟合与正则化初体验 \\
特征值几何、特征分解、Jordan 链与矩阵幂 & MIT 18.06 / 18.C06 & 第 3 章 & 动力系统稳定性、RNN/梯度迭代发散判定 \\
实对称谱定理、PSD 矩阵与条件数 $\kappa$ & 18.C06 / Berkeley EECS127 & 第 3 章 & 协方差阵、核矩阵、Hessian 凸地形与收敛快慢 \\
行随机矩阵与图拉普拉斯矩阵直觉 & MIT 18.06 / 北大通班 & 第 3 章 & 马尔可夫链、MDP、Softmax 注意力、谱聚类 \\
SVD、最优低秩近似、伪逆与 PCA & 18.06 / CS229 / 清华 & 第 4 章 & 数据降维、图像压缩、过参数化最小范数、LoRA \\
方向导数、Jacobian 链式法则与 Hessian & Stanford Math 51 / CS229 & 第 5 章 & 第 5 卷反向传播、二阶优化与局部极值判定 \\
拉格朗日乘子与矩阵微分迹技巧 & Berkeley EECS127 / CS229 & 第 5 章 & PCA 约束优化、第 4 卷对偶/KKT、矩阵目标求导 \\
二次型梯度下降步长阈值与有限差分 & CS229 / CS231n & 第 6 章 & 特征标准化调参、手写反向传播正确性自检 \\
\bottomrule
\end{tabularx}
\end{center}

\section*{0.3\quad 本卷知识阶梯与四层习题使用法}
\addcontentsline{toc}{section}{0.3\quad 本卷知识阶梯与四层习题使用法}

本卷严格按\textbf{零超前依赖}的有向无环阶梯展开（先内积投影，后谱定理与正定性，再 SVD 与微积分，最后落脚梯度下降动力学）：

\begin{center}
\resizebox{0.98\linewidth}{!}{%
\begin{tikzpicture}[
  node distance=6.5mm and 7.5mm,
  ch/.style={
    rectangle, rounded corners=1.8mm, draw=teal, thick, fill=tealLight,
    text width=34mm, minimum height=14.5mm, align=center, font=\sffamily\footnotesize, text=ink, inner sep=4pt
  },
  goal/.style={
    rectangle, rounded corners=1.8mm, draw=ochre, very thick, fill=ochreLight,
    text width=36mm, minimum height=14.5mm, align=center, font=\sffamily\footnotesize, text=ink, inner sep=4pt
  },
  arr/.style={-{Stealth[length=2.2mm]}, thick, teal, preaction={draw=white, line width=2.5pt, arrows=-}}
]
  \node[ch] (c1) {\textbf{第 1 章}\\Python/NumPy 与\\矩阵乘法四视角};
  \node[ch, right=of c1] (c2) {\textbf{第 2 章}\\范数、高维几何与\\正交投影最小二乘};
  \node[ch, right=of c2] (c3) {\textbf{第 3 章}\\特征分解、谱定理\\与半正定矩阵};

  \node[ch, below=9mm of c1] (c4) {\textbf{第 4 章}\\SVD、伪逆与\\PCA 降维实战};
  \node[ch, below=9mm of c2] (c5) {\textbf{第 5 章}\\AI 多元微积分\\与矩阵微分};
  \node[goal, below=9mm of c3] (c6) {\textbf{\color{ochre}第 6 章（通向第 2 卷）}\\二次型梯度下降与\\有限差分梯度校验};

  \draw[arr] (c1) -- (c2);
  \draw[arr] (c2) -- (c3);
  \draw[arr] ([xshift=-10mm]c2.south) -- ([xshift=4mm]c4.north);
  \draw[arr] ([xshift=-12mm]c3.south) -- ([xshift=12mm]c4.north);
  \draw[arr] (c1.south) -- (c4.north);
  \draw[arr] (c2.south) -- (c5.north);
  \draw[arr] ([xshift=-4mm]c3.south) -- ([xshift=10mm]c5.north);
  \draw[arr] (c3.south) -- (c6.north);
  \draw[arr] (c5) -- (c6);
\end{tikzpicture}}
\end{center}

\begin{insightBox}{阅读与做题建议（零压力、真掌握）}
\begin{itemize}
  \item \textbf{区分主干与选读}：正文全部为后续必用主干，顺着“具体小例 $\to$ 严格推导 $\to$ 代码验证”即可流畅读懂；紫色边框的\textbf{【选读拓展】}初读可直接跳过，二刷再看。
  \item \textbf{善用每章 8 道四层习题}：A 类（概念反例）和 B 类（手算小例）务必合上书独立完成；C 类（实用证明）卡住时先看题末【提示】；D 类（编程实验）务必亲手跑出断言通过再对照附录。
\end{itemize}
\end{insightBox}

\chapter[Python/NumPy 与矩阵乘法四视角]{Python/NumPy 数值编程\\与矩阵乘法的四种视角}\label{ch:numpy-matmul}

\chapEpigraph{程序必须写给人读，只是顺便让机器执行。}{哈罗德·阿贝尔森 \& 杰拉尔德·萨斯曼，《计算机程序的构造和解释》（SICP）}

\begin{chapIntro}{本章导学：前置基础、学习目标与下游坐标}
\textbf{前置基础}：大一 C 语言基础（循环、数组、内存地址概念）与线性代数基本矩阵乘法定义。

\textbf{学习目标}：完成从“逐元素标量循环”到“数组向量化”的思维切换；彻底扫清 NumPy 工程排错中的三大高频雷区（\texttt{view/copy}、\texttt{axis} 规约方向与广播静默维度膨胀）；掌握矩阵乘法的四种互补视角（尤其是列/行线性组合与外积秩 1 求和），并能流畅读写高维批处理张量与 \texttt{np.einsum} 表达式。

\textbf{下游坐标}：本章是本卷第 2--6 章与后续第 2--6 卷所有代码实现、Transformer 注意力计算、SVD/LoRA 低秩分解和神经网络反向传播推导的底层通用语言。
\end{chapIntro}

在 C 语言里，你想把两个 $100 \times 100$ 的矩阵相乘，第一反应是写三层 \texttt{for} 循环，把每个标量元素一个个算出来。但在人工智能与科学计算的世界里，\textbf{矩阵和张量本身就是基本运算单元}：手写 Python 三层循环不仅冗长难读，还会比底层向量化慢上数千倍。本章的目标不是把 Python 语言从头讲一遍，而是站在你已有的 C 语言与基础线代起点上，直接打通 AI 编程与读论文必需的核心表达习惯。

\section{从 C 思维到 Python 工程惯用法}\label{sec:c-to-python}

\begin{recallBox}{C 语言与 Python 的核心心智差异}
C 语言是静态编译型语言：变量必须先声明类型、数组是一段连续裸内存、循环几乎没有解释器开销。Python 是动态解释型语言：变量名只是指向堆上对象的引用，\texttt{for} 循环每步都要做动态类型检查；因此 Python 的黄金法则是\textbf{用简洁的高层语法组织控制流，把密集数值循环全部下沉到 C/NumPy 内核}。
\end{recallBox}

表~\ref{tab:c-vs-py} 浓缩了从 C 过渡到 AI 方向 Python 编程时最关键的六组差异对照。

\begin{table}[htbp]
\centering
\caption{C 语言基础到 AI 方向 Python/NumPy 工程惯用法对照}\label{tab:c-vs-py}
{\small
\begin{tabularx}{\linewidth}{@{}p{24mm}p{52mm}X@{}}
\toprule
\textbf{编程维度} & \textbf{C 语言习惯} & \textbf{Python/NumPy 工程惯用法} \\
\midrule
变量与类型 & 编译期静态声明 \texttt{double x = 1.0;} & 运行期动态绑定，函数签名用类型标注 \texttt{x: float} 提升可读性 \\
数组与切片 & 裸指针 + 长度参数，手动管理越界 & 左闭右开切片 \texttt{a[start:stop:step]}，支持负索引 \texttt{a[-1]} \\
构造序列 & 先 \texttt{malloc} 再 \texttt{for} 循环逐项赋值 & 列表/字典推导式 \texttt{[k*k for k in range(5)]}，一行表达意图 \\
算法封装 & 结构体 \texttt{struct} + 一组外部函数 & 轻量类（\texttt{\_\_init\_\_} / \texttt{fit} / \texttt{predict}），状态与方法统一 \\
防错检查 & \texttt{if (n != m) return -1;} & 关键节点写 \texttt{assert X.shape == (N, D)}，尽早暴露形状错位 \\
随机数生成 & 全局状态 \texttt{srand(42); rand();} & 独立生成器 \texttt{rng = np.random.default\_rng(42)}，隔离实验状态 \\
\bottomrule
\end{tabularx}}
\end{table}

在进入 NumPy 数组之前，只需牢记以下四条贯穿本套七卷本讲义的 Python 工程惯用法：

\begin{enumerate}
  \item \textbf{切片语义与推导式}：切片 \texttt{arr[i:j]} 永远遵循\textbf{左闭右开} $[i, j)$，长度恰好为 $j - i$；负索引 \texttt{-1} 表示最后一个元素，\texttt{arr[::-1]} 表示逆序视图。推导式如 \texttt{\{k: v for k, v in items if v > 0\}} 用于轻量数据整理，但\textbf{切勿用列表推导式替代大规模矩阵运算}。
  \item \textbf{类型标注与形状断言（Shape Assertions）}：AI 代码最常见的 Bug 不是语法报错，而是数组维度静默对错。养成在函数入口写类型标注与 \texttt{assert} 的习惯——例如 \texttt{assert X.ndim == 2}、\texttt{assert X.shape[1] == w.shape[0]}——能在错误扩散到损失函数前当场捕获。
  \item \textbf{统一模型类接口}：从第 2 卷开始，无论是线性回归、逻辑回归还是神经网络模块，我们都沿用 \texttt{fit(X, y)}（拟合参数）与 \texttt{predict(X)}（前向预测）的轻量类封装，便于无缝衔接 scikit-learn 与 PyTorch。
  \item \textbf{现代可复现随机数}：统一调用 \texttt{np.random.default\_rng(seed)} 创建独立生成器 \texttt{rng}，告别旧教材中的全局函数 \texttt{np.random.seed(42)}。这样即使多个函数交替采样，也不会互相污染随机流，保证实验结果百分百可复现。
\end{enumerate}

代码~\ref{lst:py-idiom} 给出了一个 20 行左右的迷你特征标准化器，完整演示了上述四条惯用法。

\begin{lstlisting}[caption={融合类型标注、形状断言、类封装与独立随机生成器的 Python 惯用法}, label={lst:py-idiom}]
import numpy as np

class SimpleScaler:
    """将特征矩阵按列标准化为零均值、单位标准差。"""
    def __init__(self, eps: float = 1e-8) -> None:
        self.eps = eps
        self.mean_: np.ndarray | None = None
        self.std_: np.ndarray | None = None

    def fit(self, X: np.ndarray) -> "SimpleScaler":
        assert X.ndim == 2, f"期望二维矩阵 (N, D)，实际维度 {X.ndim}"
        self.mean_ = X.mean(axis=0, keepdims=True)   # 形状 (1, D)
        self.std_ = X.std(axis=0, keepdims=True)     # 形状 (1, D)
        return self

    def transform(self, X: np.ndarray) -> np.ndarray:
        assert self.mean_ is not None and self.std_ is not None
        assert X.shape[1] == self.mean_.shape[1], "特征维数不匹配"
        return (X - self.mean_) / (self.std_ + self.eps)

rng = np.random.default_rng(42)
X_demo = rng.normal(loc=5.0, scale=2.0, size=(100, 3))
X_norm = SimpleScaler().fit(X_demo).transform(X_demo)
assert np.allclose(X_norm.mean(axis=0), 0.0, atol=1e-7)
\end{lstlisting}

\section{NumPy 核心：\texttt{ndarray} 内存布局、\texttt{axis} 与广播三步规则}\label{sec:numpy-core}

\subsection{\texttt{ndarray} 的底层结构与 \texttt{view/copy} 陷阱}

为什么 Python 原生列表 \texttt{list} 做数值计算很慢？因为 Python 的 \texttt{list} 底层存的并不是连续浮点数，而是一排指向堆内存中分散 \texttt{PyObject} 对象的指针；每次加法都要解引用指针、检查对象类型、拆箱计算再重新装箱。

相比之下，NumPy 的核心数据结构 \texttt{ndarray} 非常贴近 C 语言数组：它在底层维护一块\textbf{连续的同类型原始内存缓冲区（Data Buffer）}，外加一个轻量头部，记录三个核心元信息：
\begin{itemize}
  \item \texttt{shape}（形状）：各轴长度元组，例如 $(N, D)$ 表示 $N$ 个样本、$D$ 个特征；维数记作 \texttt{ndim = len(shape)}。
  \item \texttt{dtype}（元素类型）：决定每个元素占多少字节。科学计算与本卷梯度校验默认使用双精度 \texttt{float64}（8 字节，约 15--17 位有效数字）；深度学习大模型训练常使用单精度 \texttt{float32}（4 字节，约 7 位有效数字）乃至半精度 \texttt{bfloat16} 以减半显存占用。
  \item \texttt{strides}（步幅，单位为字节）：沿每个轴前进 1 步在内存中跳过的字节数。NumPy 默认采用 \textbf{C 连续（行优先，row-major）}布局：对一个形状为 $(N, D)$ 的 \texttt{float64} 数组，第 $(i, j)$ 个元素距离首地址的字节偏移为 $i \cdot (8D) + j \cdot 8$，即 \texttt{strides == (8*D, 8)}。因此按行读取连续命中 CPU 缓存，而跨大步幅按列读取更容易触发缓存未命中。
\end{itemize}

正是因为 \texttt{ndarray} 通过“头部元信息 + 共享数据缓冲区”表示数组，NumPy 在做普通切片和转置（\texttt{X.T}）时\textbf{不会复制底层数据}，而只是返回一个修改了 \texttt{shape}、\texttt{strides} 与起始指针的\textbf{视图（view）}。这带来了零拷贝的高效，但也埋下了初学者最常踩的隐蔽陷阱。

\begin{failBox}{\texttt{view} 与 \texttt{copy} 陷阱：切片修改静默污染原数组}
下面两段看起来非常相似的代码，行为却截然相反：
\begin{lstlisting}[numbers=none]
A = np.array([[1.0, 2.0, 3.0],
              [4.0, 5.0, 6.0]])
s = A[:, 0]          # 基础切片 -> 返回 view（与 A 共享底层内存！）
s[:] = -99.0         # 原地修改 s，A[:, 0] 同时被改成了 [-99., -99.]！

B = np.array([[1.0, 2.0, 3.0],
              [4.0, 5.0, 6.0]])
f = B[[0, 1], 0]     # 花式索引（整数列表/布尔掩码）-> 返回独立 copy
f[:] = -99.0         # B 毫发无损，B[0, 0] 依然是 1.0
\end{lstlisting}
\textbf{排错守则}：（1）基础切片（\texttt{A[1:3, ::2]}）、\texttt{A.T} 与大多数 \texttt{reshape} 返回 \textbf{view}；花式索引（整数数组索引、布尔掩码 \texttt{A[A > 0]}）总是返回 \textbf{copy}。（2）只要你打算原地修改子数组而不影响原数据，\textbf{务必显式调用 \texttt{.copy()}}，不确定时可用 \texttt{np.shares\_memory(a, b)} 直接核实。
\end{failBox}

\subsection{沿轴规约（\texttt{axis}）与 \texttt{keepdims=True}}

在机器学习中，数据表几乎总是按“\textbf{行 = 样本、列 = 特征}”排布成矩阵 $X \in \R^{N \times D}$：第 0 轴（\texttt{axis=0}）沿样本方向（竖着走），第 1 轴（\texttt{axis=1}）沿特征方向（横着走）。理解 \texttt{axis=k} 只有一个绝不记混的口诀：\textbf{\texttt{axis=k} 表示“把第 $k$ 轴压缩（折叠）掉”}。

先看一个 $N=3$ 个样本、$D=2$ 个特征的小例子：
\[
X = \begin{pmatrix} 1 & 10 \\ 3 & 20 \\ 5 & 30 \end{pmatrix} \in \R^{3 \times 2}.
\]
\begin{itemize}
  \item \texttt{X.mean(axis=0)}：压缩第 0 轴（把 3 个样本合为 1 个均值），输出形状从 \texttt{(3, 2)} 变为 \texttt{(2,)}，结果为 \texttt{[3.0, 20.0]}——即\textbf{每个特征的均值}。
  \item \texttt{X.mean(axis=1)}：压缩第 1 轴（把 2 个特征合为 1 个均值），输出形状从 \texttt{(3, 2)} 变为 \texttt{(3,)}，结果为 \texttt{[5.5, 11.5, 17.5]}——即\textbf{每个样本自身特征的均值}。
  \item \texttt{X.mean(axis=1, keepdims=True)}：把第 1 轴长度压成 1 而\textbf{保留轴位}，输出形状为 \texttt{(3, 1)}。这样再写 \texttt{X - X.mean(axis=1, keepdims=True)} 时，\texttt{(3, 2)} 与 \texttt{(3, 1)} 完美对齐；如果不加 \texttt{keepdims=True}，\texttt{(3, 2) - (3,)} 会直接报错！
\end{itemize}

\begin{table}[htbp]
\centering
\caption{数据矩阵 $X \in \R^{N \times D}$ 沿不同 \texttt{axis} 规约的形状与 AI 语义}\label{tab:axis-summary}
{\small
\begin{tabularx}{\linewidth}{@{}p{44mm}p{24mm}X@{}}
\toprule
\textbf{表达式} & \textbf{输出形状} & \textbf{统计与 AI 含义} \\
\midrule
\texttt{X.mean(axis=0)} & \texttt{(D,)} & 跨样本求每个特征的均值（特征标准化） \\
\texttt{X.mean(axis=0, keepdims=True)} & \texttt{(1, D)} & 同上，保留二维形状以便与 $X$ 直接相减 \\
\texttt{X.sum(axis=1, keepdims=True)} & \texttt{(N, 1)} & 对每个样本的所有类别得分求和（如 Softmax 分母） \\
\texttt{X.max(axis=1, keepdims=True)} & \texttt{(N, 1)} & 每个样本的最大 logit（数值稳定 Softmax 减最大值） \\
\texttt{X.mean(axis=None)} & 标量 \texttt{()} & 全部 $N \times D$ 个元素的全局均值 \\
\bottomrule
\end{tabularx}}
\end{table}

\subsection{广播（Broadcasting）三步规则与静默维度灾难}

在实际代码里，我们经常把形状为 \texttt{(100, 3)} 的矩阵加上形状为 \texttt{(3,)} 的偏置向量，或者把 \texttt{(4, 1)} 的列向量加上 \texttt{(1, 5)} 的行向量得到 \texttt{(4, 5)} 的网格。NumPy 允许不同形状的数组做逐元素运算，其背后的严格机制称为\textbf{广播（Broadcasting）}。

\begin{defBox}
\begin{definition}[NumPy 广播三步规则]\label{def:broadcasting}
给定两个参与逐元素运算的数组 $A$ 和 $B$，其形状分别为 $s_A$ 和 $s_B$，NumPy 按以下三步判定并执行广播：
\begin{enumerate}
  \item \textbf{左侧补 1 右对齐}：若两数组维数 \texttt{ndim} 不同，在维数较小的形状元组\textbf{左侧}补 \texttt{1}，直到两者维数相等。
  \item \textbf{逐维兼容性检查}：从左到右逐维比较两形状的轴长 $(d_A, d_B)$。每一维必须满足 $d_A = d_B$，或者其中至少一个为 $1$；否则立即抛出 \texttt{ValueError}。
  \item \textbf{长度为 1 的轴虚拟扩展}：输出数组每一维的长度为 $\max(d_A, d_B)$。凡是原长度为 1 的轴，逻辑上沿该轴重复展开至目标长度——底层实现通过将该轴步幅设为 \texttt{stride = 0} 完成，\textbf{不消耗额外内存复制}。
\end{enumerate}
\end{definition}
\end{defBox}

表~\ref{tab:broadcast-examples} 逐行演练了四组常见形状对齐过程，帮你把三步规则刻进肌肉记忆。

\begin{table}[htbp]
\centering
\caption{广播三步规则的形状对齐演练}\label{tab:broadcast-examples}
{\small
\begin{tabularx}{\linewidth}{@{}p{26mm}p{22mm}p{26mm}p{24mm}X@{}}
\toprule
\textbf{形状 $s_A$} & \textbf{形状 $s_B$} & \textbf{补 1 对齐后} & \textbf{输出形状} & \textbf{判定与典型场景} \\
\midrule
\texttt{(100, 3)} & \texttt{(3,)} & \texttt{(100, 3)} vs \texttt{(1, 3)} & \texttt{(100, 3)} & 合法：给每个样本加特征偏置 \\
\texttt{(4, 1)} & \texttt{(1, 5)} & \texttt{(4, 1)} vs \texttt{(1, 5)} & \texttt{(4, 5)} & 合法：生成两两组合网格/外和 \\
\texttt{(32, 10, 64)} & \texttt{(64,)} & \texttt{(32, 10, 64)} vs \texttt{(1, 1, 64)} & \texttt{(32, 10, 64)} & 合法：Transformer 层偏置广播 \\
\texttt{(100, 3)} & \texttt{(100,)} & \texttt{(100, 3)} vs \texttt{(1, 100)} & \textcolor{ochre}{\textbf{报错}} & 非法：末维 $3 \neq 100$ 且均不为 1 \\
\bottomrule
\end{tabularx}}
\end{table}

广播机制极其优雅，但正因为它“只要满足规则就默默执行”，一旦一维向量 \texttt{(N,)} 与二维列向量 \texttt{(N, 1)} 混用，就会制造出 AI 初学者最致命、最难排查的\textbf{静默 Bug}。

\begin{failBox}{经典静默维度灾难：\texttt{(100, 1) - (100,)} 膨胀成 \texttt{(100, 100)}}
假设你用模型预测了 100 个样本的标签，模型输出 \texttt{y\_pred} 的形状是二维列向量 \texttt{(100, 1)}，而真实标签 \texttt{y\_true} 的形状是一维数组 \texttt{(100,)}。当你写下均方误差：
\begin{lstlisting}[numbers=none]
loss = np.mean((y_pred - y_true) ** 2)   # 看似无害，实则大错特错！
\end{lstlisting}
按定义~\ref{def:broadcasting} 走一遍：\texttt{(100,)} 先左补 1 变成 \texttt{(1, 100)}，再与 \texttt{(100, 1)} 广播——两维都兼容！结果 \texttt{y\_pred - y\_true} 没有报错，而是静默生成了一个 \texttt{(100, 100)} 的成对差值矩阵（内存瞬间膨胀 100 倍），随后 \texttt{np.mean} 把全部 10,000 个跨样本错配误差平均成一个\textbf{完全错误的标量损失}！

\textbf{根治方法}：在计算损失前永远加一句 \texttt{assert y\_pred.shape == y\_true.shape}，或者统一用 \texttt{y\_true.reshape(-1)} 规范为一维向量。
\end{failBox}

\subsection{向量化性能实测：为何 NumPy 比 Python 循环快 5000 倍}

口说无凭，我们在本卷配套实验脚本 \texttt{code/vol1\_experiments.py}（函数 \texttt{exp1\_matmul\_speed}）中，对两个随机双精度矩阵 $A, B \in \R^{120 \times 120}$ 分别运行“纯 Python 三层循环”与“NumPy 矩阵乘法 \texttt{A @ B}”，在本机实测结果如下：

\begin{center}
{\small
\begin{tabular}{@{}lccc@{}}
\toprule
\textbf{实现方式（$120 \times 120$ 双精度矩阵乘）} & \textbf{单次耗时} & \textbf{相对加速比} & \textbf{结果最大绝对偏差} \\
\midrule
纯 Python 三层循环（列表推导式） & $135.7\text{ ms}$ & $1\times$ & \multirow{2}{*}{$1.78 \times 10^{-14}$} \\
NumPy 向量化矩阵乘法 \texttt{A @ B} & $0.027\text{ ms}$ & \textbf{$5{,}021\times$} & \\
\bottomrule
\end{tabular}}
\end{center}

两者计算结果的最大绝对偏差仅为 $1.78 \times 10^{-14}$（源于浮点加法结合律舍入顺序差异，属于双精度机器精度量级），但速度相差\textbf{五千余倍}！从你学过的 C 语言与计算机组成角度看，这 5000 倍加速来自三层红利叠加：
\begin{enumerate}
  \item \textbf{消除解释器开销}：Python 循环执行 $120^3 = 1.728 \times 10^6$ 次乘加时，每次都要做字节码分发、动态类型检查与对象拆装箱；NumPy 一次函数调用直接切入底层编译好的 C/Fortran 例程。
  \item \textbf{缓存分块与连续内存}：底层 BLAS 库的 \texttt{dgemm} 内核将大矩阵拆成贴合 CPU L1/L2 缓存大小的小块计算，大幅减少主存搬运。
  \item \textbf{SIMD 指令与多核并行}：单条 CPU 向量指令（如 AVX/NEON）可在同一时钟周期内同时完成 4--8 组双精度乘加（FMA），大矩阵还会自动启用多核线程并行。
\end{enumerate}

\begin{insightBox}{下游坐标：从 CPU 向量化到 GPU 批处理}
你在本章养成的“消灭 Python 循环、尽量表达成大块连续矩阵运算”的习惯，不仅适用于本卷第 2--6 章与第 2 卷的纯 NumPy 建模，在第 5--6 卷切换到 PyTorch 与 GPU 时依然是同一套底层逻辑——GPU 的 Tensor Core 本质上就是把“大块矩阵乘”做到极致的硬件加速器。
\end{insightBox}

\section{矩阵乘法的四种视角（AI 读论文的母语）}\label{sec:matmul-four-views}

\begin{recallBox}{矩阵乘法的标准定义}
设 $A \in \R^{m \times p}$，$B \in \R^{p \times n}$，则乘积 $C = AB \in \R^{m \times n}$ 的第 $(i, j)$ 个元素定义为
\[
c_{ij} = (AB)_{ij} = \sum_{k=1}^{p} a_{ik} b_{kj}, \qquad i=1,\dots,m,\; j=1,\dots,n.
\]
\end{recallBox}

大一线性代数通常只把上面的公式当作计算口诀来用。然而在阅读人工智能论文和写模型代码时，如果你脑子里只有“第 $i$ 行乘第 $j$ 列”，看到线性层 $Wx$、注意力加权 $PV$、奇异值分解 $\sum \sigma_i u_i v_i^T$ 或反向传播梯度 $\delta x^T$ 时，往往会觉得它们是毫无关系的神秘符号。

事实上，它们只是\textbf{同一个矩阵乘法 $AB$ 的四种等价分解视角}。我们先用图~\ref{fig:matmul-four-views} 一眼看清这四种视角的几何/结构画面，再用同一个 $2 \times 2$ 数值小例逐一验算。

\begin{figure}[htbp]
\centering
\resizebox{\linewidth}{!}{%
\begin{tikzpicture}[
  font=\sffamily\footnotesize,
  mat/.style={draw=ink, thick, fill=white, rounded corners=0.6mm},
  hiA/.style={fill=tealLight, draw=teal, thick},
  hiB/.style={fill=navyLight, draw=navy, thick},
  hiC/.style={fill=ochreLight, draw=ochre, thick},
  panel/.style={draw=tealLine, rounded corners=2mm, fill=white, line width=0.6pt}
]

% ---------- Panel 1: Dot Product ----------
\begin{scope}[shift={(0, 0)}]
  \draw[panel] (0,0) rectangle (7.4, 4.45);
  \node[anchor=north west, font=\sffamily\bfseries\small, text=tealDark] at (0.22, 4.25) {(1) 点积视角：逐元素内积填充};
  % Matrix A
  \draw[mat] (0.45, 1.48) rectangle (1.95, 3.28);
  \draw[hiA] (0.45, 2.48) rectangle (1.95, 2.83);
  \node[text=tealDark] at (1.20, 2.655) {$a_{i\cdot}$};
  \node[below=2pt, text=slate] at (1.20, 1.44) {$A\,(m\times p)$};
  % Matrix B
  \draw[mat] (2.35, 1.48) rectangle (3.85, 3.28);
  \draw[hiB] (3.00, 1.48) rectangle (3.35, 3.28);
  \node[text=navy] at (3.175, 2.38) {$b_{\cdot j}$};
  \node[below=2pt, text=slate] at (3.10, 1.44) {$B\,(p\times n)$};
  % =
  \node[font=\bfseries] at (4.25, 2.38) {$=$};
  % Matrix C
  \draw[mat] (4.65, 1.48) rectangle (6.15, 3.28);
  \draw[hiC] (5.30, 2.48) rectangle (5.65, 2.83);
  \node[text=ochre] at (5.475, 2.655) {\tiny $c_{ij}$};
  \node[below=2pt, text=slate] at (5.40, 1.44) {$C\,(m\times n)$};
  % Formula
  \node[anchor=south, text=ink] at (3.7, 0.16) {$c_{ij} = \inner{a_{i\cdot}}{b_{\cdot j}} = \sum_{k=1}^p a_{ik} b_{kj}$};
\end{scope}

% ---------- Panel 2: Column Linear Combination ----------
\begin{scope}[shift={(7.9, 0)}]
  \draw[panel] (0,0) rectangle (7.4, 4.45);
  \node[anchor=north west, font=\sffamily\bfseries\small, text=tealDark] at (0.22, 4.25) {(2) 列线性组合视角：右乘向量混合列};
  % Matrix A with columns
  \draw[mat] (0.45, 1.48) rectangle (2.25, 3.28);
  \draw[hiA] (0.55, 1.56) rectangle (1.00, 3.20);
  \draw[hiA] (1.13, 1.56) rectangle (1.58, 3.20);
  \draw[hiA] (1.71, 1.56) rectangle (2.16, 3.20);
  \node[text=tealDark] at (0.775, 2.38) {\scriptsize $a_{\cdot 1}$};
  \node[text=tealDark] at (1.355, 2.38) {\scriptsize $a_{\cdot 2}$};
  \node[text=tealDark] at (1.935, 2.38) {\scriptsize $a_{\cdot p}$};
  \node[below=2pt, text=slate] at (1.35, 1.44) {$A$ 的列字典};
  % Vector x
  \draw[hiB] (2.60, 1.68) rectangle (3.10, 3.08);
  \node[text=navy] at (2.85, 2.78) {\scriptsize $x_1$};
  \node[text=navy] at (2.85, 2.38) {\scriptsize $\vdots$};
  \node[text=navy] at (2.85, 1.95) {\scriptsize $x_p$};
  \node[below=2pt, text=slate] at (2.85, 1.44) {$x$};
  % =
  \node[font=\bfseries] at (3.48, 2.38) {$=$};
  % Output
  \node[anchor=west, text=ink] at (3.78, 2.38) {$x_1 \textcolor{teal}{a_{\cdot 1}} + \dots + x_p \textcolor{teal}{a_{\cdot p}}$};
  \node[anchor=south, text=ink] at (3.7, 0.16) {$Ax \in \col(A)$；$AB$ 的第 $j$ 列 $= A b_{\cdot j}$};
\end{scope}

% ---------- Panel 3: Row Linear Combination ----------
\begin{scope}[shift={(0, -4.90)}]
  \draw[panel] (0,0) rectangle (7.4, 4.45);
  \node[anchor=north west, font=\sffamily\bfseries\small, text=tealDark] at (0.22, 4.25) {(3) 行线性组合视角：左乘向量加权行};
  % Row vector x^T
  \draw[hiB] (0.40, 2.18) rectangle (1.95, 2.63);
  \node[text=navy] at (1.175, 2.405) {\scriptsize $x_1 \;\cdots\; x_m$};
  \node[below=2pt, text=slate] at (1.175, 2.02) {$x^T$};
  % Matrix B with rows
  \draw[mat] (2.25, 1.48) rectangle (3.95, 3.28);
  \draw[hiA] (2.33, 2.77) rectangle (3.87, 3.19);
  \draw[hiA] (2.33, 2.17) rectangle (3.87, 2.59);
  \draw[hiA] (2.33, 1.57) rectangle (3.87, 1.99);
  \node[text=tealDark] at (3.10, 2.98) {\scriptsize $b_{1\cdot}$};
  \node[text=tealDark] at (3.10, 2.38) {\scriptsize $b_{2\cdot}$};
  \node[text=tealDark] at (3.10, 1.78) {\scriptsize $b_{m\cdot}$};
  \node[below=2pt, text=slate] at (3.10, 1.44) {$B$ 的行样本};
  % =
  \node[font=\bfseries] at (4.30, 2.38) {$=$};
  % Output
  \node[anchor=west, text=ink, font=\scriptsize] at (4.49, 2.38) {$x_1 \textcolor{teal}{b_{1\cdot}} + \dots + x_m \textcolor{teal}{b_{m\cdot}}$};
  \node[anchor=south, text=ink] at (3.7, 0.16) {$x^T B \in \row(B)$（Attention 权重加权 Value 行）};
\end{scope}

% ---------- Panel 4: Outer Product Sum ----------
\begin{scope}[shift={(7.9, -4.90)}]
  \draw[panel] (0,0) rectangle (7.4, 4.45);
  \node[anchor=north west, font=\sffamily\bfseries\small, text=tealDark] at (0.22, 4.25) {(4) 外积秩 1 求和视角：低秩结构分解};
  % Sum of rank-1
  \node[anchor=west, text=ink] at (0.35, 2.38) {$AB = \displaystyle\sum_{k=1}^{p}$};
  % Col vector
  \draw[hiA] (1.95, 1.58) rectangle (2.40, 3.18);
  \node[text=tealDark] at (2.175, 2.38) {\scriptsize $a_{\cdot k}$};
  % Row vector
  \draw[hiB] (2.55, 2.78) rectangle (4.05, 3.18);
  \node[text=navy] at (3.30, 2.98) {\scriptsize $b_{k\cdot}$};
  % Equals rank-1 block
  \node[font=\bfseries] at (4.42, 2.38) {$=$};
  \draw[hiC] (4.78, 1.63) rectangle (6.12, 3.13);
  \node[align=center, text=ochre] at (5.45, 2.38) {\scriptsize 秩 $\le 1$\\\scriptsize 矩阵块};
  \node[right=1pt, text=ink] at (6.16, 2.38) {$+\cdots$};
  \node[anchor=south, text=ink] at (3.7, 0.16) {$AB = \sum_{k=1}^p a_{\cdot k} b_{k\cdot}$（SVD、LoRA、反传外积梯度）};
\end{scope}

\end{tikzpicture}}
\caption{矩阵乘法的四种互补视角：同一种运算，在 AI 中对应四种截然不同的建模语义}\label{fig:matmul-four-views}
\end{figure}

下面我们用同一对 $2 \times 2$ 矩阵手算验证这四种视角，记
\[
A = \begin{pmatrix} 2 & 1 \\ 0 & 3 \end{pmatrix}, \qquad
B = \begin{pmatrix} 1 & 4 \\ 2 & -1 \end{pmatrix}.
\]

\begin{enumerate}
  \item \textbf{视角一：点积视角（逐元素内积）}。第 $(i,j)$ 个元素等于 $A$ 的第 $i$ 行与 $B$ 的第 $j$ 列做点积：
  \[
  AB = \begin{pmatrix}
    2\cdot 1 + 1\cdot 2 & 2\cdot 4 + 1\cdot(-1) \\
    0\cdot 1 + 3\cdot 2 & 0\cdot 4 + 3\cdot(-1)
  \end{pmatrix}
  = \begin{pmatrix} 4 & 7 \\ 6 & -3 \end{pmatrix}.
  \]
  在 AI 中，这一视角用于理解\textbf{相似度打分矩阵}：若 $Q$ 的每一行是一个查询向量、$K$ 的每一行是一个键向量，则 $QK^T$ 的第 $(i,j)$ 个元素就是第 $i$ 个查询与第 $j$ 个键的点积相似度。

  \item \textbf{视角二：列线性组合视角（右乘列向量混合 $A$ 的列）}。把 $A$ 写成列向量组 $A = [a_{\cdot 1},\, a_{\cdot 2}]$，则 $AB$ 的第 1 列和第 2 列分别是 $A$ 各列的线性组合，组合系数由 $B$ 对应列给出：
  \[
  (AB)_{\cdot 1} = 1\begin{pmatrix}2\\0\end{pmatrix} + 2\begin{pmatrix}1\\3\end{pmatrix} = \begin{pmatrix}4\\6\end{pmatrix}, \qquad
  (AB)_{\cdot 2} = 4\begin{pmatrix}2\\0\end{pmatrix} + (-1)\begin{pmatrix}1\\3\end{pmatrix} = \begin{pmatrix}7\\-3\end{pmatrix}.
  \]
  这一视角直接告诉你一个深刻事实：\textbf{对任意向量 $x$，$Ax$ 永远落在 $A$ 的列空间 $\col(A)$ 中}。第 2 章做线性最小二乘 $\min_w \|Aw - b\|_2^2$ 时，正是靠这一视角把它看成“在列空间 $\col(A)$ 中找离观测向量 $b$ 最近的点”！

  \item \textbf{视角三：行线性组合视角（左乘行向量加权 $B$ 的行）}。把 $B$ 写成行向量组 $B = \begin{psmallmatrix} b_{1\cdot} \\ b_{2\cdot} \end{psmallmatrix}$，则 $AB$ 的第 1 行和第 2 行分别是 $B$ 各行的线性组合，权重由 $A$ 对应行给出：
  \[
  (AB)_{1\cdot} = 2\begin{pmatrix}1 & 4\end{pmatrix} + 1\begin{pmatrix}2 & -1\end{pmatrix} = \begin{pmatrix}4 & 7\end{pmatrix}, \qquad
  (AB)_{2\cdot} = 0\begin{pmatrix}1 & 4\end{pmatrix} + 3\begin{pmatrix}2 & -1\end{pmatrix} = \begin{pmatrix}6 & -3\end{pmatrix}.
  \]
  在 Transformer 自注意力机制 $\operatorname{Attention}(Q,K,V) = P V$ 中（第 5 卷），$V \in \R^{T \times d_v}$ 的每一行是一个词元的值向量，$P \in \R^{T \times T}$ 的第 $i$ 行是经过 Softmax 归一化的注意力概率分布。按行组合视角看，\textbf{输出的第 $i$ 行就是用第 $i$ 行的注意力权重对所有词元值向量做加权平均}！

  \item \textbf{视角四：外积秩 1 求和视角（列乘行分解）}。把乘积拆成 $A$ 的第 $k$ 列与 $B$ 的第 $k$ 行的\textbf{外积（Outer Product）}之和：
  \[
  AB = a_{\cdot 1} b_{1\cdot} + a_{\cdot 2} b_{2\cdot}
  = \begin{pmatrix}2\\0\end{pmatrix}\begin{pmatrix}1 & 4\end{pmatrix}
  + \begin{pmatrix}1\\3\end{pmatrix}\begin{pmatrix}2 & -1\end{pmatrix}
  = \begin{pmatrix}2 & 8\\0 & 0\end{pmatrix}
  + \begin{pmatrix}2 & -1\\6 & -3\end{pmatrix}
  = \begin{pmatrix}4 & 7\\6 & -3\end{pmatrix}.
  \]
  注意观察等式右边的每一项 $u v^T$（$u \in \R^m, v \in \R^n$）：它的每一列都是 $u$ 的标量倍，每一行都是 $v^T$ 的标量倍，因此任何非零外积 $uv^T$ 的秩严格等于 $1$。由此顺带看清一个重要结论：$m \times p$ 矩阵乘以 $p \times n$ 矩阵，本质上是 $p$ 个秩 1 矩阵的叠加，故 $\rank(AB) \le p$。
\end{enumerate}

\begin{insightBox}{为什么外积秩 1 视角是 AI 最值钱的视角？}
至少有三大核心算法直接建立在外积视角之上：
\begin{itemize}
  \item \textbf{奇异值分解与图像压缩（本卷第 4 章）}：把大矩阵写成 $A = \sum_{i=1}^r \sigma_i u_i v_i^T$，只保留前 $k$ 个最大奇异值对应的秩 1 块即可高精度压缩。
  \item \textbf{反向传播梯度（本卷第 5 章与第 5 卷）}：线性层 $y = Wx$ 的权重梯度恰好是外积 $\frac{\partial L}{\partial W} = \delta x^T$；整个小批次 $B$ 个样本的总梯度则是 $\Delta^T X = \sum_{n=1}^B \delta_n x_n^T$，一步矩阵乘汇总所有样本的外积！
  \item \textbf{大模型 LoRA 微调（第 6 卷）}：不直接更新巨大的 $d \times d$ 权重增量 $\Delta W$，而是参数化为低秩乘积 $\Delta W = B A$（$B \in \R^{d \times r},\, A \in \R^{r \times d},\, r \ll d$），本质上是用 $r$ 个秩 1 块逼近参数更新，可训练参数量从 $d^2$ 骤降至 $2dr$。
\end{itemize}
\end{insightBox}

除了上述标准矩阵乘法外，AI 中还频繁使用\textbf{哈达玛积（Hadamard Product，逐元素乘积）}，记作 $A \had B$。对同形状矩阵 $A, B \in \R^{m \times n}$，定义
\[
(A \had B)_{ij} = a_{ij} b_{ij}.
\]
它在深度学习中无处不在：Dropout 掩码 $h \had m$（把随机置 0 的神经元屏蔽掉）、现代大模型（如 LLaMA）的 SwiGLU 门控层 $\operatorname{SiLU}(x W_1) \had (x W_2)$、以及逐元素特征缩放。请牢记 NumPy 运算符的铁律：\textbf{\texttt{A * B} 是逐元素哈达玛积 $A \had B$，而 \texttt{A @ B} 才是矩阵乘法 $AB$}，二者绝不可混淆！

\section{批处理高维张量乘、\texttt{np.einsum} 与基础可视化}\label{sec:batch-einsum}

\subsection{批处理高维张量与 \texttt{@} 的高维规则}

真实 AI 任务很少一次只算单个样本或单个注意力头。为了喂饱现代硬件，我们会把样本批量、时间步、注意力头打包成三维或四维张量（多维 \texttt{ndarray}）：
\begin{itemize}
  \item \textbf{序列批处理张量} \texttt{(B, T, D)}：$B$ 个样本（Batch）、每个样本 $T$ 个词元（Tokens）、每个词元用 $D$ 维向量表示。
  \item \textbf{多头注意力张量} \texttt{(B, H, T, D\_h)}：把 $D$ 维拆成 $H$ 个独立注意力头（Heads），每头维度 $D_h = D/H$。
\end{itemize}

面对高维张量，NumPy 的 \texttt{A @ B} 遵循一条异常清晰的规则：\textbf{把最后两维 \texttt{(-2, -1)} 视作待相乘的二维矩阵，前面的所有批次维度 \texttt{[:-2]} 按广播三步规则（定义~\ref{def:broadcasting}）对齐}。

看两个深度学习中最常见的例子：
\begin{enumerate}
  \item \textbf{批量线性映射}：输入 \texttt{X} 形状为 \texttt{(B, T, D)}，共享权重 \texttt{W} 形状为 \texttt{(D, M)}。执行 \texttt{X @ W} 时，\texttt{W} 的前置维度自动广播为 \texttt{(1, D, M)}，与每个样本的 \texttt{(T, D)} 矩阵相乘，得到 \texttt{(B, T, M)}——等价于对每个样本中的每个词元独立施加同一线性变换 $W$。
  \item \textbf{多头注意力打分矩阵}：查询 \texttt{Q} 形状为 \texttt{(B, H, T, D\_h)}，键 \texttt{K} 形状为 \texttt{(B, H, S, D\_h)}。先令 \texttt{Kt = K.swapaxes(-1, -2)} 把最后两维转置成 \texttt{(B, H, D\_h, S)}，再计算 \texttt{Q @ Kt}，便一次性并行算完所有 $B \times H$ 组注意力打分矩阵，输出形状为 \texttt{(B, H, T, S)}！
\end{enumerate}

\subsection{爱因斯坦求和约定与 \texttt{np.einsum}}

当张量超过三维、涉及多处转置和求和时，连续写 \texttt{swapaxes} 和 \texttt{@} 往往让人眼花缭乱。NumPy 提供的 \texttt{np.einsum}（爱因斯坦求和约定）允许你\textbf{直接把下标字母写出来}，代码几乎就是数学公式本身。

阅读与书写 \texttt{np.einsum("输入下标->输出下标", ...)} 只需掌握两条规则：
\begin{enumerate}
  \item \textbf{出现在箭头左边、但未出现在箭头右边的字母}：表示沿该下标求和（缩并，contraction）。
  \item \textbf{同时出现在箭头左右两边的字母}：表示保留在输出张量中，并按箭头右边的顺序排列轴位。
\end{enumerate}

表~\ref{tab:einsum-cheatsheet} 整理了从本卷到第 6 卷最高频的六种 \texttt{np.einsum} 写法。

\begin{table}[htbp]
\centering
\caption{AI 高频张量运算的 \texttt{np.einsum} 速查表}\label{tab:einsum-cheatsheet}
{\small
\begin{tabularx}{\linewidth}{@{}p{28mm}p{38mm}p{32mm}X@{}}
\toprule
\textbf{运算名称} & \textbf{\texttt{np.einsum} 写法} & \textbf{数学表达式} & \textbf{典型用途与说明} \\
\midrule
向量点积 & \texttt{'i,i->', u, v} & $\sum_i u_i v_i$ & 内积、余弦相似度分子 \\
向量外积 & \texttt{'i,j->ij', u, v} & $(uv^T)_{ij} = u_i v_j$ & 秩 1 矩阵构造、单样本梯度 \\
标准矩阵乘 & \texttt{'ik,kj->ij', A, B} & $\sum_k A_{ik} B_{kj}$ & 等价于 \texttt{A @ B} \\
Frobenius 内积 & \texttt{'ij,ij->', A, B} & $\sum_{i,j} A_{ij} B_{ij}$ & 等于 $\tr(A^T B)$（第 5 章矩阵微分） \\
批量二次型 & \texttt{'ni,ij,nj->n', X, Q, X} & $x_n^T Q x_n$ & 一次性计算 $N$ 个样本的二次损失 \\
多头注意力打分 & \texttt{'bhtd,bhsd->bhts', Q, K} & $\sum_d Q_{b,h,t,d} K_{b,h,s,d}$ & 无需手动 \texttt{swapaxes}，可读性极佳 \\
\bottomrule
\end{tabularx}}
\end{table}

\begin{remark}[多张量缩并的性能提示]
当 \texttt{np.einsum} 同时缩并三个或更多张量时（例如批量二次型表达式 \texttt{'ni,ij,nj->n'}），默认从左到右逐对计算可能产生不必要的巨大中间张量。养成加上参数 \texttt{optimize=True} 的习惯，NumPy 会自动搜索最优加括号顺序，大幅降低计算量与内存峰值。
\end{remark}

\subsection{完整可运行代码范例与基础可视化模板}

代码~\ref{lst:ch1-demo} 整合了本章全部核心技法：校验矩阵乘法四视角等价、防静默广播残差检查、对比 \texttt{@} 与 \texttt{np.einsum} 计算多头注意力打分，并给出本系列通用的最小 \texttt{matplotlib} 绘图骨架。

\begin{lstlisting}[caption={矩阵乘法四视角、多头注意力 \texttt{einsum} 验证与最小可视化模板}, label={lst:ch1-demo}]
import numpy as np
import matplotlib.pyplot as plt

# 1. 矩阵乘法四视角数值一致性校验
A = np.array([[2.0, 1.0], [0.0, 3.0]])
B = np.array([[1.0, 4.0], [2.0, -1.0]])
C_dot = A @ B
C_outer = sum(np.outer(A[:, k], B[k, :]) for k in range(A.shape[1]))
assert np.allclose(C_dot, C_outer, atol=1e-12)

# 2. 防止静默广播的安全均方误差函数
def safe_mse(y_pred: np.ndarray, y_true: np.ndarray) -> float:
    assert y_pred.shape == y_true.shape, (
        f"形状不一致将触发静默广播：{y_pred.shape} vs {y_true.shape}"
    )
    return float(np.mean((y_pred - y_true) ** 2))

# 3. 批处理多头注意力打分：@ 与 np.einsum 完全等价
rng = np.random.default_rng(2026)
B_sz, H, T, S, Dh = 2, 4, 8, 6, 16
Q = rng.normal(size=(B_sz, H, T, Dh))
K = rng.normal(size=(B_sz, H, S, Dh))
scores_at = (Q @ K.swapaxes(-1, -2)) / np.sqrt(Dh)
scores_ein = np.einsum("bhtd,bhsd->bhts", Q, K) / np.sqrt(Dh)
assert scores_at.shape == (B_sz, H, T, S)
assert np.allclose(scores_at, scores_ein, atol=1e-12)

# 4. 最小科研绘图模板（面向对象 Axes 接口）
x = np.linspace(0.0, 2.0 * np.pi, 200)
fig, ax = plt.subplots(figsize=(5.5, 3.2), dpi=150)
ax.plot(x, np.sin(x), color="#0F5C5B", linewidth=2.0, label=r"$\sin(x)$")
ax.plot(x, np.cos(x), color="#9A4624", linestyle="--", label=r"$\cos(x)$")
ax.set_xlabel("x"); ax.set_ylabel("y"); ax.grid(alpha=0.25); ax.legend()
fig.tight_layout()
# fig.savefig("trig_demo.png")
\end{lstlisting}

\begin{extBox}{步幅 \texttt{strides} 零拷贝滑窗直觉（\texttt{sliding\_window\_view}）}
前面提到，广播底层靠将步幅 \texttt{stride} 设为 0 实现零内存拷贝。更进一步，NumPy 提供了 \texttt{np.lib.stride\_tricks.sliding\_window\_view(x, window\_shape)}：对一维时间序列或二维图像，只需巧妙调整头部 \texttt{shape} 与 \texttt{strides}，就能零拷贝生成所有滑动窗口视图——这正是第 5 卷卷积神经网络（CNN）底层 \texttt{im2col} 向量化加速的思想雏形。初读只需知道其存在，二刷或写手写卷积时再深入尝试即可。
\end{extBox}

\section{本章习题}\label{sec:ch01-exercises}

\exTier{A. 概念辨析与反例查错}{先独立思考并定位陷阱，切忌直接翻答案}

\begin{enumerate}[label=\textbf{1.\arabic*}]
  \item \textbf{广播规则与静默维度错误诊断}：
  （1）根据广播三步规则，逐一判断以下四组形状能否合法广播；若合法，写出输出形状，若非法，说明在哪一维冲突：
  \begin{center}
  (a) \texttt{(8, 1, 5)} 与 \texttt{(4, 5)}；\quad
  (b) \texttt{(8, 3, 5)} 与 \texttt{(3, 1)}；\quad
  (c) \texttt{(10, 4)} 与 \texttt{(10,)}；\quad
  (d) \texttt{(64, 1)} 与 \texttt{(64,)}。
  \end{center}
  （2）设某回归模型在 64 个样本上的预测值为 \texttt{y\_pred}（形状 \texttt{(64, 1)}），真实标签为 \texttt{y\_true}（形状 \texttt{(64,)}）。解释直接执行 \texttt{np.mean((y\_pred - y\_true)**2)} 会发生什么数值错误，并给出两种安全的修复写法。

  \item \textbf{\texttt{view/copy} 与矩阵乘法运算符排错}：阅读下面这段试图实现“归一化副本并计算矩阵乘积 $AB$”的 6 行代码，在不运行机器的前提下，逐行写出 \texttt{A[0, 0]} 的最终值与 \texttt{C} 的计算结果，指出其中隐藏的\textbf{两处 Bug} 并改正：
  \begin{lstlisting}[numbers=none]
A = np.array([[2.0, 4.0], [1.0, 3.0]])
B = np.array([[1.0, 0.0], [0.0, 2.0]])
A_scaled = A[:]          # Bug 1?
A_scaled /= 2.0
C = A_scaled * B         # Bug 2?
print(A[0, 0], C)
  \end{lstlisting}
\end{enumerate}

\exTier{B. 手算小例与基础演练}{亲手算一遍，把概念变成直觉}

\begin{enumerate}[label=\textbf{1.\arabic*}, resume]
  \item \textbf{矩阵乘法四视角手算验证}：设
  \[
  A = \begin{pmatrix} 1 & 2 \\ 0 & -1 \\ 3 & 1 \end{pmatrix}, \qquad
  B = \begin{pmatrix} 2 & -1 \\ 1 & 3 \end{pmatrix}.
  \]
  请分别使用（1）\textbf{列线性组合视角}、（2）\textbf{行线性组合视角}和（3）\textbf{外积秩 1 求和视角}计算乘积 $AB$，验证三者结果完全一致，并说明外积展开中每一项的矩阵秩是多少。

  \item \textbf{\texttt{axis} 规约与 \texttt{np.einsum} 翻译}：
  （1）设数据矩阵 $X = \begin{psmallmatrix} 2 & 10 & 4 \\ 6 & 20 & 4 \end{psmallmatrix}$，手算 \texttt{X.mean(axis=0, keepdims=True)}、\texttt{X.std(axis=0, keepdims=True)} 以及按列中心化结果 \texttt{X - X.mean(axis=0, keepdims=True)}。
  （2）将以下三个 \texttt{np.einsum} 表达式翻译为等价的数学求和式和常规 NumPy 写法（使用 \texttt{@}、\texttt{*} 或 \texttt{sum}），并写出输出形状（设 $A:(m,n)$，$x:(n,)$，$U,V:(N,D)$，$H:(B,T,D)$，$W:(D,M)$）：
  \begin{center}
  (a) \texttt{np.einsum('ij,j->i', A, x)}；\quad
  (b) \texttt{np.einsum('ni,ni->n', U, V)}；\quad
  (c) \texttt{np.einsum('btd,dm->btm', H, W)}。
  \end{center}
\end{enumerate}

\exTier{C. 实用证明与理论推导}{紧扣 AI 核心性质，无冷门偏题}

\begin{enumerate}[label=\textbf{1.\arabic*}, resume]
  \item \textbf{秩 1 外积矩阵的基本性质}：设非零向量 $u \in \R^m$，$v \in \R^n$，令 $M = u v^T \in \R^{m \times n}$。证明：
  （1）$\rank(M) = 1$；
  （2）当 $m = n$ 时，矩阵的迹满足 $\tr(u v^T) = u^T v$；
  （3）对任意矩阵 $A \in \R^{m \times p}$ 与 $B \in \R^{p \times n}$，利用外积求和视角证明 $\rank(AB) \le p$。
  \exHint{第（1）问观察 $M$ 的每一列均为 $u$ 的倍数且至少有一列非零；第（2）问直接展开对角元 $M_{ii} = u_i v_i$ 求和；第（3）问利用秩的次可加性 $\rank(X + Y) \le \rank(X) + \rank(Y)$。}

  \item \textbf{低秩分解的乘法加括号顺序与 LoRA 计算量}：设大型权重更新矩阵采用低秩参数化 $\Delta W = B A$，其中 $B \in \R^{d \times r}$、$A \in \R^{r \times d}$ 且 $r \ll d$。给定单个输入向量 $x \in \R^d$，比较两种计算顺序：方案一先显式合成 $M = BA$ 再算 $Mx$；方案二先算 $z = Ax$ 再算 $Bz$。分别计算两种方案所需的标量乘法次数，并在 $d = 4096,\, r = 16$ 时给出具体加速比。
  \exHint{一个 $p \times q$ 矩阵乘以 $q \times s$ 矩阵恰好需要 $p q s$ 次标量乘法；注意向量可看作 $d \times 1$ 矩阵。}
\end{enumerate}

\exTier{D. 数值编程与可自检实验}{独立跑通并让全部 \texttt{assert} 通过}

\begin{enumerate}[label=\textbf{1.\arabic*}, resume]
  \item \textbf{零循环实现成对欧氏距离矩阵}：给定两组点集 $X \in \R^{N \times D}$ 与 $Y \in \R^{M \times D}$，要求计算成对欧氏距离矩阵 $D \in \R^{N \times M}$，其中 $D_{ij} = \|x_i - y_j\|_2$。请分别实现以下两种零循环方法：
  \begin{itemize}
    \item 方法一（三维广播法）：通过 \texttt{X[:, None, :] - Y[None, :, :]} 直接广播求差再沿末轴求范数；
    \item 方法二（展开恒等式法）：利用 $\|x_i - y_j\|_2^2 = \|x_i\|_2^2 + \|y_j\|_2^2 - 2 x_i^T y_j$，借助矩阵乘法 \texttt{X @ Y.T} 与广播一次性算出平方距离，再用 \texttt{np.maximum(..., 0.0)} 截断浮点舍入产生的微小负数后开方。
  \end{itemize}
  \textbf{自检标准}：使用 \texttt{rng = np.random.default\_rng(0)} 生成 $N=120,\, M=80,\, D=32$ 的标准正态矩阵，将两种方法与双重循环基准结果对比，验证最大绝对误差均满足 \texttt{max\_abs\_err < 1e-12}，并比较两种方法在 $N=M=1000,\, D=64$ 时的耗时差异。
  \exHint{方法二中令 \texttt{x2 = np.sum(X**2, axis=1, keepdims=True)}（形状 \texttt{(N, 1)}）与 \texttt{y2 = np.sum(Y**2, axis=1)[None, :]}（形状 \texttt{(1, M)}），即可安全广播相加。}

  \item \textbf{多头缩放点积注意力的纯 NumPy 实现与校验}：编写函数 \texttt{multihead\_attention(Q, K, V)}，其中输入 \texttt{Q, K, V} 形状均为 \texttt{(B, H, T, D\_h)}，按以下公式计算输出 $O \in \R^{B \times H \times T \times D_h}$：
  \[
  S = \frac{Q K^T}{\sqrt{D_h}}, \qquad
  P = \operatorname{softmax}_{\text{行}}(S), \qquad
  O = P V.
  \]
  要求：（1）行 Softmax 必须先减去每行最大值（\texttt{keepdims=True}）防止 \texttt{np.exp} 溢出；（2）分别用高维 \texttt{@} 与 \texttt{np.einsum} 各实现一次。
  \textbf{自检标准}：使用 \texttt{rng = np.random.default\_rng(7)} 在 $B=2,\, H=4,\, T=16,\, D_h=8$ 上测试，验证：（a）输出形状严格为 \texttt{(2, 4, 16, 8)}；（b）两种实现的最大绝对误差 \texttt{< 1e-12}；（c）注意力权重矩阵 $P$ 每一行之和与 $1$ 的最大偏差 \texttt{< 1e-12}，且将 $S$ 整体加上常数 $1000$ 后 $P$ 保持不变。
  \exHint{利用 \texttt{np.einsum} 时，打分步用 \texttt{'bhtd,bhsd->bhts'}，加权求和步用 \texttt{'bhts,bhsd->bhtd'}。}
\end{enumerate}

\chapter[范数、高维几何与线性最小二乘]{向量范数、高维几何、\\正交投影与线性最小二乘}\label{ch:norms-proj-ls}

\chapEpigraph{数学是赋予不同事物以相同名称的艺术。}{昂利·庞加莱（Henri Poincar\'e），《科学与方法》}

\begin{chapIntro}{本章导学：前置衔接、学习目标与首个真实模型}
\textbf{前置衔接}：衔接第 1 章的矩阵乘法列空间视角；本章开篇内置\textbf{零障碍线性代数前置精华}（向量空间、张成、基、线性变换、秩与零空间、行列式体积直觉），完全自洽且\textbf{不依赖第 3 章的特征值与谱定理}。

\textbf{学习目标}：建立清晰的空间、基、线性变换与行列式几何直觉；掌握 $\ell_1/\ell_2/\ell_\infty$ 范数几何与高维近正交现象；严格掌握正交分解、线性代数四大基本子空间与“超定 vs 欠定”二象性；从三点手算直线出发，完整证明正交投影最优逼近定理、正规方程可逆判据与 Gram-Schmidt/QR 数值求解，并亲手做出你在本套讲义中的\textbf{第一个真实机器学习模型}——多项式最小二乘与岭回归！

\textbf{下游坐标}：本章是本卷第 3--4 章谱定理与 SVD/PCA 的几何底座，也是第 2 卷高斯极大似然与线性回归、第 4 卷 $L_1/L_2$ 正则化与偏差-方差权衡、第 5--6 卷向量检索（Embedding）与梯度裁剪的直接源头。
\end{chapIntro}

第 1 章把数据与参数放进了矩阵，但要真正看懂高维空间里的数据结构与模型求解，我们还需要两套底层直觉：先看清“向量张成的空间、基、线性变换、秩与行列式”的几何本质（第~\ref{sec:linalg-bridge}~节），再拿起衡量距离与夹角的尺子（范数与内积），一路走向高维近正交、正交投影，最终完成首个数据拟合与抗过拟合实验闭环。

\section{零障碍线性代数前置精华：空间、变换与行列式的几何直觉}\label{sec:linalg-bridge}

如果你此前尚未系统修完线性代数，或只记得“解方程、算行列式”的代数步骤，本节将用四段紧凑而直观的几何画面，帮你一次打通后续全书反复依赖的全部线性代数底层语言；若你已熟悉这些概念，也可把本节当作快速校准“AI 几何直觉”的热身。

\subsection{向量空间、线性组合与张成空间（Span）}

在人工智能与几何的双重视野里，一个 $n$ 维实向量 $v = (v_1, v_2, \dots, v_n)^T \in \R^n$ 同时扮演两个角色：在代码中，它是长度为 $n$ 的特征数组；在几何中，它是\textbf{从原点出发、指向点 $(v_1, \dots, v_n)$ 的一支有向箭头}。向量空间 $\R^n$ 中最基本的动作只有两个——\textbf{数乘伸缩}（$\alpha v$）与\textbf{首尾相加}（$u + v$）。把这两种动作合在一起，就得到了线性代数最核心的构造：

\begin{defBox}
\begin{definition}[线性组合、张成空间与线性子空间]\label{def:span-subspace}
给定一组向量 $v_1, v_2, \dots, v_k \in \R^n$：
\begin{enumerate}
  \item 形如 $c_1 v_1 + c_2 v_2 + \dots + c_k v_k$（其中标量 $c_1, \dots, c_k \in \R$ 任意取值）的向量称为它们的\textbf{线性组合（Linear Combination）}。
  \item 所有可能线性组合汇成的集合称为这组向量的\textbf{张成空间（Span）}：
  \[
    \spanop\{v_1, \dots, v_k\} \coloneqq \left\{\sum_{i=1}^{k} c_i v_i \;:\; c_1, \dots, c_k \in \R\right\}.
  \]
  \item 若非空子集 $U \subseteq \R^n$ 包含零向量 $0$，且对加法（$u, w \in U \Rightarrow u+w \in U$）与数乘（$u \in U, \alpha \in \R \Rightarrow \alpha u \in U$）均封闭，则称 $U$ 为 $\R^n$ 的一个\textbf{线性子空间（Subspace）}。
\end{enumerate}
\end{definition}
\end{defBox}

把“张成空间”想象成你的\textbf{步幅工具箱}：手里每多一支向量，就多了一种可以任意伸缩行走的方向。
\begin{itemize}
  \item 在平面 $\R^2$ 中，若 $v_1 = (2, 1)^T$ 与 $v_2 = (0.5, 1.5)^T$ 不共线，任意调节 $c_1 v_1 + c_2 v_2$ 就能编织出一张铺满整个平面的斜网格，$\spanop\{v_1, v_2\} = \R^2$；但如果两支向量恰好共线（例如 $u_2 = 2u_1$），无论怎么伸缩相加都逃不出同一条过原点的直线！
  \item 在三维空间 $\R^3$ 中，两支不共线的向量 $v_1, v_2$ 张成一张\textbf{穿过原点的倾斜平面}（二维子空间），如图~\ref{fig:ch2-span-geom} 所示；而任何不穿过原点的平面都不是子空间（因为它不含原点，且把向量乘以 $2$ 就会跳出该平面）。
\end{itemize}

\begin{figure}[htbp]
\centering
\resizebox{0.96\linewidth}{!}{%
\begin{tikzpicture}[font=\sffamily\footnotesize]
% Panel 1: 2D Span Grid vs Collinear Collapse
\begin{scope}[shift={(0,0)}]
  \node[font=\sffamily\bfseries\small, text=tealDark] at (1.5, 3.20) {二维张成：两支不共线向量铺满 $\R^2$};
  % Background oblique span grid of v1=(1.2, 0.35), v2=(0.45, 1.05)
  \begin{scope}[opacity=0.28]
    \foreach \i in {-1,0,1,2} {
      \draw[tealLine, thin] ({-0.9 + \i*1.2}, {-1.05 + \i*0.35}) -- ({1.35 + \i*1.2}, {2.10 + \i*0.35});
    }
    \foreach \j in {-1,0,1,2} {
      \draw[tealLine, thin] ({-1.2 + \j*0.45}, {-0.35 + \j*1.05}) -- ({3.6 + \j*0.45}, {1.05 + \j*1.05});
    }
  \end{scope}
  % Standard axes
  \draw[->, slate!60, thin] (-0.9, 0) -- (3.8, 0) node[right, text=slate] {$x_1$};
  \draw[->, slate!60, thin] (0, -0.7) -- (0, 2.50) node[left=1pt, text=slate] {$x_2$};
  % Span vectors v1, v2 and linear combination w = 1.5 v1 + 1.2 v2
  \draw[-{Stealth[length=2.3mm]}, teal, very thick] (0,0) -- (1.2, 0.35) node[below right=-1pt, text=tealDark] {$v_1$};
  \draw[-{Stealth[length=2.3mm]}, navy, very thick] (0,0) -- (0.45, 1.05) node[above left=-1pt, text=navy] {$v_2$};
  % Combination parallelogram
  \draw[dashed, gold, thick] (1.8, 0.525) -- (2.34, 1.785) -- (0.54, 1.26);
  \draw[-{Stealth[length=2.3mm]}, ochre, very thick] (0,0) -- (2.34, 1.785)
    node[above right=-1pt, text=ochre] {$c_1 v_1 + c_2 v_2$};
  \fill[tealDark] (0,0) circle (1.8pt);
  \node[below, text=slate, align=center] at (1.5, -0.95) {\scriptsize 任意平面点都有唯一组合系数 $(c_1, c_2)$\\[-1pt]\scriptsize 若 $v_2 \parallel v_1$，张成空间将塌缩为一条直线};
\end{scope}

% Panel 2: 3D Plane spanned by two vectors in R^3
\begin{scope}[shift={(7.6,0.1)}]
  \node[font=\sffamily\bfseries\small, text=navy] at (1.6, 3.10) {三维张成：两支独立向量张成过原点平面};
  % 3D axes (x_3 up, x_2 right, x_1 front-left)
  \coordinate (O) at (1.5, 0.75);
  \draw[->, slate!65, thin] (O) -- ++(2.55, 0) node[right, text=slate] {$x_2$};
  \draw[->, slate!65, thin] (O) -- ++(0, 1.78) node[left=1pt, text=slate] {$x_3$};
  \draw[->, slate!65, thin] (O) -- ++(-1.55, -0.95) node[below left, text=slate] {$x_1$};

  % Tilted 2D plane spanned by u=(1.35, 0.45) and w=(-0.85, 0.72) in screen coords
  \coordinate (P1) at ($(O) + (-1.35,-0.45) + (0.85,-0.72)$);
  \coordinate (P2) at ($(O) + (1.55,0.52) + (0.85,-0.72)$);
  \coordinate (P3) at ($(O) + (1.55,0.52) + (-0.95,0.80)$);
  \coordinate (P4) at ($(O) + (-1.35,-0.45) + (-0.95,0.80)$);
  \filldraw[fill=tealLight, fill opacity=0.72, draw=tealLine, thick] (P1) -- (P2) -- (P3) -- (P4) -- cycle;

  % Faint grid lines on the tilted plane
  \foreach \t in {-0.5, 0, 0.5} {
    \draw[tealLine!80, thin]
      ($(O) + \t*(1.45,0.48) + (0.85,-0.72)$) -- ($(O) + \t*(1.45,0.48) + (-0.95,0.80)$);
    \draw[tealLine!80, thin]
      ($(O) + (-1.35,-0.45) + \t*(-0.90,0.76)$) -- ($(O) + (1.55,0.52) + \t*(-0.90,0.76)$);
  }

  % Basis vectors on the plane
  \draw[-{Stealth[length=2.3mm]}, teal, very thick] (O) -- ++(1.25, 0.42) node[right, text=tealDark] {$v_1$};
  \draw[-{Stealth[length=2.3mm]}, navy, very thick] (O) -- ++(-0.78, 0.66) node[above left=-1pt, text=navy] {$v_2$};
  % A vector outside the plane (for intuition)
  \draw[-{Stealth[length=2.1mm]}, ochre, thick, dashed] (O) -- ++(0.35, 1.48) node[right, text=ochre] {$b \notin \spanop\{v_1,v_2\}$};
  \fill[tealDark] (O) circle (1.9pt);
  \node[below, text=slate, align=center] at (1.6, -1.05) {\scriptsize $\spanop\{v_1,v_2\}$ 是 $\R^3$ 中的二维切面\\[-1pt]\scriptsize 含噪目标 $b$ 往往落在平面外（引出正交投影）};
\end{scope}
\end{tikzpicture}}
\caption{张成空间（Span）的 2D 与 3D 几何直觉：线性组合的全体构成穿过原点的平坦子空间}\label{fig:ch2-span-geom}
\end{figure}

\subsection{线性无关、基（Basis）与维数（Dimension）}

如果工具箱里的某支向量 $v_k$ 已经可以由其他向量组合出来（即 $v_k \in \spanop\{v_1, \dots, v_{k-1}\}$），那么去掉 $v_k$ 后张成空间\textbf{一点都不会变小}——这支向量就是多余的。把所有冗余剔除干净，剩下的“极简无废话坐标系”就是\textbf{基}。

\begin{defBox}
\begin{definition}[线性无关、基与维数]\label{def:indep-basis-dim}
设 $v_1, \dots, v_k \in \R^n$，$U \subseteq \R^n$ 是线性子空间：
\begin{enumerate}
  \item \textbf{线性无关（Linear Independence）}：方程 $c_1 v_1 + \dots + c_k v_k = 0$ 仅有唯一零解 $c_1 = \dots = c_k = 0$（等价地，没有任何一支向量能写成其余向量的线性组合）。否则称其\textbf{线性相关}。
  \item \textbf{基（Basis）}：若向量组 $\{b_1, \dots, b_r\} \subseteq U$ 同时满足（i）\textbf{线性无关}（无冗余）且（ii）\textbf{张成整个 $U$}（$\spanop\{b_1, \dots, b_r\} = U$，无遗漏），则称它为 $U$ 的一组基。
  \item \textbf{维数（Dimension）}：可以证明，同一子空间 $U$ 的任意两组基所含向量个数必然相等，这个唯一确定的非负整数称为 $U$ 的\textbf{维数}，记作 $\dim(U) = r$（约定 $\dim\{0\} = 0$）。
\end{enumerate}
\end{definition}
\end{defBox}

为什么“基”如此珍贵？因为\textbf{一旦选定子空间 $U$ 的一组基 $\{b_1, \dots, b_r\}$，$U$ 中每个向量 $x$ 都可以唯一地写成坐标展开 $x = c_1 b_1 + \dots + c_r b_r$}！（若有两种写法 $\sum c_i b_i = \sum c'_i b_i$，相减得 $\sum (c_i - c'_i) b_i = 0$，由线性无关立刻迫使 $c_i = c'_i$。）

在机器学习的数据表 $X \in \R^{m \times n}$ 中，每一列就是一个特征向量。如果你把“身高（厘米）”和“身高（米）”同时放进特征列，这两列就\textbf{线性相关}（$x_{\text{cm}} = 100\, x_{\text{m}}$）——模型多了一个毫无新信息的冗余参数，线性回归权重将立刻失去唯一性！

\subsection{线性变换、矩阵的秩（Rank）与零空间（Null Space）}

现在迎来自学线性代数最关键的一次视角跃迁：\textbf{不要只把矩阵 $A \in \R^{m \times n}$ 看成一张静态数字表格，而要把它看成一台把输入空间 $\R^n$ 搬运到输出空间 $\R^m$ 的“线性变形机” $T(x) = Ax$}。

所谓“线性”，几何上只有两条铁律：（1）原点原地不动（$A0 = 0$）；（2）所有直线变换后依然是直线，且平行等距的网格线依然保持平行等距（旋转、拉伸、剪切、投影都是线性变换，但弯曲扭折不是）。结合第 1 章的列视角，只需盯住标准基向量 $e_1, \dots, e_n$ 的去向，就能完全锁定整台变形机的行为：
\[
  A e_j = a_{\cdot j} \quad (\text{矩阵 } A \text{ 的第 } j \text{ 列就是第 } j \text{ 支标准基向量变换后的落点！}).
\]
由此，任意输入 $x = \sum_{j=1}^n x_j e_j$ 经 $A$ 作用后的落点自然是 $Ax = \sum_{j=1}^n x_j a_{\cdot j}$。所有可能输出汇成的集合，正是第 1 章提到的\textbf{列空间} $\col(A) = \spanop\{a_{\cdot 1}, \dots, a_{\cdot n}\}$。于是矩阵的两个核心不变量立刻有了血肉：
\begin{itemize}
  \item \textbf{矩阵的秩（Rank）}：$\rank(A) \coloneqq \dim\col(A)$，即\textbf{经过变换 $A$ 之后，输出世界真正剩下的维数}。秩越高，保留的信息自由度越完整；若 $\rank(A) < n$，说明空间在变换中发生了“维数塌缩”（压扁）。
  \item \textbf{零空间（Null Space / 核空间）}：$\nul(A) \coloneqq \{x \in \R^n : Ax = 0\}$，即\textbf{所有被 $A$ 彻底压扁归零的输入方向全体}。显然 $\nul(A)$ 是 $\R^n$ 的线性子空间：如果 $x_n \in \nul(A)$，那么把任何输入 $x$ 加上 $x_n$ 后输出毫无变化：$A(x + x_n) = Ax + 0 = Ax$。
\end{itemize}

输入空间原本有 $n$ 个自由度，经过 $A$ 作用后，一部分自由度活了下来（撑起 $\rank(A)$ 维输出空间），剩下的自由度被彻底压缩进了原点（构成 $\dim\nul(A)$ 维零空间）——不多不少，恰好守恒！这就是线性代数最核心的维数守恒定律：
\begin{equation}\label{eq:rank-nullity}
  \rank(A) + \dim\nul(A) = n \qquad (\text{秩-零化度定理, Rank--Nullity Theorem}).
\end{equation}

用两个 $2 \times 2$ 矩阵立刻对比感受：
\[
  A_1 = \begin{pmatrix} 2 & 1 \\ 0 & 1.5 \end{pmatrix}, \qquad
  A_2 = \begin{pmatrix} 2 & 1 \\ 1 & 0.5 \end{pmatrix}.
\]
对 $A_1$，两列 $(2,0)^T$ 与 $(1,1.5)^T$ 不共线，故 $\col(A_1) = \R^2$、$\rank(A_1) = 2$，没有任何非零向量被压扁，$\nul(A_1) = \{0\}$（$2 + 0 = 2$）。而对 $A_2$，第 2 列恰好是第 1 列的一半，整个二维平面被彻底压扁到直线 $\spanop\{(2, 1)^T\}$ 上（$\rank(A_2) = 1$）；同时，整条直线 $x_2 = -2x_1$（即 $\spanop\{(1, -2)^T\}$）上的所有点都会被 $A_2$ 送到原点（$\dim\nul(A_2) = 1$，$1 + 1 = 2$）。

\subsection{行列式 $\det(A)$ 的几何本质：有向体积缩放因子}

传统教材往往一上来就背诵繁琐的逆序数或余子式展开来定义方阵 $A \in \R^{n \times n}$ 的行列式 $\det(A)$。但从线性变换的几何眼睛望去，行列式的本质惊人地单纯：

\begin{defBox}
\begin{definition}[行列式的几何意义]\label{def:det-geom}
对方阵 $A \in \R^{n \times n}$，考察标准基向量 $e_1, \dots, e_n$ 围成的单位超立方体（2D 为面积 $=1$ 的单位正方形，3D 为体积 $=1$ 的单位立方体）。在变换 $x \mapsto Ax$ 下，它变成由列向量 $a_{\cdot 1}, \dots, a_{\cdot n}$ 围成的平行多面体：
\begin{itemize}
  \item \textbf{绝对值 $|\det(A)|$}：等于该平行多面体的\textbf{面积/体积缩放倍数}（例如 2D 方阵 $\begin{psmallmatrix}a & b \\ c & d\end{psmallmatrix}$ 变换后的平行四边形有向面积恰好为 $ad - bc$）；
  \item \textbf{正负号}：$\det(A) > 0$ 表示空间保持原有定向（如旋转、拉伸、剪切）；$\det(A) < 0$ 表示空间发生了镜像翻转（如左右对调）；$\det(A) = 0$ 表示空间被\textbf{彻底压扁至低维，体积归零}。
\end{itemize}
\end{definition}
\end{defBox}

图~\ref{fig:ch2-det-geom} 用刚才的 $A_1$ 与 $A_2$ 将这一几何图像画得一目了然。

\begin{figure}[htbp]
\centering
\resizebox{0.97\linewidth}{!}{%
\begin{tikzpicture}[font=\sffamily\footnotesize]
% Panel 1: Unit square (det = 1)
\begin{scope}[shift={(0,0)}]
  \node[font=\sffamily\bfseries\small, text=ink] at (1.15, 2.75) {(1) 原始单位正方形};
  \draw[step=0.65cm, codeborder, very thin] (-0.35,-0.35) grid (2.6, 2.15);
  \draw[->, slate!75, thin] (-0.35, 0) -- (2.75, 0) node[right] {$x_1$};
  \draw[->, slate!75, thin] (0, -0.35) -- (0, 2.25) node[left=1pt] {$x_2$};
  \filldraw[fill=navyLight, fill opacity=0.85, draw=navy, thick] (0,0) rectangle (0.65, 0.65);
  \draw[-{Stealth[length=2.1mm]}, teal, very thick] (0,0) -- (0.65, 0) node[below, text=tealDark] {$e_1$};
  \draw[-{Stealth[length=2.1mm]}, navy, very thick] (0,0) -- (0, 0.65) node[left, text=navy] {$e_2$};
  \node[text=navy, font=\scriptsize\bfseries] at (0.325, 0.325) {$1$};
  \node[below, text=slate] at (1.15, -0.58) {面积 $= 1$};
\end{scope}

% Arrow 1 -> 2
\draw[-{Stealth[length=2.4mm]}, teal, thick] (3.05, 0.95) -- (4.05, 0.95)
  node[midway, above, text=tealDark, font=\scriptsize] {$A_1$};

% Panel 2: A1 = [[2,1],[0,1.5]], det = 3 (unit = 0.65cm)
\begin{scope}[shift={(4.45,0)}]
  \node[font=\sffamily\bfseries\small, text=tealDark] at (1.45, 2.75) {(2) 可逆变换：$\det(A_1)=3$};
  % Transformed grid lines (e1->(1.3,0), e2->(0.65,0.975))
  \begin{scope}[opacity=0.28]
    \foreach \j in {0,1,2} {
      \draw[tealLine, thin] ({\j*0.65}, {\j*0.975}) -- ({2.6+\j*0.65}, {\j*0.975});
    }
    \foreach \i in {0,1,2} {
      \draw[tealLine, thin] ({\i*1.3}, 0) -- ({\i*1.3+1.3}, 1.95);
    }
  \end{scope}
  \draw[->, slate!75, thin] (-0.25, 0) -- (3.35, 0) node[right] {$x_1$};
  \draw[->, slate!75, thin] (0, -0.35) -- (0, 2.25) node[left=1pt] {$x_2$};
  % Parallelogram (0,0) -> (1.3,0) -> (1.95,0.975) -> (0.65,0.975)
  \filldraw[fill=tealLight, fill opacity=0.90, draw=teal, thick]
    (0,0) -- (1.3,0) -- (1.95,0.975) -- (0.65,0.975) -- cycle;
  \draw[-{Stealth[length=2.1mm]}, teal, very thick] (0,0) -- (1.3, 0) node[below, text=tealDark] {$A_1 e_1$};
  \draw[-{Stealth[length=2.1mm]}, navy, very thick] (0,0) -- (0.65, 0.975) node[above left=-1pt, text=navy] {$A_1 e_2$};
  \node[text=tealDark, font=\scriptsize\bfseries] at (0.98, 0.48) {面积 $= 3$};
  \node[below, text=slate] at (1.45, -0.58) {$\rank=2,\;\nul=\{0\}$（可逆）};
\end{scope}

% Arrow 2 -> 3 (from Panel 1 conceptually, label A2)
\draw[-{Stealth[length=2.4mm]}, ochre, thick] (8.15, 0.95) -- (9.15, 0.95)
  node[midway, above, text=ochre, font=\scriptsize] {对比 $A_2$};

% Panel 3: A2 = [[2,1],[1,0.5]], det = 0
\begin{scope}[shift={(9.55,0)}]
  \node[font=\sffamily\bfseries\small, text=ochre] at (1.40, 2.75) {(3) 奇异压扁：$\det(A_2)=0$};
  \draw[->, slate!75, thin] (-0.25, 0) -- (3.15, 0) node[right] {$x_1$};
  \draw[->, slate!75, thin] (0, -0.35) -- (0, 2.25) node[left=1pt] {$x_2$};
  % Collapsed 1D column space line y = 0.5 x
  \draw[ochre!35, line width=4.5pt, line cap=round] (-0.2, -0.1) -- (2.9, 1.45);
  \draw[ochre, thick, dashed] (-0.2, -0.1) -- (2.9, 1.45) node[above left=2pt and -2pt, text=ochre, font=\scriptsize] {$\col(A_2)$};
  % A2 e1 = (2,1) -> (1.3, 0.65), A2 e2 = (1,0.5) -> (0.65, 0.325)
  \draw[-{Stealth[length=2.1mm]}, teal, very thick] (0,0) -- (1.3, 0.65) node[above left=-1pt, text=tealDark] {$A_2 e_1$};
  \draw[-{Stealth[length=2.1mm]}, navy, very thick] (0,0) -- (0.65, 0.325) node[below right=-1pt, text=navy] {$A_2 e_2$};
  % Null space line y = -2x hint
  \draw[slate, dotted, thick] (-0.25, 0.50) -- (0.20, -0.40);
  \node[text=slate, font=\scriptsize] at (1.05, 2.02) {平面被压扁为直线};
  \node[below, text=ochre] at (1.40, -0.58) {面积 $= 0$（不可逆、丢失信息）};
\end{scope}
\end{tikzpicture}}
\caption{行列式 $\det(A)$ 的几何本质：单位正方形经线性变换后的有向面积缩放倍数；一旦 $\det(A)=0$，整个平面不可逆地塌缩到一维直线}\label{fig:ch2-det-geom}
\end{figure}

带着这幅“体积缩放器”的画面，线性代数中两组最核心的结论几乎不证自明：
\begin{enumerate}
  \item \textbf{乘积法则}：$\det(AB) = \det(A)\det(B)$——先后做两次体积缩放，总缩放倍数就是两次倍数的乘积！
  \item \textbf{奇异性总等价链}：对方阵 $A \in \R^{n \times n}$，以下五句话说的是\textbf{同一件几何事实}：
  \begin{equation}\label{eq:singular-equiv}
    \det(A) = 0
    \;\iff\; \text{空间被压扁}
    \;\iff\; \rank(A) < n
    \;\iff\; \nul(A) \neq \{0\}
    \;\iff\; A \text{ 不可逆（奇异）}.
  \end{equation}
\end{enumerate}

\begin{insightBox}{为什么这两条直觉贯穿整套 AI 讲义？}
一旦理解了“$\det(A)=0 \iff$ 存在非零向量被压扁到 $0$”，第 3 章的\textbf{特征方程} $\det(A - \lambda I) = 0$ 就不再是神秘的代数咒语；而第 2 卷多元高斯分布的概率密度归一化常数 $\frac{1}{\sqrt{\det(\Sigma)}}$，以及第 5 卷标准化流（Normalizing Flows）的雅可比行列式 $|\det J|$，本质上全都是\textbf{“空间局部体积被拉伸了多少倍”}这一同一几何思想的化身！
\end{insightBox}

\section{向量范数、余弦相似度与高维几何直觉}\label{sec:norms-highdim}

\subsection{向量范数与 $\ell_1, \ell_2, \ell_\infty$ 单位球几何}

衡量一个向量 $x \in \R^n$ 的“大小”不止有欧氏长度一种方式。在人工智能里，换一把尺子往往会改变整个模型的解结构（例如让参数自动稀疏化）。抽象来看，任何合理的“长度度量”都只需满足三条几何公理。

\begin{defBox}
\begin{definition}[向量范数与 $\ell_p$ 范数]\label{def:vector-norm}
映射 $\norm{\cdot}: \R^n \to \R$ 若对任意 $x, y \in \R^n$ 及 $\alpha \in \R$ 满足：
\begin{enumerate}
  \item \textbf{非负性与正定性}：$\norm{x} \ge 0$，且 $\norm{x} = 0 \iff x = 0$；
  \item \textbf{绝对齐次性}：$\norm{\alpha x} = \abs{\alpha}\,\norm{x}$；
  \item \textbf{三角不等式}：$\norm{x + y} \le \norm{x} + \norm{y}$，
\end{enumerate}
则称 $\norm{\cdot}$ 为 $\R^n$ 上的一个\textbf{范数（Norm）}。AI 中最常用的三类 $\ell_p$ 范数为：
\[
\norm{x}_1 = \sum_{i=1}^{n} \abs{x_i}, \qquad
\norm{x}_2 = \sqrt{\sum_{i=1}^{n} x_i^2}, \qquad
\norm{x}_\infty = \max_{1 \le i \le n} \abs{x_i}.
\]
\end{definition}
\end{defBox}

先拿二维小向量 $x = (3, -4)^T \in \R^2$ 感受三把尺子的差别：
\[
\norm{x}_\infty = \max(3, 4) = 4
\;\le\;
\norm{x}_2 = \sqrt{3^2 + (-4)^2} = 5
\;\le\;
\norm{x}_1 = 3 + 4 = 7.
\]
这个不等式链 $\norm{x}_\infty \le \norm{x}_2 \le \norm{x}_1 \le \sqrt{n}\norm{x}_2$ 对所有 $x \in \R^n$ 恒成立（习题~2.5 给出完整证明）。更有启发的是画出这三种范数的\textbf{单位球} $\{x \in \R^2 : \norm{x}_p \le 1\}$，如图~\ref{fig:unit-balls} 所示。

\begin{figure}[htbp]
\centering
\resizebox{0.96\linewidth}{!}{%
\begin{tikzpicture}[font=\sffamily\footnotesize]
% Panel 1: L1 ball
\begin{scope}[shift={(0,0)}]
  \draw[->, slate] (-1.55, 0) -- (1.65, 0) node[right] {$x_1$};
  \draw[->, slate] (0, -1.55) -- (0, 1.65) node[above] {$x_2$};
  \filldraw[fill=tealLight, draw=teal, very thick] (1,0) -- (0,1) -- (-1,0) -- (0,-1) -- cycle;
  \foreach \p in {(1,0), (0,1), (-1,0), (0,-1)} {
    \fill[ochre] \p circle (2.2pt);
  }
  \node[above right=1pt, text=ochre] at (0,1) {\scriptsize 尖角在坐标轴上};
  \node[below, font=\sffamily\bfseries\small, text=tealDark] at (0, -1.72) {$\ell_1$ 单位球（菱形）};
  \node[below, text=slate] at (0, -2.18) {$\abs{x_1} + \abs{x_2} \le 1$};
\end{scope}

% Panel 2: L2 ball
\begin{scope}[shift={(5.3,0)}]
  \draw[->, slate] (-1.55, 0) -- (1.65, 0) node[right] {$x_1$};
  \draw[->, slate] (0, -1.55) -- (0, 1.65) node[above] {$x_2$};
  \filldraw[fill=navyLight, draw=navy, very thick] (0,0) circle (1.0);
  \node[above right=2pt, text=navy] at (0.71, 0.71) {\scriptsize 各向同性光滑};
  \node[below, font=\sffamily\bfseries\small, text=navy] at (0, -1.72) {$\ell_2$ 单位球（圆盘）};
  \node[below, text=slate] at (0, -2.18) {$x_1^2 + x_2^2 \le 1$};
\end{scope}

% Panel 3: L_inf ball
\begin{scope}[shift={(10.6,0)}]
  \draw[->, slate] (-1.55, 0) -- (1.65, 0) node[right] {$x_1$};
  \draw[->, slate] (0, -1.55) -- (0, 1.65) node[above] {$x_2$};
  \filldraw[fill=ochreLight, draw=ochre, very thick] (-1,-1) rectangle (1,1);
  \node[above right=1pt, text=ochre] at (0.12, 1.0) {\scriptsize 逐分量上限约束};
  \node[below, font=\sffamily\bfseries\small, text=ochre] at (0, -1.72) {$\ell_\infty$ 单位球（正方形）};
  \node[below, text=slate] at (0, -2.18) {$\max(\abs{x_1}, \abs{x_2}) \le 1$};
\end{scope}
\end{tikzpicture}}
\caption{$\R^2$ 中 $\ell_1$、$\ell_2$ 与 $\ell_\infty$ 单位球的几何形状差异}\label{fig:unit-balls}
\end{figure}

\begin{insightBox}{三把尺子在后续卷次中的分工}
\begin{itemize}
  \item \textbf{$\ell_1$ 菱形与稀疏性（第 4 卷 Lasso）}：$\ell_1$ 单位球在坐标轴上存在尖锐顶点；当损失函数的等高线扩张并与菱形相切时，切点大概率落在坐标轴顶点处，使若干权重精确变为 $0$，从而自动完成特征筛选。
  \item \textbf{$\ell_2$ 圆盘与光滑收缩（第~\ref{sec:poly-ridge}~节岭回归、第 5 卷 Weight Decay）}：$\ell_2$ 球处处光滑且旋转不变，惩罚 $\norm{w}_2^2$ 会温和均匀地缩小所有过大权重，防止过拟合，同时处处可导。
  \item \textbf{$\ell_\infty$ 正方形与对抗扰动（第 7 卷安全鲁棒性）}：图像像素级对抗攻击常用 $\norm{\delta}_\infty \le \varepsilon$ 约束每个像素的肉眼不可见最大改变量。
\end{itemize}
\end{insightBox}

\subsection{实内积、柯西-施瓦茨不等式与余弦相似度}

\begin{recallBox}{实欧氏内积与正交}
在 $\R^n$ 上，标准实内积定义为 $\inner{u}{v} = u^T v = \sum_{i=1}^n u_i v_i$。它满足四条公理：（1）对称性 $\inner{u}{v} = \inner{v}{u}$；（2）双线性 $\inner{\alpha u_1 + \beta u_2}{v} = \alpha\inner{u_1}{v} + \beta\inner{u_2}{v}$；（3）正定性 $\inner{u}{u} \ge 0$ 且等号当且仅当 $u=0$；（4）诱导欧氏范数 $\norm{u}_2 = \sqrt{\inner{u}{u}}$。若 $\inner{u}{v} = 0$，称 $u$ 与 $v$ \textbf{正交}，记作 $u \perp v$。
\end{recallBox}

下面这条不等式是整个内积空间几何、相关性分析与数值优化的基石。

\begin{thmBox}
\begin{theorem}[柯西-施瓦茨不等式、三角不等式与勾股定理]\label{thm:cauchy-schwarz}
对任意 $u, v \in \R^n$，恒有
\begin{equation}\label{eq:cs}
\abs{\inner{u}{v}} \le \norm{u}_2 \norm{v}_2,
\end{equation}
等号成立当且仅当 $u$ 与 $v$ 线性相关（共线）。由此立得：
\begin{enumerate}
  \item \textbf{三角不等式}：$\norm{u + v}_2 \le \norm{u}_2 + \norm{v}_2$；
  \item \textbf{勾股定理}：若 $u \perp v$，则 $\norm{u + v}_2^2 = \norm{u}_2^2 + \norm{v}_2^2$。
\end{enumerate}
\end{theorem}
\end{thmBox}

\begin{proof}
若 $v = 0$，不等式~\eqref{eq:cs} 两端均为 $0$，且 $\{u, 0\}$ 显然线性相关。下设 $v \neq 0$，对任意实数 $t \in \R$ 构造非负二次函数：
\[
q(t) = \norm{u + t v}_2^2 = \inner{u + tv}{u + tv} = \norm{v}_2^2 t^2 + 2\inner{u}{v} t + \norm{u}_2^2 \ge 0.
\]
由于 $\norm{v}_2^2 > 0$ 且抛物线 $q(t)$ 恒在横轴上方（或相切），其判别式必非正：
\[
\Delta = 4\inner{u}{v}^2 - 4\norm{u}_2^2 \norm{v}_2^2 \le 0
\implies \abs{\inner{u}{v}} \le \norm{u}_2 \norm{v}_2.
\]
等号成立当且仅当 $\Delta = 0$，即存在唯一实根 $t_0$ 使 $q(t_0) = \norm{u + t_0 v}_2^2 = 0$，由范数正定性等价于 $u = -t_0 v$，即 $u, v$ 线性相关。

进而展开 $\norm{u + v}_2^2 = \norm{u}_2^2 + 2\inner{u}{v} + \norm{v}_2^2$：
\begin{itemize}
  \item 代入 $\inner{u}{v} \le \norm{u}_2\norm{v}_2$ 得 $\norm{u+v}_2^2 \le (\norm{u}_2 + \norm{v}_2)^2$，开方即得三角不等式；
  \item 若 $u \perp v$，交叉项 $2\inner{u}{v} = 0$，立即得到勾股定理 $\norm{u+v}_2^2 = \norm{u}_2^2 + \norm{v}_2^2$。
\end{itemize}
\end{proof}

有了柯西-施瓦茨不等式保证比值落在 $[-1, 1]$ 内，我们就可以在任意高维空间 $\R^d$ 中合理定义两个非零向量的夹角余弦——即 AI 向量检索与表征学习的核心指标\textbf{余弦相似度（Cosine Similarity）}：
\begin{equation}\label{eq:cosine-sim}
\cos\theta(u, v) = \frac{\inner{u}{v}}{\norm{u}_2 \norm{v}_2} \in [-1, 1] \qquad (u, v \neq 0).
\end{equation}
余弦相似度最重要的工程特性是\textbf{尺度不变性}：对任意正数 $\alpha, \beta > 0$，$\cos\theta(\alpha u, \beta v) = \cos\theta(u, v)$。举个具体例子：假设文档 A 的词频向量为 $u = (2, 1, 0)^T$，文档 B 把文档 A 原封不动复制了 10 遍，词频向量变成 $v = (20, 10, 0)^T$。两篇文档主题完全相同，但欧氏距离 $\norm{u - v}_2 = 9\sqrt{5} \approx 20.1$ 极大；而余弦相似度 $\cos\theta(u, v) = 1.0$，一眼识别出它们方向完全一致！

\subsection{高维随机向量的“几乎必然正交”现象}

三维世界培养的几何直觉到了高维空间（如大语言模型常见的 $d = 768$ 或 $4096$ 维 Embedding 空间）会剧烈失效。其中一个对 AI 至关重要的反直觉现象是：\textbf{在超高维空间中，任意独立随机抽取的两个向量几乎总是彼此垂直（近正交）}！

先看简单的直觉估算：设 $u, v \in \R^d$ 的各分量独立服从标准正态分布 $\mathcal{N}(0, 1)$，归一化后得到球面上的均匀随机单位向量 $\hat{u} = u/\norm{u}_2$ 与 $\hat{v} = v/\norm{v}_2$。每个坐标大约是 $\pm 1/\sqrt{d}$ 量级，内积 $\cos\theta = \sum_{i=1}^d \hat{u}_i \hat{v}_i$ 是 $d$ 个正负随机抵消的微小乘积之和。大数抵消的结果是：$\cos\theta$ 的均值为 $0$，而其标准差严格按 $1/\sqrt{d}$ 衰减！

我们在配套实验脚本的 \texttt{exp2\_highdim\_ortho} 函数（位于 \texttt{code/vol1\_experiments.py}）中，对每个维度 $d \in \{2, 5, 10, \dots, 1000\}$ 独立采样 $5{,}000$ 对随机单位向量，实测结果见表~\ref{tab:highdim-cos} 与图~\ref{fig:highdim-ortho}。

\begin{table}[htbp]
\centering
\caption{高维随机单位向量余弦相似度实测统计（每维 $5{,}000$ 对样本，源自 \texttt{data/highdim\_cos.csv}）}\label{tab:highdim-cos}
{\small
\begin{tabular}{@{}rcccc@{}}
\toprule
\textbf{维度 $d$} & \textbf{实测标准差 $\operatorname{Std}(\cos\theta)$} & \textbf{理论值 $1/\sqrt{d}$} & \textbf{$P(\abs{\cos\theta} < 0.1)$} & \textbf{最大 $\abs{\cos\theta}$} \\
\midrule
2 & 0.7118 & 0.7071 & 6.50\% & 1.0000 \\
10 & 0.3150 & 0.3162 & 22.72\% & 0.9101 \\
50 & 0.1392 & 0.1414 & 51.74\% & 0.5156 \\
100 & 0.1002 & 0.1000 & 68.16\% & 0.3651 \\
500 & 0.0455 & 0.0447 & 96.82\% & 0.1724 \\
1000 & 0.0317 & 0.0316 & \textbf{99.86\%} & 0.1086 \\
\bottomrule
\end{tabular}}
\end{table}

\begin{figure}[htbp]
\centering
\begin{tikzpicture}
\begin{axis}[
  width=0.47\linewidth,
  height=5.2cm,
  xmode=log, ymode=log,
  xlabel={维度 $d$（对数刻度）},
  ylabel={余弦相似度标准差},
  grid=both,
  grid style={line width=0.25pt, draw=codeborder},
  legend pos=south west,
  legend style={font=\scriptsize, fill=white, draw=codeborder},
  tick label style={font=\scriptsize},
  label style={font=\footnotesize\sffamily}
]
  \addplot[only marks, mark=*, mark size=2.2pt, color=teal]
    table[x=d, y=emp_std, col sep=comma]{data/highdim_cos.csv};
  \addlegendentry{实测标准差}
  \addplot[thick, dashed, color=ochre]
    table[x=d, y=theory_std, col sep=comma]{data/highdim_cos.csv};
  \addlegendentry{理论曲线 $1/\sqrt{d}$}
\end{axis}
\end{tikzpicture}
\hfill
\begin{tikzpicture}
\begin{axis}[
  width=0.47\linewidth,
  height=5.2cm,
  xmode=log,
  ymin=0, ymax=105,
  xlabel={维度 $d$（对数刻度）},
  ylabel={落入近正交区比例 $P(\abs{\cos\theta}<0.1)$ (\%)},
  grid=both,
  grid style={line width=0.25pt, draw=codeborder},
  tick label style={font=\scriptsize},
  label style={font=\footnotesize\sffamily}
]
  \addplot[thick, mark=square*, mark size=2pt, color=navy]
    table[x=d, y=pct_abs_lt_01, col sep=comma]{data/highdim_cos.csv};
  \node[anchor=west, font=\scriptsize\sffamily, text=ochre, fill=white, inner sep=1.5pt]
    at (axis cs:2.2, 87) {$d=1000$ 时达 $99.86\%$};
\end{axis}
\end{tikzpicture}
\caption{高维随机向量的近正交集中现象（数据直接读取自 \texttt{data/highdim\_cos.csv}）：左图实测标准差与 $1/\sqrt{d}$ 严丝合缝，右图近正交比例随维度升至 $99.86\%$}\label{fig:highdim-ortho}
\end{figure}

\begin{insightBox}{高维近正交对 AI 是一把“双刃剑”}
\begin{itemize}
  \item \textbf{表征学习的红利}：在平面 $\R^2$ 里最多只能摆 2 个两两垂直的向量；但在 $d=1000$ 维空间里，只要容忍 $\abs{\cos\theta} < 0.1$ 的微小夹角偏差，就能塞进\textbf{远超 $d$ 个（甚至指数级数量）几乎互不干扰的语义方向}。这正是大语言模型用几千维向量表征几十万词义与概念的几何底气。
  \item \textbf{朴素近邻算法的维度灾难}：如果直接把原始高维噪声特征丢进 $k$-近邻（$k$-NN）分类器，任意两点几乎都互相垂直、距离趋同（$\norm{\hat{u}-\hat{v}}_2 = \sqrt{2 - 2\cos\theta} \approx \sqrt{2}$），导致“最近邻”与“最远邻”几乎没有区分度——除非数据本身集中在低维流形上（第 4 章 PCA）或经过深度网络学习出紧凑表征。
\end{itemize}
\end{insightBox}

\section{正交补、四大基本子空间与“超定 vs 欠定”二象性}\label{sec:four-subspaces}

\subsection{标准正交基与正交分解定理}

设 $U \subseteq \R^n$ 是线性子空间。如果向量组 $\{q_1, \dots, q_k\}$ 是 $U$ 的一组基，且满足两两正交、长度为 1，即
\[
\inner{q_i}{q_j} = \delta_{ij} = \begin{cases} 1, & i = j, \\ 0, & i \neq j, \end{cases}
\]
则称 $\{q_1, \dots, q_k\}$ 为 $U$ 的一组\textbf{标准正交基（Orthonormal Basis）}。与 $U$ 中所有向量都垂直的向量全体组成 $U$ 的\textbf{正交补（Orthogonal Complement）}：
\[
U^\perp = \{w \in \R^n : \inner{w}{u} = 0,\; \forall u \in U\}.
\]
先看三维直觉：在 $\R^3$ 中，若 $U$ 是水平 $xy$-平面，则 $U^\perp$ 就是竖直 $z$-轴；空间中任意一点 $(x, y, z)^T$ 都可以唯一地拆成“落在平面内的影子”$(x, y, 0)^T \in U$ 与“垂直于平面的高度”$(0, 0, z)^T \in U^\perp$ 之和。下面的定理把这一几何直觉严格推广到任意维度。

\begin{thmBox}
\begin{theorem}[正交分解定理]\label{thm:ortho-decomp}
设 $U$ 是 $\R^n$ 的任意线性子空间，则 $\R^n$ 可唯一分解为 $U$ 与其正交补 $U^\perp$ 的直和：
\[
\R^n = U \oplus U^\perp.
\]
也就是说，对任意 $v \in \R^n$，存在唯一的 $u \in U$ 与 $w \in U^\perp$ 使得 $v = u + w$。进而有 $\dim U + \dim U^\perp = n$ 且 $(U^\perp)^\perp = U$。
\end{theorem}
\end{thmBox}

\begin{proof}
若 $U = \{0\}$，取 $u = 0,\, w = v$ 显然成立。下设 $\dim U = k \ge 1$，取 $U$ 的一组标准正交基 $\{q_1, \dots, q_k\}$（下一节的 Gram-Schmidt 过程保证任何有限维非零子空间都存在标准正交基）。

\textbf{（1）存在性}：对任意 $v \in \R^n$，构造
\[
u = \sum_{i=1}^{k} \inner{v}{q_i} q_i \in U, \qquad w = v - u.
\]
验证 $w \in U^\perp$：对每个基向量 $q_j$（$1 \le j \le k$），利用内积双线性与标准正交性：
\[
\inner{w}{q_j} = \inner{v}{q_j} - \sum_{i=1}^{k} \inner{v}{q_i} \inner{q_i}{q_j}
= \inner{v}{q_j} - \inner{v}{q_j} = 0.
\]
由于 $U$ 中任意向量都是 $q_1, \dots, q_k$ 的线性组合，故 $w$ 与 $U$ 中所有向量正交，即 $w \in U^\perp$，且 $v = u + w$。

\textbf{（2）唯一性}：若另有一组分解 $v = u' + w'$（$u' \in U,\, w' \in U^\perp$），则令 $z = u - u' = w' - w$。此时 $z$ 同时属于 $U$ 与 $U^\perp$，必有 $\inner{z}{z} = 0$，由内积正定性得 $z = 0$，故 $u = u'$ 且 $w = w'$。

由直和维数公式立即得到 $\dim U + \dim U^\perp = n$。再由定义显然有 $U \subseteq (U^\perp)^\perp$，而 $\dim(U^\perp)^\perp = n - \dim U^\perp = \dim U$，故 $(U^\perp)^\perp = U$。
\end{proof}

\subsection{线性代数基本定理（四大基本子空间）}

\begin{recallBox}{矩阵 $A \in \R^{m \times n}$ 的四个基本子空间与维数}
设矩阵 $A \in \R^{m \times n}$ 的秩为 $r = \rank(A)$，则它自然关联着四个线性子空间：
\begin{itemize}
  \item \textbf{列空间} $\col(A) = \{Ax : x \in \R^n\} \subseteq \R^m$，维数为 $r$；
  \item \textbf{行空间} $\row(A) = \col(A^T) = \{A^T y : y \in \R^m\} \subseteq \R^n$，维数为 $r$；
  \item \textbf{零空间（核空间）} $\nul(A) = \{x \in \R^n : Ax = 0\} \subseteq \R^n$，维数为 $n - r$；
  \item \textbf{左零空间} $\nul(A^T) = \{y \in \R^m : A^T y = 0\} \subseteq \R^m$，维数为 $m - r$。
\end{itemize}
\end{recallBox}

这四个子空间并不是零散孤立的，而是通过正交补结成严丝合缝的两对直和分解。

\begin{thmBox}
\begin{theorem}[线性代数基本定理：四大子空间正交对偶]\label{thm:four-subspaces}
对任意矩阵 $A \in \R^{m \times n}$，恒有以下两组正交补关系与空间直和分解：
\[
\nul(A) = \row(A)^\perp, \qquad \nul(A^T) = \col(A)^\perp,
\]
\[
\R^n = \row(A) \oplus \nul(A), \qquad \R^m = \col(A) \oplus \nul(A^T).
\]
\end{theorem}
\end{thmBox}

\begin{proof}
记 $A$ 的各行为 $a_{1\cdot}^T, \dots, a_{m\cdot}^T \in \R^n$。由矩阵乘法点积视角，向量 $x \in \nul(A)$ 当且仅当 $Ax = 0$，即对每个 $i = 1, \dots, m$ 均有 $\inner{a_{i\cdot}}{x} = 0$。这等价于 $x$ 与 $A$ 的所有行的任意线性组合都正交，即 $x \in \row(A)^\perp$。因此 $\nul(A) = \row(A)^\perp$。

把上述结论直接应用到转置矩阵 $A^T$，立刻得到 $\nul(A^T) = \row(A^T)^\perp = \col(A)^\perp$。再结合正交分解定理（定理~\ref{thm:ortho-decomp}），即得两条直和分解。
\end{proof}

图~\ref{fig:four-subspaces} 用一张经典全景图刻画了线性变换 $A: \R^n \to \R^m$ 的全部秘密：任何输入 $x \in \R^n$ 都唯一拆成 $x = x_r + x_n$（$x_r \in \row(A),\, x_n \in \nul(A)$）；乘以 $A$ 时，零空间分量 $x_n$ 被彻底压缩为 $0$，而行空间 $\row(A)$ 被 $A$ \textbf{一一满射（双射）}地映到列空间 $\col(A)$ 上！

\begin{figure}[htbp]
\centering
\resizebox{0.96\linewidth}{!}{%
\begin{tikzpicture}[
  font=\sffamily\footnotesize,
  subR/.style={draw=teal, very thick, fill=tealLight, rounded corners=1.2mm},
  subN/.style={draw=slate, thick, fill=codebg, rounded corners=1.2mm},
  subC/.style={draw=navy, very thick, fill=navyLight, rounded corners=1.2mm},
  subLN/.style={draw=ochre, thick, fill=ochreLight, rounded corners=1.2mm},
  arr/.style={-{Stealth[length=2.5mm]}, very thick}
]
% Left domain R^n
\node[font=\sffamily\bfseries\normalsize, text=ink] at (2.1, 4.95) {输入空间 $\R^n$};
% Row space
\draw[subR] (0.2, 2.35) rectangle (4.0, 4.45);
\node[font=\sffamily\bfseries, text=tealDark] at (2.1, 3.85) {行空间 $\row(A)$};
\node[text=slate] at (2.1, 3.30) {$\dim = r$};
\fill[tealDark] (2.85, 2.75) circle (2pt) node[left=2pt, text=tealDark] {$x_r$};
% Null space
\draw[subN] (0.2, 0.25) rectangle (4.0, 2.35);
\node[font=\sffamily\bfseries, text=ink] at (2.1, 1.65) {零空间 $\nul(A)$};
\node[text=slate] at (2.1, 1.10) {$\dim = n - r$};
\fill[slate] (2.85, 0.65) circle (2pt) node[left=2pt, text=slate] {$x_n$};
% Right angle symbol between row(A) and null(A)
\draw[ink, thick] (0.2, 2.35) ++(0.32, 0) -- ++(0, 0.32) -- ++(-0.32, 0);
\node[anchor=west, text=tealDark] at (0.62, 2.58) {\scriptsize $\perp$};

% Right codomain R^m
\node[font=\sffamily\bfseries\normalsize, text=ink] at (11.1, 4.95) {输出空间 $\R^m$};
% Column space
\draw[subC] (9.2, 2.35) rectangle (13.0, 4.45);
\node[font=\sffamily\bfseries, text=navy] at (11.1, 3.85) {列空间 $\col(A)$};
\node[text=slate] at (11.1, 3.30) {$\dim = r$};
\fill[navy] (10.35, 2.75) circle (2pt) node[right=2pt, text=navy] {$b = A x_r$};
% Left null space
\draw[subLN] (9.2, 0.25) rectangle (13.0, 2.35);
\node[font=\sffamily\bfseries, text=ochre] at (11.1, 1.65) {左零空间 $\nul(A^T)$};
\node[text=slate] at (11.1, 1.10) {$\dim = m - r$};
\fill[ochre] (10.35, 2.35) circle (2.2pt) node[below right=1pt, text=ochre] {$0$};
% Right angle symbol
\draw[ink, thick] (13.0, 2.35) ++(-0.32, 0) -- ++(0, 0.32) -- ++(0.32, 0);
\node[anchor=east, text=navy] at (12.58, 2.58) {\scriptsize $\perp$};

% Arrows between spaces
\draw[arr, teal] (4.15, 3.80) to[bend left=10] node[above=2pt, text=tealDark] {$A$（一一对应双射）} (9.05, 3.80);
\draw[arr, navy, dashed] (9.05, 3.20) to[bend left=8] node[below=2pt, text=navy] {$A^T$（反向双射）} (4.15, 3.20);
\draw[arr, slate] (4.15, 0.90) -- node[below=4pt, sloped, text=slate] {$A x_n = 0$（压缩至原点）} (10.18, 2.28);
\end{tikzpicture}}
\caption{线性代数基本定理：矩阵 $A \in \R^{m \times n}$ 的四大基本子空间及其正交补对偶}\label{fig:four-subspaces}
\end{figure}

\subsection{AI 中的“超定 vs 欠定”二象性}

把数据矩阵 $A \in \R^{m \times n}$（$m$ 个训练样本、$n$ 个待估参数）与标签向量 $b \in \R^m$ 放进图~\ref{fig:four-subspaces}，人工智能历史上前后两个时代的核心建模矛盾立刻一览无余（见表~\ref{tab:over-under}）：

\begin{table}[htbp]
\centering
\caption{线性方程 $Aw = b$ 在人工智能中的“超定 vs 欠定”二象性}\label{tab:over-under}
{\small
\begin{tabularx}{\linewidth}{@{}p{24mm}p{56mm}X@{}}
\toprule
\textbf{维度对比} & \textbf{超定系统（$m > n$，典型 $\rank(A) = n$）} & \textbf{欠定系统（$m < n$，典型 $\rank(A) = m$）} \\
\midrule
AI 时代对应 & 传统统计学习：样本数远多于特征数 & 现代过参数化深度网络：参数量远大于样本数 \\
子空间几何 & $\col(A)$ 只是 $\R^m$ 的低维切片，$\nul(A)=\{0\}$ & $\col(A) = \R^m$ 铺满空间，$\nul(A)$ 维数高达 $n-m$ \\
精确解存在性 & 含噪观测 $b \notin \col(A)$，\textbf{通常无精确解} & 满足 $Aw=b$ 的解构成仿射空间 $w_0 + \nul(A)$，\textbf{无穷多解} \\
核心求解策略 & 在 $\col(A)$ 中做\textbf{正交投影}（最小二乘，第~\ref{sec:proj-ls}~节） & 引入归纳偏置挑选零空间分量为 $0$ 的\textbf{最小范数解}（第 4 章伪逆） \\
\bottomrule
\end{tabularx}}
\end{table}

\section{正交投影与线性最小二乘}\label{sec:proj-ls}

\subsection{从三点拟合直线出发}

现在我们直面超定系统。假设你测得了 3 个实验数据点 $(0, 1),\, (1, 2),\, (2, 4)$，想用一条直线 $y = w_0 + w_1 x$ 去拟合它们。把三组坐标依次代入，得到超定线性方程组 $Aw = b$：
\[
\begin{cases}
w_0 + 0\cdot w_1 = 1, \\
w_0 + 1\cdot w_1 = 2, \\
w_0 + 2\cdot w_1 = 4,
\end{cases}
\quad \Longleftrightarrow \quad
\underbrace{\begin{pmatrix} 1 & 0 \\ 1 & 1 \\ 1 & 2 \end{pmatrix}}_{A}
\underbrace{\begin{pmatrix} w_0 \\ w_1 \end{pmatrix}}_{w}
=
\underbrace{\begin{pmatrix} 1 \\ 2 \\ 4 \end{pmatrix}}_{b}.
\]
一眼可见它\textbf{无精确解}：前两式已经唯一钉死了 $w_0 = 1,\, w_1 = 1$，但代入第三式只能得到 $1 + 2 = 3 \neq 4$。用第 1 章的列线性组合视角翻译：无论你怎么调 $w_0, w_1$，向量 $Aw = w_0 (1,1,1)^T + w_1 (0,1,2)^T$ 永远被锁死在 $\R^3$ 的二维平面 $\col(A)$ 内，而观测点 $b = (1, 2, 4)^T$ 漂浮在这个平面之外！

既然碰不到 $b$，最自然的退而求其次便是寻找 $\hat{w} \in \R^n$，使平面上的预测向量 $A\hat{w}$ 离目标 $b$ 的欧氏距离平方最小——这就是著名的\textbf{线性最小二乘问题（Linear Least Squares）}：
\begin{equation}\label{eq:ls-obj}
\hat{w} = \argmin_{w \in \R^n} \norm{Aw - b}_2^2.
\end{equation}

\subsection{正交投影最优逼近定理与投影矩阵刻画}

三维空间里，从平面外一点 $b$ 到平面 $\col(A)$ 最短的线段是什么？初中几何就告诉你：\textbf{向平面作垂线}！垂足 $\hat{b} = A\hat{w}$ 就是最佳逼近点，而残差向量 $e = b - A\hat{w}$ 必须垂直于整个平面 $\col(A)$（见图~\ref{fig:ortho-proj}）。下面我们用勾股定理将这一直观变成任意高维下的严格定理。

\begin{figure}[htbp]
\centering
\resizebox{0.82\linewidth}{!}{%
\begin{tikzpicture}[font=\sffamily\footnotesize]
  % Subspace plane col(A)
  \fill[tealLight, opacity=0.85] (-0.8, -0.75) -- (6.2, -0.75) -- (7.8, 2.25) -- (0.8, 2.25) -- cycle;
  \draw[teal, thick] (-0.8, -0.75) -- (6.2, -0.75) -- (7.8, 2.25) -- (0.8, 2.25) -- cycle;
  \node[anchor=north east, font=\sffamily\bfseries\small, text=tealDark] at (7.35, 2.05) {列空间 $\col(A)$};

  % Origin and points (clean triangle separation)
  \coordinate (O) at (0.8, 0.35);
  \coordinate (Bhat) at (4.7, 1.35);
  \coordinate (B) at (4.7, 4.25);
  \coordinate (Aw) at (3.2, -0.30);

  % Vectors
  \draw[-{Stealth[length=2.5mm]}, very thick, navy] (O) -- (B) node[midway, above left=3pt, text=navy] {$b$（观测向量）};
  \draw[-{Stealth[length=2.5mm]}, very thick, teal] (O) -- (Bhat) node[pos=0.52, above=3pt, sloped, text=tealDark] {$\hat{b} = A\hat{w} = Pb$};
  \draw[-{Stealth[length=2.5mm]}, very thick, ochre] (Bhat) -- (B) node[pos=0.58, right=5pt, text=ochre] {$e = b - A\hat{w} \perp \col(A)$};

  % Arbitrary point Aw and right triangle
  \draw[-{Stealth[length=2.2mm]}, thick, slate, dashed] (O) -- (Aw);
  \draw[thick, slate, dashed] (Aw) -- (Bhat) node[midway, below right=2pt, text=slate] {\scriptsize $\hat{b} - Aw$};
  \draw[thick, slate, dotted] (Aw) -- (B) node[pos=0.61, left=2pt, text=slate] {\scriptsize $b - Aw$};

  % Right angle mark at Bhat
  \draw[ochre, thick] (Bhat) ++(0, 0.32) -- ++(-0.32, -0.08) -- ++(0, -0.32);

  % Dots & labels
  \fill[ink] (O) circle (2pt) node[below left=1pt] {$0$};
  \fill[tealDark] (Bhat) circle (2.2pt);
  \fill[navy] (B) circle (2.2pt);
  \fill[slate] (Aw) circle (2pt) node[below=2pt, text=slate] {$Aw$};
\end{tikzpicture}}
\caption{线性最小二乘的几何本质：最优拟合 $\hat{b} = A\hat{w}$ 是 $b$ 在列空间 $\col(A)$ 上的正交投影，对任意 $Aw \in \col(A)$ 构成直角三角形}\label{fig:ortho-proj}
\end{figure}

\begin{thmBox}
\begin{theorem}[正交投影最优逼近定理]\label{thm:best-approx}
设 $U \subseteq \R^m$ 为线性子空间，给定 $b \in \R^m$，按定理~\ref{thm:ortho-decomp} 将其唯一分解为 $b = \hat{b} + e$，其中 $\hat{b} \in U$、$e \in U^\perp$。则对任意 $u \in U$，恒有
\[
\norm{b - u}_2^2 = \norm{e}_2^2 + \norm{\hat{b} - u}_2^2 \ge \norm{e}_2^2,
\]
等号成立当且仅当 $u = \hat{b}$。因此 $\hat{b}$ 是 $U$ 中离 $b$ 唯一的欧氏最近点，称为 $b$ 在 $U$ 上的\textbf{正交投影}。
\end{theorem}
\end{thmBox}

\begin{proof}
对任意 $u \in U$，由于 $\hat{b} \in U$ 且 $U$ 是子空间，有 $\hat{b} - u \in U$；又 $e \in U^\perp$，故 $e \perp (\hat{b} - u)$。将差向量改写为 $b - u = (b - \hat{b}) + (\hat{b} - u) = e + (\hat{b} - u)$，直接应用勾股定理（定理~\ref{thm:cauchy-schwarz}）即得
\[
\norm{b - u}_2^2 = \norm{e}_2^2 + \norm{\hat{b} - u}_2^2 \ge \norm{e}_2^2.
\]
等号成立当且仅当 $\norm{\hat{b} - u}_2^2 = 0$，即 $u = \hat{b}$。
\end{proof}

若将每个向量 $b \in \R^m$ 映到它在 $U$ 上的正交投影 $\hat{b}$，这一映射是线性的，对应唯一的 $m \times m$ 矩阵 $P$（即 $\hat{b} = Pb$）。如何一眼认出一个矩阵是不是正交投影矩阵？只需检验两条代数性质。

\begin{thmBox}
\begin{theorem}[正交投影矩阵的充要刻画]\label{thm:proj-matrix}
方阵 $P \in \R^{m \times m}$ 是某个子空间 $U \subseteq \R^m$ 上的正交投影矩阵，当且仅当它同时满足\textbf{幂等性}与\textbf{对称性}：
\[
P^2 = P \qquad \text{且} \qquad P^T = P.
\]
\end{theorem}
\end{thmBox}

\begin{proof}
\textbf{（必要性 $\Rightarrow$）}：设 $P$ 是到 $U$ 的正交投影。对任意 $b \in \R^m$，向量 $Pb$ 已经落在 $U$ 中，再投影一次仍等于自身，即 $P(Pb) = Pb$，故 $P^2 = P$。其次，对任意 $x, y \in \R^m$，拆分为 $x = Px + (x - Px)$ 与 $y = Py + (y - Py)$，由于 $Px, Py \in U$ 且 $x - Px, y - Py \in U^\perp$，交叉内积为零，于是
\[
\inner{Px}{y} = \inner{Px}{Py + (y - Py)} = \inner{Px}{Py} = \inner{Px + (x - Px)}{Py} = \inner{x}{Py}.
\]
写成矩阵形式即 $x^T P^T y = x^T P y$ 对一切 $x, y$ 成立，故 $P^T = P$。

\textbf{（充分性 $\Leftarrow$）}：设 $P^2 = P = P^T$，令子空间 $U = \col(P)$。对任意 $b \in \R^m$，分解 $b = Pb + (b - Pb)$。显然 $Pb \in U$；对 $U$ 中任意向量 $Pu$，有
\[
\inner{b - Pb}{Pu} = \inner{P^T(b - Pb)}{u} = \inner{Pb - P^2 b}{u} = \inner{0}{u} = 0.
\]
故 $b - Pb \in U^\perp$。由正交分解的唯一性，$Pb$ 正是 $b$ 在 $U$ 上的正交投影。
\end{proof}

\subsection{正规方程、可逆性判据与三点手算验证}

把定理~\ref{thm:best-approx} 应用到子空间 $U = \col(A)$：最优参数 $\hat{w}$ 必须使残差 $e = b - A\hat{w}$ 落入 $\col(A)^\perp$。回顾线性代数基本定理（定理~\ref{thm:four-subspaces}），$\col(A)^\perp = \nul(A^T)$！于是几何正交条件立刻翻译成代数方程：
\begin{equation}\label{eq:normal-eq}
A^T (b - A\hat{w}) = 0 \quad \Longleftrightarrow \quad A^T A \hat{w} = A^T b.
\end{equation}
方程~\eqref{eq:normal-eq} 称为\textbf{正规方程（Normal Equations）}。下面这个定理不借用任何特征值工具，直接回答了最小二乘解何时唯一存在。

\begin{thmBox}
\begin{theorem}[正规方程的可逆性判据]\label{thm:ata-inv}
对任意矩阵 $A \in \R^{m \times n}$，恒有 $\nul(A^T A) = \nul(A)$。因此方阵 $A^T A \in \R^{n \times n}$ 可逆当且仅当 $A$ \textbf{列满秩}（即 $\rank(A) = n$）。此时最小二乘问题~\eqref{eq:ls-obj} 有唯一解
\[
\hat{w} = (A^T A)^{-1} A^T b,
\]
且向 $\col(A)$ 的正交投影矩阵为 $P = A (A^T A)^{-1} A^T$。
\end{theorem}
\end{thmBox}

\begin{proof}
先证 $\nul(A^T A) = \nul(A)$：若 $Av = 0$，两边左乘 $A^T$ 立得 $A^T A v = 0$，故 $\nul(A) \subseteq \nul(A^T A)$。反过来，若 $A^T A v = 0$，两边左乘 $v^T$ 得
\[
0 = v^T A^T A v = (Av)^T (Av) = \norm{Av}_2^2.
\]
由范数正定性必有 $Av = 0$，即 $v \in \nul(A)$。两向包含证毕。

方阵 $A^T A$ 可逆 $\iff \nul(A^T A) = \{0\} \iff \nul(A) = \{0\} \iff \rank(A) = n$。当 $A^T A$ 可逆时，正规方程给出唯一解 $\hat{w} = (A^T A)^{-1} A^T b$，相应投影为 $\hat{b} = A\hat{w} = A(A^T A)^{-1} A^T b = Pb$。直接检验可知 $P^T = P$ 且 $P^2 = A(A^T A)^{-1}(A^T A)(A^T A)^{-1} A^T = P$。
\end{proof}

\begin{example}[三点拟合直线的手算闭环]\label{ex:3pt-ls}
回到本节开头的三点 $(0,1),\, (1,2),\, (2,4)$，设计矩阵 $A = \begin{psmallmatrix} 1 & 0 \\ 1 & 1 \\ 1 & 2 \end{psmallmatrix}$ 显然列满秩（$\rank(A)=2$）。计算正规方程两端：
\[
A^T A = \begin{pmatrix} 1 & 1 & 1 \\ 0 & 1 & 2 \end{pmatrix}\begin{pmatrix} 1 & 0 \\ 1 & 1 \\ 1 & 2 \end{pmatrix}
= \begin{pmatrix} 3 & 3 \\ 3 & 5 \end{pmatrix}, \qquad
A^T b = \begin{pmatrix} 1 & 1 & 1 \\ 0 & 1 & 2 \end{pmatrix}\begin{pmatrix} 1 \\ 2 \\ 4 \end{pmatrix}
= \begin{pmatrix} 7 \\ 10 \end{pmatrix}.
\]
由于 $\det(A^T A) = 3\times 5 - 3\times 3 = 6$，直接求逆得
\[
\hat{w} = \frac{1}{6}\begin{pmatrix} 5 & -3 \\ -3 & 3 \end{pmatrix}\begin{pmatrix} 7 \\ 10 \end{pmatrix}
= \frac{1}{6}\begin{pmatrix} 5 \\ 9 \end{pmatrix}
= \begin{pmatrix} 5/6 \\ 3/2 \end{pmatrix}.
\]
于是最佳拟合直线为 $\hat{y} = \frac{5}{6} + \frac{3}{2} x$。验证残差正交性：预测值 $\hat{b} = A\hat{w} = (\frac{5}{6},\, \frac{14}{6},\, \frac{23}{6})^T$，残差 $e = b - \hat{b} = (\frac{1}{6},\, -\frac{2}{6},\, \frac{1}{6})^T$，检查内积：
\[
A^T e = \begin{pmatrix} \frac{1}{6} - \frac{2}{6} + \frac{1}{6} \\[2pt] 0\cdot\frac{1}{6} - 1\cdot\frac{2}{6} + 2\cdot\frac{1}{6} \end{pmatrix}
= \begin{pmatrix} 0 \\ 0 \end{pmatrix}, \qquad
\norm{e}_2^2 = \frac{1 + 4 + 1}{36} = \frac{1}{6}.
\]
\end{example}

\subsection{Gram-Schmidt 正交化与 QR 分解：为什么工程上不直接求 $(A^T A)^{-1}$？}

如果矩阵 $A$ 的各列本身就构成标准正交向量组（记作 $Q \in \R^{m \times n}$，满足 $Q^T Q = I_n$），正规方程会出现奇迹：$Q^T Q \hat{w} = Q^T b$ 直接坍缩为 $\hat{w} = Q^T b$，连求逆都省了！

自然地，我们可以用\textbf{格拉姆-施密特（Gram-Schmidt）正交化}，把任意列满秩矩阵 $A = [a_1, \dots, a_n]$ 的列逐一“掰直”为标准正交列 $q_1, \dots, q_n$：第 $k$ 步从 $a_k$ 中减掉它在前面所有已得正交基上的投影，再归一化：
\begin{equation}\label{eq:gs}
\tilde{q}_k = a_k - \sum_{j=1}^{k-1} \underbrace{\inner{q_j}{a_k}}_{r_{jk}} q_j, \qquad
r_{kk} = \norm{\tilde{q}_k}_2 > 0, \qquad
q_k = \frac{\tilde{q}_k}{r_{kk}}.
\end{equation}
把式~\eqref{eq:gs} 反写出来：$a_k = \sum_{j=1}^k r_{jk} q_j$。按第 1 章的列线性组合视角打包成矩阵乘法，就是大名鼎鼎的\textbf{精简 QR 分解（Reduced QR Factorization）}：
\[
A = Q R, \qquad Q \in \R^{m \times n}\;(Q^T Q = I_n), \qquad R \in \R^{n \times n}\;\text{（对角元全正的上三角阵）}.
\]
将 $A = QR$ 代入正规方程 $A^T A \hat{w} = A^T b$：
\[
R^T (Q^T Q) R \hat{w} = R^T Q^T b \implies R^T R \hat{w} = R^T Q^T b.
\]
因为 $R$ 的对角元全正，故 $R^T$ 可逆，两端消去 $R^T$ 立得极简的\textbf{上三角回代方程}：
\begin{equation}\label{eq:qr-ls}
R \hat{w} = Q^T b.
\end{equation}

\begin{example}[对三点拟合矩阵手算 QR 分解并复核]\label{ex:3pt-qr}
仍取例~\ref{ex:3pt-ls} 的 $A = \begin{psmallmatrix} 1 & 0 \\ 1 & 1 \\ 1 & 2 \end{psmallmatrix}$ 与 $b = (1, 2, 4)^T$：
\begin{itemize}
  \item 第 1 列 $a_1 = (1, 1, 1)^T$：$r_{11} = \norm{a_1}_2 = \sqrt{3}$，故 $q_1 = \frac{1}{\sqrt{3}}(1, 1, 1)^T$。
  \item 第 2 列 $a_2 = (0, 1, 2)^T$：$r_{12} = \inner{q_1}{a_2} = \frac{3}{\sqrt{3}} = \sqrt{3}$，减去投影得
  \[
  \tilde{q}_2 = (0, 1, 2)^T - \sqrt{3}\cdot\frac{1}{\sqrt{3}}(1, 1, 1)^T = (-1, 0, 1)^T,
  \]
  长度 $r_{22} = \sqrt{2}$，故 $q_2 = \frac{1}{\sqrt{2}}(-1, 0, 1)^T$。
  \item 计算 $Q^T b = \begin{psmallmatrix} 7/\sqrt{3} \\ 3/\sqrt{2} \end{psmallmatrix}$，解上三角方程 $R\hat{w} = Q^T b$：
  \[
  \begin{pmatrix} \sqrt{3} & \sqrt{3} \\ 0 & \sqrt{2} \end{pmatrix}\begin{pmatrix} w_0 \\ w_1 \end{pmatrix}
  = \begin{pmatrix} 7/\sqrt{3} \\ 3/\sqrt{2} \end{pmatrix}
  \implies w_1 = \frac{3}{2}, \quad w_0 = \frac{7}{3} - \frac{3}{2} = \frac{5}{6},
  \]
  与正规方程手算结果一模一样！
\end{itemize}
\end{example}

既然正规方程与 QR 在精确算术下结果完全相同，为什么所有现代数值库（包括 \texttt{np.linalg.lstsq}）都强烈警告\textbf{绝不要显式写 \texttt{np.linalg.inv(A.T @ A) @ A.T @ b}}？

\begin{failBox}{数值精度杀手：显式构造 $A^T A$ 会把敏感度“平方放大”}
考虑一个极简的 $3 \times 2$ 矩阵（Läuchli 矩阵），取微小正数 $\varepsilon = 10^{-8}$：
\[
A = \begin{pmatrix} 1 & 1 \\ \varepsilon & 0 \\ 0 & \varepsilon \end{pmatrix}, \qquad
A^T A = \begin{pmatrix} 1 + \varepsilon^2 & 1 \\ 1 & 1 + \varepsilon^2 \end{pmatrix}.
\]
在双精度浮点数 \texttt{float64} 中（机器精度 $\varepsilon_{\text{mach}} \approx 2.22 \times 10^{-16}$），$A$ 本身两列清楚地线性无关；可是一旦乘出 $A^T A$，由于 $\varepsilon^2 = 10^{-16} < \varepsilon_{\text{mach}}/2$，计算机里 $1 + \varepsilon^2$ 会直接舍入成 $1.0$！原本列满秩的问题瞬间退化成奇异矩阵 $\begin{psmallmatrix}1 & 1 \\ 1 & 1\end{psmallmatrix}$，信息荡然无存！

用第 3 章将严格定义的语言来说，形成 $A^T A$ 会使矩阵条件数平方放大：$\cond(A^T A) = \cond(A)^2$。\textbf{工程铁律}：求解最小二乘时，永远直接调用基于正交变换（QR/SVD）的 \texttt{np.linalg.lstsq(A, b, rcond=None)}。
\end{failBox}

\begin{extBox}{改进 Gram-Schmidt 与 Householder 镜面反射直觉}
经典 Gram-Schmidt（CGS）在有限精度下容易因舍入误差让新向量逐渐丧失与前面基向量的正交性；稍微调整顺序、每减去一个分量就立刻更新 $\tilde{q}_k$ 的\textbf{改进 Gram-Schmidt（MGS）}会稳健不少。而 \texttt{np.linalg.qr} 底层真正采用的是精度更高的\textbf{Householder 镜面反射}：用形如 $H = I - 2 u u^T$（$\norm{u}_2 = 1$，满足 $H^T = H$ 与 $H^T H = I$）的保长正交反射矩阵，像折纸一样逐列把对角线下方元素反射清零，即使在 $\varepsilon = 10^{-8}$ 的 Läuchli 矩阵上也能保持 $10^{-15}$ 级机器精度（见习题~2.7 实测）。初读了解“工业界靠正交变换保精度”即可，不必手推反射公式。
\end{extBox}

\section{首个实战闭环：多项式回归、过拟合与岭回归初探}\label{sec:poly-ridge}

\subsection{特征映射：用线性最小二乘拟合非线性曲线}

现在我们把第~\ref{sec:proj-ls}~节的工具升级为真正的机器学习模型。读者常有一个误解：线性回归只能画直线。其实不然——\textbf{“线性模型”里的“线性”，指的是对未知参数 $w$ 呈线性，而对原始输入 $x$ 可以高度非线性}！

给定标量输入 $x \in \R$，我们构造一个 $d$ 次\textbf{多项式特征映射（Feature Map）}：
\[
\phi(x) = \begin{pmatrix} 1 & x & x^2 & \cdots & x^d \end{pmatrix}^T \in \R^{d+1}.
\]
模型预测值 $\hat{y}(x) = w^T \phi(x) = w_0 + w_1 x + w_2 x^2 + \dots + w_d x^d$ 虽然关于 $x$ 是蜿蜒曲折的 $d$ 次曲线，但关于参数 $w$ 却是纯粹的线性组合！把 $m$ 个训练样本 $(x_1, y_1), \dots, (x_m, y_m)$ 堆叠成\textbf{范德蒙德（Vandermonde）设计矩阵}：
\[
A = \begin{pmatrix}
1 & x_1 & x_1^2 & \cdots & x_1^d \\
1 & x_2 & x_2^2 & \cdots & x_2^d \\
\vdots & \vdots & \vdots & \ddots & \vdots \\
1 & x_m & x_m^2 & \cdots & x_m^d
\end{pmatrix} \in \R^{m \times (d+1)},
\]
拟合非线性曲线的问题就原封不动地回到了线性最小二乘 $\min_{w} \norm{Aw - y}_2^2$！

\subsection{过拟合危机与岭回归（Ridge Regression）}

是不是多项式次数 $d$ 越高越好呢？当我们只有 $m = 15$ 个带噪声的训练点、却硬塞进 $d = 9$ 甚至 $d = 11$ 次多项式时，灾难发生了：一方面，$x^k$ 与 $x^{k+1}$ 在区间 $[0, 1]$ 上长得极像，导致设计矩阵 $A$ 的各列近乎共线；另一方面，为了死命穿过每一个带噪声的训练点，普通最小二乘（OLS）会把多项式系数推高到惊人的正负百万级，靠巨额正负抵消在训练点间剧烈震荡——这就是机器学习的核心宿敌：\textbf{过拟合（Overfitting）}！

怎么驯服发疯的系数？最经典的方法是在损失函数后面加上一项 $\ell_2$ 惩罚（约束参数不能太大），得到\textbf{岭回归（Ridge Regression）}：
\begin{equation}\label{eq:ridge-obj}
\min_{w \in \R^{d+1}} \; \norm{Aw - y}_2^2 + \lambda \norm{w}_2^2 \qquad (\lambda > 0),
\end{equation}
其中超参数 $\lambda > 0$ 控制“贴合训练数据”与“保持权重简洁平滑”之间的平衡。

绝妙的是，我们\textbf{不需要微积分，也不需要第 3 章的谱定理}，用一个小技巧就能把岭回归直接变回刚学过的普通最小二乘！把两项平方和拼成一个\textbf{增广向量的欧氏范数平方}：
\begin{equation}\label{eq:ridge-aug}
\norm{Aw - y}_2^2 + \lambda \norm{w}_2^2
= \norm{\underbrace{\begin{pmatrix} A \\ \sqrt{\lambda} I \end{pmatrix}}_{\tilde{A}} w - \underbrace{\begin{pmatrix} y \\ 0 \end{pmatrix}}_{\tilde{y}}}_2^2.
\end{equation}
立刻套用第~\ref{sec:proj-ls}~节的正规方程 $\tilde{A}^T \tilde{A} \hat{w}_\lambda = \tilde{A}^T \tilde{y}$，展开即得岭回归解析解方程：
\begin{equation}\label{eq:ridge-normal}
(A^T A + \lambda I) \hat{w}_\lambda = A^T y.
\end{equation}
而且，这个加上“山脊”$\lambda I$ 的矩阵永远不会奇异——甚至在样本数 $m < n$ 或特征完全重复时都如此！

\begin{thmBox}
\begin{theorem}[岭回归矩阵恒可逆定理]\label{thm:ridge-inv}
对任意形状、任意秩的矩阵 $A \in \R^{m \times n}$ 及任意正则化系数 $\lambda > 0$，方阵 $A^T A + \lambda I_n$ 必然可逆，因此岭回归问题~\eqref{eq:ridge-obj} 存在唯一解析解 $\hat{w}_\lambda = (A^T A + \lambda I_n)^{-1} A^T y$。
\end{theorem}
\end{thmBox}

\begin{proof}
设向量 $v \in \R^n$ 满足 $(A^T A + \lambda I_n) v = 0$。两边左乘 $v^T$ 得
\[
0 = v^T (A^T A + \lambda I_n) v = \norm{Av}_2^2 + \lambda \norm{v}_2^2.
\]
由于 $\norm{Av}_2^2 \ge 0$ 且 $\lambda > 0$，必有 $\lambda \norm{v}_2^2 = 0 \implies \norm{v}_2 = 0 \implies v = 0$。由此 $\nul(A^T A + \lambda I_n) = \{0\}$，故方阵 $A^T A + \lambda I_n$ 恒可逆。
\end{proof}

\begin{remark}[工程细节：为何常豁免截距项 $w_0$？]
在具体实现中，截距项 $w_0$ 仅负责把拟合曲线整体上下平移，不控制曲线的弯曲剧烈程度；压缩 $w_0$ 反而会引入无谓的整体偏差。因此工程上常令惩罚矩阵 $P = \diag(0, 1, \dots, 1)$，仅惩罚非常数项权重 $w_1, \dots, w_d$，求解 $(A^T A + \lambda P)\hat{w}_\lambda = A^T y$。只要设计矩阵第 1 列是全 $1$ 向量 $\mathbf{1} \neq 0$，用与定理~\ref{thm:ridge-inv} 完全相同的能量论证可知：$(A^T A + \lambda P)v = 0 \implies \norm{Av}_2^2 + \lambda\sum_{j=1}^d v_j^2 = 0 \implies v_1 = \dots = v_d = 0$，继而 $\norm{v_0 \mathbf{1}}_2^2 = m v_0^2 = 0 \implies v_0 = 0$，矩阵依然恒可逆！
\end{remark}

\subsection{实测数据、U 型泛化曲线与可运行代码}

在配套脚本 \texttt{code/vol1\_experiments.py}（函数 \texttt{exp2\_poly\_ridge}）中，我们设定真实规律为正弦曲线 $y = \sin(2\pi x)$（$x \in [0, 1]$），仅抽取 $m = 15$ 个叠加标准差 $\sigma = 0.25$ 高斯噪声的训练样本，并在 $200$ 个无噪测试点上评估均方误差（MSE）。表~\ref{tab:poly-ridge} 与图~\ref{fig:poly-ridge} 完整展示了从欠拟合、甜蜜点到严重过拟合、再到岭回归救场的真实实测曲线。

\begin{table}[htbp]
\centering
\caption{多项式回归（OLS vs 岭回归 $\lambda = 10^{-3}$）实测对比（源自 \texttt{data/poly\_ridge.csv}）}\label{tab:poly-ridge}
{\small
\begin{tabular}{@{}rcccc@{}}
\toprule
\textbf{多项式次数 $d$} & \textbf{OLS 训练 MSE} & \textbf{OLS 测试 MSE} & \textbf{岭回归测试 MSE} & \textbf{OLS 权重范数 $\norm{w_{\text{OLS}}}_2$} \\
\midrule
1 & 0.2472 & 0.2288 & 0.2289 & 1.44 \\
2 & 0.2421 & 0.2330 & 0.2329 & 2.43 \\
3 & 0.0403 & 0.0160 & 0.0514 & 36.49 \\
5 & 0.0347 & \textbf{0.0145} & 0.0296 & 123.05 \\
7 & 0.0305 & 0.0211 & 0.0257 & 5,324.09 \\
9 & 0.0111 & \textcolor{ochre}{\textbf{0.0613}} & 0.0218 & 402,933.89 \\
11 & \textbf{0.0048} & \textcolor{ochre}{0.0484} & \textbf{0.0216} & \textcolor{ochre}{\textbf{9,940,472.74}} \\
\bottomrule
\end{tabular}}
\end{table}

\begin{figure}[htbp]
\centering
\begin{tikzpicture}
\begin{axis}[
  width=0.84\linewidth,
  height=5.8cm,
  xlabel={多项式次数 $d$},
  ylabel={均方误差（MSE）},
  xtick={1,2,3,5,7,9,11},
  ymin=0, ymax=0.27,
  grid=both,
  grid style={line width=0.25pt, draw=codeborder},
  legend pos=north east,
  legend style={font=\footnotesize, fill=white, draw=codeborder},
  tick label style={font=\scriptsize},
  label style={font=\small\sffamily}
]
  \addplot[thick, dashed, mark=triangle*, mark size=2.3pt, color=slate]
    table[x=deg, y=train_mse, col sep=comma]{data/poly_ridge.csv};
  \addlegendentry{OLS 训练 MSE（单调逼近 0）}
  \addplot[very thick, mark=*, mark size=2.4pt, color=ochre]
    table[x=deg, y=test_mse, col sep=comma]{data/poly_ridge.csv};
  \addlegendentry{OLS 测试 MSE（经典 U 型过拟合）}
  \addplot[very thick, mark=square*, mark size=2.2pt, color=teal]
    table[x=deg, y=ridge_test_mse, col sep=comma]{data/poly_ridge.csv};
  \addlegendentry{岭回归测试 MSE（$\lambda = 10^{-3}$，稳健抑制）}
  % Annotations
  \node[anchor=center, font=\scriptsize\sffamily, text=slate, fill=white, inner sep=1.5pt] at (axis cs:1.45, 0.165) {欠拟合区};
  \node[anchor=south, font=\scriptsize\sffamily, text=tealDark] at (axis cs:5, 0.044) {甜蜜点 $d=5$};
  \node[anchor=south, font=\scriptsize\sffamily, text=ochre] at (axis cs:9.1, 0.074) {过拟合抬头（$\|w\|_2 \sim 10^7$）};
\end{axis}
\end{tikzpicture}
\caption{多项式回归的 U 型泛化曲线与岭回归抗过拟合实测（数据直接读取自 \texttt{data/poly\_ridge.csv}）}\label{fig:poly-ridge}
\end{figure}

仔细品读表~\ref{tab:poly-ridge} 与图~\ref{fig:poly-ridge} 的三个阶段：
\begin{enumerate}
  \item \textbf{欠拟合（$d = 1, 2$）}：直线或抛物线根本弯不出一个完整正弦波，训练 MSE 与测试 MSE 双双高达 $0.23$ 以上。
  \item \textbf{甜蜜点（$d = 3, 5$）}：模型容量刚好匹配真实曲线，OLS 测试 MSE 骤降至 $0.0145$。
  \item \textbf{过拟合与岭回归救场（$d = 9, 11$）}：OLS 训练 MSE 继续跌到 $0.0048$，但测试 MSE 反而反弹回升至 $0.0613$，权重范数爆炸到 $9.94 \times 10^6$！加上极小的 $\lambda = 10^{-3}$ 岭惩罚后，即使在 $d = 11$ 阶高次多项式上，测试 MSE 也平稳压制在 $0.0216$，权重范数暴降数十万倍（仅约 $14.7$）。
\end{enumerate}

代码~\ref{lst:poly-ridge} 给出了完整的纯 NumPy 实战代码，建议你亲手运行并试着改动样本量 $m$ 与 $\lambda$ 观察曲线变化。

\begin{lstlisting}[caption={多项式 OLS 与岭回归实战闭环（复现表~\ref{tab:poly-ridge} 实测数据）}, label={lst:poly-ridge}]
import numpy as np

rng = np.random.default_rng(42)
x_tr = np.linspace(0.0, 1.0, 15)
y_tr = np.sin(2.0 * np.pi * x_tr) + rng.normal(0.0, 0.25, size=x_tr.shape)
x_te = np.linspace(0.0, 1.0, 200)
y_te = np.sin(2.0 * np.pi * x_te)

def vandermonde(x: np.ndarray, deg: int) -> np.ndarray:
    return np.column_stack([x ** k for k in range(deg + 1)])

def fit_ridge_augmented(A: np.ndarray, y: np.ndarray, lam: float) -> np.ndarray:
    """用增广矩阵 + lstsq 稳健求解岭回归（豁免截距项 w_0）。"""
    n_feat = A.shape[1]
    reg = np.sqrt(lam) * np.eye(n_feat)
    reg[0, 0] = 0.0
    A_aug = np.vstack([A, reg])
    y_aug = np.concatenate([y, np.zeros(n_feat)])
    w, *_ = np.linalg.lstsq(A_aug, y_aug, rcond=None)
    return w

for d in [1, 3, 5, 9, 11]:
    A_tr, A_te = vandermonde(x_tr, d), vandermonde(x_te, d)
    w_ols, *_ = np.linalg.lstsq(A_tr, y_tr, rcond=None)
    w_rid = fit_ridge_augmented(A_tr, y_tr, lam=1e-3)
    te_ols = float(np.mean((A_te @ w_ols - y_te) ** 2))
    te_rid = float(np.mean((A_te @ w_rid - y_te) ** 2))
    print(f"deg={d:2d} | OLS te={te_ols:.4f} (||w||={np.linalg.norm(w_ols):.1f})"
          f" | Ridge te={te_rid:.4f} (||w||={np.linalg.norm(w_rid):.1f})")
\end{lstlisting}

\begin{insightBox}{下游坐标：你刚刚摸到了机器学习的两大支柱}
在本章我们纯靠“正交投影 + $\ell_2$ 权重惩罚”推出了岭回归；到第 2 卷学完概率论与贝叶斯推断后，你会惊喜地发现：\textbf{最小二乘等价于高斯噪声下的极大似然估计（MLE），而岭回归恰好等价于给权重加上零均值高斯先验后的最大后验估计（MAP）}！到第 4 卷，图~\ref{fig:poly-ridge} 的 U 型曲线还将由“偏差-方差分解（Bias-Variance Tradeoff）”给出精确的统计解释。
\end{insightBox}

\section{本章习题}\label{sec:ch02-exercises}

\exTier{A. 概念辨析与反例查错}{抓清边界条件，避免似是而非}

\begin{enumerate}[label=\textbf{2.\arabic*}]
  \item \textbf{范数公理、斜投影与余弦相似度辨析}：
  （1）定义 $\R^2$ 上的函数 $f(x) = (\sqrt{\abs{x_1}} + \sqrt{\abs{x_2}})^2$（即形式上的“$\ell_{1/2}$ 范数”）。它是否构成真正的范数？请用具体的二维向量 $u, v \in \R^2$ 验证它是否违反三角不等式。
  （2）构造一个具体的 $2 \times 2$ 实矩阵 $P$，使其满足幂等性 $P^2 = P$ 但\textbf{不对称}（$P^T \neq P$），并用一个具体向量 $b$ 说明 $Pb$ 不是正交投影（即 $b - Pb$ 不垂直于 $Pb$）。
  （3）判断正误并修正：“两个非零向量 $u, v$ 满足 $\cos\theta(u, v) = 1$ 当且仅当 $u = v$。”

  \item \textbf{线性最小二乘与过拟合误区诊断}：
  （1）设超定问题 $\min_w \norm{Aw - b}_2^2$ 中的设计矩阵 $A \in \R^{m \times n}$ 有两列完全相同（从而列不满秩）。请问：最小二乘最优参数向量 $\hat{w}$ 是否唯一？最优拟合预测向量 $\hat{b} = A\hat{w}$ 是否唯一？请结合列空间正交投影说明理由。
  （2）在表~\ref{tab:poly-ridge} 的多项式实验中，某同学见 $d = 11$ 次普通最小二乘的训练 MSE 已经降到全场最低的 $0.0048$，便认为 $d=11$ 是最优模型。请指出该判断的两处关键疏漏，并解释为何 $\norm{w_{\text{OLS}}}_2 \approx 9.94 \times 10^6$ 是明显的过拟合危险信号。
\end{enumerate}

\exTier{B. 手算小例与基础演练}{把四大子空间与 QR 最小二乘彻底算透}

\begin{enumerate}[label=\textbf{2.\arabic*}, resume]
  \item \textbf{四大基本子空间与正交性手算}：给定矩阵
  \[
  A = \begin{pmatrix} 1 & 0 & 1 \\ 0 & 1 & 1 \\ 1 & 1 & 2 \end{pmatrix}.
  \]
  求 $\row(A)$、$\nul(A)$、$\col(A)$ 与 $\nul(A^T)$ 各自的一组基，验证它们的维数满足 $\dim\row(A) + \dim\nul(A) = 3$ 与 $\dim\col(A) + \dim\nul(A^T) = 3$，并逐对检验基向量之间的正交关系 $\row(A) \perp \nul(A)$ 与 $\col(A) \perp \nul(A^T)$。

  \item \textbf{正规方程与 Gram-Schmidt QR 手算对比}：给定 4 个数据点 $(-1, 0),\, (0, 1),\, (1, 1),\, (2, 3)$，拟用直线 $y = w_0 + w_1 x$ 进行线性最小二乘拟合。
  （1）写出设计矩阵 $A$ 与标签向量 $b$，用正规方程 $A^T A \hat{w} = A^T b$ 手算最优参数 $\hat{w}$ 与残差 $e = b - A\hat{w}$，并验证 $A^T e = 0$。
  （2）对 $A$ 手算精简 QR 分解 $A = QR$，再解上三角方程 $R\hat{w} = Q^T b$，确认结果与第（1）问完全一致。
\end{enumerate}

\exTier{C. 实用证明与理论推导}{纯靠内积与最优性完成严格证明，零超前依赖}

\begin{enumerate}[label=\textbf{2.\arabic*}, resume]
  \item \textbf{范数等价常数与互补正交投影}：
  （1）利用柯西-施瓦茨不等式证明：对任意 $x \in \R^n$，恒有
  \[
  \norm{x}_\infty \le \norm{x}_2 \le \norm{x}_1 \le \sqrt{n}\norm{x}_2,
  \]
  并分别举出使每个不等号取等号的非零向量 $x \in \R^n$。
  （2）设 $P \in \R^{m \times m}$ 是正交投影矩阵（$P^2 = P = P^T$），证明 $I_m - P$ 也是正交投影矩阵，且对任意 $b \in \R^m$ 恒有能量分解等式 $\norm{b}_2^2 = \norm{Pb}_2^2 + \norm{(I_m - P)b}_2^2$。
  \exHint{第（1）问证 $\norm{x}_1 \le \sqrt{n}\norm{x}_2$ 时，把 $\norm{x}_1$ 写成向量 $(\abs{x_1}, \dots, \abs{x_n})^T$ 与全 $1$ 向量 $(1, \dots, 1)^T$ 的内积；第（2）问直接验证 $(I-P)^2 = I-P = (I-P)^T$ 并检查 $\inner{Pb}{(I-P)b} = 0$。}

  \item \textbf{岭回归解的单调收缩与拟合代价权衡}：设 $A \in \R^{m \times n}$，$y \in \R^m$。对任意 $\lambda > 0$，记 $\hat{w}_\lambda$ 为岭回归目标函数 $L_\lambda(w) = \norm{Aw - y}_2^2 + \lambda \norm{w}_2^2$ 的唯一极小值点。证明：若 $0 < \lambda_1 < \lambda_2$，则必有
  \[
  \norm{\hat{w}_{\lambda_2}}_2 \le \norm{\hat{w}_{\lambda_1}}_2
  \qquad \text{且} \qquad
  \norm{A\hat{w}_{\lambda_2} - y}_2^2 \ge \norm{A\hat{w}_{\lambda_1} - y}_2^2.
  \]
  也就是说，加大正则化系数 $\lambda$ 必然单调收缩权重范数，并以增大训练残差平方和为代价。
  \exHint{全程不需要求导或特征值！利用极小值定义写出 $L_{\lambda_1}(\hat{w}_{\lambda_1}) \le L_{\lambda_1}(\hat{w}_{\lambda_2})$ 与 $L_{\lambda_2}(\hat{w}_{\lambda_2}) \le L_{\lambda_2}(\hat{w}_{\lambda_1})$，将两式相加消去残差项即可推出第一个不等式，再回代即得第二个不等式。}
\end{enumerate}

\exTier{D. 数值编程与可自检实验}{用真实代码亲眼见证数值稳定性与正则化威力}

\begin{enumerate}[label=\textbf{2.\arabic*}, resume]
  \item \textbf{Läuchli 病态矩阵上正规方程 vs 正交分解稳定性实测}：构造矩阵 $A = \begin{psmallmatrix} 1 & 1 \\ \varepsilon & 0 \\ 0 & \varepsilon \end{psmallmatrix}$，取真实参数 $w^\star = (1, 2)^T$ 及无噪观测 $b = A w^\star = (3, \varepsilon, 2\varepsilon)^T$。令 $\varepsilon$ 依次取 $10^{-3}, 10^{-5}, 10^{-7}, 10^{-8}$，分别用以下三种方法求 $\hat{w}$ 并打印误差 $\norm{\hat{w} - w^\star}_2$：
  \begin{enumerate}[label=(\alph*)]
    \item 显式正规方程：用 \texttt{np.linalg.solve(A.T @ A, A.T @ b)} 求解，并捕获可能出现的 \texttt{LinAlgError} 奇异报错；
    \item Householder QR 回代：\texttt{Q, R = np.linalg.qr(A, mode="reduced")} 后解 \texttt{R @ w = Q.T @ b}；
    \item 工业标准求解器：\texttt{np.linalg.lstsq(A, b, rcond=None)}。
  \end{enumerate}
  \textbf{自检标准}：验证在 $\varepsilon = 10^{-7}$ 时正规方程误差已劣化至约 $10^{-2}$ 量级、在 $\varepsilon = 10^{-8}$ 时触发奇异矩阵报错，而 Householder QR 与 \texttt{lstsq} 在全部四个 $\varepsilon$ 下误差均保持在 \texttt{< 1e-14}。
  \exHint{可顺便自己手写一份经典/改进 Gram-Schmidt 函数跑同一组 $\varepsilon$，观察它为何也会在 $\varepsilon = 10^{-8}$ 处失真，体会 Householder 反射的价值。}

  \item \textbf{高维近正交验证与岭回归 $\lambda$ 扫描实验}：
  （1）用 \texttt{rng = np.random.default\_rng(2026)}，在维度 $d \in \{10, 100, 1000\}$ 下各生成 $5{,}000$ 对随机单位向量，计算余弦相似度样本标准差，验证其与理论值 $1/\sqrt{d}$ 的相对误差均 \texttt{< 5\%}，且在 $d = 1000$ 时 $P(\abs{\cos\theta} < 0.1) > 99\%$。
  （2）在代码~\ref{lst:poly-ridge} 的 $m=15$ 含噪正弦数据上固定多项式次数 $d = 11$，令 $\lambda$ 扫描 $\{0, 10^{-6}, 10^{-4}, 10^{-3}, 10^{-2}, 10^{-1}, 1\}$，用\textbf{增广矩阵 + \texttt{np.linalg.lstsq}}（豁免截距 $w_0$）求解岭回归，打印每组的训练 MSE、测试 MSE 与 $\norm{w}_2$。
  \textbf{自检标准}：验证随着 $\lambda$ 从 $10^{-6}$ 增至 $1$，$\norm{w}_2$ 单调递减而训练 MSE 单调递增（印证习题~2.6）；同时验证在 $\lambda = 10^{-3}$ 处测试 MSE 降至 \texttt{< 0.025}（约 $0.0216$），$\norm{w}_2$ 较 OLS（约 $9.94 \times 10^6$）缩小超过 $5 \times 10^5$ 倍。
  \exHint{注意当 $\lambda = 0$ 时直接对原矩阵 $A_{\text{tr}}$ 调用 \texttt{np.linalg.lstsq}，避免构造零增广块。}
\end{enumerate}

\chapter[特征分解、谱定理与半正定矩阵]{特征分解、实对称阵谱定理、半正定矩阵与特殊 AI 矩阵}
\label{ch:ch3}

\chapEpigraph{对称性，无论你赋予它多宽泛或多狭窄的含义，都是人类千百年来用以理解并创造秩序、美与完美的观念。}{赫尔曼·外尔（Hermann Weyl），《对称》}

\begin{chapIntro}{本章定位与零循环路线}
如果说第 2 章用正交投影驯服了线性方程组，那么本章则要把方阵当成一台反复运转的动力机器：看清哪些特殊方向在矩阵作用下只做伸缩、绝不偏转。在人工智能中，我们\textbf{不需要}花几十页去手算高阶矩阵的复杂 Jordan 块；相反，我们将从特征值与特征向量的 2D 几何直觉出发，依托第 2 章已经建立的\textbf{行列式体积直觉、内积、正交补分解 $\R^n = U \oplus U^\perp$ 与四大基本子空间}，沿一条\textbf{零循环、零跳步}的路线推进：
\begin{enumerate}
  \item 从几何画面直观看懂\textbf{特征值、特征向量与特征方程 $\det(A-\lambda I)=0$}，并用 $2\times 2$ 剪切矩阵直观理解 Jordan 链引发的“瞬态暴涨”现象，看清为何现代数值 AI 避开一般 Jordan 分解；
  \item 从零严格证明\textbf{实对称矩阵谱定理} $S = Q\Lambda Q^T$，一步推出 Rayleigh 商界；
  \item 作为谱定理的直接推论，\textbf{无循环地}证出正定/半正定判据（$Q \succ 0 \iff \lambda_i > 0$），并解剖半正定矩阵（PSD）在 AI 中的三大核心化身（协方差阵、核矩阵、Hessian）与条件数 $\kappa$；
  \item 轻量认识两类在 AI 中随处可见的特殊矩阵：\textbf{行随机矩阵}（马尔可夫链与 Transformer Softmax 注意力）与\textbf{图拉普拉斯矩阵}（图信号平滑度与谱聚类）。
\end{enumerate}
\end{chapIntro}

% =========================================================================
\section{特征值几何直觉、特征分解与 Jordan 链}
\label{sec:ch3-eigen-jordan}

\subsection{从几何直觉理解特征值、特征向量与特征多项式}

回忆第 2 章第~\ref{sec:linalg-bridge}~节的线性变换视角：当我们用方阵 $A \in \R^{n\times n}$ 作用在空间里的向量 $x$ 上时，绝大多数向量都会被推离原来的直线（既旋转又伸缩）。然而，在整片空间翻腾变形之际，往往藏着几条极为特殊的方向——落在这些直线上的非零向量被 $A$ 作用后\textbf{依然留在自己原来的直线上}，仅仅被拉长、缩短或反向翻转！

\begin{defBox}
\begin{definition}[特征值与特征向量]\label{def:ch3-eigen}
设 $A \in \R^{n\times n}$ 为方阵。若存在标量 $\lambda$（可能为实数或复数）与\textbf{非零}向量 $v \neq 0$，使得
\begin{equation}\label{eq:ch3-eig-def}
  A v = \lambda v,
\end{equation}
则称 $\lambda$ 为 $A$ 的\textbf{特征值（Eigenvalue）}，$v$ 为 $A$ 属于 $\lambda$ 的\textbf{特征向量（Eigenvector）}。几何上，$v$ 是变换 $A$ 下方向不变的“宿命轴线”，$\lambda$ 是沿该轴线的纯粹伸缩倍数（$\lambda < 0$ 表示反向翻转，$\lambda = 0$ 表示被压扁归零）。
\end{definition}
\end{defBox}

以矩阵 $A = \begin{pmatrix} 3 & 1 \\ 0 & 2 \end{pmatrix}$ 为例，图~\ref{fig:ch3-eig-geom} 直观展示了它对平面上三支向量的作用：
\begin{itemize}
  \item $v_1 = (1, 0)^T \mapsto A v_1 = (3, 0)^T = 3v_1$：稳稳留在水平轴上，纯粹拉长 $3$ 倍；
  \item $v_2 = (-1, 1)^T \mapsto A v_2 = (-2, 2)^T = 2v_2$：稳稳留在 $-45^\circ$ 对角线上，纯粹拉长 $2$ 倍；
  \item 普通向量 $u = (0, 1)^T \mapsto A u = (1, 2)^T$：被彻底推离了竖直轴，不是特征向量。
\end{itemize}

\begin{figure}[htbp]
\centering
\resizebox{0.96\linewidth}{!}{%
\begin{tikzpicture}[font=\sffamily\footnotesize]
% Left Panel: Before transformation
\begin{scope}[shift={(0,0)}]
  \node[font=\sffamily\bfseries\small, text=ink] at (0.35, 2.75) {(1) 变换前：两条不变直线与普通向量};
  \draw[step=0.65cm, codeborder!75, very thin] (-1.6,-0.55) grid (2.3, 2.05);
  \draw[->, slate!75, thin] (-1.65, 0) -- (2.45, 0) node[right] {$x_1$};
  \draw[->, slate!75, thin] (0, -0.60) -- (0, 2.15) node[right=2pt] {$x_2$};
  % Span lines of eigenvectors (unit = 0.65cm)
  \draw[tealLine, thick, dashed] (-1.55, 0) -- (2.25, 0);
  \draw[navy!40, thick, dashed] (-1.55, 1.55) -- (0.45, -0.45);
  \draw[ochre!45, thin, dotted] (0, -0.45) -- (0, 1.95);
  % Vectors: v1=(1,0), v2=(-1,1), u=(0,1)
  \draw[-{Stealth[length=2.2mm]}, teal, very thick] (0,0) -- (0.65, 0) node[below right=1pt and -4pt, text=tealDark] {$v_1=(1,0)^T$};
  \draw[-{Stealth[length=2.2mm]}, navy, very thick] (0,0) -- (-0.65, 0.65) node[below left=-1pt, text=navy] {$v_2=(-1,1)^T$};
  \draw[-{Stealth[length=2.2mm]}, ochre, very thick] (0,0) -- (0, 0.65) node[above right=-1pt, text=ochre] {$u=(0,1)^T$};
  \fill[ink] (0,0) circle (1.7pt);
\end{scope}

% Transition Arrow
\draw[-{Stealth[length=2.6mm]}, teal, very thick] (2.85, 0.85) -- (4.35, 0.85)
  node[midway, above=2pt, text=tealDark, font=\footnotesize\bfseries] {$x \mapsto Ax$}
  node[midway, below=2pt, text=slate, font=\scriptsize] {$A=\begin{psmallmatrix}3&1\\0&2\end{psmallmatrix}$};

% Right Panel: After transformation
\begin{scope}[shift={(6.45,0)}]
  \node[font=\sffamily\bfseries\small, text=tealDark] at (0.35, 2.75) {(2) 变换后：$v_1,v_2$ 方向不变，$u$ 偏离原轴};
  \draw[step=0.65cm, codeborder!75, very thin] (-1.6,-0.55) grid (2.3, 2.05);
  \draw[->, slate!75, thin] (-1.65, 0) -- (2.45, 0) node[right] {$x_1$};
  \draw[->, slate!75, thin] (0, -0.60) -- (0, 2.15) node[right=2pt] {$x_2$};
  % Same invariant lines
  \draw[tealLine, thick, dashed] (-1.55, 0) -- (2.25, 0);
  \draw[navy!40, thick, dashed] (-1.55, 1.55) -- (0.45, -0.45);
  \draw[ochre!45, thin, dotted] (0, -0.45) -- (0, 1.95);
  % Ghost originals
  \draw[-{Stealth[length=1.8mm]}, teal!40, thick] (0,0) -- (0.65, 0);
  \draw[-{Stealth[length=1.8mm]}, navy!35, thick] (0,0) -- (-0.65, 0.65);
  \draw[-{Stealth[length=1.8mm]}, ochre!40, thick] (0,0) -- (0, 0.65);
  % Transformed vectors: Av1=(3,0)->(1.95,0), Av2=(-2,2)->(-1.30,1.30), Au=(1,2)->(0.65,1.30)
  \draw[-{Stealth[length=2.3mm]}, teal, very thick] (0,0) -- (1.95, 0)
    node[below=2pt, text=tealDark] {$Av_1 = 3v_1$};
  \draw[-{Stealth[length=2.3mm]}, navy, very thick] (0,0) -- (-1.30, 1.30)
    node[above=1pt, text=navy] {$Av_2 = 2v_2$};
  \draw[-{Stealth[length=2.3mm]}, ochre, very thick] (0,0) -- (0.65, 1.30)
    node[above right=-1pt, text=ochre] {$Au=(1,2)^T$};
  % Curved deflection arrow from vertical axis to Au
  \draw[-{Stealth[length=1.8mm]}, ochre, dashed] (0.04, 1.25) to[bend left=18] (0.58, 1.25);
  \node[text=ochre, font=\scriptsize] at (1.42, 0.85) {被推离原轴！};
  \fill[ink] (0,0) circle (1.7pt);
\end{scope}
\end{tikzpicture}}
\caption{特征值与特征向量的 2D 几何画面：经 $A=\begin{psmallmatrix}3&1\\0&2\end{psmallmatrix}$ 变换后，只有落在特征直线上的 $v_1, v_2$ 保持方向不变、只做倍数伸缩，普通向量 $u$ 则被推离原直线}\label{fig:ch3-eig-geom}
\end{figure}

那么，给定任意方阵 $A$，怎么系统找出所有特征值 $\lambda$？把定义式 $Av = \lambda v$ 改写成齐次方程：
\[
  (A - \lambda I) v = 0 \qquad (v \neq 0).
\]
此刻，第 2 章第~\ref{sec:linalg-bridge}~节建立的\textbf{行列式压扁直觉}（式~\eqref{eq:singular-equiv}）立即发威：矩阵 $A - \lambda I$ 会把某个非零向量 $v$ 压扁归零（即 $\nul(A - \lambda I) \neq \{0\}$），\textbf{当且仅当它把整个空间的体积压扁为零}！于是立刻得到\textbf{特征方程}：
\begin{equation}\label{eq:ch3-charpoly}
  \det(A - \lambda I) = 0 \quad \Longleftrightarrow \quad p_A(\lambda) \coloneqq \det(\lambda I - A) = 0.
\end{equation}
拿刚才的 $A = \begin{psmallmatrix}3 & 1 \\ 0 & 2\end{psmallmatrix}$ 手算验证一遍：
\[
  p_A(\lambda) = \det\begin{pmatrix} \lambda - 3 & -1 \\ 0 & \lambda - 2 \end{pmatrix} = (\lambda - 3)(\lambda - 2) = 0 \implies \lambda_1 = 3,\; \lambda_2 = 2.
\]
代回求解零空间：$(A - 3I)v = \begin{psmallmatrix}0 & 1 \\ 0 & -1\end{psmallmatrix}v = 0 \Rightarrow v_1 = (1, 0)^T$；$(A - 2I)v = \begin{psmallmatrix}1 & 1 \\ 0 & 0\end{psmallmatrix}v = 0 \Rightarrow v_2 = (-1, 1)^T$。同时验证两条极其好用的速查恒等式：$\tr(A) = 3 + 2 = 5 = \lambda_1 + \lambda_2$，且 $\det(A) = 3 \times 2 = 6 = \lambda_1 \lambda_2$！

\begin{recallBox}{特征多项式、相似不变性与重数速览}
一般地，对 $n$ 阶矩阵 $A \in \R^{n\times n}$，特征多项式 $p_A(\lambda) = \det(\lambda I - A)$ 是关于 $\lambda$ 的 $n$ 次首一多项式，在复数域 $\C$ 上有 $n$ 个根（计重数），且恒满足 $\tr(A) = \sum_{i=1}^n \lambda_i$ 与 $\det(A) = \prod_{i=1}^n \lambda_i$。若存在可逆阵 $P$ 使 $B = P^{-1}AP$（几何上相当于换了一组坐标系看同一线性变换），则称 $A$ 与 $B$ \textbf{相似}，它们具有完全相同的特征多项式、迹、行列式与特征值。对任意特征值 $\lambda$，它在 $p_A$ 中的重根次数称为\textbf{代数重数} $m_a(\lambda)$，特征子空间维数 $m_g(\lambda) \coloneqq \dim\nul(A - \lambda I)$ 称为\textbf{几何重数}，恒有 $1 \le m_g(\lambda) \le m_a(\lambda)$。
\end{recallBox}

\subsection{特征分解、矩阵幂与离散动力稳定性}

在人工智能的迭代算法中（例如第 6 章的梯度下降 $e^{(t)} = (I - \eta Q)^t e^{(0)}$，以及第 5 卷循环神经网络 RNN 的隐状态演化 $h^{(t)} = W h^{(t-1)}$），我们反复面对同一个问题：\textbf{把矩阵 $A$ 连乘 $k$ 次，向量 $A^k x^{(0)}$ 最终会趋于零、保持稳定，还是指数爆炸？}特征分解给出了最清晰的坐标解耦视角。我们先补全一条短小精悍的基础定理。

\begin{thmBox}
\begin{theorem}[互异特征值对应特征向量的线性无关性]
\label{thm:ch3-distinct-indep}
设 $v_1, \dots, v_m$ 分别是方阵 $A$ 对应于\textbf{互不相同}的特征值 $\lambda_1, \dots, \lambda_m$ 的特征向量，则 $v_1, \dots, v_m$ 线性无关。
\end{theorem}
\end{thmBox}

\begin{proof}
对向量个数 $m$ 施数学归纳法。当 $m=1$ 时，特征向量非零，自然线性无关。假设对任意 $m-1$ 个互异特征值的特征向量结论成立，现设存在标量 $c_1, \dots, c_m$ 使得
\begin{equation}
  c_1 v_1 + c_2 v_2 + \dots + c_m v_m = 0. \label{eq:ch3-lincomb}
\end{equation}
将 \eqref{eq:ch3-lincomb} 两边左乘 $A$，利用 $Av_i = \lambda_i v_i$ 得
\[
  c_1 \lambda_1 v_1 + c_2 \lambda_2 v_2 + \dots + c_m \lambda_m v_m = 0.
\]
再将 \eqref{eq:ch3-lincomb} 两边乘以 $\lambda_m$ 后与上式相减，消去最后一项：
\[
  c_1(\lambda_1 - \lambda_m) v_1 + c_2(\lambda_2 - \lambda_m) v_2 + \dots + c_{m-1}(\lambda_{m-1} - \lambda_m) v_{m-1} = 0.
\]
由归纳假设，$v_1, \dots, v_{m-1}$ 线性无关，故对每个 $i \le m-1$ 均有 $c_i(\lambda_i - \lambda_m) = 0$。因特征值互异（$\lambda_i \neq \lambda_m$），必有 $c_1 = \dots = c_{m-1} = 0$。代回 \eqref{eq:ch3-lincomb} 得 $c_m v_m = 0$，又 $v_m \neq 0$，从而 $c_m = 0$。
\end{proof}

由此立即得到可对角化判据与矩阵幂的显式公式。

\begin{proposition}[可对角化充要条件与解耦矩阵幂]
\label{prop:ch3-diag}
$n$ 阶矩阵 $A$ 可相似对角化为 $A = P \Lambda P^{-1}$（其中 $\Lambda = \diag(\lambda_1,\dots,\lambda_n)$，$P = [v_1,\dots,v_n]$ 可逆），当且仅当 $A$ 存在 $n$ 个线性无关的特征向量（等价于在 $\C$ 上每个特征值都满足 $m_g(\lambda) = m_a(\lambda)$）。特别地，若 $A$ 有 $n$ 个互异特征值，则 $A$ 必可对角化。此时对任意正整数 $k$，
\begin{equation}
  A^k = (P\Lambda P^{-1})^k = P \Lambda^k P^{-1}, \qquad A^k x^{(0)} = \sum_{i=1}^n c_i \lambda_i^k v_i, \label{eq:ch3-matpower}
\end{equation}
其中坐标向量 $c = P^{-1} x^{(0)}$。
\end{proposition}

式 \eqref{eq:ch3-matpower} 把耦合的 $n$ 维矩阵连乘彻底拆解成了 $n$ 条互不干扰的标量等比数列 $c_i \lambda_i^k$。定义 $A$ 的\textbf{谱半径}（spectral radius）为 $\rho(A) := \max_i |\lambda_i|$。对于可对角化矩阵，只需一眼看 $\rho(A)$：若 $\rho(A) < 1$，每个分量 $\lambda_i^k \to 0$，系统全局指数衰减；若 $\rho(A) > 1$，沿主特征方向的分量将指数爆炸。

\subsection{不可对角化时发生了什么：Jordan 链与剪切瞬态增长}

然而，并非所有方阵都可对角化——当某个重特征值的几何重数严格小于代数重数（$m_g(\lambda) < m_a(\lambda)$）时，特征向量不够凑满一组基。最简单的例子就是 $2\times 2$ 的\textbf{Jordan 块}：
\[
  J_2(\lambda) = \begin{pmatrix} \lambda & 1 \\ 0 & \lambda \end{pmatrix}.
\]
显然它唯一的特征值是 $\lambda$（$m_a(\lambda)=2$），但解 $(J_2(\lambda)-\lambda I)v = \begin{psmallmatrix}0&1\\0&0\end{psmallmatrix}v = 0$ 仅得一维特征子空间 $\nul(J_2(\lambda)-\lambda I) = \spanop\{e_1\}$（$m_g(\lambda)=1$）。

缺失的第二个基向量去哪儿了？观察 $e_2 = (0,1)^T$ 在 $N := J_2(\lambda) - \lambda I = \begin{psmallmatrix}0&1\\0&0\end{psmallmatrix}$ 下的作用：
\[
  e_2 \xrightarrow{\;J_2(\lambda)-\lambda I\;} e_1 \xrightarrow{\;J_2(\lambda)-\lambda I\;} 0,
\]
向量序列 $(e_2, e_1)$ 构成一条长度为 $2$ 的\textbf{Jordan 链}（Jordan chain），它们依次张成一条严格递增的广义特征子空间序列（核的递增链）：
\[
  \nul\bigl(J_2(\lambda)-\lambda I\bigr) = \spanop\{e_1\} \;\subsetneq\; \nul\bigl((J_2(\lambda)-\lambda I)^2\bigr) = \R^2.
\]
从几何上看，$J_2(\lambda) = \lambda I + N$ 在把空间按 $\lambda$ 缩放的同时，还沿 $e_1$ 方向附加了一个\textbf{剪切}（shear）操作 $N$。由于 $\lambda I$ 与幂零矩阵 $N$（满足 $N^2 = 0$）可交换，直接套用二项式定理即可算出它的精确 $k$ 次幂：
\begin{equation}
  J_2(\lambda)^k = (\lambda I + N)^k = \lambda^k I + k \lambda^{k-1} N = \begin{pmatrix} \lambda^k & k \lambda^{k-1} \\ 0 & \lambda^k \end{pmatrix}. \label{eq:ch3-jordan-power}
\end{equation}
一般地，对任意 $m$ 阶 Jordan 块 $J_m(\lambda) = \lambda I + N$（次对角线全为 $1$，满足 $N^m = 0$），同样有 $J_m(\lambda)^k = \sum_{j=0}^{m-1} \binom{k}{j}\lambda^{k-j} N^j$。

\begin{example}[$|\lambda|<1$ 时的瞬态暴涨实测]
\label{ex:ch3-jordan-transient}
取 $\lambda = 0.9 < 1$，考察 $J_2(0.9) = \begin{psmallmatrix}0.9 & 1 \\ 0 & 0.9\end{psmallmatrix}$。虽然谱半径 $\rho(J_2(0.9)) = 0.9 < 1$ 保证了 $k \to \infty$ 时 $J_2(0.9)^k \to 0$，但右上角的剪切项 $k \cdot 0.9^{k-1}$ 包含一个线性增长因子 $k$！在衰减真正占据上风之前，它会先经历剧烈的\textbf{瞬态增长}（transient growth）：由实测数据（\texttt{data/jordan\_transient.csv}），右上角元素 $k\cdot 0.9^{k-1}$ 在 $k=9$ 与 $k=10$ 处达到并列峰值 $9\times 0.9^8 = 10\times 0.9^9 \approx 3.8742$，矩阵 Frobenius 范数 $\|J_2(0.9)^k\|_F$ 在 $k=9$ 处冲高至 $3.9128$（初始值的 $2.77$ 倍），直到 $k=30$ 之后才回落到初始范数以下（图 \ref{fig:ch3-jordan-transient}）。
\end{example}

\begin{figure}[htbp]
  \centering
  \begin{tikzpicture}
    \begin{axis}[
      width=0.82\linewidth,
      height=5.6cm,
      xlabel={迭代步数 $k$},
      ylabel={范数 / 元素幅值},
      xmin=0, xmax=60,
      ymin=0, ymax=4.7,
      grid=major,
      grid style={dashed, gray!25},
      legend pos=north east,
      legend style={font=\footnotesize, fill=white, draw=codeborder, rounded corners=1pt},
      tick label style={font=\footnotesize},
      label style={font=\small\sffamily}
    ]
      \addplot[color=teal, very thick] table[x=k, y=fro_norm, col sep=comma] {data/jordan_transient.csv};
      \addlegendentry{$\|J_2(0.9)^k\|_F$（总范数）}
      \addplot[color=ochre, thick, dashed] table[x=k, y=upper_right, col sep=comma] {data/jordan_transient.csv};
      \addlegendentry{剪切项 $k\cdot 0.9^{k-1}$}
      \addplot[color=navy, thick, dotted] table[x=k, y=diag, col sep=comma] {data/jordan_transient.csv};
      \addlegendentry{对角衰减项 $0.9^k$}

      % Peak marker
      \addplot[only marks, mark=*, mark size=2.2pt, color=ochre] coordinates {(9, 3.912755) (10, 3.874205)};
      \node[anchor=west, font=\footnotesize\sffamily, color=ochre, align=left, fill=white, fill opacity=0.85, text opacity=1, inner sep=1.5pt] at (axis cs:14, 4.05)
        {瞬态峰值：$\|J_2^9\|_F \approx 3.913$\\右上角 $k=9,10$ 达 $3.874$};
      \draw[-{Stealth[length=2mm]}, color=ochre, thin] (axis cs:13.7, 4.05) -- (axis cs:9.6, 3.92);
    \end{axis}
  \end{tikzpicture}
  \caption{不可对角化矩阵 $J_2(0.9)$ 幂次的瞬态增长曲线（由 \texttt{code/vol1\_experiments.py} 实测生成）。尽管 $|\lambda|=0.9<1$ 保证了最终收敛到 $0$，非正交的剪切结构仍使范数在前期翻倍暴涨；在深度循环网络中，这种非正规（non-normal）瞬态放大正是梯度不稳定的重要根源。}
  \label{fig:ch3-jordan-transient}
\end{figure}

\begin{failBox}{为什么现代数值计算与 AI 几乎从不求一般 Jordan 标准形？}
Jordan 标准形在纯代数上是优美的存在性定理，但在浮点数世界里却是\textbf{极度病态}的。考虑微扰后的 $2\times 2$ Jordan 块：
\[
  J_2(0) = \begin{pmatrix} 0 & 1 \\ 0 & 0 \end{pmatrix}
  \quad\xrightarrow{\;\text{左下角加 } \varepsilon = 10^{-16}\;}\quad
  A_\varepsilon = \begin{pmatrix} 0 & 1 \\ \varepsilon & 0 \end{pmatrix}.
\]
原矩阵不可对角化，特征值为二重根 $0$；仅加上一个机器精度级别的舍入误差 $\varepsilon = 10^{-16}$，特征方程立刻变成 $\lambda^2 - \varepsilon = 0$，特征值跳变为 $\pm\sqrt{\varepsilon} = \pm 10^{-8}$（误差被放大了整整 $10^8$ 倍！），矩阵也瞬间变成了可对角化矩阵。也就是说，\textbf{Jordan 标准形对扰动是不连续的}。因此，现代人工智能和数值线性代数不会去计算脆弱的一般 Jordan 分解，而是坚定转向两把\textbf{数值极其稳定、基于正交变换的武器}：本章的\textbf{实对称阵谱定理}与第 4 章的\textbf{奇异值分解（SVD）}。
\end{failBox}

% =========================================================================
\section{实对称矩阵谱定理的完整严格证明}
\label{sec:ch3-spectral}

实对称矩阵（$S^T = S$）是线性代数中最“完美”的矩阵家族：它们\textbf{永远不会出现复特征值，永远不会出现 Jordan 块（必可对角化），而且特征向量永远可以选得两两正交}。在进入一般 $n$ 维证明之前，我们先亲手解剖一个 $2\times 2$ 实例，建立清晰的几何画面。

\begin{example}[$2\times 2$ 实对称阵手算解剖与外积和]
\label{ex:ch3-sym2x2}
考察实对称矩阵
\[
  S = \begin{pmatrix} 5 & 3 \\ 3 & 5 \end{pmatrix}.
\]
\begin{enumerate}
  \item \textbf{求特征值}：特征多项式 $\det(\lambda I - S) = (\lambda - 5)^2 - 9 = (\lambda - 8)(\lambda - 2) = 0$，得两个实特征值 $\lambda_1 = 8,\; \lambda_2 = 2$。
  \item \textbf{求单位特征向量}：
    \begin{itemize}
      \item 对 $\lambda_1 = 8$，解 $(S - 8I)v = \begin{psmallmatrix}-3 & 3 \\ 3 & -3\end{psmallmatrix}v = 0$，单位化得 $q_1 = \frac{1}{\sqrt{2}}\begin{pmatrix}1\\1\end{pmatrix}$。
      \item 对 $\lambda_2 = 2$，解 $(S - 2I)v = \begin{psmallmatrix}3 & 3 \\ 3 & 3\end{psmallmatrix}v = 0$，单位化得 $q_2 = \frac{1}{\sqrt{2}}\begin{pmatrix}-1\\1\end{pmatrix}$。
    \end{itemize}
  \item \textbf{验证正交性与秩 1 外积和}：立刻验证 $q_1^T q_2 = \frac{1}{2}(-1+1) = 0$！拼成正交矩阵 $Q = [q_1, q_2]$，它满足 $Q^{-1} = Q^T$（无需矩阵求逆）。套用第 1 章矩阵乘法的第 4 视角（列乘行外积和），有
  \[
    S = Q \Lambda Q^T
      = 8\, q_1 q_1^T + 2\, q_2 q_2^T
      = 8 \begin{pmatrix} 1/2 & 1/2 \\ 1/2 & 1/2 \end{pmatrix}
      + 2 \begin{pmatrix} 1/2 & -1/2 \\ -1/2 & 1/2 \end{pmatrix}
      = \begin{pmatrix} 5 & 3 \\ 3 & 5 \end{pmatrix}.
  \]
\end{enumerate}
\textbf{几何直觉}：对任意向量 $x\in\R^2$，计算 $Sx = Q(\Lambda(Q^T x))$ 分三步完成——先正交旋转到主轴坐标系 ($Q^T x$)，再沿互相垂直的 $q_1, q_2$ 轴分别纯粹拉伸 $8$ 倍和 $2$ 倍 ($\Lambda$)，最后旋转回原坐标系 ($Q$)，\textbf{完全没有剪切}！
\end{example}

现在，借助第 2 章已经证明的正交补定理 $\R^n = U \oplus U^\perp$ 与 Gram--Schmidt 正交化，我们用三条极其干净的引理加上归纳法，\textbf{零跳步地}证明一般 $n$ 维实对称矩阵的谱定理。

\begin{lemma}[实对称阵的特征值全为实数，且存在实特征向量]
\label{lem:ch3-real-eig}
设 $S \in \R^{n\times n}$ 满足 $S^T = S$。若存在复标量 $\lambda \in \C$ 与非零复向量 $v \in \C^n$ 使得 $Sv = \lambda v$，则必有 $\lambda \in \R$，且存在非零实向量 $x \in \R^n$ 满足 $Sx = \lambda x$。
\end{lemma}

\begin{proof}
记 $v^* := \bar{v}^T$ 为 $v$ 的共轭转置。在 $Sv = \lambda v$ 两边左乘 $v^*$ 得
\[
  v^* S v = \lambda (v^* v) = \lambda \sum_{j=1}^n |v_j|^2.
\]
由于 $v^* S v$ 是一个 $1\times 1$ 标量，取其复共轭等价于取共轭转置。利用 $S$ 为实对称矩阵（故 $S^* = \bar{S}^T = S^T = S$），有
\[
  \overline{v^* S v} = (v^* S v)^* = v^* S^* (v^*)^* = v^* S v.
\]
一个复数等于自身的共轭，说明 $v^* S v \in \R$。又 $v \neq 0$ 保证了 $\sum_{j=1}^n |v_j|^2 > 0$，故 $\lambda = (v^*Sv)/\sum_j |v_j|^2 \in \R$。一旦 $\lambda \in \R$，矩阵 $S - \lambda I$ 便是奇异的\textbf{实矩阵}，齐次线性方程组 $(S - \lambda I)x = 0$ 必有非零实解 $x \in \R^n$（也可直接取 $v$ 的实部或虚部中的非零者）。
\end{proof}

\begin{lemma}[互异特征值对应的特征向量正交]
\label{lem:ch3-ortho-eig}
设 $S \in \R^{n\times n}$ 为实对称矩阵，$Sv_1 = \lambda_1 v_1$，$Sv_2 = \lambda_2 v_2$，且 $\lambda_1 \neq \lambda_2$，则 $v_1^T v_2 = 0$。
\end{lemma}

\begin{proof}
利用对称性 $S^T = S$，从左右两个方向计算标量 $v_1^T S v_2$：
\[
  \lambda_1 (v_1^T v_2) = (\lambda_1 v_1)^T v_2 = (Sv_1)^T v_2 = v_1^T S^T v_2 = v_1^T S v_2 = v_1^T (\lambda_2 v_2) = \lambda_2 (v_1^T v_2).
\]
移项得 $(\lambda_1 - \lambda_2)(v_1^T v_2) = 0$。因 $\lambda_1 \neq \lambda_2$，故 $v_1^T v_2 = 0$。
\end{proof}

\begin{lemma}[不变子空间的正交补仍是不变子空间]
\label{lem:ch3-perp-inv}
设 $S \in \R^{n\times n}$ 为实对称矩阵，$U \subseteq \R^n$ 是 $S$-不变子空间（即 $\forall u \in U,\; Su \in U$）。则其正交补 $U^\perp$ 也是 $S$-不变子空间（即 $\forall y \in U^\perp,\; Sy \in U^\perp$）。
\end{lemma}

\begin{proof}
任取 $y \in U^\perp$ 与 $u \in U$。由 $S^T = S$ 与 $Su \in U$，有
\[
  \inner{Sy}{u} = (Sy)^T u = y^T S^T u = y^T (Su) = \inner{y}{Su} = 0.
\]
由于这对一切 $u \in U$ 均成立，按正交补定义即得 $Sy \in U^\perp$。
\end{proof}

引理 \ref{lem:ch3-perp-inv} 正是第 2 章正交补理论的绝妙回报：它让我们可以在不假设特征值互异（允许任意重根！）的情况下，一步剥离出一个一维特征子空间，并对余下的 $(n-1)$ 维正交补空间实施归纳。

\begin{thmBox}
\begin{theorem}[实对称矩阵谱定理, Real Symmetric Spectral Theorem]
\label{thm:ch3-spectral}
对任意 $n$ 阶实对称矩阵 $S \in \R^{n\times n}$（$S^T = S$），存在正交矩阵 $Q = [q_1,\dots,q_n] \in \R^{n\times n}$（满足 $Q^T Q = Q Q^T = I_n$）以及实对角矩阵 $\Lambda = \diag(\lambda_1,\dots,\lambda_n)$，使得
\begin{equation}
  S = Q \Lambda Q^T = \sum_{i=1}^n \lambda_i q_i q_i^T. \label{eq:ch3-spectral}
\end{equation}
\end{theorem}
\end{thmBox}

\begin{proof}
对矩阵阶数 $n$ 施数学归纳法。

\textbf{基础步（$n=1$）}：$S = [s_{11}]$，取 $Q = [1]$、$\Lambda = [s_{11}]$，显然成立。

\textbf{归纳步}：假设对任意 $(n-1)$ 阶实对称矩阵定理成立，现考虑 $n$ 阶实对称阵 $S$。
\begin{enumerate}
  \item 由代数基本定理，特征方程 $\det(\lambda I - S) = 0$ 在 $\C$ 上至少有一个根 $\lambda_1$；再由引理 \ref{lem:ch3-real-eig}，$\lambda_1 \in \R$ 且存在单位实特征向量 $q_1 \in \R^n$（$\|q_1\|_2 = 1$）满足 $Sq_1 = \lambda_1 q_1$。
  \item 记一维子空间 $U := \spanop\{q_1\}$，显然 $SU \subseteq U$。由第 2 章的正交补分解定理，$\R^n = U \oplus U^\perp$ 且 $\dim(U^\perp) = n-1$。利用 Gram--Schmidt 正交化在 $U^\perp$ 中取一组标准正交基，按列拼成矩阵 $Q_2 \in \R^{n\times(n-1)}$（满足 $Q_2^T Q_2 = I_{n-1}$ 且 $q_1^T Q_2 = 0^T$）。于是 $W := [q_1,\; Q_2] \in \R^{n\times n}$ 是正交矩阵。
  \item 计算正交相似变换 $W^T S W$：
  \[
    W^T S W
    = \begin{pmatrix} q_1^T \\ Q_2^T \end{pmatrix} S \begin{pmatrix} q_1 & Q_2 \end{pmatrix}
    = \begin{pmatrix} q_1^T S q_1 & q_1^T S Q_2 \\ Q_2^T S q_1 & Q_2^T S Q_2 \end{pmatrix}.
  \]
  其中 $q_1^T S q_1 = \lambda_1 q_1^T q_1 = \lambda_1$，$Q_2^T S q_1 = \lambda_1 Q_2^T q_1 = 0$；而由引理 \ref{lem:ch3-perp-inv}（或直接取转置 $q_1^T S Q_2 = (Q_2^T S q_1)^T = 0^T$），右上块同样为零向量！记 $S_2 := Q_2^T S Q_2 \in \R^{(n-1)\times(n-1)}$，则
  \[
    W^T S W = \begin{pmatrix} \lambda_1 & 0^T \\ 0 & S_2 \end{pmatrix}.
  \]
  \item 验证对称性：$S_2^T = (Q_2^T S Q_2)^T = Q_2^T S^T Q_2 = Q_2^T S Q_2 = S_2$，故 $S_2$ 是 $(n-1)$ 阶实对称矩阵。由归纳假设，存在 $(n-1)$ 阶正交矩阵 $\widetilde{Q}_2$ 与实对角阵 $\Lambda_2 = \diag(\lambda_2,\dots,\lambda_n)$ 使得 $\widetilde{Q}_2^T S_2 \widetilde{Q}_2 = \Lambda_2$。
  \item 令分块正交阵 $Q := W \begin{psmallmatrix} 1 & 0^T \\ 0 & \widetilde{Q}_2 \end{psmallmatrix} = [q_1,\; Q_2 \widetilde{Q}_2]$。显然 $Q^T Q = I_n$，且
  \[
    Q^T S Q
    = \begin{pmatrix} 1 & 0^T \\ 0 & \widetilde{Q}_2^T \end{pmatrix}
      \begin{pmatrix} \lambda_1 & 0^T \\ 0 & S_2 \end{pmatrix}
      \begin{pmatrix} 1 & 0^T \\ 0 & \widetilde{Q}_2 \end{pmatrix}
    = \begin{pmatrix} \lambda_1 & 0^T \\ 0 & \Lambda_2 \end{pmatrix} = \Lambda.
  \]
  左乘 $Q$、右乘 $Q^T$ 即得 $S = Q\Lambda Q^T = \sum_{i=1}^n \lambda_i q_i q_i^T$。
\end{enumerate}
\end{proof}

有了谱定理，下面这条在 PCA、优化与图算法中频繁亮相的 \textbf{Rayleigh 商界}只需 4 行便水到渠成。

\begin{thmBox}
\begin{corollary}[Rayleigh 商的上下界与极值点]
\label{cor:ch3-rayleigh}
设实对称矩阵 $S \in \R^{n\times n}$ 的特征值按降序排列为 $\lambda_{\max} = \lambda_1 \ge \lambda_2 \ge \dots \ge \lambda_n = \lambda_{\min}$，对应单位正交特征向量为 $q_1, \dots, q_n$。则对任意非零向量 $x \in \R^n$，\textbf{Rayleigh 商} $R_S(x) := \dfrac{x^T S x}{x^T x}$ 恒满足
\begin{equation}
  \lambda_{\min} \;\le\; \frac{x^T S x}{x^T x} \;\le\; \lambda_{\max}, \label{eq:ch3-rayleigh}
\end{equation}
且最大值在 $x = q_1$ 处取到，最小值在 $x = q_n$ 处取到。
\end{corollary}
\end{thmBox}

\begin{proof}
由谱定理 $S = Q\Lambda Q^T$，令正交坐标变换 $y := Q^T x \in \R^n$。因 $Q$ 正交，有 $y^T y = x^T Q Q^T x = x^T x > 0$ 且 $x^T S x = y^T \Lambda y = \sum_{i=1}^n \lambda_i y_i^2$。因此
\[
  \lambda_{\min} = \frac{\sum_{i=1}^n \lambda_{\min} y_i^2}{\sum_{i=1}^n y_i^2}
  \;\le\; \frac{\sum_{i=1}^n \lambda_i y_i^2}{\sum_{i=1}^n y_i^2}
  = \frac{x^T S x}{x^T x}
  \;\le\; \frac{\sum_{i=1}^n \lambda_{\max} y_i^2}{\sum_{i=1}^n y_i^2} = \lambda_{\max}.
\]
代入 $x = q_1$ 得 $q_1^T S q_1 / q_1^T q_1 = \lambda_1 = \lambda_{\max}$，代入 $x = q_n$ 得 $\lambda_{\min}$。
\end{proof}

\begin{insightBox}{谱定理在全书下游的四大核心落点}
\begin{itemize}
  \item \textbf{本卷第 4 章（SVD 与 PCA）}：对一般长方形矩阵 $A$，对称阵 $A^T A$ 的谱分解直接构造出 SVD；而样本协方差阵 $\Sigma_{\mathrm{cov}}$ 的 Rayleigh 商最大值恰好给出第一主成分方向！
  \item \textbf{本卷第 6 章（梯度下降收敛率）}：二次目标 $f(w) = \frac{1}{2}w^T Q w$ 的梯度下降误差满足 $e^{(t)} = (I - \eta Q)^t e^{(0)}$，正交变换 $Q = U\Lambda U^T$ 将其解耦为 $n$ 条独立标量迭代，直接导出步长红线 $0 < \eta < 2/\lambda_{\max}$。
  \item \textbf{第 2 卷（多元高斯分布）}：协方差阵 $\Sigma = Q\Lambda Q^T$ 决定了高斯概率等高椭球的主轴方向与半轴长。
  \item \textbf{第 4 卷（凸优化）}：目标函数 Hessian 的最大特征值 $\lambda_{\max}$ 对应 $L$-光滑常数，最小特征值 $\lambda_{\min}$ 对应 $\mu$-强凸常数。
\end{itemize}
\end{insightBox}

% =========================================================================
\section{二次型、半正定矩阵（PSD）的三大 AI 化身与条件数 $\kappa$}
\label{sec:ch3-psd-cond}

\subsection{二次型对称化与半正定判据（零循环推论）}

在机器学习的损失函数与泰勒展开中，我们随处可见形如 $q(x) = x^T A x$ 的\textbf{二次型}（quadratic form）。在讨论它的正负性之前，必须牢记第一条铁律：\textbf{必须先把矩阵对称化}。

由于 $x^T A x$ 是一个标量，它等于自身的转置 $(x^T A x)^T = x^T A^T x$，从而对任意方阵 $A \in \R^{n\times n}$ 恒有
\[
  x^T A x = x^T \!\left(\frac{A + A^T}{2}\right) x.
\]
记 $S := \frac{1}{2}(A + A^T)$，它是唯一的对称阵，决定着二次型在所有方向上的取值。直接看非对称阵 $A$ 的特征值来判断 $x^T A x$ 的正负，会导致致命误判！

\begin{failBox}{非对称矩阵的“全正特征值”陷阱}
考虑非对称上三角矩阵 $A = \begin{pmatrix} 1 & 4 \\ 0 & 1 \end{pmatrix}$。它的特征值显然是 $\lambda_1 = \lambda_2 = 1 > 0$，全部严格为正。但取向量 $x = (1, -1)^T$ 代入二次型：
\[
  x^T A x = \begin{pmatrix} 1 & -1 \end{pmatrix} \begin{pmatrix} 1 & 4 \\ 0 & 1 \end{pmatrix} \begin{pmatrix} 1 \\ -1 \end{pmatrix} = 1 - 4 + 1 = -2 < 0!
\]
原因一目了然：对称化后的真实核心矩阵是 $S = \frac{1}{2}(A + A^T) = \begin{pmatrix} 1 & 2 \\ 2 & 1 \end{pmatrix}$，它的特征值是 $3$ 和 $-1$（存在负特征值，实际是马鞍面而非碗状）！因此，\textbf{定义正定性与半正定性时，前提永远是实对称矩阵}。
\end{failBox}

\begin{defBox}
\begin{definition}[正定、半正定与不定矩阵]
\label{def:ch3-psd}
设 $Q \in \R^{n\times n}$ 为\textbf{实对称矩阵}（$Q^T = Q$）：
\begin{itemize}
  \item 若对任意非零向量 $x \in \R^n$ 均有 $x^T Q x > 0$，则称 $Q$ 为\textbf{正定矩阵}（Positive Definite, PD），记作 $Q \succ 0$；
  \item 若对任意向量 $x \in \R^n$ 均有 $x^T Q x \ge 0$，则称 $Q$ 为\textbf{半正定矩阵}（Positive Semi-Definite, PSD），记作 $Q \succeq 0$；
  \item 若存在 $x, y \in \R^n$ 使得 $x^T Q x > 0$ 且 $y^T Q y < 0$，则称 $Q$ 为\textbf{不定矩阵}（Indefinite）。
\end{itemize}
\end{definition}
\end{defBox}

现在，因为我们在第 \ref{sec:ch3-spectral} 节已经严格证完谱定理与 Rayleigh 商界，正定与半正定的特征值充要条件成为一个\textbf{3 行即毕、绝无逻辑循环}的推论！

\begin{thmBox}
\begin{theorem}[正定与半正定矩阵的特征值充要判据]
\label{thm:ch3-psd-eig}
设 $Q \in \R^{n\times n}$ 为实对称矩阵，其特征值为 $\lambda_1, \dots, \lambda_n \in \R$。则：
\begin{enumerate}
  \item $Q \succ 0 \iff \forall i \in \{1,\dots,n\},\; \lambda_i > 0$（等价于 $\lambda_{\min}(Q) > 0$）；
  \item $Q \succeq 0 \iff \forall i \in \{1,\dots,n\},\; \lambda_i \ge 0$（等价于 $\lambda_{\min}(Q) \ge 0$）。
\end{enumerate}
特别地，正定矩阵必可逆（因 $\det(Q) = \prod_i \lambda_i > 0$），且其逆矩阵 $Q^{-1} = U\Lambda^{-1}U^T$ 仍为正定矩阵。
\end{theorem}
\end{thmBox}

\begin{proof}
由推论 \ref{cor:ch3-rayleigh}，二次型在单位球面 $\|x\|_2 = 1$ 上的最小值恰好等于 $\lambda_{\min}(Q)$。因此：
\begin{itemize}
  \item 若 $\lambda_{\min}(Q) > 0$，则对一切 $x \neq 0$ 有 $x^T Q x \ge \lambda_{\min}(Q)\|x\|_2^2 > 0$，故 $Q \succ 0$；反之若 $Q \succ 0$，取最小特征值对应的单位特征向量 $u_n$ 即得 $\lambda_{\min}(Q) = u_n^T Q u_n > 0$。
  \item 将严格不等号 $>$ 替换为 $\ge$，完全同理得 $Q \succeq 0 \iff \lambda_{\min}(Q) \ge 0$。
\end{itemize}
\end{proof}

从几何上看，二维二次型 $f(x_1, x_2) = x^T Q x$ 的等高线清晰地揭示了三种曲面的本质差异（图 \ref{fig:ch3-quad-contours}）：正定阵对应封闭同心椭圆构成的\textbf{严格凸碗}（有唯一全局极小点）；半正定退化阵对应平行直线构成的\textbf{平底山谷}（沿零特征方向存在无穷多个极小点）；不定阵对应双曲线构成的\textbf{马鞍面}（沿一轴上升、沿另一轴下降，原点既非极大也非极小）。

\begin{figure}[htbp]
  \centering
  \resizebox{0.97\linewidth}{!}{%
  \begin{tikzpicture}[>=Stealth]
    % Panel (a): Positive Definite Bowl x1^2 + 4 x2^2
    \begin{scope}[shift={(0,0)}]
      \draw[->, slate, thin] (-2.2,0) -- (2.2,0) node[right, font=\footnotesize] {$x_1$};
      \draw[->, slate, thin] (0,-1.8) -- (0,1.8) node[above, font=\footnotesize] {$x_2$};
      \foreach \r/\op in {0.55/0.35, 1.05/0.55, 1.55/0.75, 1.95/0.95} {
        \draw[teal, thick, opacity=\op] (0,0) ellipse ({\r} and {\r/2});
      }
      \fill[ochre] (0,0) circle (2.2pt);
      \node[below right, font=\scriptsize\sffamily, color=ochre, fill=white, fill opacity=0.88, text opacity=1, inner sep=1pt] at (0.08,-0.06) {唯一极小点};
      \node[below, font=\small\sffamily\bfseries, color=tealDark] at (0,-2.05) {(a) 正定碗：$x_1^2 + 4x_2^2$};
      \node[below, font=\scriptsize, color=slate] at (0,-2.5) {$\lambda_1=1>0,\;\lambda_2=4>0$};
    \end{scope}

    % Panel (b): Positive Semi-Definite Valley x1^2
    \begin{scope}[shift={(5.8,0)}]
      \draw[->, slate, thin] (-2.2,0) -- (2.2,0) node[right, font=\footnotesize] {$x_1$};
      \draw[->, slate, thin] (0,-1.8) -- (0,1.8) node[above, font=\footnotesize] {$x_2$};
      \foreach \x/\op in {0.45/0.4, 0.9/0.6, 1.35/0.8, 1.8/0.95} {
        \draw[navy, thick, opacity=\op] (\x,-1.5) -- (\x,1.5);
        \draw[navy, thick, opacity=\op] (-\x,-1.5) -- (-\x,1.5);
      }
      \draw[ochre, very thick] (0,-1.55) -- (0,1.55);
      \node[right, font=\scriptsize\sffamily, color=ochre, fill=white, inner sep=1pt] at (0.12, 1.15) {平底谷线 $f=0$};
      \node[below, font=\small\sffamily\bfseries, color=navy] at (0,-2.05) {(b) 半正定山谷：$x_1^2$};
      \node[below, font=\scriptsize, color=slate] at (0,-2.5) {$\lambda_1=1>0,\;\lambda_2=0$};
    \end{scope}

    % Panel (c): Indefinite Saddle x1^2 - x2^2
    \begin{scope}[shift={(11.6,0)}]
      \draw[->, slate, thin] (-2.2,0) -- (2.2,0) node[right, font=\footnotesize] {$x_1$};
      \draw[->, slate, thin] (0,-1.8) -- (0,1.8) node[above, font=\footnotesize] {$x_2$};
      % Asymptotes
      \draw[slate, dashed, thin] (-1.5,-1.5) -- (1.5,1.5);
      \draw[slate, dashed, thin] (-1.5,1.5) -- (1.5,-1.5);
      % Hyperbolas x1^2 - x2^2 = c > 0 (teal, ascent)
      \foreach \c in {0.4, 1.0} {
        \draw[teal, thick, domain=-1.15:1.15, samples=40] plot ({sqrt(\c + \x*\x)}, {\x});
        \draw[teal, thick, domain=-1.15:1.15, samples=40] plot ({-sqrt(\c + \x*\x)}, {\x});
      }
      % Hyperbolas x2^2 - x1^2 = c > 0 (ochre, descent)
      \foreach \c in {0.4, 1.0} {
        \draw[ochre, thick, domain=-1.15:1.15, samples=40] plot ({\x}, {sqrt(\c + \x*\x)});
        \draw[ochre, thick, domain=-1.15:1.15, samples=40] plot ({\x}, {-sqrt(\c + \x*\x)});
      }
      \fill[ink] (0,0) circle (2.2pt);
      \node[font=\scriptsize\sffamily, color=teal] at (1.72, 0.32) {$f>0$};
      \node[font=\scriptsize\sffamily, color=ochre] at (0.48, 1.52) {$f<0$};
      \node[below, font=\small\sffamily\bfseries, color=ochre] at (0,-2.05) {(c) 不定马鞍：$x_1^2 - x_2^2$};
      \node[below, font=\scriptsize, color=slate] at (0,-2.5) {$\lambda_1=1>0,\;\lambda_2=-1<0$};
    \end{scope}
  \end{tikzpicture}}
  \caption{二维二次型等高线的三种典型几何形态。正定性保证唯一全局极小点，半正定性允许沿零特征子空间的平底谷，而不定性产生高维优化中常见的马鞍面。}
  \label{fig:ch3-quad-contours}
\end{figure}

\subsection{半正定矩阵（PSD）在人工智能中的三大化身}

为什么人工智能论文和教科书里动不动就出现“$M \succeq 0$”？因为 AI 中最核心的三类矩阵天然都具有 $B^T B$ 的平方结构，而\textbf{任意实矩阵的格拉姆型乘积 $B^T B$ 必半正定}：对一切 $v$，$v^T (B^T B) v = \|Bv\|_2^2 \ge 0$！表 \ref{tab:ch3-psd-three} 汇总了这三大化身：

\begin{table}[htbp]
  \centering
  \caption{半正定矩阵（PSD）在人工智能中的三大核心化身与下游坐标}
  \label{tab:ch3-psd-three}
  \small
  \begin{tabularx}{\linewidth}{@{}p{26mm}p{35mm}p{35mm}X@{}}
    \toprule
    \textbf{AI 化身} & \textbf{构造表达式} & \textbf{半正定性验证} & \textbf{几何/统计意义与下游卷次} \\
    \midrule
    \textbf{1. 样本协方差阵} & $\Sigma = \dfrac{1}{n}\widetilde{X}^T\widetilde{X} \in \R^{d\times d}$ & $v^T \Sigma v = \dfrac{1}{n}\|\widetilde{X}v\|_2^2 \ge 0$ & 投影到方向 $v$ 上的样本方差非负；高维小样本（$n<d$）时必退化（$\Sigma \succeq 0$ 但非正定）。\textbf{（第 4 章 PCA、第 2 卷统计）} \\[3pt]
    \textbf{2. 核 / Gram 矩阵} & $K_{ij} = \inner{\phi(x_i)}{\phi(x_j)}$，即 $K = \Phi^T \Phi \in \R^{n\times n}$ & $c^T K c = \norm{\sum_{i=1}^n c_i \phi(x_i)}_2^2 \ge 0$ & 刻画样本在高维特征空间的相似度几何；合法核函数必须对任意样本集生成 PSD 矩阵（Mercer 定理）。\textbf{（第 4 卷 SVM 与核方法）} \\[3pt]
    \textbf{3. Hessian 矩阵} & $H(w) = \nabla^2 f(w) \in \R^{d\times d}$ & $\Delta w^T H \Delta w \ge 0$（曲率非负） & 局部二次近似向上弯曲；$H(w) \succeq 0$ 处处成立 $\iff$ 目标函数为凸函数，无伪局部极小值。\textbf{（第 5--6 章、第 4 卷凸优化）} \\
    \bottomrule
  \end{tabularx}
\end{table}

顺便一提，定理 \ref{thm:ch3-psd-eig} 还彻底解释了我们在第 2 章遗留的一个问题：\textbf{为什么岭回归只要在 $X^T X$ 对角线上加上 $\lambda I$（$\lambda > 0$），原本可能奇异的矩阵就一定可逆？}因为 $X^T X \succeq 0$ 的全部特征值 $\mu_i \ge 0$，加上 $\lambda I$ 后，对每个特征向量 $u_i$ 都有 $(X^T X + \lambda I)u_i = (\mu_i + \lambda)u_i$，于是新矩阵的全部特征值都平移到 $\mu_i + \lambda \ge \lambda > 0$，瞬间变成\textbf{严格正定矩阵}，必然可逆！

\subsection{条件数 $\kappa(Q)$：椭球扁平度与病态峡谷}

对于正定矩阵 $Q = U\Lambda U^T \succ 0$（特征值 $\lambda_{\max} = \lambda_1 \ge \dots \ge \lambda_n = \lambda_{\min} > 0$），在特征坐标 $y = U^T x$ 下，等高面 $x^T Q x = c$（$c > 0$）化为标准椭球方程：
\[
  \sum_{i=1}^n \lambda_i y_i^2 = c \;\Longleftrightarrow\; \sum_{i=1}^n \frac{y_i^2}{\bigl(\sqrt{c/\lambda_i}\bigr)^2} = 1.
\]
沿着特征向量 $u_i$ 方向的半轴长正是 $\sqrt{c/\lambda_i}$！因此，椭球\textbf{最长半轴与最短半轴之比}为
\[
  \frac{\text{最长半轴}}{\text{最短半轴}} = \frac{\sqrt{c/\lambda_{\min}}}{\sqrt{c/\lambda_{\max}}} = \sqrt{\frac{\lambda_{\max}}{\lambda_{\min}}} = \sqrt{\cond(Q)},
\]
其中核心指标
\begin{equation}
  \cond(Q) := \frac{\lambda_{\max}(Q)}{\lambda_{\min}(Q)} \ge 1 \label{eq:ch3-cond}
\end{equation}
称为正定矩阵 $Q$ 的\textbf{（谱）条件数}（condition number）。
\begin{itemize}
  \item 当 $\cond(Q) = 1$ 时（$Q = \lambda I$），等高面是完美正球面，在任何一点处负梯度方向都直指碗底最优解，梯度下降一步即可到位；
  \item 当 $\cond(Q) \gg 1$ 时（例如 $\cond(Q) = 10^6$，长短轴比高达 $1000:1$），等高面被压成极其狭长的\textbf{病态峡谷}：负梯度几乎垂直于谷底走向，我们在第 6 章将看到，梯度下降会被迫在峡谷两侧剧烈锯齿震荡、步履维艰；同时求解线性方程组 $Qx = b$ 时，右端噪声 $\Delta b$ 的相对误差最高会被放大 $\cond(Q)$ 倍（$\|\Delta x\|_2/\|x\|_2 \le \cond(Q)\|\Delta b\|_2/\|b\|_2$）。
\end{itemize}

\begin{extBox}{Sylvester 顺序主子式判据与 Cholesky 分解}
在手算低阶矩阵或编写高性能数值库时，无需算出全部特征值也能判定正定性：
\begin{enumerate}
  \item \textbf{Sylvester 判据}：实对称阵 $Q \succ 0$ 当且仅当其全部 $n$ 个\textbf{顺序主子式}（左上角 $k\times k$ 子块行列式 $\Delta_k$，$k=1,\dots,n$）都严格大于 $0$。例如对 $2\times 2$ 实对称阵 $\begin{psmallmatrix}a&b\\b&c\end{psmallmatrix}$，只需检查 $a > 0$ 且 $ac - b^2 > 0$。但务必小心：\textbf{判定半正定（$Q \succeq 0$）时，“顺序主子式 $\ge 0$”并不充分！}例如 $Q = \diag(0, -1)$ 满足 $\Delta_1 = 0, \Delta_2 = 0$，却显然不是半正定阵（须检验\textbf{所有}主子式均 $\ge 0$）。
  \item \textbf{Cholesky 分解}：实对称阵 $Q \succ 0$ 当且仅当存在唯一的对角元全正的下三角矩阵 $L$，使得 $Q = L L^T$。它的计算量仅为 $O(n^3/3)$（LU 分解的一半），因此 \texttt{np.linalg.cholesky(Q)} 是工程中验证严格正定性、求解正定线性方程组以及在第 2 卷采样多元高斯向量（$x = \mu + L z,\; z\sim\mathcal{N}(0,I)$）的首选工具。
\end{enumerate}
\end{extBox}

% =========================================================================
\section{轻量专题：AI 中高频出现的两类特殊矩阵}
\label{sec:ch3-special-matrices}

在步入第 4 章的 SVD 之前，我们用一节轻快而扎实的应用专题，把本章的谱理论立刻对接到人工智能中两类出镜率极高的结构化矩阵：\textbf{行随机矩阵}与\textbf{图拉普拉斯矩阵}。

\subsection{行随机矩阵：马尔可夫链与 Transformer 注意力权重}

\begin{defBox}
\begin{definition}[行随机矩阵, Row-Stochastic Matrix]
\label{def:ch3-stochastic}
若方阵 $P \in \R^{n\times n}$ 满足元素全非负（$P_{ij} \ge 0$）且每一行元素之和均为 $1$（即 $\sum_{j=1}^n P_{ij} = 1$，等价于矩阵方程 $P\mathbf{1} = \mathbf{1}$，其中 $\mathbf{1} = (1,\dots,1)^T$），则称 $P$ 为\textbf{行随机矩阵}（或马尔可夫转移矩阵）。
\end{definition}
\end{defBox}

行随机矩阵的每行都是一组概率分布。它的谱结构有一个极其简洁优美的普遍规律：

\begin{thmBox}
\begin{proposition}[行随机矩阵的谱半径恒为 $1$]
\label{prop:ch3-stoch-spec}
对任意 $n\times n$ 行随机矩阵 $P$：
\begin{enumerate}
  \item $\lambda = 1$ 必是 $P$ 的特征值（对应右特征向量 $\mathbf{1}$）；
  \item $P$ 的任意复特征值 $\lambda \in \C$ 均满足 $|\lambda| \le 1$，故谱半径 $\rho(P) = 1$。
\end{enumerate}
\end{proposition}
\end{thmBox}

\begin{proof}
第 1 条直接由定义 $P\mathbf{1} = 1\cdot\mathbf{1}$ 得出。对第 2 条，设 $Pv = \lambda v$（$v \in \C^n \setminus \{0\}$）。在 $v$ 的 $n$ 个分量中选取绝对值最大的一个下标 $m$，使得 $|v_m| = \|v\|_\infty > 0$。考察方程 $Pv = \lambda v$ 的第 $m$ 行并取绝对值，利用三角不等式与 $P_{mj} \ge 0$：
\[
  |\lambda|\,|v_m| = \abs{(Pv)_m} = \abs{\sum_{j=1}^n P_{mj} v_j}
  \;\le\; \sum_{j=1}^n P_{mj} |v_j|
  \;\le\; \left(\sum_{j=1}^n P_{mj}\right)|v_m| = 1\cdot |v_m|.
\]
两边消去正数 $|v_m|$，立刻得到 $|\lambda| \le 1$。
\end{proof}

在有限状态马尔可夫链中，若第 $k$ 步的状态概率分布记为\textbf{行向量} $\pi^{(k)} \in \R^{1\times n}$，则下一步的分布按 $\pi^{(k+1)} = \pi^{(k)} P$ 演化，即 $\pi^{(k)} = \pi^{(0)} P^k$。若存在概率分布 $\pi^*$ 满足 $\pi^* P = \pi^*$（即 $\pi^{*T}$ 是 $P^T$ 对应特征值 $\lambda = 1$ 的特征向量），则称 $\pi^*$ 为\textbf{平稳分布}（stationary distribution）。当除 $\lambda_1 = 1$ 外的其余特征值满足 $|\lambda_2| \ge \dots \ge |\lambda_n|$ 且 $|\lambda_2| < 1$ 时，任意初始分布都将按 $|\lambda_2|^k$ 的指数速率收敛到唯一的平稳分布 $\pi^*$！

\begin{example}[$3\times 3$ 马尔可夫链收敛实测]
\label{ex:ch3-markov-exp}
考察 \texttt{code/vol1\_experiments.py} 中的 $3\times 3$ 行随机矩阵：
\[
  P = \begin{pmatrix} 0.7 & 0.2 & 0.1 \\ 0.1 & 0.6 & 0.3 \\ 0.2 & 0.3 & 0.5 \end{pmatrix},
  \qquad \pi^{(0)} = (1,\; 0,\; 0).
\]
解线性方程组 $(P^T - I)\pi^{*T} = 0$ 并令各分量之和为 $1$，可得精确平稳分布
\[
  \pi^* = \left(\frac{11}{34},\; \frac{13}{34},\; \frac{10}{34}\right) \approx (0.3235,\; 0.3824,\; 0.2941).
\]
该矩阵的三个特征值为 $\lambda_1 = 1$、$\lambda_2 = 0.4 + \frac{\sqrt{2}}{10} \approx 0.5414$ 与 $\lambda_3 = 0.4 - \frac{\sqrt{2}}{10} \approx 0.2586$。表 \ref{tab:ch3-markov} 展示了实测迭代收敛过程：经过 20 步后 $L_1$ 误差已降至 $6.68 \times 10^{-6}$，与次大特征值衰减律 $|\lambda_2|^{20} \approx 4.6 \times 10^{-6}$ 高度吻合！
\end{example}

\begin{table}[htbp]
  \centering
  \caption{$3\times 3$ 行随机矩阵从 $\pi^{(0)}=(1,0,0)$ 出发收敛到平稳分布 $\pi^*$ 的实测记录}
  \label{tab:ch3-markov}
  \small
  \begin{tabular}{@{}ccc@{}}
    \toprule
    \textbf{迭代步数 $k$} & \textbf{当前分布 $\pi^{(k)} = \pi^{(0)} P^k$} & \textbf{$L_1$ 误差 $\|\pi^{(k)} - \pi^*\|_1$} \\
    \midrule
    $0$  & $(1.0000,\; 0.0000,\; 0.0000)$ & $1.35\times 10^{0}$ \\
    $1$  & $(0.7000,\; 0.2000,\; 0.1000)$ & $7.53\times 10^{-1}$ \\
    $5$  & $(0.3566,\; 0.3685,\; 0.2748)$ & $6.62\times 10^{-2}$ \\
    $10$ & $(0.3251,\; 0.3817,\; 0.2932)$ & $3.08\times 10^{-3}$ \\
    $20$ & $(0.3235,\; 0.3824,\; 0.2941)$ & $6.68\times 10^{-6}$ \\
    \midrule
    平稳分布 $\pi^*$ & $(0.3235,\; 0.3824,\; 0.2941)$ & $0$ \\
    \bottomrule
  \end{tabular}
\end{table}

\begin{insightBox}{Transformer 里的每一个 Softmax 注意力矩阵都是行随机矩阵！}
除了第 3 卷的隐马尔可夫模型（HMM）与马尔可夫决策过程（MDP），行随机矩阵在当代大模型里更是无处不在：在第 5 卷将要学习的 Transformer 自注意力机制中，注意力权重矩阵
\[
  A = \operatorname{softmax}\!\left(\frac{Q K^T}{\sqrt{d_k}}\right) \in \R^{T\times T}
\]
经过逐行 Softmax 归一化后，元素全正且每行之和严格等于 $1$——\textbf{每一个注意力矩阵都是一个行随机矩阵}！由命题 \ref{prop:ch3-stoch-spec}，它的谱半径恒为 $1$，意味着纯注意力加权 $AV$ 只是对 Value 向量的凸组合（不会把词元范数放大爆炸）；反过来，如果把几十层纯注意力矩阵连乘而不加残差连接，所有词元表示就会像马尔可夫链收敛到平稳分布一样趋向同一条秩 1 向量——这正是深度学习中著名的\textbf{深层注意力秩坍塌（rank collapse / oversmoothing）}现象的线性代数根源！
\end{insightBox}

\subsection{图拉普拉斯矩阵：图信号平滑度与谱聚类直觉}

许多现实数据并非规整表格，而是由节点和边构成的\textbf{图}（社交网络、分子结构、知识图谱、像素邻域）。给定 $n$ 节点无向带权图 $G = (V, E, W)$，其\textbf{邻接矩阵} $W \in \R^{n\times n}$ 满足对称非负（$w_{ij} = w_{ji} \ge 0,\; w_{ii} = 0$），\textbf{度矩阵}为对角阵 $D = \diag(d_1,\dots,d_n)$（其中 $d_i = \sum_{j=1}^n w_{ij}$）。定义\textbf{（非归一化）图拉普拉斯矩阵}（Graph Laplacian）为
\[
  L := D - W.
\]

\begin{thmBox}
\begin{proposition}[图拉普拉斯矩阵的半正定性与平滑度二次型]
\label{prop:ch3-laplacian}
对任意无向非负加权图的拉普拉斯矩阵 $L = D - W$：
\begin{enumerate}
  \item $L$ 是实对称矩阵，且 $L\mathbf{1} = 0$（即 $\lambda_1 = 0$ 恒为其特征值）；
  \item 对任意定义在节点上的信号向量 $x = (x_1,\dots,x_n)^T \in \R^n$，恒有
  \begin{equation}
    x^T L x = \frac{1}{2}\sum_{i=1}^n \sum_{j=1}^n w_{ij}(x_i - x_j)^2 = \sum_{\{i,j\}\in E} w_{ij}(x_i - x_j)^2 \;\ge\; 0, \label{eq:ch3-lap-quad}
  \end{equation}
  因此 $L \succeq 0$ 是半正定矩阵；
  \item 特征值 $0$ 的重数恰好等于图 $G$ 的\textbf{连通分量个数}。
\end{enumerate}
\end{proposition}
\end{thmBox}

\begin{proof}
对称性显然；第 $i$ 行元素之和为 $(L\mathbf{1})_i = d_i - \sum_j w_{ij} = 0$，故 $L\mathbf{1}=0$。展开式 \eqref{eq:ch3-lap-quad} 右侧并利用对称性 $w_{ij} = w_{ji}$：
\begin{align*}
  \frac{1}{2}\sum_{i=1}^n \sum_{j=1}^n w_{ij}(x_i - x_j)^2
  &= \frac{1}{2}\sum_{i,j} w_{ij} x_i^2 - \sum_{i,j} w_{ij} x_i x_j + \frac{1}{2}\sum_{i,j} w_{ij} x_j^2 \\
  &= \sum_{i=1}^n d_i x_i^2 - \sum_{i,j} w_{ij} x_i x_j
   = x^T D x - x^T W x = x^T L x.
\end{align*}
因 $w_{ij} \ge 0$，故 $x^T L x \ge 0$，即 $L \succeq 0$。进一步，由于 $L \succeq 0$，由谱定理 $Lx = 0 \iff x^T L x = 0 \iff$ 对每条正权边 $\{i,j\}\in E$ 均有 $x_i = x_j$，亦即 $x$ 在每个连通分量内部必须取常数，从而 $\dim\nul(L)$ 等于连通分量个数。
\end{proof}

式 \eqref{eq:ch3-lap-quad} 的物理直觉无比清晰：$x^T L x$ 度量了节点信号 $x$ 在图上的\textbf{总变差（粗糙度）}——相连且边权大的节点若取值接近，则 $x^T L x$ 很小（信号平滑）；若相邻节点正负剧烈震荡，则 $x^T L x$ 很大！

\begin{figure}[htbp]
  \centering
  \resizebox{0.96\linewidth}{!}{%
  \begin{tikzpicture}[
    v/.style={circle, draw=tealDark, thick, fill=tealLight, minimum size=8.5mm, font=\sffamily\bfseries\small, text=ink},
    strong/.style={line width=1.6pt, teal},
    weak/.style={line width=1.1pt, ochre, dashed}
  ]
    % Left: 4-node graph
    \node[v] (n1) at (0, 0) {1};
    \node[v] (n2) at (2.1, 0) {2};
    \node[v] (n3) at (5.1, 0) {3};
    \node[v] (n4) at (7.2, 0) {4};

    \draw[strong] (n1) -- node[above, font=\footnotesize\sffamily, color=tealDark] {$w=1.0$} (n2);
    \draw[weak]   (n2) -- node[above, font=\footnotesize\sffamily, color=ochre] {$w=0.2$} node[below, font=\scriptsize\sffamily, color=ochre] {（弱桥）} (n3);
    \draw[strong] (n3) -- node[above, font=\footnotesize\sffamily, color=tealDark] {$w=1.0$} (n4);

    % Cluster brackets below
    \draw[decorate, decoration={brace, amplitude=5pt, mirror}, teal] (-0.4,-0.75) -- (2.5,-0.75)
      node[midway, below=5pt, font=\footnotesize\sffamily, color=tealDark] {紧耦合簇 A（$q_2 > 0$）};
    \draw[decorate, decoration={brace, amplitude=5pt, mirror}, navy] (4.7,-0.75) -- (7.6,-0.75)
      node[midway, below=5pt, font=\footnotesize\sffamily, color=navy] {紧耦合簇 B（$q_2 < 0$）};

    % Right: compact matrix & smoothness comparison box
    \node[anchor=west, align=left, font=\footnotesize, draw=codeborder, fill=codebg, rounded corners=1.5mm, inner sep=6pt] at (8.4, -0.2) {%
      $L = D - W = \begin{pmatrix}
        1 & -1 & 0 & 0 \\
        -1 & 1.2 & -0.2 & 0 \\
        0 & -0.2 & 1.2 & -1 \\
        0 & 0 & -1 & 1
      \end{pmatrix}$\\[6pt]
      \textbf{平滑信号} $x_{\mathrm{s}}=(1,1,-1,-1)^T$：\;$x_{\mathrm{s}}^T L x_{\mathrm{s}} = \mathbf{0.8}$\\
      \textbf{震荡信号} $x_{\mathrm{o}}=(1,-1,1,-1)^T$：\;$x_{\mathrm{o}}^T L x_{\mathrm{o}} = \mathbf{8.8}$\\
      \textbf{Fiedler 向量}（$\lambda_2 \approx 0.180$）：$q_2 \approx (0.55, 0.45, -0.45, -0.55)^T$
    };
  \end{tikzpicture}}
  \caption{4 节点双簇弱桥图的拉普拉斯矩阵与信号平滑度对比。第二小特征值 $\lambda_2$ 对应的特征向量（Fiedler 向量）在保证与 $\mathbf{1}$ 正交的前提下让 $x^T L x$ 最小，因此天然切开弱桥、把两个紧耦合社群按正负号分开——这正是第 4 卷\textbf{谱聚类}与图神经网络（GNN）平滑正则化的核心思想。}
  \label{fig:ch3-graph-lap}
\end{figure}

\subsection{NumPy 代码实战：谱分解、马尔可夫迭代与图拉普拉斯}

在工程中处理实对称矩阵时，务必牢记一条 NumPy 黄金法则：\textbf{永远优先调用 \texttt{np.linalg.eigh}（专为 Hermitian/对称阵设计的算法），而不是通用的 \texttt{np.linalg.eig}}。\texttt{eigh} 不仅速度快约 2--3 倍，而且保证返回的特征值严格为升序实数、特征向量严格机器精度正交。

\begin{lstlisting}[caption={实对称谱分解（\texttt{eigh}）、行随机阵平稳分布与图拉普拉斯平滑度验证},label={lst:ch3-verify}]
import numpy as np

# 1. Real symmetric spectral decomposition via eigh (guaranteed real + orthogonal)
S = np.array([[5.0, 3.0], [3.0, 5.0]])
eigvals, Q = np.linalg.eigh(S)  # ascending: [2.0, 8.0]
S_recon = Q @ np.diag(eigvals) @ Q.T
assert np.allclose(S, S_recon) and np.allclose(Q.T @ Q, np.eye(2))
kappa_S = eigvals[-1] / eigvals[0]
print(f"eigs={eigvals}, cond={kappa_S:.2f}")  # [2. 8.], cond=4.00

# 2. Row-stochastic Markov power iteration
P = np.array([[0.7, 0.2, 0.1], [0.1, 0.6, 0.3], [0.2, 0.3, 0.5]])
assert np.allclose(P.sum(axis=1), 1.0)
pi_20 = np.array([1.0, 0.0, 0.0]) @ np.linalg.matrix_power(P, 20)
pi_star = np.array([11.0, 13.0, 10.0]) / 34.0
print(f"pi_20={np.round(pi_20, 4)}, L1_err={np.linalg.norm(pi_20 - pi_star, 1):.2e}")

# 3. 4-node Graph Laplacian quadratic smoothness & Fiedler vector
W = np.array([
    [0.0, 1.0, 0.0, 0.0],
    [1.0, 0.0, 0.2, 0.0],
    [0.0, 0.2, 0.0, 1.0],
    [0.0, 0.0, 1.0, 0.0],
])
L = np.diag(W.sum(axis=1)) - W
l_vals, l_vecs = np.linalg.eigh(L)
x_smooth = np.array([1.0, 1.0, -1.0, -1.0])
quad_mat = float(x_smooth @ L @ x_smooth)
quad_sum = 0.5 * float(np.sum(W * (x_smooth[:, None] - x_smooth[None, :]) ** 2))
assert np.isclose(quad_mat, quad_sum)
print(f"L eigs={np.round(l_vals, 4)}, Fiedler q2={np.round(l_vecs[:, 1], 3)}")
\end{lstlisting}

% =========================================================================
\section{本章习题}
\label{sec:ch3-exercises}

\exTier{A. 概念辨析与反例查错}{合上书独立作答，精准定位定义边界}

\begin{enumerate}[label=\textbf{3.\arabic*}, leftmargin=2.4em]
  \item \textbf{（对称化与正定判据反例）} 某同学声称：“若实方阵 $A$ 的所有特征值都严格大于 $0$，则二次型 $x^T A x$ 对一切非零向量 $x$ 必严格大于 $0$。”请构造一个 $2\times 2$ 反例推翻该说法，并通过对称化矩阵 $S = \frac{1}{2}(A + A^T)$ 解释其错误根源；进一步举出一个 $2\times 2$ 实对称阵反例，说明“所有顺序主子式 $\ge 0$”并不足以保证矩阵半正定。

  \item \textbf{（特征分解、瞬态增长与行随机阵辨析）} 判断以下四个命题正误；对错误命题给出最小反例，对正确命题简述理由：
  \begin{enumerate}[label=(\alph*), leftmargin=1.8em]
    \item 若两个实方阵 $A, B$ 均可对角化，则 $A + B$ 必可对角化；
    \item 实对称矩阵属于不同特征值的特征向量必相互正交；
    \item 若实方阵 $A$ 的全部特征值均满足 $|\lambda_i| < 1$，则序列 $\|A^k\|_F$（$k=1,2,\dots$）必单调递减；
    \item 两个 $n\times n$ 行随机矩阵的乘积 $P_1 P_2$ 仍是行随机矩阵。
  \end{enumerate}
\end{enumerate}

\exTier{B. 小例手算与数值直觉}{亲手打通特征向量、椭球与拉普拉斯计算}

\begin{enumerate}[label=\textbf{3.\arabic*}, leftmargin=2.4em, resume]
  \item \textbf{（2D 实对称谱分解与等高椭圆）} 给定实对称矩阵 $S = \begin{pmatrix} 3 & -1 \\ -1 & 3 \end{pmatrix}$：
  \begin{enumerate}[label=(\arabic*), leftmargin=1.8em]
    \item 手算 $S$ 的全部特征值与单位正交特征向量，写出正交分解 $S = Q\Lambda Q^T$ 及秩 1 外积和形式；
    \item 求条件数 $\cond(S)$，指出椭圆等高线 $x^T S x = 8$ 的两条主轴方向及对应的半轴长度；
    \item 不直接乘矩阵，利用谱分解口算 $S^{10}$ 的最大特征值并写出 $S^{-1}$。
  \end{enumerate}

  \item \textbf{（3 节点图拉普拉斯与随机游走平稳分布）} 考虑 3 节点无向链图 $1 - 2 - 3$（边 $(1,2)$ 与 $(2,3)$ 权重均为 $1$）：
  \begin{enumerate}[label=(\arabic*), leftmargin=1.8em]
    \item 写出邻接矩阵 $W$、度矩阵 $D$ 与图拉普拉斯矩阵 $L = D - W$；
    \item 手算 $L$ 的全部特征值与单位正交特征向量，验证 $L \succeq 0$ 且特征值 $0$ 仅为一重根；
    \item 构造随机游走转移矩阵 $P = D^{-1}W$，验证 $P$ 为行随机矩阵，并求出平稳分布 $\pi^*$（满足 $\pi^* P = \pi^*,\; \sum_i \pi_i^* = 1$）。
  \end{enumerate}
\end{enumerate}

\exTier{C. 实用推导与性质证明}{AI 理论与后续卷次反复用到的核心工具}

\begin{enumerate}[label=\textbf{3.\arabic*}, leftmargin=2.4em, resume]
  \item \textbf{（PSD 凸锥性质、岭回归正定化与矩阵平方根）} 设实对称矩阵 $A, B \in \R^{n\times n}$ 满足 $A \succeq 0,\; B \succeq 0$，常数 $\alpha, \beta \ge 0$。证明：
  \begin{enumerate}[label=(\arabic*), leftmargin=1.8em]
    \item $\alpha A + \beta B \succeq 0$（即半正定矩阵集合构成一个凸锥）；
    \item 对任意设计矩阵 $X \in \R^{m\times n}$ 与 $\lambda > 0$，岭回归矩阵 $H = X^T X + \lambda I_n$ 严格正定，并用 $X^T X$ 的最大/最小特征值 $\mu_{\max}, \mu_{\min}$ 给出条件数 $\cond(H)$ 的精确表达式（说明增大 $\lambda$ 如何改善条件数）；
    \item 对任意 $A \succeq 0$，存在半正定对称矩阵 $A^{1/2} \succeq 0$ 满足 $(A^{1/2})^2 = A$（并简证该半正定平方根唯一）。
  \end{enumerate}
  \exHint{第 (3) 问可直接令 $A^{1/2} = Q \Lambda^{1/2} Q^T$；证唯一性时，设 $B \succeq 0$ 亦满足 $B^2 = A$，先观察 $B$ 与 $A$ 的特征向量关系：对任意 $B$ 的特征对 $Bv = \mu v$（$\mu \ge 0$），有 $Av = B^2 v = \mu^2 v$，故 $v$ 必落入 $A$ 对应 $\mu^2$ 的特征子空间中。}

  \item \textbf{（Rayleigh 商的约束极小值与 Fiedler 向量）} 设实对称矩阵 $S \in \R^{n\times n}$ 的特征值按升序排列为 $\lambda_1 \le \lambda_2 \le \dots \le \lambda_n$，对应单位正交特征向量为 $q_1, \dots, q_n$。证明：
  \[
    \lambda_2 = \min_{\substack{x \neq 0 \\ x \perp q_1}} \frac{x^T S x}{x^T x}.
  \]
  并据此说明：对于连通无向图的拉普拉斯矩阵 $L$（其 $q_1 = \frac{1}{\sqrt{n}}\mathbf{1}$），在约束 $\sum_{i=1}^n x_i = 0$（消除常数平移）且 $\|x\|_2 = 1$ 下，使平滑度能量 $\sum_{\{i,j\}\in E} w_{ij}(x_i - x_j)^2$ 最小的非平凡信号恰为 Fiedler 向量 $q_2$。
  \exHint{在标准正交基 $\{q_1, \dots, q_n\}$ 下展开 $x = \sum_{i=1}^n y_i q_i$，正交约束 $x \perp q_1$ 恰好等价于第一个坐标 $y_1 = 0$。}
\end{enumerate}

\exTier{D. 编程实战与可自检实验}{用 NumPy 亲手跑出数值现象并通过断言检验}

\begin{enumerate}[label=\textbf{3.\arabic*}, leftmargin=2.4em, resume]
  \item \textbf{（高阶 Jordan 块瞬态增长与 \texttt{eigh} vs \texttt{eig} 对比实验）} 编写 Python 脚本完成以下两项实验：
  \begin{enumerate}[label=(\arabic*), leftmargin=1.8em]
    \item 构造 $m$ 阶 Jordan 块 $J_m(0.9)$（$m \in \{2, 3, 4\}$），计算 $k = 0, 1, \dots, 120$ 时的 $\|J_m(0.9)^k\|_F$，输出各阶对应的峰值和达峰步数；
    \item 随机生成 $200\times 200$ 实对称阵 $S = \frac{1}{2}(A + A^T)$，分别调用 \texttt{np.linalg.eig} 与 \texttt{np.linalg.eigh}，对比二者运行耗时、特征值虚部的最大绝对值，以及特征向量正交误差 $\|Q^T Q - I\|_F$。
  \end{enumerate}
  \par\smallskip\noindent{\footnotesize\sffamily\color{tealDark}\textbf{【自检标准】}\,当 $m=2$ 时峰值须与例 \ref{ex:ch3-jordan-transient} 一致（$k=9$ 时 $\|J_2(0.9)^9\|_F \approx 3.9128$），$m=4$ 时峰值须超过 $100$；\texttt{eigh} 所得特征值虚部严格为 $0$，且正交误差 $\|Q^T Q - I\|_F < 10^{-12}$。}

  \item \textbf{（双社群图的拉普拉斯谱聚类实验）} 构造一个 20 节点无向图：将节点 $0\text{--}9$ 与 $10\text{--}19$ 分为两个社群，社群内部以全连接（或每对以概率 $0.8$ 连边）赋予边权 $1.0$，两社群之间仅连 3 条权重为 $0.2$ 的弱桥边。编写代码构造 $L = D - W$ 并调用 \texttt{np.linalg.eigh(L)}：
  \begin{enumerate}[label=(\arabic*), leftmargin=1.8em]
    \item 验证最小特征值 $|\lambda_1| < 10^{-12}$ 且 $\lambda_2 > 0$（图连通）；
    \item 提取 Fiedler 向量 $q_2$，用 $\operatorname{sign}(q_2)$ 给 20 个节点分类；
    \item 用随机置换打乱 20 个节点的原始编号后重新运行，计算社群恢复准确率。
  \end{enumerate}
  \par\smallskip\noindent{\footnotesize\sffamily\color{tealDark}\textbf{【自检标准】}\,断言二次型恒等式 $\abs{q_2^T L q_2 - \frac{1}{2}\sum_{i,j} w_{ij}(q_{2,i}-q_{2,j})^2} < 10^{-12}$ 通过；无论节点如何随机打乱，按 $q_2$ 符号（允许整体反号）划分两社群的准确率均应达到 $100\%$。}
\end{enumerate}

\chapter[奇异值分解、伪逆与 PCA 降维]{奇异值分解（SVD）、最优低秩近似、伪逆与 PCA 降维}
\label{ch:ch4}

\chapEpigraph{完美不是无以复加，而是无可删减。}{安托万·德·圣-埃克苏佩里（Antoine de Saint-Exup\'ery），《风沙星辰》}

\begin{chapIntro}{本章定位：AI 线性代数的王冠}
第 3 章的实对称阵谱定理美轮美奂，但它有一个严格门槛：矩阵必须\textbf{既是方阵又对称}（$S^T = S$）。然而，人工智能中绝大多数真实矩阵——样本特征表 $X \in \R^{n\times d}$、灰度图像 $I \in \R^{H\times W}$、神经网络线性层权重 $W \in \R^{d_{\mathrm{out}}\times d_{\mathrm{in}}}$——几乎都是\textbf{长方形且非对称}的。

\textbf{奇异值分解（Singular Value Decomposition, SVD）}把谱定理的威力解放到了\textbf{任意实矩阵 $A \in \R^{m\times n}$}上。本章将以零跳步的方式完成四件大事：
\begin{enumerate}
  \item 直接用第 3 章的谱定理构造出 SVD，看清它对四大基本子空间与三步旋转拉伸几何的精确刻画；
  \item 用 10 行极简证明拿下\textbf{Eckart--Young 最优低秩近似定理}，直连第 6 卷大模型 \textbf{LoRA} 低秩微调；
  \item 定义 \textbf{Moore--Penrose 伪逆}，回答第 2 章留下的核心悬念：欠定（过参数化）系统如何选出唯一最小范数解；
  \item 证明\textbf{主成分分析（PCA）}两大视角的严格等价，并用 $64\times 64$ 图像压缩实测完成理论与代码的闭环。
\end{enumerate}
\end{chapIntro}

% =========================================================================
\section{从谱定理到奇异值分解（SVD）的完整构造与三步几何}
\label{sec:ch4-svd-construction}

\subsection{为什么从 $A^T A$ 出发？四步构造完整证明}

给定任意长方形矩阵 $A \in \R^{m\times n}$，我们无法直接对 $A$ 谈特征值。但只需构造它的格拉姆矩阵 $A^T A \in \R^{n\times n}$，奇迹便立刻发生：
\begin{itemize}
  \item $A^T A$ 是 $n\times n$ \textbf{实对称方阵}：$(A^T A)^T = A^T A$；
  \item $A^T A$ 是\textbf{半正定矩阵}（第 3 章）：对任意 $v \in \R^n$，$v^T (A^T A) v = \|Av\|_2^2 \ge 0$；
  \item $\rank(A^T A) = \rank(A) =: r$（第 2 章核心结论：$\nul(A^T A) = \nul(A)$）。
\end{itemize}
因此，第 3 章的谱定理立刻对 $A^T A$ 生效——这就是构造 SVD 的全部钥匙。

\begin{thmBox}
\begin{theorem}[奇异值分解存在性定理, Singular Value Decomposition]
\label{thm:ch4-svd}
对任意实矩阵 $A \in \R^{m\times n}$（设 $\rank(A) = r$），必存在正交矩阵 $U = [u_1,\dots,u_m] \in \R^{m\times m}$、正交矩阵 $V = [v_1,\dots,v_n] \in \R^{n\times n}$，以及长方形对角矩阵 $\Sigma \in \R^{m\times n}$（其前 $r$ 个主对角元为正实数 $\sigma_1 \ge \sigma_2 \ge \dots \ge \sigma_r > 0$，其余元素全为 $0$），使得
\begin{equation}
  A = U \Sigma V^T = U_r \Sigma_r V_r^T = \sum_{i=1}^r \sigma_i u_i v_i^T, \label{eq:ch4-svd}
\end{equation}
其中 $U_r = [u_1,\dots,u_r] \in \R^{m\times r}$，$V_r = [v_1,\dots,v_r] \in \R^{n\times r}$，$\Sigma_r = \diag(\sigma_1,\dots,\sigma_r)$。标量 $\sigma_i$ 称为 $A$ 的\textbf{奇异值}，$u_i$ 与 $v_i$ 分别称为\textbf{左奇异向量}与\textbf{右奇异向量}。
\end{theorem}
\end{thmBox}

\begin{proof}
分四步直接构造：
\begin{enumerate}
  \item \textbf{对 $A^T A$ 施实对称谱定理}：因 $A^T A \in \R^{n\times n}$ 对称半正定且秩为 $r$，由定理 \ref{thm:ch3-spectral} 与定理 \ref{thm:ch3-psd-eig}，存在正交矩阵 $V = [v_1,\dots,v_n] \in \R^{n\times n}$ 及按降序排列的非负特征值
  \[
    \lambda_1 \ge \lambda_2 \ge \dots \ge \lambda_r > 0 = \lambda_{r+1} = \dots = \lambda_n,
  \]
  满足 $A^T A v_i = \lambda_i v_i$（$i=1,\dots,n$）。定义正奇异值 $\sigma_i := \sqrt{\lambda_i} > 0$（$1 \le i \le r$）。
  \item \textbf{构造前 $r$ 个左奇异向量并验证标准正交性}：对每个 $1 \le i \le r$，定义
  \[
    u_i := \frac{1}{\sigma_i} A v_i \in \R^m.
  \]
  直接计算任意两个 $u_i, u_j$（$1 \le i, j \le r$）的内积：
  \[
    u_i^T u_j
    = \frac{1}{\sigma_i \sigma_j} (A v_i)^T (A v_j)
    = \frac{1}{\sigma_i \sigma_j} v_i^T (A^T A v_j)
    = \frac{\sigma_j^2}{\sigma_i \sigma_j} v_i^T v_j
    = \frac{\sigma_j}{\sigma_i} \delta_{ij}
    = \delta_{ij}.
  \]
  故 $\{u_1,\dots,u_r\}$ 在 $\R^m$ 中天然标准正交！此外，对任意 $j > r$，因 $\|A v_j\|_2^2 = v_j^T (A^T A) v_j = 0$，必有 $A v_j = 0$。
  \item \textbf{补全 $\R^m$ 的标准正交基}：利用第 2 章的 Gram--Schmidt 正交化，将 $\{u_1,\dots,u_r\}$ 扩充为 $\R^m$ 的一组标准正交基 $\{u_1,\dots,u_r, u_{r+1},\dots,u_m\}$，拼成正交矩阵 $U \in \R^{m\times m}$。
  \item \textbf{验证 $AV = U\Sigma$}：逐列对照两个 $m\times n$ 矩阵。前 $r$ 列中，$(AV)_{:,i} = A v_i = \sigma_i u_i = (U\Sigma)_{:,i}$；后 $n-r$ 列中，$(AV)_{:,j} = A v_j = 0 = (U\Sigma)_{:,j}$。因此 $AV = U\Sigma$。两边右乘 $V^T$（利用 $V V^T = I_n$）并按第 1 章外积视角展开，即得
  \[
    A = U \Sigma V^T = U_r \Sigma_r V_r^T = \sum_{i=1}^r \sigma_i u_i v_i^T.
  \]
\end{enumerate}
\end{proof}

\textbf{三种尺寸写法速辨}：实际编码中（记 $p = \min(m, n)$），除上述\textbf{完整 SVD}（$U_{m\times m},\;\Sigma_{m\times n},\;V^T_{n\times n}$）外，更常用节省内存的\textbf{精简 SVD}（Thin SVD，对应 \texttt{full\_matrices=False}：$U_{m\times p},\;\Sigma_{p\times p},\;V^T_{p\times n}$）以及仅保留非零奇异值的\textbf{紧凑 SVD}（$U_r \Sigma_r V_r^T$）。三者相乘还原的都是同一个矩阵 $A$。

\subsection{四大基本子空间、矩阵范数与一般条件数}

SVD 一次性给出了第 2 章\textbf{全部四大基本子空间的标准正交基}，毫无冗余（表 \ref{tab:ch4-subspaces}）：

\begin{table}[htbp]
  \centering
  \caption{SVD 对一般矩阵 $A \in \R^{m\times n}$（秩为 $r$）四大基本子空间的精确正交基刻画}
  \label{tab:ch4-subspaces}
  \small
  \begin{tabularx}{\linewidth}{@{}p{26mm}p{38mm}p{16mm}X@{}}
    \toprule
    \textbf{基本子空间} & \textbf{标准正交基} & \textbf{维数} & \textbf{SVD 几何角色} \\
    \midrule
    列空间 $\col(A)$ & $\{u_1, \dots, u_r\}$ & $r$ & $A$ 的输出值域；每列 $Ax = \sum_{i=1}^r \sigma_i(v_i^T x)u_i$ \\
    左零空间 $\nul(A^T)$ & $\{u_{r+1}, \dots, u_m\}$ & $m - r$ & $\col(A)^\perp$；回归残差 $b - \hat{b}$ 所在空间 \\
    行空间 $\row(A)$ & $\{v_1, \dots, v_r\}$ & $r$ & 真正被 $A$ 映射到非零输出的“有效输入方向” \\
    零空间 $\nul(A)$ & $\{v_{r+1}, \dots, v_n\}$ & $n - r$ & $\row(A)^\perp$；被 $A$ 完全压扁为零的输入方向 \\
    \bottomrule
  \end{tabularx}
\end{table}

此外，SVD 让两个最常用的矩阵范数以及一般矩阵的条件数变得一目了然。

\begin{proposition}[SVD 刻画算子 2-范数、Frobenius 范数与一般条件数]
\label{prop:ch4-norms}
设非零矩阵 $A \in \R^{m\times n}$ 的非零奇异值为 $\sigma_1 \ge \dots \ge \sigma_r > 0$，则：
\begin{enumerate}
  \item \textbf{算子 2-范数（谱范数）}：$\|A\|_2 := \max_{\|x\|_2 = 1} \|Ax\|_2 = \sigma_1$（在 $x = v_1$ 处取到）；
  \item \textbf{Frobenius 范数}：$\|A\|_F := \sqrt{\sum_{i,j} A_{ij}^2} = \sqrt{\tr(A^T A)} = \sqrt{\sum_{i=1}^r \sigma_i^2}$；
  \item \textbf{一般 2-范数条件数}：对列满秩（或方阵可逆）矩阵，$\cond_2(A) := \dfrac{\sigma_1}{\sigma_r}$（一般非零矩阵亦称有效条件数）；特别地，当 $A \succ 0$ 时 $\sigma_i = \lambda_i$，与第 3 章的 $\cond(A) = \lambda_{\max}/\lambda_{\min}$ 完全一致！
\end{enumerate}
\end{proposition}

\begin{proof}
对任意单位向量 $x \in \R^n$，令 $y = V^T x$（因 $V$ 正交，$\|y\|_2 = \|x\|_2 = 1$）。再由 $U$ 正交保长：
\[
  \|Ax\|_2^2 = \|U\Sigma V^T x\|_2^2 = \|\Sigma y\|_2^2 = \sum_{i=1}^r \sigma_i^2 y_i^2 \;\le\; \sigma_1^2 \sum_{i=1}^r y_i^2 \;\le\; \sigma_1^2,
\]
当 $x = v_1$（即 $y = e_1$）时取等号，故 $\|A\|_2 = \sigma_1$。对 Frobenius 范数，利用迹的循环不变性与正交性，
\[
  \|A\|_F^2 = \tr(A^T A) = \tr(V \Sigma^T \Sigma V^T) = \tr(\Sigma^T \Sigma V^T V) = \tr(\Sigma^T \Sigma) = \sum_{i=1}^r \sigma_i^2.
\]
\end{proof}

\subsection{手算两例与三步几何可视化}

\begin{example}[$2\times 2$ 非对称阵的 SVD 与三步几何]
\label{ex:ch4-svd-2x2}
考察非对称矩阵（它的特征值甚至是复数 $\frac{3 \pm i\sqrt{7}}{2}$，在实数域上根本无法特征对角化！）：
\[
  M = \begin{pmatrix} 2 & 2 \\ -1 & 1 \end{pmatrix}.
\]
\begin{enumerate}
  \item \textbf{构造对称阵 $M^T M$}：
  \[
    M^T M = \begin{pmatrix} 2 & -1 \\ 2 & 1 \end{pmatrix}\begin{pmatrix} 2 & 2 \\ -1 & 1 \end{pmatrix} = \begin{pmatrix} 5 & 3 \\ 3 & 5 \end{pmatrix}.
  \]
  它正好是第 3 章例 \ref{ex:ch3-sym2x2} 中手算过的对称矩阵！直接复用结果：特征值 $\lambda_1 = 8,\; \lambda_2 = 2$，因此奇异值为 $\sigma_1 = 2\sqrt{2},\; \sigma_2 = \sqrt{2}$，右奇异向量为
  \[
    v_1 = \frac{1}{\sqrt{2}}\begin{pmatrix} 1 \\ 1 \end{pmatrix}, \qquad
    v_2 = \frac{1}{\sqrt{2}}\begin{pmatrix} -1 \\ 1 \end{pmatrix}.
  \]
  \item \textbf{计算左奇异向量}：
  \[
    u_1 = \frac{1}{\sigma_1} M v_1 = \frac{1}{2\sqrt{2}}\cdot \frac{1}{\sqrt{2}}\begin{pmatrix} 4 \\ 0 \end{pmatrix} = \begin{pmatrix} 1 \\ 0 \end{pmatrix}, \qquad
    u_2 = \frac{1}{\sigma_2} M v_2 = \frac{1}{\sqrt{2}}\cdot \frac{1}{\sqrt{2}}\begin{pmatrix} 0 \\ 2 \end{pmatrix} = \begin{pmatrix} 0 \\ 1 \end{pmatrix}.
  \]
  \item \textbf{拼装与验算}：$U = I_2$，$\Sigma = \diag(2\sqrt{2}, \sqrt{2})$，$V^T = \frac{1}{\sqrt{2}}\begin{psmallmatrix}1 & 1 \\ -1 & 1\end{psmallmatrix}$，立刻验证 $U\Sigma V^T = M$，且 $\|M\|_2 = 2\sqrt{2},\; \|M\|_F = \sqrt{8+2} = \sqrt{10},\; \cond_2(M) = 2$。
\end{enumerate}
\end{example}

例 \ref{ex:ch4-svd-2x2} 揭示了任意线性变换的\textbf{普适几何本质}（图 \ref{fig:ch4-svd-geo}）：无论矩阵 $A$ 看起来多么歪斜甚至不可对角化，它对单位球面的作用永远可以拆解为三步——先用 $V^T$ \textbf{正交旋转/反射}对齐标准轴，再用 $\Sigma$ 沿各坐标轴做纯\textbf{轴向拉伸/压扁}，最后用 $U$ 做一次\textbf{正交旋转/反射}摆放到输出空间！

\begin{figure}[htbp]
  \centering
  \resizebox{0.97\linewidth}{!}{%
  \begin{tikzpicture}[>=Stealth]
    % Panel 1: Input unit circle with v1, v2
    \begin{scope}[shift={(0,0)}]
      \draw[->, slate, thin] (-1.65,0) -- (1.65,0) node[right, font=\scriptsize] {$x_1$};
      \draw[->, slate, thin] (0,-1.65) -- (0,1.65) node[above, font=\scriptsize] {$x_2$};
      \draw[teal, thick, fill=tealLight, fill opacity=0.45] (0,0) circle (1.15);
      \draw[->, ochre, very thick] (0,0) -- (0.813, 0.813) node[above right, font=\footnotesize\bfseries] {$v_1$};
      \draw[->, navy, very thick]  (0,0) -- (-0.813, 0.813) node[above left, font=\footnotesize\bfseries] {$v_2$};
      \draw[slate, thin] (0.14, 0.14) -- (0, 0.28) -- (-0.14, 0.14); % right angle mark
      \node[below, font=\small\sffamily\bfseries, color=tealDark] at (0,-1.85) {(1) 输入单位圆};
      \node[below, font=\scriptsize, color=slate] at (0,-2.3) {正交基 $v_1, v_2$（$45^\circ$）};
    \end{scope}

    % Arrow 1 -> 2
    \draw[->, very thick, teal] (2.05, 0) -- (3.35, 0)
      node[midway, above, font=\footnotesize\sffamily\bfseries, color=tealDark] {$V^T$}
      node[midway, below, font=\scriptsize\sffamily, color=slate] {旋转对齐};

    % Panel 2: Rotated unit circle aligned with e1, e2
    \begin{scope}[shift={(5.4,0)}]
      \draw[->, slate, thin] (-1.65,0) -- (1.65,0) node[right, font=\scriptsize] {$y_1$};
      \draw[->, slate, thin] (0,-1.65) -- (0,1.65) node[above, font=\scriptsize] {$y_2$};
      \draw[teal, thick, fill=tealLight, fill opacity=0.45] (0,0) circle (1.15);
      \draw[->, ochre, very thick] (0,0) -- (1.15, 0) node[below right, font=\footnotesize\bfseries] {$e_1$};
      \draw[->, navy, very thick]  (0,0) -- (0, 1.15) node[above right, font=\footnotesize\bfseries] {$e_2$};
      \node[below, font=\small\sffamily\bfseries, color=tealDark] at (0,-1.85) {(2) 主轴对齐};
      \node[below, font=\scriptsize, color=slate] at (0,-2.3) {$V^T v_1=e_1,\;V^T v_2=e_2$};
    \end{scope}

    % Arrow 2 -> 3
    \draw[->, very thick, teal] (7.45, 0) -- (8.85, 0)
      node[midway, above, font=\footnotesize\sffamily\bfseries, color=tealDark] {$U\Sigma$}
      node[midway, below, font=\scriptsize\sffamily, color=slate] {拉伸与摆放};

    % Panel 3: Output ellipse with semi-axes sigma1*u1 and sigma2*u2
    \begin{scope}[shift={(11.85,0)}]
      \draw[->, slate, thin] (-2.65,0) -- (2.65,0) node[right, font=\scriptsize] {$z_1$};
      \draw[->, slate, thin] (0,-1.65) -- (0,1.65) node[above, font=\scriptsize] {$z_2$};
      \draw[teal, thick, fill=tealLight, fill opacity=0.45] (0,0) ellipse (2.25 and 1.125);
      \draw[->, ochre, very thick] (0,0) -- (2.25, 0) node[above left, font=\footnotesize\bfseries, yshift=1pt] {$\sigma_1 u_1$};
      \draw[->, navy, very thick]  (0,0) -- (0, 1.125) node[above right, font=\footnotesize\bfseries] {$\sigma_2 u_2$};
      \node[below, font=\small\sffamily\bfseries, color=tealDark] at (0,-1.85) {(3) 输出超椭球};
      \node[below, font=\scriptsize, color=slate] at (0,-2.3) {半轴长 $\sigma_1=2\sqrt{2},\;\sigma_2=\sqrt{2}$};
    \end{scope}
  \end{tikzpicture}}
  \caption{非对称矩阵 $M = \begin{psmallmatrix}2&2\\-1&1\end{psmallmatrix}$ 的 SVD 三步几何动作：先正交旋转 $V^T$ 将右奇异向量对齐坐标轴，再按奇异值 $\sigma_1,\sigma_2$ 轴向拉伸并由 $U$ 旋转摆放到输出空间（此处 $U=I_2$）。}
  \label{fig:ch4-svd-geo}
\end{figure}

\begin{example}[$3\times 2$ 长方形矩阵的完整与紧凑 SVD]
\label{ex:ch4-svd-3x2}
考察 $A = \begin{pmatrix} 1 & 0 \\ 0 & 1 \\ 1 & 0 \end{pmatrix} \in \R^{3\times 2}$。计算 $A^T A = \begin{pmatrix} 2 & 0 \\ 0 & 1 \end{pmatrix}$，得奇异值 $\sigma_1 = \sqrt{2},\; \sigma_2 = 1$ 及右奇异向量 $v_1 = (1,0)^T,\; v_2 = (0,1)^T$。于是
\[
  u_1 = \frac{1}{\sqrt{2}} A v_1 = \frac{1}{\sqrt{2}}\begin{pmatrix} 1 \\ 0 \\ 1 \end{pmatrix}, \qquad
  u_2 = \frac{1}{1} A v_2 = \begin{pmatrix} 0 \\ 1 \\ 0 \end{pmatrix}.
\]
这就已经给出了紧凑 SVD：$A = \sqrt{2}\,u_1 v_1^T + 1\,u_2 v_2^T$。若要补成完整 SVD，只需在 $\col(A)^\perp = \nul(A^T)$ 中找一个单位向量与 $u_1, u_2$ 正交，立刻看出 $u_3 = \frac{1}{\sqrt{2}}(1,0,-1)^T$，于是
\[
  A = \begin{pmatrix}
    1/\sqrt{2} & 0 & 1/\sqrt{2} \\
    0 & 1 & 0 \\
    1/\sqrt{2} & 0 & -1/\sqrt{2}
  \end{pmatrix}
  \begin{pmatrix}
    \sqrt{2} & 0 \\
    0 & 1 \\
    0 & 0
  \end{pmatrix}
  \begin{pmatrix} 1 & 0 \\ 0 & 1 \end{pmatrix}.
\]
\end{example}

% =========================================================================
\section{最优低秩近似（Eckart--Young 短证）与大模型 LoRA 直觉}
\label{sec:ch4-eckart-young}

观察 SVD 的外积展开式 $A = \sum_{i=1}^r \sigma_i u_i v_i^T$。因为奇异值按降序排列 $\sigma_1 \ge \sigma_2 \ge \dots \ge \sigma_r > 0$，且每个秩 1 矩阵 $u_i v_i^T$ 的 2-范数和 Frobenius 范数都恰好等于 $1$，所以这个求和天然是\textbf{按“能量/重要性”从大到小排列的}！

如果为了压缩存储或去噪，我们只保留前 $k$ 项（$1 \le k < r$），得到\textbf{秩-$k$ 截断 SVD}（truncated SVD）：
\[
  A_k := \sum_{i=1}^k \sigma_i u_i v_i^T = U_k \Sigma_k V_k^T.
\]
显然 $\rank(A_k) = k$，而且截断误差矩阵 $A - A_k = \sum_{i=k+1}^r \sigma_i u_i v_i^T$ 的奇异值就是被丢弃的尾巴 $\sigma_{k+1}, \dots, \sigma_r$，由命题 \ref{prop:ch4-norms} 立刻有
\[
  \|A - A_k\|_2 = \sigma_{k+1}, \qquad \|A - A_k\|_F = \sqrt{\sum_{i=k+1}^r \sigma_i^2}.
\]
令人震撼的是：在\textbf{所有}秩不超过 $k$ 的矩阵中，没有任何其他构造能比直接截断 SVD 更接近原矩阵 $A$！

\begin{thmBox}
\begin{theorem}[Eckart--Young--Mirsky 最优低秩近似定理]
\label{thm:ch4-eym}
设 $A = \sum_{i=1}^r \sigma_i u_i v_i^T \in \R^{m\times n}$，且整数 $1 \le k < r$。则对\textbf{任意}满足 $\rank(B) \le k$ 的矩阵 $B \in \R^{m\times n}$，同时有
\begin{equation}
  \|A - B\|_2 \;\ge\; \|A - A_k\|_2 = \sigma_{k+1}, \qquad
  \|A - B\|_F \;\ge\; \|A - A_k\|_F = \sqrt{\sum_{i=k+1}^r \sigma_i^2}. \label{eq:ch4-eym}
\end{equation}
\end{theorem}
\end{thmBox}

\begin{proof}[谱范数版本的 10 行精简证明]
我们要证：对任意满足 $\rank(B) \le k$ 的矩阵 $B$，总能找到一个单位向量 $z \in \R^n$ 使得 $\|(A-B)z\|_2 \ge \sigma_{k+1}$。
\begin{enumerate}
  \item 由第 2 章秩-零度定理，$\rank(B) \le k \implies \dim\nul(B) = n - \rank(B) \ge n - k$。
  \item 记前 $k+1$ 个右奇异向量张成的子空间为 $W := \spanop\{v_1, \dots, v_{k+1}\} \subseteq \R^n$，显然 $\dim(W) = k+1$。
  \item 由子空间维数公式（因 $\nul(B) + W \subseteq \R^n$ 维数至多为 $n$）：
  \[
    \dim\bigl(\nul(B) \cap W\bigr) = \dim\nul(B) + \dim(W) - \dim\bigl(\nul(B) + W\bigr) \ge (n - k) + (k + 1) - n = 1.
  \]
  \item 因此交空间 $\nul(B) \cap W$ 中至少存在一个\textbf{单位向量} $z$（$\|z\|_2 = 1$）。既然 $z \in W$，可写成 $z = \sum_{i=1}^{k+1} c_i v_i$，且系数满足 $\sum_{i=1}^{k+1} c_i^2 = \|z\|_2^2 = 1$。
  \item 又因 $z \in \nul(B)$ 即 $Bz = 0$，利用 $\{u_1,\dots,u_{k+1}\}$ 的标准正交性与 $\sigma_1 \ge \dots \ge \sigma_{k+1}$：
  \[
    \|(A - B)z\|_2^2 = \|Az\|_2^2
    = \norm{\sum_{i=1}^{k+1} \sigma_i c_i u_i}_2^2
    = \sum_{i=1}^{k+1} \sigma_i^2 c_i^2
    \;\ge\; \sigma_{k+1}^2 \sum_{i=1}^{k+1} c_i^2
    = \sigma_{k+1}^2.
  \]
\end{enumerate}
两边开方并由算子 2-范数定义，得 $\|A - B\|_2 \ge \|(A - B)z\|_2 \ge \sigma_{k+1} = \|A - A_k\|_2$。
\end{proof}

\begin{extBox}{Frobenius 范数版本的证明思路}
Frobenius 范数版本 $\|A - B\|_F^2 \ge \sum_{i=k+1}^r \sigma_i^2$ 同样不需要繁重的行列式技巧，只需结合第 2 章的正交投影即可证明：固定 $\col(B)$（其维数 $\le k$），在所有列空间含于 $\col(B)$ 的矩阵中，最小化 $\|A - B\|_F^2 = \sum_{j=1}^n \|a_j - b_j\|_2^2$ 的唯一最优解正是把 $A$ 的每一列正交投影到 $\col(B)$ 上，即 $B^* = P_k A$（$P_k$ 为秩 $\le k$ 的正交投影阵）。于是 $\|A - B\|_F^2 \ge \|(I - P_k)A\|_F^2 = \tr\bigl((I - P_k)A A^T\bigr)$，归结为在 $(m-k)$ 维子空间上极小化对称阵 $AA^T$ 的迹（由谱定理/Ky Fan 极值原理，最小值恰为最小 $m-k$ 个特征值之和 $\sum_{i=k+1}^r \sigma_i^2$）。初读直接记住并使用式 \eqref{eq:ch4-eym} 即可。
\end{extBox}

\begin{insightBox}{从截断 SVD 到第 6 卷大模型 LoRA 微调（$\Delta W = B A$）}
在微调十亿至千亿参数的大语言模型时，直接更新每一层的全尺寸权重矩阵 $W_0 \in \R^{d\times d}$（例如 $d = 4096$ 时单个矩阵就含 $1677$ 万个参数）代价极高。研究发现，微调产生的权重增量 $\Delta W = W_{\mathrm{ft}} - W_0$ 虽然尺寸庞大，但其奇异值衰减极快，具有很低的\textbf{内在秩}（intrinsic rank）。

\textbf{LoRA（Low-Rank Adaptation，第 6 卷核心技术）}正是 Eckart--Young 思想的直接工程化：冻结原始权重 $W_0$，把增量强制参数化为两个瘦长矩阵的乘积（秩 $k \ll d$）：
\[
  W = W_0 + \Delta W = W_0 + B A, \qquad B \in \R^{d\times k},\quad A \in \R^{k\times d}.
\]
当 $d = 4096,\; k = 8$ 时，可训练参数量从 $4096^2 = 16\,777\,216$ 暴降至 $8 \times (4096 + 4096) = 65\,536$——整整\textbf{压缩了 256 倍（仅占 0.39\%）}，却能保留 $\Delta W$ 绝大部分的主导低秩能量！
\end{insightBox}

% =========================================================================
\section{Moore--Penrose 伪逆与欠定系统的最小范数解}
\label{sec:ch4-pinv}

回到第 2 章末尾留下的开放问题：当线性回归设计矩阵 $A \in \R^{m\times n}$ \textbf{列满秩且超定}（$m > n,\; r = n$）时，正规方程给出的最小二乘解 $x^* = (A^T A)^{-1} A^T b$ 是唯一的。但如果 $A$ \textbf{秩亏}（例如特征共线）或者处于\textbf{欠定/过参数化}状态（样本数少于参数数 $m < n$，如同当代神经网络），$A^T A$ 不可逆，使 $\|Ax - b\|_2^2$ 最小的解有无穷多个（构成仿射子空间 $x_0 + \nul(A)$）。面对无穷多个同样能把训练误差压到最低的解，我们应当挑哪一个？SVD 给出了最自然、最唯一的答案：\textbf{Moore--Penrose 伪逆}。

\begin{defBox}
\begin{definition}[Moore--Penrose 伪逆, Pseudoinverse]
\label{def:ch4-pinv}
设矩阵 $A \in \R^{m\times n}$ 的 SVD 为 $A = U\Sigma V^T = U_r \Sigma_r V_r^T$（$\rank(A) = r$）。定义 $A$ 的 \textbf{Moore--Penrose 伪逆}为 $n\times m$ 矩阵：
\begin{equation}
  A^+ := V \Sigma^+ U^T = V_r \Sigma_r^{-1} U_r^T = \sum_{i=1}^r \frac{1}{\sigma_i} v_i u_i^T \;\in\; \R^{n\times m}, \label{eq:ch4-pinv}
\end{equation}
其中 $\Sigma^+ \in \R^{n\times m}$ 由将 $\Sigma$ 转置并把所有非零奇异值 $\sigma_i$ 替换为倒数 $1/\sigma_i$、零元保持为 $0$ 得到。
\end{definition}
\end{defBox}

立刻核对伪逆的几个关键特例与投影恒等式：
\begin{itemize}
  \item \textbf{可逆方阵}（$r = m = n$）：$A^+ = V \Sigma^{-1} U^T = A^{-1}$，退化为普通逆矩阵；
  \item \textbf{列满秩}（$r = n \le m$）：$A^+ = (A^T A)^{-1} A^T$（左逆，正好还原第 2 章的最小二乘解析解）；
  \item \textbf{行满秩}（$r = m \le n$）：$A^+ = A^T (A A^T)^{-1}$（右逆，满足 $A A^+ = I_m$）；
  \item \textbf{两大正交投影恒等式}（对照表 \ref{tab:ch4-subspaces} 的标准正交基！）：
  \begin{equation}
    A A^+ = U_r U_r^T = P_{\col(A)}, \qquad A^+ A = V_r V_r^T = P_{\row(A)}. \label{eq:ch4-pinv-proj}
  \end{equation}
  由此顺手验证全部四条 Moore--Penrose 代数公理：$A A^+ A = A$、$A^+ A A^+ = A^+$、$(AA^+)^T = AA^+$、$(A^+A)^T = A^+A$。
\end{itemize}

下面用 SVD 坐标变换一步证明伪逆在人工智能与优化中的核心地位。

\begin{thmBox}
\begin{theorem}[伪逆给出最小二乘集合中的唯一最小范数解]
\label{thm:ch4-min-norm}
对\textbf{任意}矩阵 $A \in \R^{m\times n}$（不限形状与秩）与任意向量 $b \in \R^m$，记 $x^+ := A^+ b$。则：
\begin{enumerate}
  \item $x^+$ 是最小二乘问题 $\min_{x\in\R^n} \|Ax - b\|_2^2$ 的全局极小解；
  \item 在所有满足最小二乘极小的解 $x^*$ 中，$x^+$ 具有\textbf{严格最小的欧氏范数}：$\|x^+\|_2 \le \|x^*\|_2$，等号成立当且仅当 $x^* = x^+$。
\end{enumerate}
\end{theorem}
\end{thmBox}

\begin{proof}
代入 $A = U\Sigma V^T$，令正交变量代换 $y := V^T x \in \R^n$ 与 $c := U^T b \in \R^m$。由于正交矩阵保欧氏范数，有 $\|x\|_2^2 = \|y\|_2^2$ 且
\begin{equation}
  \|Ax - b\|_2^2 = \|U(\Sigma V^T x - U^T b)\|_2^2 = \|\Sigma y - c\|_2^2
  = \sum_{i=1}^r (\sigma_i y_i - c_i)^2 + \sum_{i=r+1}^m c_i^2. \label{eq:ch4-svd-ls}
\end{equation}
观察式 \eqref{eq:ch4-svd-ls} 的两项：第二项 $\sum_{i=r+1}^m c_i^2$ 与变量 $y$ 完全无关（它是 $b$ 落在 $\nul(A^T)$ 中的不可消残差）；第一项是非负平方和，其最小值为 $0$，且\textbf{当且仅当}
\[
  y_i = \frac{c_i}{\sigma_i} \qquad (i = 1, \dots, r)
\]
时取到！而余下的坐标 $y_{r+1}, \dots, y_n$（对应零空间 $\nul(A)$ 的分量）可以任意取值而不改变残差。

现在，在所有最小二乘解中进一步极小化欧氏范数平方：
\[
  \|x\|_2^2 = \|y\|_2^2 = \sum_{i=1}^r \left(\frac{c_i}{\sigma_i}\right)^2 + \sum_{i=r+1}^n y_i^2.
\]
显然，此式取得严格唯一最小值的充要条件是令所有自由零空间分量为零：$y_{r+1} = \dots = y_n = 0$！将最优坐标 $y^+$ 换回原空间，即得
\[
  x^+ = V y^+ = \sum_{i=1}^r \frac{c_i}{\sigma_i} v_i = \sum_{i=1}^r \frac{u_i^T b}{\sigma_i} v_i = V_r \Sigma_r^{-1} U_r^T b = A^+ b.
\]
\end{proof}

\begin{example}[二维欠定方程的最小范数解]
\label{ex:ch4-underdet}
解单方程双未知数欠定系统 $x_1 + 2x_2 = 5$，即 $A = \begin{pmatrix}1 & 2\end{pmatrix} \in \R^{1\times 2},\; b = (5)$。因 $A$ 行满秩，直接算得 $A^+ = A^T(AA^T)^{-1} = \frac{1}{5}\begin{pmatrix}1\\2\end{pmatrix}$，于是最小范数解为 $x^+ = A^+ b = \begin{pmatrix}1\\2\end{pmatrix}$（范数 $\|x^+\|_2 = \sqrt{5}$）。从几何上看，所有解构成直线 $x_1 + 2x_2 = 5$，而 $x^+ = (1,2)^T \in \row(A)$ 恰好是\textbf{从原点向该直线所作垂线的垂足}——离原点最近的点！
\end{example}

\begin{insightBox}{过参数化神经网络的隐式偏置（Implicit Bias）}
现代深度学习模型往往参数量远超样本数（$n \gg m$），能够以零训练误差插值整个数据集。为什么它们常常没有彻底过拟合？一个关键理论起点正来自定理 \ref{thm:ch4-min-norm}：对欠定线性回归 $Aw = b$，若从零初始化 $w^{(0)} = 0$ 出发运行梯度下降 $w^{(t+1)} = w^{(t)} - \eta A^T(Aw^{(t)} - b)$，由于每一步更新向量 $A^T(\cdot)$ 永远落在行空间 $\row(A) = \nul(A)^\perp$ 内，迭代终点 $w^{(\infty)}$ 必然在 $\nul(A)$ 上零分量——也就是说，\textbf{哪怕不加任何显式正则项，零初始化梯度下降也会自动收敛到最小范数伪逆解 $w^+ = A^+ b$}！这正是第 4--5 卷“梯度下降隐式正则化”的雏形（习题 4.7 将带你亲手跑通这一惊人现象）。
\end{insightBox}

\begin{failBox}{微小奇异值放大噪声与截断阈值 \texttt{rcond}}
观察公式 $x^+ = \sum_{i=1}^r \frac{u_i^T b}{\sigma_i} v_i$：如果数据矩阵接近秩亏，某个奇异值小到 $\sigma_r = 10^{-12}$，那么观测向量 $b$ 中即便只有 $10^{-6}$ 的微小测量噪声，除以 $\sigma_r$ 后也会被瞬间放大 $10^{12}$ 倍，彻底摧毁解的精度！因此，NumPy 的 \texttt{np.linalg.pinv(A, rcond=1e-15)} 默认把所有小于 $\text{rcond}\cdot \sigma_1$ 的极小奇异值视作零而丢弃；而在机器学习回归中，更平滑的替代方案是第 2 章的岭回归（习题 4.6 将证明岭回归把生硬的 $1/\sigma_i$ 平滑阻尼为 $\frac{\sigma_i}{\sigma_i^2 + \lambda}$）。
\end{failBox}

% =========================================================================
\section{主成分分析（PCA）的双重视角等价与 SVD 图像压缩实战}
\label{sec:ch4-pca-compression}

\subsection{从二维倾斜点云出发：降维的两种等价愿望}

设想我们收集了 $n$ 名学生的“数学成绩”与“物理成绩”二维数据，画在平面上是一团沿着 $45^\circ$ 方向拉长的倾斜椭圆点云（中心在样本均值 $\mu$）。现在要把这批二维点压缩到穿过均值 $\mu$ 的\textbf{一条一维直线}上，该挑哪条直线？

直觉上有两种看似不同的标准：
\begin{itemize}
  \item \textbf{视角 1（统计视角 · 最大投影方差）}：希望投影到直线上的坐标尽可能拉开差距（方差最大），从而保留最多的信息区分度；
  \item \textbf{视角 2（几何视角 · 最小正交重构误差）}：希望每个样本点到这条直线的垂直距离平方和尽可能小，从而用一维坐标还原二维点时失真最小。
\end{itemize}
在任意 $d$ 维空间中把数据降到 $k$ 维子空间时，这两种愿望竟是\textbf{完全等价}的，而且其精确解正是去中心化数据矩阵的 SVD！

\subsection{PCA 双重视角等价定理的完整证明}

设原始数据矩阵为 $X \in \R^{n\times d}$（第 $i$ 行为样本 $x_i^T \in \R^{1\times d}$）。第一步\textbf{必须去中心化}：计算样本均值 $\mu = \frac{1}{n}\sum_{i=1}^n x_i = \frac{1}{n}X^T \mathbf{1} \in \R^d$，构造中心化矩阵
\[
  \widetilde{X} := X - \mathbf{1}\mu^T \in \R^{n\times d}, \qquad \text{第 } i \text{ 行为 } \tilde{x}_i^T = (x_i - \mu)^T.
\]
对应的样本协方差矩阵为 $\Sigma_{\mathrm{cov}} := \frac{1}{n}\widetilde{X}^T \widetilde{X} \in \R^{d\times d}$（第 3 章已证 $\Sigma_{\mathrm{cov}} \succeq 0$）。设我们寻找 $k$ 个两两正交的单位主轴向量，按列拼成矩阵 $W \in \R^{d\times k}$（满足 $W^T W = I_k$）。样本 $\tilde{x}_i$ 在该 $k$ 维子空间上的低维坐标为 $z_i = W^T \tilde{x}_i \in \R^k$，对应的正交重构点为 $\hat{x}_i = W z_i = W W^T \tilde{x}_i \in \R^d$。

\begin{thmBox}
\begin{theorem}[PCA 双重视角等价定理与 SVD 解]
\label{thm:ch4-pca}
设中心化数据矩阵 $\widetilde{X} \in \R^{n\times d}$ 的 SVD 为 $\widetilde{X} = U \Sigma V^T$，非零奇异值 $\sigma_1 \ge \dots \ge \sigma_r > 0$。对任意目标维数 $1 \le k \le r$，在正交约束 $W^T W = I_k$ 下：
\begin{enumerate}
  \item \textbf{最大投影方差问题} $\displaystyle\max_{W^T W = I_k} \tr(W^T \Sigma_{\mathrm{cov}} W) = \max_{W^T W = I_k} \frac{1}{n}\|\widetilde{X} W\|_F^2$，与
  \item \textbf{最小正交重构误差问题} $\displaystyle\min_{W^T W = I_k} \frac{1}{n}\sum_{i=1}^n \|\tilde{x}_i - W W^T \tilde{x}_i\|_2^2 = \min_{W^T W = I_k} \frac{1}{n}\|\widetilde{X} - \widetilde{X} W W^T\|_F^2$
\end{enumerate}
\textbf{严格等价}，且共同最优主轴矩阵由 $\widetilde{X}$ 的前 $k$ 个右奇异向量给出：$W^* = V_k = [v_1,\dots,v_k]$。降维后的主成分坐标矩阵为
\begin{equation}
  Z_k = \widetilde{X} V_k = U_k \Sigma_k \;\in\; \R^{n\times k}, \label{eq:ch4-pca-scores}
\end{equation}
前 $k$ 个主成分的\textbf{累计解释方差比}（explained variance ratio）为 $\dfrac{\sum_{i=1}^k \sigma_i^2}{\sum_{i=1}^r \sigma_i^2}$。
\end{theorem}
\end{thmBox}

\begin{proof}
\textbf{第一步：用勾股定理证双视角严格等价。}
因 $W^T W = I_k$，矩阵 $P := W W^T$ 是对称正交投影阵（$P^2 = P = P^T$），故对每个去中心化样本 $\tilde{x}_i$，重构误差向量 $\tilde{x}_i - P\tilde{x}_i$ 与投影向量 $P\tilde{x}_i$ 严格正交。由勾股定理，
\[
  \|\tilde{x}_i\|_2^2 = \|P\tilde{x}_i\|_2^2 + \|\tilde{x}_i - P\tilde{x}_i\|_2^2
  = \|W^T \tilde{x}_i\|_2^2 + \|\tilde{x}_i - W W^T \tilde{x}_i\|_2^2,
\]
其中用到 $\|W(W^T\tilde{x}_i)\|_2^2 = (W^T\tilde{x}_i)^T (W^T W)(W^T\tilde{x}_i) = \|W^T\tilde{x}_i\|_2^2$。对 $i = 1, \dots, n$ 求和并除以 $n$，得矩阵恒等式：
\begin{equation}
  \underbrace{\frac{1}{n}\|\widetilde{X}\|_F^2}_{\text{常数（总方差）}}
  = \underbrace{\frac{1}{n}\|\widetilde{X} W\|_F^2}_{\text{投影总方差}}
  + \underbrace{\frac{1}{n}\|\widetilde{X} - \widetilde{X} W W^T\|_F^2}_{\text{正交重构误差}}. \label{eq:ch4-pca-pythag}
\end{equation}
等号左侧与 $W$ 无关！因此最大化投影总方差与最小化正交重构误差是\textbf{同一件事}。

\textbf{第二步：求最优 $W^*$ 与坐标 $Z_k$。}
\begin{itemize}
  \item 当 $k = 1$ 时，投影方差 $\frac{1}{n}\|\widetilde{X}w\|_2^2 = w^T \Sigma_{\mathrm{cov}} w$ 正是协方差矩阵的 Rayleigh 商！因 $\Sigma_{\mathrm{cov}} = \frac{1}{n} V \Sigma^T \Sigma V^T$ 的最大特征值为 $\sigma_1^2/n$，由第 3 章推论 \ref{cor:ch3-rayleigh}，最优方向就是 $w^* = v_1$。
  \item 对一般 $k \ge 1$，注意到重构矩阵 $\widehat{X} := \widetilde{X} W W^T$ 的秩至多为 $k$。由 Eckart--Young--Mirsky 定理（定理 \ref{thm:ch4-eym} 的 Frobenius 范数界），
  \[
    \frac{1}{n}\|\widetilde{X} - \widetilde{X} W W^T\|_F^2 \;\ge\; \frac{1}{n}\sum_{i=k+1}^r \sigma_i^2.
  \]
  取 $W^* = V_k$，直接算出 $\widetilde{X} V_k V_k^T = (U\Sigma V^T) V_k V_k^T = U_k \Sigma_k V_k^T = \widetilde{X}_k$，其重构误差恰好等于理论下界 $\frac{1}{n}\sum_{i=k+1}^r \sigma_i^2$！此时投影方差为 $\frac{1}{n}\sum_{i=1}^k \sigma_i^2$，总方差为 $\frac{1}{n}\sum_{i=1}^r \sigma_i^2$，降维坐标满足 $Z_k = \widetilde{X} V_k = U\Sigma V^T V_k = U_k \Sigma_k$。
\end{itemize}
\end{proof}

\begin{failBox}{PCA 工程实践的两大经典翻车反例}
\begin{enumerate}
  \item \textbf{忘记去中心化（直接对原始 $X$ 做 SVD）}：设二维样本紧凑分布在 $(100, 100)$ 附近，真实数据沿 $(1, -1)$ 方向伸展（标准差为 $2$），沿 $(1,1)$ 方向几乎无变化。若忘记减去均值 $\mu = (100,100)^T$，原点到点云中心的距离平方高达 $20\,000$，彻底淹没真实方差 $4$——未中心化的“第一主成分”会盲目指向均值方向 $(1,1)/\sqrt{2}$，与真正方差最大的 $(1,-1)/\sqrt{2}$ 垂直！
  \item \textbf{异质特征未做量纲标准化}：设房价数据包含“面积（单位 $\mathrm{cm}^2$，量级 $10^6$）”与“卧室数量（量级 $3$）”。前者原始方差比后者大 $10^{11}$ 倍，直接做 PCA 会让第一主成分完全退化为“纯面积轴”，彻底无视卧室数与其他相关性。当各特征量纲不同时，务必先做 Z-score 标准化（减均值再除以各自标准差）再执行 PCA！
\end{enumerate}
\end{failBox}

\subsection{实测实战：$64\times 64$ 灰度图像的 SVD 压缩与 NumPy 实现}

为了直观感受低秩近似的压缩威力，我们在 \texttt{code/vol1\_experiments.py} 中构造了一幅非可分的 $64\times 64$ 合成灰度图像（叠加圆盘、圆环、对角高斯光带与二维正弦条纹，总计 $4096$ 个浮点数），对其做截断 SVD $A_k = U_k \Sigma_k V_k^T$。存储 $U_k \in \R^{64\times k}$、$\sigma_{1:k} \in \R^k$ 与 $V_k^T \in \R^{k\times 64}$ 共需 $k(64 + 64 + 1) = 129k$ 个浮点数，存储占比为 $129k / 4096$。表 \ref{tab:ch4-svd-compress} 与图 \ref{fig:ch4-svd-compress} 展示了实测数据（\texttt{data/svd\_compress.csv}）：

\begin{table}[htbp]
  \centering
  \caption{$64\times 64$ 非可分灰度图像截断 SVD 压缩实测结果（\texttt{data/svd\_compress.csv}）}
  \label{tab:ch4-svd-compress}
  \small
  \begin{tabular}{@{}cccc@{}}
    \toprule
    \textbf{保留秩 $k$} & \textbf{相对误差 $\|A-A_k\|_F/\|A\|_F$} & \textbf{累计能量占比 $\sum_{i\le k}\sigma_i^2/\sum\sigma_i^2$} & \textbf{存储占比 $129k/4096$} \\
    \midrule
    $1$  & $0.5657$ & $68.00\%$ & $3.15\%$ \\
    $2$  & $0.4526$ & $79.52\%$ & $6.30\%$ \\
    $4$  & $0.2617$ & $\mathbf{93.15\%}$ & $\mathbf{12.60\%}$ \\
    $8$  & $0.1320$ & $\mathbf{98.26\%}$ & $\mathbf{25.20\%}$ \\
    $12$ & $0.0874$ & $99.24\%$ & $37.79\%$ \\
    $16$ & $0.0542$ & $\mathbf{99.71\%}$ & $50.39\%$ \\
    $24$ & $0.0100$ & $99.99\%$ & $75.59\%$ \\
    \bottomrule
  \end{tabular}
\end{table}

请注意表中每行数据都严格满足勾股能量守恒恒等式：$\text{相对误差}^2 + \text{累计能量比} = 1$（例如 $k=4$ 时 $0.26167^2 + 0.93153 = 1.0000$）。仅保留 $k=4$ 个秩 1 图层，就能用 \textbf{12.60\%} 的存储体积还原 \textbf{93.15\%} 的图像总能量；保留 $k=8$ 时仅需约四分之一存储即可恢复 \textbf{98.26\%} 的能量！

\begin{figure}[htbp]
  \centering
  \begin{tikzpicture}
    \begin{axis}[
      name=leftplot,
      width=0.47\linewidth,
      height=5.2cm,
      xlabel={保留秩 $k$},
      ylabel={百分比（\%）},
      xmin=0, xmax=25,
      ymin=0, ymax=106,
      grid=major,
      grid style={dashed, gray!25},
      legend pos=south east,
      legend style={font=\scriptsize, fill=white, draw=codeborder},
      tick label style={font=\footnotesize},
      label style={font=\small\sffamily},
      title={\small\sffamily\bfseries 累计能量 vs 存储占比}
    ]
      \addplot[color=teal, very thick, mark=*, mark size=1.8pt] table[x=k, y=energy_pct, col sep=comma] {data/svd_compress.csv};
      \addlegendentry{累计能量占比}
      \addplot[color=ochre, thick, dashed, mark=square*, mark size=1.6pt] table[x=k, y=ratio_pct, col sep=comma] {data/svd_compress.csv};
      \addlegendentry{存储占比}
      \node[anchor=west, font=\scriptsize\sffamily, color=tealDark, fill=white, fill opacity=0.85, text opacity=1, inner sep=1pt] at (axis cs:6.8, 76) {$k=4$: $93.15\%$};
      \draw[-{Stealth[length=1.8mm]}, color=tealDark, thin] (axis cs:6.7, 79) -- (axis cs:4.3, 91.5);
    \end{axis}

    \begin{axis}[
      at={(leftplot.east)},
      anchor=west,
      xshift=1.4cm,
      width=0.47\linewidth,
      height=5.2cm,
      xlabel={保留秩 $k$},
      ylabel={相对 Frobenius 误差},
      xmin=0, xmax=25,
      ymin=0, ymax=0.63,
      grid=major,
      grid style={dashed, gray!25},
      tick label style={font=\footnotesize},
      label style={font=\small\sffamily},
      title={\small\sffamily\bfseries 重构相对误差衰减曲线}
    ]
      \addplot[color=navy, very thick, mark=triangle*, mark size=2pt] table[x=k, y=rel_err, col sep=comma] {data/svd_compress.csv};
      \node[anchor=west, font=\scriptsize\sffamily, color=navy] at (axis cs:5.5, 0.30) {$k=4\;(0.262)$};
      \draw[-{Stealth[length=1.6mm]}, color=navy, thin] (axis cs:5.4, 0.295) -- (axis cs:4.25, 0.268);
      \node[anchor=west, font=\scriptsize\sffamily, color=navy] at (axis cs:10.0, 0.22) {$k=8\;(0.132)$};
      \draw[-{Stealth[length=1.6mm]}, color=navy, thin] (axis cs:10.0, 0.205) -- (axis cs:8.25, 0.142);
      \node[anchor=west, font=\scriptsize\sffamily, color=navy] at (axis cs:15.5, 0.13) {$k=16\;(0.054)$};
      \draw[-{Stealth[length=1.6mm]}, color=navy, thin] (axis cs:17.2, 0.112) -- (axis cs:16.1, 0.064);
    \end{axis}
  \end{tikzpicture}
  \caption{$64\times 64$ 灰度图像 SVD 压缩的实测权衡曲线（数据源：\texttt{data/svd\_compress.csv}）。左图显示极低秩（$k=4, 8$）就能捕获绝大多数能量；右图显示相对重构误差随 $k$ 迅速衰减。}
  \label{fig:ch4-svd-compress}
\end{figure}

最后，给出工业界标准的纯 NumPy 实现。\textbf{数值工程要点}：实现 PCA 时应当直接对中心化矩阵 $\widetilde{X}$ 调用 \texttt{np.linalg.svd}，而\textbf{不要}先算格拉姆矩阵 $\widetilde{X}^T\widetilde{X}$ 再调 \texttt{eigh}——因为形成 $\widetilde{X}^T\widetilde{X}$ 会将条件数平方（$\cond(\widetilde{X}^T\widetilde{X}) = \cond(\widetilde{X})^2$），在浮点运算下可能彻底淹没较小的奇异值（见习题 4.8）。

\begin{lstlisting}[caption={工业级 \texttt{pca\_via\_svd} 与 \texttt{svd\_compress} 纯 NumPy 实现},label={lst:ch4-pca-svd}]
import numpy as np

def pca_via_svd(X: np.ndarray, k: int):
    """Exact PCA via economy SVD on centered data (avoids squaring cond(X))."""
    mu = X.mean(axis=0, keepdims=True)         # (1, d)
    Xc = X - mu                                # mandatory centering
    U, s, Vt = np.linalg.svd(Xc, full_matrices=False)
    V_k = Vt[:k].T                             # (d, k) principal axes
    Z_k = U[:, :k] * s[:k]                     # (n, k) == Xc @ V_k
    X_hat = mu + Z_k @ V_k.T                   # (n, d) reconstruction
    explained_ratio = (s[:k] ** 2) / np.sum(s ** 2)
    return Z_k, V_k, X_hat, explained_ratio

def svd_compress(img: np.ndarray, k: int):
    """Rank-k optimal Frobenius approximation via truncated SVD."""
    U, s, Vt = np.linalg.svd(img, full_matrices=False)
    img_k = (U[:, :k] * s[:k]) @ Vt[:k, :]     # broadcasting avoids diag(s)
    rel_err = np.linalg.norm(img - img_k, "fro") / np.linalg.norm(img, "fro")
    energy = float(np.sum(s[:k] ** 2) / np.sum(s ** 2))
    assert np.isclose(rel_err ** 2 + energy, 1.0)
    return img_k, rel_err, energy

# Quick self-check on a 2D tilted cloud
X_demo = np.array([[1.0, 4.0], [3.0, 2.0], [5.0, 8.0], [7.0, 6.0]])
Z1, V1, X1_hat, ratio = pca_via_svd(X_demo, k=1)
print(f"PC1 ratio={ratio[0]:.2%}, axis={np.round(V1[:, 0], 4)}")  # 80.00%
\end{lstlisting}

% =========================================================================
\section{本章习题}
\label{sec:ch4-exercises}

\exTier{A. 概念辨析与反例查错}{区分 SVD 与特征分解、看破伪逆与 PCA 的常见误区}

\begin{enumerate}[label=\textbf{4.\arabic*}, leftmargin=2.4em]
  \item \textbf{（SVD、PCA 与伪逆核心概念辨析）} 判断以下四个命题正误；对错误命题给出反例，对正确命题简述理由：
  \begin{enumerate}[label=(\alph*), leftmargin=1.8em]
    \item 对任意实方阵 $A \in \R^{n\times n}$，其奇异值必然等于其特征值的绝对值 $|\lambda_i|$；
    \item 对样本均值显著偏离原点的紧凑点云，若忘记去中心化直接做 SVD，得到的第一右奇异向量通常会偏向样本均值方向而非点云真实的伸展主方向；
    \item 用 \texttt{eigh(Xc.T @ Xc)} 与直接用 \texttt{svd(Xc)} 求 PCA 主轴在精确算术下等价，但在浮点计算中后者数值更稳健，因为前者会将矩阵条件数平方；
    \item 对欠定线性系统 $Ax = b$，只要方程组有解（即 $b \in \col(A)$），$x^+ = A^+ b$ 就必然精确满足 $Ax^+ = b$ 且具有最小欧氏范数。
  \end{enumerate}

  \item \textbf{（伪逆的代数反例）} 对于可逆方阵，恒有 $(AB)^{-1} = B^{-1}A^{-1}$ 与 $AA^{-1} = A^{-1}A = I$。请构造一对简单的 $2\times 2$ 矩阵 $A, B$ 作为反例，验证一般情形下 $(AB)^+ \neq B^+ A^+$；并举出一个简单反例说明 $A A^+$ 一般不等于 $A^+ A$。
\end{enumerate}

\exTier{B. 小例手算与数值直觉}{亲手算出长方形 SVD、最小范数解与完整 2D PCA}

\begin{enumerate}[label=\textbf{4.\arabic*}, leftmargin=2.4em, resume]
  \item \textbf{（$2\times 3$ 长方形矩阵手算 SVD 与伪逆）} 给定矩阵 $A = \begin{pmatrix} 1 & 1 & 0 \\ 0 & 0 & 2 \end{pmatrix} \in \R^{2\times 3}$：
  \begin{enumerate}[label=(\arabic*), leftmargin=1.8em]
    \item 手算 $A$ 的完整 SVD（$U, \Sigma, V^T$）与紧凑 SVD，写出秩 1 最佳近似矩阵 $A_1$ 及其相对 Frobenius 误差 $\|A - A_1\|_F / \|A\|_F$；
    \item 计算 Moore--Penrose 伪逆 $A^+$，并求出欠定方程组 $Ax = \begin{pmatrix}2 \\ 4\end{pmatrix}$ 的最小 $\ell_2$ 范数解 $x^+$；
    \item 写出 $\nul(A)$ 的一组基，直接验证 $x^+ \perp \nul(A)$。
  \end{enumerate}

  \item \textbf{（2D 小样本手算 PCA 全流程）} 给定 4 个二维样本点：$(1,4),\; (3,2),\; (5,8),\; (7,6)$：
  \begin{enumerate}[label=(\arabic*), leftmargin=1.8em]
    \item 计算样本均值 $\mu$ 与去中心化数据矩阵 $\widetilde{X} \in \R^{4\times 2}$；
    \item 计算样本协方差矩阵 $\Sigma_{\mathrm{cov}} = \frac{1}{4}\widetilde{X}^T\widetilde{X}$，手算其特征值与单位正交特征向量；
    \item 将样本降到 1 维，写出主成分坐标 $Z_1 = \widetilde{X} v_1$、重构后的二维坐标 $\widehat{X}_1 = \mathbf{1}\mu^T + Z_1 v_1^T$，以及第一主成分的解释方差比。
  \end{enumerate}
\end{enumerate}

\exTier{C. 实用推导与性质证明}{理解极分解的几何美感与岭回归的谱过滤本质}

\begin{enumerate}[label=\textbf{4.\arabic*}, leftmargin=2.4em, resume]
  \item \textbf{（SVD 的正交不变性与矩阵极分解）} 设任意实矩阵 $A \in \R^{m\times n}$ 的 SVD 为 $A = U\Sigma V^T$。
  \begin{enumerate}[label=(\arabic*), leftmargin=1.8em]
    \item 证明对任意正交矩阵 $Q_1 \in \R^{m\times m}$ 与 $Q_2 \in \R^{n\times n}$，矩阵 $Q_1 A Q_2$ 与 $A$ 具有完全相同的奇异值、算子 2-范数与 Frobenius 范数；
    \item 当 $A \in \R^{n\times n}$ 为方阵时，证明 $A$ 可分解为 $A = Q S$，其中 $Q = U V^T$ 为正交矩阵，$S = V\Sigma V^T \succeq 0$ 为半正定对称矩阵（称为\textbf{极分解}，是复数极坐标 $z = e^{i\theta}r$ 在高维矩阵上的精确推广）。
  \end{enumerate}
  \exHint{第 (1) 问考察 $(Q_1 A Q_2)^T(Q_1 A Q_2)$ 与 $A^T A$ 的正交相似关系；第 (2) 问在 $U\Sigma V^T$ 中间插入 $V^T V = I_n$。}

  \item \textbf{（岭回归的谱过滤与伪逆极限）} 设任意实矩阵 $A = U\Sigma V^T \in \R^{m\times n}$（$\rank(A) = r$）及目标向量 $b \in \R^m$。对任意正则参数 $\lambda > 0$，岭回归目标 $J_\lambda(x) = \|Ax - b\|_2^2 + \lambda\|x\|_2^2$ 的唯一解析解为 $x_\lambda = (A^T A + \lambda I_n)^{-1} A^T b$。利用 SVD 证明：
  \[
    x_\lambda = \sum_{i=1}^r \frac{\sigma_i}{\sigma_i^2 + \lambda}(u_i^T b) v_i,
    \qquad \lim_{\lambda \to 0^+} x_\lambda = A^+ b.
  \]
  并从滤波器因子 $f_\lambda(\sigma_i) = \dfrac{\sigma_i^2}{\sigma_i^2 + \lambda}$ 的角度，解释岭回归如何选择性地抑制极小奇异值上的噪声放大。
  \exHint{将 $A^T A + \lambda I_n = V(\Sigma^T \Sigma + \lambda I_n)V^T$ 代入求逆，再右乘 $A^T = V\Sigma^T U^T$。}
\end{enumerate}

\exTier{D. 编程实战与可自检实验}{用代码亲眼验证隐式偏置与条件数平方精度损失}

\begin{enumerate}[label=\textbf{4.\arabic*}, leftmargin=2.4em, resume]
  \item \textbf{（欠定系统梯度下降隐式偏置实验）} 用固定随机种子生成欠定线性方程组 $A \in \R^{30\times 100}$、$b \in \R^{30}$：
  \begin{enumerate}[label=(\arabic*), leftmargin=1.8em]
    \item 用 \texttt{np.linalg.pinv(A) @ b} 计算最小范数解 $x^+$；
    \item 从零初始化 $x^{(0)} = 0$ 出发，运行梯度下降 $x^{(t+1)} = x^{(t)} - \eta A^T(A x^{(t)} - b)$（取步长 $\eta = 1/\sigma_1(A)^2$，迭代 $3000$ 步）得到 $x_{\mathrm{GD}}$；
    \item 从标准正态随机初值 $x^{(0)} \sim \mathcal{N}(0, I_{100})$ 出发运行相同梯度下降得到 $x_{\mathrm{GD,rand}}$。比较三者的残差 $\|Ax - b\|_2$、范数 $\|x\|_2$，以及 $\|x_{\mathrm{GD}} - x^+\|_2$。
  \end{enumerate}
  \par\smallskip\noindent{\footnotesize\sffamily\color{tealDark}\textbf{【自检标准】}\,三者的残差 $\|Ax - b\|_2$ 均应 $< 10^{-8}$；零初始化解须满足 $\|x_{\mathrm{GD}} - x^+\|_2 < 10^{-7}$；而随机初始化解的范数应严格大于 $\|x^+\|_2$，且差向量 $x_{\mathrm{GD,rand}} - x^+$ 在行空间 $\row(A)$ 上的投影范数 $\|A^+ A(x_{\mathrm{GD,rand}} - x^+)\|_2 < 10^{-7}$。}

  \item \textbf{（图像压缩能量守恒与 \texttt{svd} vs \texttt{eigh} 精度对比）} 编写 Python 脚本完成以下两项验证：
  \begin{enumerate}[label=(\arabic*), leftmargin=1.8em]
    \item 复现本节 $64\times 64$ 合成灰度图像（参考 \texttt{code/vol1\_experiments.py}），对 $k \in \{1,2,4,8,12,16,24\}$ 调用 \texttt{svd\_compress}，验证每组 $k$ 均满足 $|\text{rel\_err}^2 + \text{energy} - 1| < 10^{-12}$；
    \item 构造经典病态矩阵 $X = \begin{pmatrix} 1 & 1 \\ \varepsilon & 0 \\ 0 & \varepsilon \end{pmatrix}$（取 $\varepsilon = 10^{-8}$，其精确奇异值为 $\sigma_1 = \sqrt{2 + \varepsilon^2},\; \sigma_2 = \varepsilon = 10^{-8}$），对比直接调用 \texttt{np.linalg.svd(X)} 与先计算格拉姆矩阵再调用 \texttt{np.linalg.eigh(X.T @ X)} 后对非负特征值开方所求得的最小奇异值 $\sigma_2$。
  \end{enumerate}
  \par\smallskip\noindent{\footnotesize\sffamily\color{tealDark}\textbf{【自检标准】}\,$k=4$ 时累计能量须为 $93.15\%$、相对误差 $0.2617$；对 $\varepsilon = 10^{-8}$ 的 $3\times 2$ 矩阵，\texttt{np.linalg.svd} 求得的 $\sigma_2$ 相对误差应 $< 10^{-14}$，而 \texttt{eigh(X.T @ X)} 会因浮点吸收 $1 + 10^{-16} = 1$ 彻底丢失 $\sigma_2$（输出 $0$ 或严重失真）。}
\end{enumerate}

\chapter[AI 多元微积分与矩阵微分]{AI 多元微积分与矩阵微分}
\label{ch:multivar-matrix-calc}

\chapEpigraph{优秀的符号以最简短的形式直指事物的本质，从而极大地减轻思维的负担。}{戈特弗里德·威廉·莱布尼茨（Gottfried Wilhelm Leibniz）}

\begin{chapIntro}{本章定位}
大一微积分下册花费大量篇幅讲授曲线积分、曲面积分以及格林、高斯、斯托克斯三大矢量场公式；但在现代人工智能中，你几乎永远不会去算一个曲面上的通量。AI 模型的世界是另一套完全不同的微积分图景：模型把数百万乃至数百亿个参数向量与权重矩阵，经过层层线性变换与非线性激活，最终映射为一个\textbf{标量损失函数} $L(\theta)\in\R$。我们需要每天反复使用的是\textbf{梯度与最速下降方向、Jacobian 矩阵与反向模式链式法则、二阶 Taylor 与 Hessian 曲率判别、拉格朗日乘子法}，以及专门为矩阵参数打造的系统化求导武器——\textbf{全微分与迹技巧（Trace Trick）及 Hadamard 逐元素求导}。本章全部工具都是第 6 章梯度下降、第 2 卷逻辑/Softmax 回归以及第 5 卷手写自动微分的直接底座。
\end{chapIntro}

% ====================================================================
\section{偏导、方向导数、梯度与最速下降方向}
\label{sec:ch5-grad}

\begin{recallBox}{一元导数的本质是“局部线性逼近”}
对一元可微函数 $f:\R\to\R$，在点 $x$ 处加微小扰动 $h$ 时，导数 $f'(x)$ 刻画了函数的一阶线性增量：
\[
  f(x + h) = f(x) + f'(x)\,h + o(|h|), \qquad h \to 0.
\]
当 $h$ 足够小时，误差项 $o(|h|)$ 相比 $|h|$ 可以忽略，因此只要 $f'(x)\neq 0$，沿着 $-\operatorname{sign}(f'(x))$ 走一小步就必然让函数值下降。
\end{recallBox}

在人工智能中，参数不再是单个实数 $x\in\R$，而是向量 $x = (x_1,\dots,x_n)^T \in \R^n$（例如线性回归的权重向量 $w\in\R^d$）。如何把“一阶线性逼近”自然升级到高维空间？

\subsection{梯度与方向导数的定义}

\begin{defBox}
\begin{definition}[可微性、梯度与方向导数]\label{def:ch5-grad}
设多元标量函数 $f:\R^n\to\R$ 在点 $x\in\R^n$ 的某个邻域内有定义。
\begin{enumerate}
  \item \textbf{梯度（Gradient）}：将 $f$ 对各个坐标分量的偏导数排成\textbf{列向量}，记作
  \[
    \nabla f(x) \coloneqq
    \begin{bmatrix}
      \dfrac{\partial f}{\partial x_1}(x) &
      \dfrac{\partial f}{\partial x_2}(x) &
      \cdots &
      \dfrac{\partial f}{\partial x_n}(x)
    \end{bmatrix}^T \in \R^n.
  \]
  \item \textbf{全可微（一阶线性逼近）}：若当扰动向量 $h\in\R^n$ 满足 $\norm{h}_2\to 0$ 时，有
  \begin{equation}\label{eq:ch5-first-order}
    f(x + h) = f(x) + \inner{\nabla f(x)}{h} + o(\norm{h}_2),
  \end{equation}
  则称 $f$ 在 $x$ 处\textbf{可微}。由连续偏导数定理，只要各偏导数在 $x$ 处连续，式~\eqref{eq:ch5-first-order} 就自动成立。
  \item \textbf{方向导数（Directional Derivative）}：给定任意单位方向向量 $u\in\R^n$（$\norm{u}_2=1$），函数沿方向 $u$ 的瞬时变化率定义为
  \[
    D_u f(x) \coloneqq \lim_{t\to 0} \frac{f(x + t u) - f(x)}{t}.
  \]
\end{enumerate}
\end{definition}
\end{defBox}

把 $h = t u$（其中 $t\in\R$，$\norm{u}_2=1$）代入式~\eqref{eq:ch5-first-order}，两边减去 $f(x)$ 再同除以 $t$，令 $t\to 0$，立即得到方向导数与梯度的核心关系：
\begin{equation}\label{eq:ch5-dir-deriv}
  D_u f(x) = \inner{\nabla f(x)}{u} = \nabla f(x)^T u.
\end{equation}
在讨论抽象定理之前，我们先通过一个具体的二维碗状函数亲手感受梯度、等高线与方向导数之间的几何关系。

\begin{example}[二维椭圆碗上的梯度与方向导数]\label{ex:ch5-2d-grad}
考虑二元函数 $f(x_1, x_2) = x_1^2 + 3x_2^2$，在点 $x_0 = (1, 1)^T$ 处函数值为 $f(x_0) = 1^2 + 3\cdot 1^2 = 4$。求偏导数得梯度向量：
\[
  \nabla f(x_1, x_2) = \begin{bmatrix} 2x_1 \\ 6x_2 \end{bmatrix}
  \implies
  \nabla f(x_0) = \begin{bmatrix} 2 \\ 6 \end{bmatrix},
  \qquad
  \norm{\nabla f(x_0)}_2 = \sqrt{2^2 + 6^2} = 2\sqrt{10} \approx 6.325.
\]
比较从 $x_0=(1,1)^T$ 出发沿四个不同单位方向 $u$ 的瞬时变化率 $D_u f(x_0) = \inner{\nabla f(x_0)}{u}$：
\begin{itemize}
  \item 沿横轴 $e_1=(1,0)^T$：$D_{e_1} f(x_0) = 2$（正是 $\partial f/\partial x_1$）；
  \item 沿纵轴 $e_2=(0,1)^T$：$D_{e_2} f(x_0) = 6$（正是 $\partial f/\partial x_2$，纵轴比横轴陡 3 倍）；
  \item 沿等高线 $x_1^2 + 3x_2^2 = 4$ 在 $(1,1)^T$ 处的单位切向量 $v = \frac{1}{\sqrt{10}}(3, -1)^T$：
  \[
    D_v f(x_0) = \frac{1}{\sqrt{10}}(2\times 3 + 6\times(-1)) = 0,
  \]
  说明沿等高线切线移动时函数值一阶不变；
  \item 沿负梯度单位方向 $u^* = -\frac{\nabla f(x_0)}{\norm{\nabla f(x_0)}_2} = -\frac{1}{\sqrt{10}}(1, 3)^T$：
  \[
    D_{u^*} f(x_0) = -\frac{1}{\sqrt{10}}(2\times 1 + 6\times 3) = -2\sqrt{10} \approx -6.325,
  \]
  比单沿 $-e_1$（变化率 $-2$）和单沿 $-e_2$（变化率 $-6$）下降得都更剧烈！
\end{itemize}
值得注意的是，点 $x_0=(1,1)^T$ 指向碗底原点 $(0,0)^T$ 的单位方向是 $-\frac{1}{\sqrt{2}}(1,1)^T$，与负梯度方向 $-\frac{1}{\sqrt{10}}(1,3)^T$ \textbf{并不重合}！正如图~\ref{fig:ch5-grad-contour} 所示，负梯度永远严格垂直于当前等高线，当等高线被拉扁成椭圆时，负梯度会偏向陡峭轴，而不是直指极小值点——这正是第 6 章“窄谷锯齿震荡”的几何根源。
\end{example}

\begin{figure}[htbp]
\centering
\begin{tikzpicture}[scale=1.18]
  % Axes
  \draw[-{Stealth[length=2.2mm]}, slate!70, thin] (-2.9,0) -- (3.2,0) node[right, font=\footnotesize, ink] {$x_1$};
  \draw[-{Stealth[length=2.2mm]}, slate!70, thin] (0,-1.85) -- (0,2.35) node[above, font=\footnotesize, ink] {$x_2$};

  % Elliptical contours x1^2 + 3 x2^2 = c
  % Semi-axes: a = sqrt(c), b = sqrt(c/3)
  \draw[tealLine, thick] (0,0) ellipse ({sqrt(1)} and {sqrt(1/3)});
  \draw[tealLine, thick] (0,0) ellipse ({sqrt(2)} and {sqrt(2/3)});
  \draw[teal, very thick] (0,0) ellipse ({2} and {2/sqrt(3)});
  \draw[tealLine, thick, dashed] (0,0) ellipse ({sqrt(7)} and {sqrt(7/3)});

  % Contour labels
  \node[font=\scriptsize, teal, fill=white, inner sep=0.8pt] at (-1.05, 0.40) {$c=1$};
  \node[font=\scriptsize, teal, fill=white, inner sep=0.8pt] at (-2.05, -0.58) {$c=4$};
  \node[font=\scriptsize, slate] at (-2.65, -0.95) {$c=7$};

  % Minimum origin
  \fill[tealDark] (0,0) circle (1.8pt) node[below left, font=\scriptsize, tealDark] {$x^*=(0,0)$};

  % Point x0 = (1, 1) on c=4
  \coordinate (P) at (1,1);
  \fill[ink] (P) circle (2.2pt);
  \node[font=\footnotesize, ink, right=3pt, yshift=4pt] at (P) {$x_0=(1,1)$};

  % Tangent line to x1^2 + 3x2^2 = 4 at (1,1): x1 + 3x2 = 4 => dir (3,-1)
  \draw[navy, dashed, thick] ($(P) + (-1.35, 0.45)$) -- ($(P) + (1.50, -0.50)$);
  \draw[-{Stealth[length=2.3mm]}, navy, very thick] (P) -- ($(P) + (0.96, -0.32)$)
    node[above right, font=\scriptsize, navy, yshift=-1pt] {切方向 $v$};

  % Gradient direction (2,6) -> unit (1/sqrt(10), 3/sqrt(10)) = (0.316, 0.949)
  \draw[-{Stealth[length=2.5mm]}, ochre, very thick] (P) -- ($(P) + (0.41, 1.23)$)
    node[right, font=\footnotesize, ochre] {$\nabla f(x_0)=(2,6)^T$（最速上升）};

  % Steepest descent direction -grad
  \draw[-{Stealth[length=2.5mm]}, tealDark, very thick] (P) -- ($(P) + (-0.41, -1.23)$)
    node[below right, font=\footnotesize, tealDark, fill=white, fill opacity=0.88, text opacity=1, inner sep=1pt, xshift=3pt, yshift=2pt] {$-\nabla f(x_0)$（最速下降）};

  % Right-angle mark between tangent (0.949, -0.316) and gradient (0.316, 0.949)
  \coordinate (A) at ($(P) + (0.18, -0.06)$);
  \coordinate (B) at ($(P) + (0.24, 0.12)$);
  \coordinate (C) at ($(P) + (0.06, 0.18)$);
  \draw[ink, thin] (A) -- (B) -- (C);
\end{tikzpicture}
\caption{函数 $f(x_1,x_2)=x_1^2+3x_2^2$ 在 $x_0=(1,1)^T$ 处的几何图景：梯度 $\nabla f(x_0)$ 严格正交于等高线 $c=4$ 的切方向 $v$ 并指向函数值增大最快的方向；负梯度 $-\nabla f(x_0)$ 是局部最速下降方向，但并不直接指向碗底原点 $(0,0)$。}
\label{fig:ch5-grad-contour}
\end{figure}

\subsection{最速下降方向与等值面正交性的严格证明}

借助第 2 章证明过的 Cauchy-Schwarz 不等式，我们可以把例~\ref{ex:ch5-2d-grad} 中的几何现象提升为任意 $n$ 维空间的普遍定理。

\begin{thmBox}
\begin{theorem}[梯度的最速变化与正交性定理]\label{thm:ch5-steepest}
设 $f:\R^n\to\R$ 在点 $x_0$ 处可微且 $\nabla f(x_0)\neq 0$，则：
\begin{enumerate}
  \item \textbf{局部最速下降与最速上升}：在所有单位方向 $u\in\R^n$（$\norm{u}_2=1$）中，方向导数 $D_u f(x_0)$ 满足
  \[
    -\norm{\nabla f(x_0)}_2 \le D_u f(x_0) \le \norm{\nabla f(x_0)}_2.
  \]
  上界在唯一单位方向 $u_{\max} = \dfrac{\nabla f(x_0)}{\norm{\nabla f(x_0)}_2}$ 处取得；下界在唯一单位方向 $u_{\min} = -\dfrac{\nabla f(x_0)}{\norm{\nabla f(x_0)}_2}$ 处取得。因此\textbf{负梯度方向是唯一的一阶最速局部下降方向}。
  \item \textbf{正交于等值面}：设 $S_c = \{x\in\R^n : f(x) = c\}$ 是过 $x_0$ 的等值面（$c = f(x_0)$）。对 $S_c$ 上任意过 $x_0$ 的光滑曲线 $\gamma:(-\delta,\delta)\to S_c$（$\gamma(0)=x_0$），梯度与曲线切向量严格正交：
  \[
    \inner{\nabla f(x_0)}{\gamma'(0)} = 0.
  \]
\end{enumerate}
\end{theorem}
\end{thmBox}

\begin{proof}
\textbf{(1)} 对任意单位向量 $\norm{u}_2=1$，由式~\eqref{eq:ch5-dir-deriv} 及第 2 章的 Cauchy-Schwarz 不等式得
\[
  \abs{D_u f(x_0)} = \abs{\inner{\nabla f(x_0)}{u}} \le \norm{\nabla f(x_0)}_2 \norm{u}_2 = \norm{\nabla f(x_0)}_2,
\]
即 $-\norm{\nabla f(x_0)}_2 \le D_u f(x_0) \le \norm{\nabla f(x_0)}_2$。回忆 Cauchy-Schwarz 等号成立当且仅当 $u$ 与 $\nabla f(x_0)$ 线性相关（共线），即存在实数 $\alpha\in\R$ 使 $u = \alpha \nabla f(x_0)$。由 $\norm{u}_2 = 1$ 得 $|\alpha| = 1/\norm{\nabla f(x_0)}_2$：
\begin{itemize}
  \item 取 $\alpha = +1/\norm{\nabla f(x_0)}_2$ 时，$D_u f(x_0) = +\norm{\nabla f(x_0)}_2$，达到唯一最大值；
  \item 取 $\alpha = -1/\norm{\nabla f(x_0)}_2$ 时，$D_u f(x_0) = -\norm{\nabla f(x_0)}_2$，达到唯一最小值。
\end{itemize}

\textbf{(2)} 由于曲线 $\gamma(t)$ 始终落在等值面 $S_c$ 上，复合函数 $h(t) \coloneqq f(\gamma(t)) \equiv c$ 为常数，故 $h'(0) = 0$。另一方面，将 $\gamma(t) = x_0 + t\gamma'(0) + o(|t|)$ 代入可微性展开式~\eqref{eq:ch5-first-order}：
\[
  0 = \lim_{t\to 0}\frac{f(\gamma(t)) - f(x_0)}{t}
    = \lim_{t\to 0}\frac{\inner{\nabla f(x_0)}{t\gamma'(0) + o(|t|)} + o(|t|)}{t}
    = \inner{\nabla f(x_0)}{\gamma'(0)}.
\]
因此 $\nabla f(x_0)$ 正交于等值面在 $x_0$ 处的每一条切向量。
\end{proof}

\begin{insightBox}{“最速下降”是局部性质，不是全局承诺}
定理~\ref{thm:ch5-steepest} 仅在步长 $\eta\to 0^+$ 的无穷小邻域内保证了 $f(x_0 - \eta \nabla f(x_0)) < f(x_0)$。如果步长 $\eta$ 迈得太大，高阶项 $o(\norm{h}_2)$ 会接管一切，导致“一步跨过山谷冲到更高的对岸”甚至发散。第 6 章我们将用特征值精确算出这个允许步长的临界上限 $\eta < 2/\lambda_{\max}$。
\end{insightBox}

% ====================================================================
\section{Jacobian 矩阵与多元链式法则：反向传播的微积分骨架}
\label{sec:ch5-jacobian}

上一节讨论的是“向量输入 $\to$ 标量输出”（$f:\R^n\to\R$）。但在神经网络中，每一层都在做\textbf{向量到向量}的变换 $f:\R^n\to\R^m$（例如线性层 $z = Wx+b$、激活层 $a = \sigma(z)$）。此时一阶线性逼近的系数不再是向量，而是一个 $m\times n$ 矩阵——\textbf{Jacobian 矩阵}。

\subsection{先看具体二维小例，再抽象定义 Jacobian}

\begin{example}[线性层与 Sigmoid 激活层的手算偏导矩阵]\label{ex:ch5-2d-jac}
考虑两层极简网络的前向计算：输入 $x = (x_1, x_2)^T \in \R^2$，权重 $W = \begin{bmatrix} 1 & -2 \\ 0.5 & 1 \end{bmatrix}$，偏置 $b = \begin{bmatrix} 0 \\ -1 \end{bmatrix}$，
\[
  z = f(x) = W x + b = \begin{bmatrix} x_1 - 2x_2 \\ 0.5 x_1 + x_2 - 1 \end{bmatrix},
  \qquad
  a = g(z) = \sigma(z) = \begin{bmatrix} \sigma(z_1) \\ \sigma(z_2) \end{bmatrix},
\]
其中 $\sigma(t) = \frac{1}{1+e^{-t}}$，其导数为 $\sigma'(t) = \sigma(t)(1-\sigma(t))$。把输出各分量对输入各分量的偏导数排成 $2\times 2$ 矩阵（第 $i$ 行对应第 $i$ 个输出，第 $j$ 列对应第 $j$ 个输入）：
\begin{enumerate}
  \item 对线性层 $z = Wx + b$：
  \[
    \frac{\partial z}{\partial x} =
    \begin{bmatrix}
      \partial z_1/\partial x_1 & \partial z_1/\partial x_2 \\
      \partial z_2/\partial x_1 & \partial z_2/\partial x_2
    \end{bmatrix}
    = \begin{bmatrix} 1 & -2 \\ 0.5 & 1 \end{bmatrix} = W.
  \]
  \item 对逐元素激活层 $a = \sigma(z)$：由于 $a_1$ 只依赖 $z_1$、$a_2$ 只依赖 $z_2$，非对角偏导全为 $0$，得到对角矩阵：
  \[
    \frac{\partial a}{\partial z} =
    \begin{bmatrix}
      \sigma'(z_1) & 0 \\
      0 & \sigma'(z_2)
    \end{bmatrix}
    = \diag(\sigma'(z_1), \sigma'(z_2)).
  \]
\end{enumerate}
代入具体点 $x = (2, 1)^T$，算得 $z = (0, 1)^T$，从而 $\sigma(0) = 0.5$，$\sigma'(0) = 0.25$；$\sigma(1)\approx 0.7311$，$\sigma'(1)\approx 0.1966$。若想知道输出 $a$ 对输入 $x$ 的敏感度，直接把两个矩阵相乘：
\[
  \frac{\partial a}{\partial x} = \frac{\partial a}{\partial z}\frac{\partial z}{\partial x}
  = \begin{bmatrix} 0.25 & 0 \\ 0 & 0.1966 \end{bmatrix}
    \begin{bmatrix} 1 & -2 \\ 0.5 & 1 \end{bmatrix}
  = \begin{bmatrix} 0.25 & -0.50 \\ 0.0983 & 0.1966 \end{bmatrix}.
\]
\end{example}

上面的矩阵乘积并非巧合，而是多元微积分最核心的\textbf{Jacobian 链式法则}。

\begin{defBox}
\begin{definition}[Jacobian 矩阵]\label{def:ch5-jacobian}
设向量值函数 $f:\R^n\to\R^m$ 的分量形式为 $f(x) = (f_1(x), \dots, f_m(x))^T$。其在点 $x\in\R^n$ 处的 \textbf{Jacobian 矩阵} $J_f(x) \in \R^{m\times n}$ 定义为
\[
  \big(J_f(x)\big)_{ij} \coloneqq \frac{\partial f_i}{\partial x_j}(x),
  \qquad i=1,\dots,m,\; j=1,\dots,n,
\]
即第 $i$ 行恰好是第 $i$ 个分量函数 $f_i$ 的梯度行向量 $\nabla f_i(x)^T$。若当 $\norm{h}_2\to 0$ 时满足
\begin{equation}\label{eq:ch5-jac-linear}
  f(x + h) = f(x) + J_f(x)\,h + o(\norm{h}_2),
\end{equation}
则称 $f$ 在 $x$ 处可微。特别地，当 $m=1$（标量函数）时，$J_f(x) = \nabla f(x)^T \in \R^{1\times n}$ 是行向量，而梯度 $\nabla f(x) \in \R^{n\times 1}$ 是列向量。
\end{definition}
\end{defBox}

\begin{thmBox}
\begin{theorem}[多元链式法则]\label{thm:ch5-chain}
设 $f:\R^n\to\R^m$ 在 $x$ 处可微，$g:\R^m\to\R^p$ 在 $y = f(x)$ 处可微，则复合映射 $h \coloneqq g \circ f : \R^n\to\R^p$ 在 $x$ 处可微，且其 Jacobian 矩阵等于两层 Jacobian 矩阵的乘积：
\begin{equation}\label{eq:ch5-chain}
  J_{g\circ f}(x) = J_g\big(f(x)\big)\, J_f(x) \in \R^{p\times n}.
\end{equation}
\end{theorem}
\end{thmBox}

\begin{proof}
记 $y = f(x)$，对扰动 $h\in\R^n$ 记 $\Delta y = f(x+h) - f(x)$。由 $f$ 的可微性~\eqref{eq:ch5-jac-linear}：
\[
  \Delta y = J_f(x) h + r_f(h), \qquad \lim_{h\to 0}\frac{\norm{r_f(h)}_2}{\norm{h}_2} = 0.
\]
由于矩阵算子范数有界，存在常数 $M>0$ 使当 $\norm{h}_2$ 充分小时 $\norm{\Delta y}_2 \le M\norm{h}_2$。再将 $\Delta y$ 代入 $g$ 在 $y$ 处的可微展开式 $g(y+\Delta y) = g(y) + J_g(y)\Delta y + r_g(\Delta y)$（其中 $\norm{r_g(\Delta y)}_2 = o(\norm{\Delta y}_2)$）：
\[
  g(f(x+h)) - g(f(x))
  = J_g(y)\big(J_f(x)h + r_f(h)\big) + r_g(\Delta y)
  = \big(J_g(y) J_f(x)\big) h + \underbrace{J_g(y) r_f(h) + r_g(\Delta y)}_{R(h)}.
\]
最后检验余项：$\norm{J_g(y) r_f(h)}_2 / \norm{h}_2 \le \norm{J_g(y)}_2 \cdot \norm{r_f(h)}_2/\norm{h}_2 \to 0$；且由 $\norm{\Delta y}_2 \le M\norm{h}_2$，亦有 $\norm{r_g(\Delta y)}_2 / \norm{h}_2 \to 0$。故 $R(h) = o(\norm{h}_2)$，对比式~\eqref{eq:ch5-jac-linear} 的唯一性即得 $J_{g\circ f}(x) = J_g(f(x)) J_f(x)$。
\end{proof}

\subsection{为什么 AI 选反向模式（VJP）而非前向模式？}

定理~\ref{thm:ch5-chain} 表面上看只是矩阵乘法，但它隐藏了深度学习最关键的计算复杂度分野：\textbf{矩阵乘法满足结合律，但不同加括号顺序的计算代价相差百万倍！}

考虑一个 3 阶段标量损失复合链条（输入维数 $n$、隐藏维数 $d$、输出标量损失维数 $1$）：
\[
  x \in \R^n \xrightarrow{\;f\;} z \in \R^d \xrightarrow{\;g\;} a \in \R^d \xrightarrow{\;h\;} L \in \R.
\]
由链式法则~\eqref{eq:ch5-chain}，总损失对输入 $x$ 的 Jacobian 是一个 $1\times n$ 行向量：
\[
  J_L(x) = \nabla_x L^T =
  \underbrace{J_h(a)}_{1\times d}\;
  \underbrace{J_g(z)}_{d\times d}\;
  \underbrace{J_f(x)}_{d\times n}.
\]
对比两种计算加括号顺序（设 $n = d = 1000$）：
\begin{enumerate}
  \item \textbf{前向模式（Forward-mode，从输入往输出推，算式从右向左结合）}：
  \[
    J_h(a)\,\Big(\underbrace{J_g(z)\,J_f(x)}_{d\times n}\Big).
  \]
  第一步就要计算 $(d\times d)(d\times n)$ 的\textbf{矩阵乘矩阵}，需要 $O(d^2 n)$ 次浮点运算（当 $n=d=1000$ 时为 $10^9$ 次乘法！）并需要显式存储 $d\times n$ 的中间大矩阵。若输入是百万级参数，前向模式彻底不可行。
  \item \textbf{反向模式（Reverse-mode / 反向传播，从输出 $L$ 往输入 $x$ 推，算式从左向右结合）}：
  \[
    \Big(\underbrace{J_h(a)\,J_g(z)}_{1\times d}\Big)\,J_f(x).
  \]
  因为最左端 $J_h(a) = \nabla_a L^T$ 是\textbf{行向量}，每一步都是行向量乘 Jacobian 矩阵——即\textbf{向量-Jacobian 积（VJP）}！第一步 $(1\times d)(d\times d)$ 仅需 $O(d^2)$（对逐元素激活甚至仅需 $O(d)$），第二步 $(1\times d)(d\times n)$ 仅需 $O(dn)$，总代价仅 $O(d^2 + dn)$（约 $2\times 10^6$ 次乘法），与一次前向传播同量级！
\end{enumerate}

\begin{figure}[htbp]
\centering
\resizebox{0.96\linewidth}{!}{%
\begin{tikzpicture}[
  node distance=16mm,
  var/.style={rectangle, rounded corners=1.5mm, draw=teal, thick, fill=tealLight,
    minimum width=16mm, minimum height=9mm, font=\sffamily\small, text=ink, align=center},
  loss/.style={rectangle, rounded corners=1.5mm, draw=ochre, very thick, fill=ochreLight,
    minimum width=16mm, minimum height=9mm, font=\sffamily\small\bfseries, text=ochre, align=center},
  fwd/.style={-{Stealth[length=2.4mm]}, thick, teal},
  bwd/.style={-{Stealth[length=2.4mm]}, very thick, ochre, dashed}
]
  \node[var] (x) {$x \in \R^n$};
  \node[var, right=of x] (z) {$z \in \R^d$};
  \node[var, right=of z] (a) {$a \in \R^d$};
  \node[loss, right=of a] (L) {$L \in \R$};

  % Forward pass
  \draw[fwd] (x) -- node[above, font=\scriptsize, tealDark] {$f$} (z);
  \draw[fwd] (z) -- node[above, font=\scriptsize, tealDark] {$g$} (a);
  \draw[fwd] (a) -- node[above, font=\scriptsize, tealDark] {$h$} (L);

  % Backward VJP pass (underneath)
  \draw[bwd] ([yshift=-2.5mm]L.west) to[bend left=28]
    node[below, font=\scriptsize, ochre, align=center] {$g_a^T = J_h(a)$\\$(1\times d)$} ([yshift=-2.5mm]a.east);
  \draw[bwd] ([yshift=-2.5mm]a.west) to[bend left=28]
    node[below, font=\scriptsize, ochre, align=center] {$g_z^T = g_a^T J_g(z)$\\VJP: $(1\times d)$} ([yshift=-2.5mm]z.east);
  \draw[bwd] ([yshift=-2.5mm]z.west) to[bend left=28]
    node[below, font=\scriptsize, ochre, align=center] {$g_x^T = g_z^T J_f(x)$\\VJP: $(1\times n)$} ([yshift=-2.5mm]x.east);

  \node[font=\footnotesize\sffamily\bfseries, teal] at ($(x)!0.5!(L) + (0, 9mm)$) {前向计算（缓存中间值 $x,z,a$）};
  \node[font=\footnotesize\sffamily\bfseries, ochre] at ($(x)!0.5!(L) + (0, -18mm)$) {反向模式 VJP（从标量 $L$ 向左回传，全程保持行向量，免去矩阵乘矩阵）};
\end{tikzpicture}}
\caption{3 阶段标量损失复合计算图与反向模式（VJP）信号回传：只要输出维度远小于参数维度（$1 \ll n$），从右向左逐层做向量-Jacobian 积就能以与前向同量级的算力一次性拿到所有参数的梯度。}
\label{fig:ch5-vjp-chain}
\end{figure}

\begin{insightBox}{直通第 5 卷自动微分（Autograd）}
图~\ref{fig:ch5-vjp-chain} 实际上已经画出了 PyTorch 的 \texttt{loss.backward()} 与第 5 卷 Micro-Autograd 的全部数学内核：前向计算时顺手把每层的局部输入缓存下来；反向传播时从 $\partial L/\partial L = 1$ 出发，每经过一个算子，只需调用该算子的\textbf{局部 VJP 规则}（把上游传回的梯度向量乘上本层局部 Jacobian），连完整的 $d\times d$ Jacobian 矩阵都无需在内存里显式构造（例如逐元素激活层只需一次逐元素乘法）。
\end{insightBox}

% ====================================================================
\section{二阶 Taylor 展开、Hessian 矩阵与拉格朗日乘子法}
\label{sec:ch5-hessian-lagrange}

一阶梯度告诉我们局部往哪走下降最快，但它无法分辨“山谷最低点”与“山脊鞍点”（两者都满足 $\nabla f(x^*) = 0$）。要看清函数在驻点附近的弯曲形状，需要把微积分推进到二阶；而当参数还必须满足等式约束（如第 4 章 PCA 中的单位范数约束 $v^T v = 1$）时，则需要拉格朗日乘子法。

\subsection{多元二阶 Taylor 展开与 Hessian 矩阵}

\begin{recallBox}{一元二阶 Taylor 展开与极值判别}
若一元函数 $f$ 二阶连续可微，则 $f(x+\Delta x) = f(x) + f'(x)\Delta x + \frac{1}{2}f''(x)(\Delta x)^2 + o((\Delta x)^2)$。在驻点 $f'(x^*)=0$ 处：$f''(x^*)>0$ 意味着向上弯曲（严格局部极小值），$f''(x^*)<0$ 意味着向下弯曲（严格局部极大值）。
\end{recallBox}

为了把一元二阶展开自然搬到多元函数 $f:\R^n\to\R$（假设二阶偏导数连续，记作 $f\in C^2$），只需固定点 $x$ 与扰动 $\Delta x$，构造沿直线的\textbf{一元切片函数} $\phi(t) \coloneqq f(x + t\Delta x)$（$t\in[0,1]$）。利用链式法则对 $t$ 求一阶导与二阶导：
\[
  \phi'(t) = \sum_{i=1}^n \frac{\partial f}{\partial x_i}(x + t\Delta x)\,\Delta x_i
           = \nabla f(x+t\Delta x)^T \Delta x,
\]
\[
  \phi''(t) = \sum_{i=1}^n \sum_{j=1}^n \frac{\partial^2 f}{\partial x_i \partial x_j}(x + t\Delta x)\,\Delta x_i \Delta x_j
            = \Delta x^T H_f(x+t\Delta x)\,\Delta x,
\]
其中 $H_f(x) \in \R^{n\times n}$ 称为 $f$ 在 $x$ 处的 \textbf{Hessian 矩阵}（即梯度映射的 Jacobian 矩阵）：
\[
  \big(H_f(x)\big)_{ij} \coloneqq \frac{\partial^2 f}{\partial x_i \partial x_j}(x), \qquad 1 \le i,j \le n.
\]
对一元函数 $\phi(t)$ 在 $t=0$ 处做二阶 Taylor 展开并取 $t=1$，便得到\textbf{多元二阶 Taylor 展开}：
\begin{equation}\label{eq:ch5-taylor2}
  f(x + \Delta x) = f(x) + \nabla f(x)^T \Delta x + \frac{1}{2}\,\Delta x^T H_f(x)\,\Delta x + o(\norm{\Delta x}_2^2).
\end{equation}

关键衔接来了：由微积分的混合偏导数对称性（Schwarz 定理），只要 $f\in C^2$，就有 $\frac{\partial^2 f}{\partial x_i \partial x_j} = \frac{\partial^2 f}{\partial x_j \partial x_i}$，因此 \textbf{Hessian 矩阵 $H_f(x)$ 必然是实对称矩阵}！这意味着第 3 章的\textbf{实对称矩阵谱定理与正定/半正定二次型几何分类}立刻无缝生效：

\begin{thmBox}
\begin{proposition}[驻点处的二阶充分条件与鞍点判别]\label{prop:ch5-second-order}
设 $f\in C^2(\R^n)$，且 $x^*$ 为驻点（$\nabla f(x^*) = 0$）。记 $H = H_f(x^*)$ 的特征值为 $\lambda_1 \ge \dots \ge \lambda_n$：
\begin{enumerate}
  \item 若 $H \succ 0$（所有 $\lambda_i > 0$，向上开口的碗），则 $x^*$ 是 $f$ 的\textbf{严格局部极小值点}；
  \item 若 $H \prec 0$（所有 $\lambda_i < 0$，向下倒扣的碗），则 $x^*$ 是 $f$ 的\textbf{严格局部极大值点}；
  \item 若 $H$ \textbf{不定}（既有 $\lambda_{\max}>0$ 又有 $\lambda_{\min}<0$），则 $x^*$ 是\textbf{鞍点}：沿特征向量 $v_1$ 方向函数局部上升，沿特征向量 $v_n$ 方向函数局部下降。
\end{enumerate}
\end{proposition}
\end{thmBox}

\begin{proof}
因 $\nabla f(x^*) = 0$，式~\eqref{eq:ch5-taylor2} 化为 $f(x^*+\Delta x) - f(x^*) = \frac{1}{2}\Delta x^T H \Delta x + o(\norm{\Delta x}_2^2)$。
\begin{itemize}
  \item 若 $H \succ 0$，由第 3 章 Rayleigh 商界得 $\Delta x^T H \Delta x \ge \lambda_{\min}\norm{\Delta x}_2^2$（$\lambda_{\min}>0$）。当 $\norm{\Delta x}_2$ 充分小时，余项 $|o(\norm{\Delta x}_2^2)| \le \frac{1}{4}\lambda_{\min}\norm{\Delta x}_2^2$，故对所有足够小的非零 $\Delta x$ 均有 $f(x^*+\Delta x) - f(x^*) \ge \frac{1}{4}\lambda_{\min}\norm{\Delta x}_2^2 > 0$，即严格局部极小值。$H\prec 0$ 同理。
  \item 若 $H$ 不定，取 $\lambda_1>0$ 对应的单位特征向量 $v_1$ 与 $\lambda_n<0$ 对应的单位特征向量 $v_n$。令 $\Delta x = t v_1$ 得 $f(x^*+t v_1) - f(x^*) = \frac{1}{2}\lambda_1 t^2 + o(t^2) > 0$；令 $\Delta x = t v_n$ 得 $f(x^*+t v_n) - f(x^*) = \frac{1}{2}\lambda_n t^2 + o(t^2) < 0$，故在任意小邻域内既非极大也非极小（鞍点）。
\end{itemize}
\end{proof}

\begin{example}[二维极小值与鞍点的手算辨析]\label{ex:ch5-min-vs-saddle}
比较两个在原点 $(0,0)^T$ 处均满足 $\nabla f(0,0) = 0$ 的二元函数：
\begin{enumerate}
  \item $f_1(x_1, x_2) = x_1^2 + 3x_2^2 + 2x_1 x_2$：
  \[
    H_1 = \begin{bmatrix} 2 & 2 \\ 2 & 6 \end{bmatrix}.
  \]
  迹 $\tr(H_1) = \lambda_1 + \lambda_2 = 8 > 0$，行列式 $\det(H_1) = \lambda_1\lambda_2 = 12 - 4 = 8 > 0$，故两特征值均严格为正（$\lambda = 4 \pm 2\sqrt{2} > 0$），$H_1 \succ 0$，原点是\textbf{严格局部（且全局）极小值}。
  \item $f_2(x_1, x_2) = x_1^2 - x_2^2 + 4x_1 x_2$：
  \[
    H_2 = \begin{bmatrix} 2 & 4 \\ 4 & -2 \end{bmatrix}.
  \]
  行列式 $\det(H_2) = -4 - 16 = -20 < 0$，两特征值一正一负（$\lambda = \pm 2\sqrt{5}$），$H_2$ 不定，原点是典型的\textbf{鞍点}。
\end{enumerate}
\end{example}

\subsection{等式约束优化与拉格朗日乘子法}

现在给优化问题加上一条等式约束：
\begin{equation}\label{eq:ch5-eq-opt}
  \min_{x\in\R^n}\; f(x) \qquad \text{subject to}\qquad g(x) = 0,
\end{equation}
其中 $f,g:\R^n\to\R$ 连续可微。在无约束优化中，极值点必须满足 $\nabla f(x^*) = 0$；但有了约束 $g(x)=0$ 后，可行点被“锁”在约束曲面 $S = \{x : g(x) = 0\}$ 上，$\nabla f(x^*)$ 通常\textbf{不再等于零}——我们只能沿着约束曲面移动，无法穿透约束墙去追寻更低的点。

那么极值点 $x^*$ 满足什么几何条件？想象你正沿着山坡上的一条狭窄栈道 $g(x)=0$ 行走，试图寻找栈道上的最低点 $x^*$：
\begin{itemize}
  \item 沿栈道经过 $x^*$ 的任意一条光滑路径 $\gamma(t)$（满足 $g(\gamma(t))\equiv 0$，$\gamma(0)=x^*$），由定理~\ref{thm:ch5-steepest}(2) 知其切向量必须垂直于约束梯度：$\inner{\nabla g(x^*)}{\gamma'(0)} = 0$；
  \item 同时，由于 $x^*$ 是栈道上的极值点，高度函数 $f(\gamma(t))$ 在 $t=0$ 处的导数也必须为零，即 $\inner{\nabla f(x^*)}{\gamma'(0)} = 0$。
\end{itemize}
既然 $\nabla f(x^*)$ 与约束曲面的\textbf{每一个切方向}都正交（没有任何沿栈道滑动的分量），那么只要 $\nabla g(x^*)\neq 0$，它必然与约束曲面的法向量 $\nabla g(x^*)$ \textbf{完全共线}（见图~\ref{fig:ch5-lagrange}）：
\begin{equation}\label{eq:ch5-lagrange-cond}
  \exists\,\lambda\in\R, \qquad \nabla f(x^*) = \lambda\,\nabla g(x^*).
\end{equation}
为了把式~\eqref{eq:ch5-lagrange-cond} 和约束 $g(x^*)=0$ 统一写成一个无约束驻点方程，我们引入\textbf{拉格朗日函数（Lagrangian）}：
\[
  \mathcal{L}(x, \lambda) \coloneqq f(x) - \lambda\, g(x),
\]
那么必要条件立刻变成极其整洁的方程组：
\[
  \nabla_x \mathcal{L}(x^*, \lambda^*) = \nabla f(x^*) - \lambda^*\nabla g(x^*) = 0,
  \qquad
  \frac{\partial \mathcal{L}}{\partial \lambda}(x^*, \lambda^*) = -g(x^*) = 0.
\]
若有 $m$ 条独立等式约束 $g_1(x)=0,\dots,g_m(x)=0$（且各 $\nabla g_k(x^*)$ 线性无关），则只需为每条约束配一个乘子 $\lambda_k$，构造 $\mathcal{L}(x, \boldsymbol{\lambda}) = f(x) - \sum_{k=1}^m \lambda_k g_k(x)$ 并令 $\nabla_x \mathcal{L}=0, \nabla_{\boldsymbol{\lambda}}\mathcal{L}=0$。

\begin{figure}[htbp]
\centering
\begin{tikzpicture}[scale=1.08]
  % Axes
  \draw[-{Stealth[length=2.2mm]}, slate!70, thin] (-0.4,0) -- (5.2,0) node[right, font=\footnotesize, ink] {$x_1$};
  \draw[-{Stealth[length=2.2mm]}, slate!70, thin] (0,-0.4) -- (0,3.6) node[above, font=\footnotesize, ink] {$x_2$};

  % Objective contours centered at (4.2, 2.7)
  \coordinate (M) at (4.2, 2.7);
  \fill[tealDark] (M) circle (1.6pt) node[above right, font=\scriptsize, tealDark] {无约束极小点};
  \draw[tealLine, thick] (M) circle (0.75);
  \draw[tealLine, thick] (M) circle (1.40);
  \draw[teal, very thick] (M) circle (2.05);
  \draw[tealLine, thick, dashed] (M) circle (2.65);

  % Tangent point x* on circle radius 2.05 at angle 220 deg from M
  % cos(220) = -0.766, sin(220) = -0.643 -> M + 2.05*(-0.766, -0.643) = (2.63, 1.38)
  \coordinate (X) at (2.63, 1.38);

  % Constraint curve g(x)=0 passing through X smoothly tangent to circle
  % Normal direction from M to X is (-0.766, -0.643); tangent is (0.643, -0.766)
  \draw[navy, very thick] (1.15, 3.30) .. controls (1.90, 2.25) and (2.25, 1.83) .. (X)
    .. controls (3.00, 0.93) and (3.65, 0.40) .. (4.70, 0.15)
    node[above right, font=\footnotesize, navy] {约束曲线 $g(x)=0$};

  % Common tangent line at X
  \draw[slate, dashed, thin] ($(X) + (-0.95, 1.13)$) -- ($(X) + (1.05, -1.25)$);

  % Point x*
  \fill[ochre] (X) circle (2.5pt);
  \node[font=\footnotesize\bfseries, ochre, left=5pt, yshift=-2pt, fill=white, inner sep=0.8pt] at (X) {$x^*$};

  % Collinear gradients along normal (-0.766, -0.643)
  \draw[-{Stealth[length=2.4mm]}, tealDark, very thick] (X) -- ($(X) + (-1.05, -0.88)$)
    node[left, font=\footnotesize, tealDark] {$\nabla f(x^*)$};
  \draw[-{Stealth[length=2.4mm]}, ochre, very thick] (X) -- ($(X) + (0.88, 0.74)$)
    node[above left, font=\footnotesize, ochre, fill=white, fill opacity=0.88, text opacity=1, inner sep=1pt, xshift=2pt] {$\nabla g(x^*)$};

  % Contour labels
  \node[font=\scriptsize, teal] at (4.85, 1.02) {$f(x)=c^*$};
  \node[font=\scriptsize, slate] at (4.2, 3.22) {$f$ 减小};
\end{tikzpicture}
\caption{拉格朗日乘子法的几何本质：在约束极值点 $x^*$ 处，目标函数等高线 $f(x)=c^*$ 与约束曲线 $g(x)=0$ 必然相切，从而两者的法向量 $\nabla f(x^*)$ 与 $\nabla g(x^*)$ 共线，即 $\nabla f(x^*) = \lambda \nabla g(x^*)$。}
\label{fig:ch5-lagrange}
\end{figure}

现在，让我们用拉格朗日乘子法给出第 4 章 PCA 第一主成分方向的“四行惊艳重证”——体会多元微积分与线性代数合流时的简洁之美！

\begin{example}[四行重推第 4 章 PCA 第一主成分方向]\label{ex:ch5-pca-lagrange}
设中心化数据的协方差矩阵为实对称半正定矩阵 $C = \frac{1}{N}X^T X \in \R^{d\times d}$。PCA 寻找投影方差最大的单位方向：
\[
  \max_{v\in\R^d}\; v^T C v \qquad \text{subject to}\qquad v^T v = 1.
\]
用拉格朗日乘子法仅需四步：
\begin{enumerate}
  \item 构造拉格朗日函数：$\mathcal{L}(v, \lambda) = v^T C v - \lambda(v^T v - 1)$；
  \item 对 $v$ 求梯度并令其为零（下节将系统证明对称阵二次型梯度 $\nabla_v(v^T C v) = 2Cv$）：
  \[
    \nabla_v \mathcal{L}(v, \lambda) = 2 C v - 2 \lambda v = 0 \implies C v = \lambda v;
  \]
  \item 把 $Cv = \lambda v$ 与约束 $v^T v = 1$ 代回目标函数：
  \[
    v^T C v = v^T (\lambda v) = \lambda (v^T v) = \lambda;
  \]
  \item 因此，约束驻点必然是 $C$ 的单位特征向量，且对应的目标值恰为特征值 $\lambda$；要使方差最大，必须取最大特征值 $\lambda_{\max}(C)$ 对应的单位特征向量！
\end{enumerate}
\end{example}

\begin{insightBox}{从等式约束到第 4 卷的 KKT 条件}
本节的拉格朗日乘子法完美处理了\textbf{等式约束} $g(x)=0$。当问题出现\textbf{不等式约束} $h(x)\le 0$（例如第 4 卷支持向量机 SVM 的最大间隔约束 $y_i(w^Tx_i+b)\ge 1$）时，拉格朗日乘子必须满足非负性 $\mu\ge 0$ 与互补松弛条件 $\mu h(x^*)=0$——这就是第 4 卷凸优化将系统展开的 \textbf{KKT（Karush-Kuhn-Tucker）条件}。
\end{insightBox}

% ====================================================================
\section{矩阵微积分的系统化方法：全微分与迹技巧}
\label{sec:ch5-matrix-calc}

面对标量损失 $f(W)$ 对整个权重矩阵 $W\in\R^{m\times n}$ 求导，许多初学者试图拆成下标 $\sum_{i,j,k}$ 逐项求偏导——不仅极易写错转置和下标，面对两层以上的矩阵乘法更是寸步难行。本节建立一套\textbf{完全免下标、几乎不可能算错}的系统化矩阵求导方法：\textbf{分母布局 + 全微分与迹技巧（Trace Trick）+ Hadamard 逐元素求导}。

\subsection{布局约定与黄金桥梁定理}

\begin{defBox}
\begin{definition}[矩阵梯度与分母布局约定]\label{def:ch5-denom-layout}
对标量函数 $f:\R^{m\times n}\to\R$，定义其对矩阵 $X\in\R^{m\times n}$ 的\textbf{梯度矩阵} $\nabla_X f(X) \in \R^{m\times n}$ 为逐元素偏导组成的同形状矩阵：
\[
  \big(\nabla_X f(X)\big)_{ij} \coloneqq \frac{\partial f}{\partial X_{ij}}(X),
  \qquad 1\le i\le m,\; 1\le j\le n.
\]
这一约定称为\textbf{分母布局（Denominator Layout）}，其最大优势是\textbf{梯度与参数永远保持完全相同的形状}：$\nabla_x f \in \R^n$ 与 $x\in\R^n$ 同形，$\nabla_W f \in \R^{m\times n}$ 与 $W\in\R^{m\times n}$ 同形，可直接写梯度下降 $W \leftarrow W - \eta \nabla_W f$，且与 PyTorch 中 \texttt{W.grad.shape == W.shape} 完全一致。
\end{definition}
\end{defBox}

如何不拆下标就能直接读出 $\nabla_X f(X)$？关键在于把全微分 $df$ 写成矩阵的 Frobenius 内积 $\inner{G}{dX}_F = \tr(G^T dX)$ 形式！

\begin{thmBox}
\begin{theorem}[矩阵求导的黄金桥梁：全微分-迹等价定理]\label{thm:ch5-golden-bridge}
设标量函数 $f:\R^{m\times n}\to\R$ 在 $X$ 处可微。若对任意微小矩阵扰动 $dX\in\R^{m\times n}$，函数的一阶全微分可化为
\begin{equation}\label{eq:ch5-bridge}
  df = \tr\!\big(G^T\, dX\big),
\end{equation}
其中矩阵 $G\in\R^{m\times n}$ 不依赖于 $dX$，则必有
\[
  \nabla_X f(X) = G.
\]
特别地，对向量 $x\in\R^n$，式~\eqref{eq:ch5-bridge} 简化为 $df = g^T dx \iff \nabla_x f(x) = g$。
\end{theorem}
\end{thmBox}

\begin{proof}
由多元函数全微分定义及矩阵迹的展开 $\tr(A^T B) = \sum_{i=1}^m\sum_{j=1}^n A_{ij} B_{ij}$：
\[
  df = \sum_{i=1}^m \sum_{j=1}^n \frac{\partial f}{\partial X_{ij}}\, dX_{ij}
     = \tr\!\Big(\big(\nabla_X f(X)\big)^T dX\Big).
\]
若同时有 $df = \tr(G^T dX)$ 对所有扰动 $dX\in\R^{m\times n}$ 成立，则 $\tr\big((\nabla_X f(X) - G)^T dX\big) = 0$ 恒成立。令 $dX = E_{ij}$（仅 $(i,j)$ 位置为 $1$、其余为 $0$ 的标准基矩阵），即得 $(\nabla_X f(X))_{ij} - G_{ij} = 0$。因 $(i,j)$ 任意，故 $\nabla_X f(X) = G$。
\end{proof}

有了定理~\ref{thm:ch5-golden-bridge}，计算任何矩阵导数都变成了一个机械、优雅的三步流水线：
\textbf{(1)} 对标量 $f$ 两边取全微分 $df$；\textbf{(2)} 套一层迹（因标量 $s = \tr(s)$）并利用迹的循环不变性把 $dX$ 搬到最右端；\textbf{(3)} 凑成 $df = \tr(G^T dX)$，直接读出 $\nabla_X f = G$！

为了流畅运行这三步，只需熟记表~\ref{tab:ch5-trace-rules} 中的 6 条基本规则（全部可一行验证）。

\begin{table}[htbp]
\centering
\caption{全微分、迹技巧与 Hadamard 逐元素求导的 6 条常用规则}
\label{tab:ch5-trace-rules}
\small
\begin{tabularx}{\linewidth}{@{}p{26mm}p{58mm}X@{}}
\toprule
\textbf{规则名称} & \textbf{公式} & \textbf{核心提醒} \\
\midrule
线性与转置 & $d(\alpha X+\beta Y)=\alpha\,dX+\beta\,dY$;\; $d(X^T)=(dX)^T$ & 微分算子 $d$ 与转置可交换 \\
乘积法则 & $d(XY) = (dX)Y + X(dY)$ & \textbf{严禁交换乘法顺序} \\
迹的线性与转置 & $\tr(A^T)=\tr(A)$;\; $d(\tr(X))=\tr(dX)$ & 可用来把 $dX^T$ 翻正为 $dX$ \\
迹的循环不变性 & $\tr(ABC) = \tr(CAB) = \tr(BCA)$ & 仅可\textbf{循环轮换}，不可任意对调两项 \\
标量套迹 & 若 $s\in\R$，则 $s = \tr(s)$ & 向量内积 $a^T b = \tr(a^T b) = \tr(b a^T)$ \\
Hadamard 规则 & $d\phi(Z) = \phi'(Z)\had dZ$;\; $\tr(A^T(B\had C)) = \tr((A\had B)^T C)$ & $\had$ 表逐元素乘，三条项对称等价于 $\sum_{i,j}A_{ij}B_{ij}C_{ij}$ \\
\bottomrule
\end{tabularx}
\end{table}

\subsection{五大 AI 必用求导实战（全流程一步不跳）}

下面五个推导覆盖了第 2 卷回归/分类模型、第 4 卷凸优化目标以及第 5 卷神经网络反向传播的绝大多数梯度公式，请务必亲手跟着推一遍。

\subsubsection{推导 1：线性函数 $f(x) = a^T x$}
取全微分：$df = a^T dx$。直接对照 $df = g^T dx$ 得
\begin{equation}\label{eq:ch5-grad-linear}
  \nabla_x (a^T x) = a.
\end{equation}

\subsubsection{推导 2：一般二次型 $f(x) = x^T Q x$（$Q\in\R^{n\times n}$ 未必对称）}
由乘积法则：
\[
  df = (dx)^T Q x + x^T Q\, dx.
\]
第一项 $(dx)^T Q x$ 是标量，取其转置值不变：$(dx)^T Q x = (x^T Q^T dx)^T = x^T Q^T dx$。合并两项得
\[
  df = x^T Q^T dx + x^T Q\, dx = \big((Q + Q^T)x\big)^T dx.
\]
由定理~\ref{thm:ch5-golden-bridge} 立即读出：
\begin{equation}\label{eq:ch5-grad-quad}
  \nabla_x (x^T Q x) = (Q + Q^T)x.
\end{equation}
仅当 $Q = Q^T$ 对称时，式~\eqref{eq:ch5-grad-quad} 才进一步化简为 $2Qx$（这正是例~\ref{ex:ch5-pca-lagrange} 中 PCA 推导用到的公式）。

\subsubsection{推导 3：线性回归均方误差（MSE）的梯度与 Hessian}
考虑第 2 章的线性回归目标函数（$X\in\R^{N\times d}$，权重 $w\in\R^d$，标签 $y\in\R^N$）：
\[
  f(w) = \frac{1}{2N}\norm{Xw - y}_2^2 = \frac{1}{2N}(Xw - y)^T(Xw - y).
\]
记残差向量 $r = Xw - y$，则 $dr = X\,dw$。取全微分：
\[
  df = \frac{1}{2N}\big((dr)^T r + r^T dr\big) = \frac{1}{N} r^T dr = \frac{1}{N}(Xw - y)^T X\, dw
     = \Big(\frac{1}{N} X^T(Xw - y)\Big)^T dw.
\]
于是梯度为
\begin{equation}\label{eq:ch5-grad-mse}
  \nabla_w f(w) = \frac{1}{N} X^T (Xw - y).
\end{equation}
令 $\nabla_w f(w) = 0$，立即重现第 2 章的正规方程 $X^T X w = X^T y$！再对梯度映射 $g(w) = \frac{1}{N}X^T X w - \frac{1}{N}X^T y$ 求 Jacobian，得线性回归的常值 Hessian 矩阵：
\[
  H_f(w) = \frac{1}{N} X^T X \succeq 0.
\]
回忆第 3 章：对任何矩阵 $X$，Gram 矩阵 $X^T X$ 必然对称半正定；若 $X$ 列满秩则严格正定，从而正规方程的解必为全局唯一极小值。

\subsubsection{推导 4：多输出矩阵最小二乘 $f(W) = \frac{1}{2}\norm{XW - Y}_F^2$}
在多标签回归或线性投影层中，权重是矩阵 $W\in\R^{d\times k}$，输入 $X\in\R^{N\times d}$，目标矩阵 $Y\in\R^{N\times k}$。记残差矩阵 $R = XW - Y$（则 $dR = X\,dW$），利用 Frobenius 范数与迹的恒等式 $\norm{R}_F^2 = \tr(R^T R)$：
\begin{align*}
  df &= \frac{1}{2}\,d\big(\tr(R^T R)\big)
      = \frac{1}{2}\tr\!\big((dR)^T R + R^T dR\big) \\
     &= \tr(R^T dR) \qquad\text{（因 $\tr((dR)^T R) = \tr(((dR)^T R)^T) = \tr(R^T dR)$）} \\
     &= \tr(R^T X\, dW) = \tr\!\big((X^T R)^T dW\big).
\end{align*}
对照 $df = \tr(G^T dW)$ 直接读出：
\begin{equation}\label{eq:ch5-grad-matls}
  \nabla_W \Big(\frac{1}{2}\norm{XW - Y}_F^2\Big) = X^T (XW - Y) \in \R^{d\times k}.
\end{equation}
注意检查形状：$X^T$ 是 $d\times N$，$XW-Y$ 是 $N\times k$，乘积恰为 $d\times k$，与 $W$ 完美吻合！

\subsubsection{推导 5：神经网络全连接 + 逐元素激活层的完整反向传播}
设一层神经网络接收批量输入矩阵 $X\in\R^{d_{\text{in}}\times B}$（每列为一个样本，共 $B$ 个样本），经权重矩阵 $W\in\R^{d_{\text{out}}\times d_{\text{in}}}$ 映射为预激活 $Z = W X \in \R^{d_{\text{out}}\times B}$，再经逐元素光滑激活函数 $\phi$（如 $\tanh$ 或 $\sigma$）输出 $A = \phi(Z) \in \R^{d_{\text{out}}\times B}$，最终接入上游标量损失 $L(A)\in\R$。

假设上游已经传回了损失对输出激活的梯度矩阵 $G_A \coloneqq \nabla_A L \in \R^{d_{\text{out}}\times B}$。我们来一步推导出 $\nabla_Z L$、$\nabla_W L$ 与 $\nabla_X L$：
\begin{enumerate}
  \item \textbf{穿过激活层（Hadamard 链式法则）}：由黄金桥梁定理与表~\ref{tab:ch5-trace-rules} 的 Hadamard 规则，$dA = \phi'(Z) \had dZ$，于是
  \[
    dL = \tr(G_A^T dA) = \tr\!\big(G_A^T (\phi'(Z) \had dZ)\big) = \tr\!\big((G_A \had \phi'(Z))^T dZ\big).
  \]
  直接读出预激活梯度：
  \begin{equation}\label{eq:ch5-backprop-z}
    G_Z \coloneqq \nabla_Z L = (\nabla_A L) \had \phi'(Z) \in \R^{d_{\text{out}}\times B}.
  \end{equation}
  \item \textbf{穿过线性层 $Z = WX$（对 $W$ 与对 $X$ 同时求导）}：由乘积法则 $dZ = (dW)X + W(dX)$，代入 $dL = \tr(G_Z^T dZ)$ 并拆成两项：
  \[
    dL = \tr\!\big(G_Z^T (dW) X\big) + \tr\!\big(G_Z^T W (dX)\big).
  \]
  对第一项使用迹的循环不变性把 $X$ 轮换到最左端：$\tr(G_Z^T (dW) X) = \tr(X G_Z^T dW) = \tr((G_Z X^T)^T dW)$；第二项直接写成 $\tr((W^T G_Z)^T dX)$。对照定理~\ref{thm:ch5-golden-bridge} 立即同时读出权重梯度与输入梯度：
  \begin{equation}\label{eq:ch5-backprop-wx}
    \nabla_W L = (\nabla_Z L)\, X^T \in \R^{d_{\text{out}}\times d_{\text{in}}},
    \qquad
    \nabla_X L = W^T (\nabla_Z L) \in \R^{d_{\text{in}}\times B}.
  \end{equation}
\end{enumerate}
请停下来欣赏一下式~\eqref{eq:ch5-backprop-z} 和~\eqref{eq:ch5-backprop-wx}：没有任何下标求和，短短四行全微分推导，就得到了第 5 卷多层感知机（MLP）反向传播的全部核心代码——且 $\nabla_W L = G_Z X^T = \sum_{b=1}^B (G_Z)_{:,b} X_{:,b}^T$ 正是第 1 章强调过的\textbf{矩阵乘法外积视角}（批量中每个样本各自贡献一个秩 1 梯度外积再累加）！

\begin{failBox}{形状核对铁律与非对称二次型暗坑}
在手推和编写矩阵梯度时，有两个最常见、最隐蔽的错误必须时刻防范：
\begin{enumerate}
  \item \textbf{形状核对铁律（Shape Checking）}：在分母布局下，无论公式多复杂，\texttt{grad\_W.shape == W.shape} 必须无条件成立。例如在 $Z = WX$（$W\in\R^{m\times n}, X\in\R^{n\times B}, G_Z\in\R^{m\times B}$）中，想用 $G_Z$ 和 $X$ 拼出 $m\times n$ 的矩阵，唯有 $G_Z X^T$ 合法（$X^T G_Z$ 或 $G_Z^T X$ 维度根本对不上）。当 $W$ 恰好是方阵时维度检查会失效，此时务必用第 6 章的有限差分梯度校验来捉虫！
  \item \textbf{把非对称二次型错写成 $2Qx$}：设 $Q = \begin{bmatrix} 1 & 2 \\ 0 & 3 \end{bmatrix}$，则 $f(x) = x^T Q x = x_1^2 + 2x_1 x_2 + 3x_2^2$。直接求偏导得 $\nabla f(x) = \begin{bmatrix} 2x_1 + 2x_2 \\ 2x_1 + 6x_2 \end{bmatrix} = (Q+Q^T)x$；若误套对称阵公式 $2Qx = \begin{bmatrix} 2x_1 + 4x_2 \\ 6x_2 \end{bmatrix}$，则会静默产生系数差 2 倍和缺项的灾难性 Bug！
\end{enumerate}
\end{failBox}

% ====================================================================
\section{本章习题}
\label{sec:ch5-exercises}

\exTier{A. 概念辨析与反例查错}{合上书独立判断正误并构造具体反例}
\begin{enumerate}[label=\textbf{5.\arabic*}]
  \item \textbf{判断正误并纠错}：
  (a) “梯度方向指向全局最大值，负梯度方向直接指向全局最小值。”
  (b) “若多元可微函数 $f$ 在 $x^*$ 处满足 $\nabla f(x^*)=0$ 且 $H_f(x^*) \succeq 0$（半正定），则 $x^*$ 一定是局部极小值点。”
  请对每一条给出明确的真假判断；若为假，请构造一个极简的 1D 或 2D 显式反例。

  \item \textbf{迹技巧与矩阵导数找茬}：某同学尝试推导 $f(W) = \tr(A W B W^T)$（$A\in\R^{m\times m}, B\in\R^{n\times n}, W\in\R^{m\times n}$，且 $A,B$ 未必对称）的梯度时写道：
  \[
    df = \tr(A (dW) B W^T + A W B (dW)^T) = 2\tr(B W^T A\, dW) \implies \nabla_W f = 2 A^T W B^T.
  \]
  请指出该推导中\textbf{唯一但致命的一处错误}，并写出对一般（非对称）矩阵 $A,B$ 正确的梯度公式 $\nabla_W f$。
\end{enumerate}

\exTier{B. 具体小例手算}{亲手算通 2D/3D 数值，建立第一手肌肉记忆}
\begin{enumerate}[label=\textbf{5.\arabic*}, resume]
  \item \textbf{方向导数、Hessian 与驻点分类}：设 $f(x_1, x_2) = x_1^3 - 3x_1 x_2 + x_2^3$。
  (a) 求在点 $(1, 0)^T$ 处沿单位方向 $u = (\frac{3}{5}, \frac{4}{5})^T$ 的方向导数 $D_u f(1,0)$；
  (b) 求出 $f$ 的全部驻点，并用 Hessian 矩阵的特征值（或行列式与迹）逐一判定每个驻点是严格局部极小值、严格局部极大值还是鞍点。

  \item \textbf{两层标量网络的前向与反向手算}：考虑小型函数 $L(x) = \frac{1}{2}\norm{a - y}_2^2$，其中 $z = Wx$，$a = \tanh(z)$，具体数值取 $W = \begin{bmatrix} 1 & -1 \\ 0 & 2 \end{bmatrix}$，$x = \begin{bmatrix} 1 \\ 0 \end{bmatrix}$，$y = \begin{bmatrix} 1 \\ 1 \end{bmatrix}$。
  (a) 计算前向中间值 $z$、$a$ 与损失值 $L$（提示：$\tanh(1)\approx 0.7616$，$\tanh(0)=0$）；
  (b) 利用式~\eqref{eq:ch5-backprop-z}--\eqref{eq:ch5-backprop-wx}，按反向模式 VJP 依次手算出 $\nabla_a L$、$\nabla_z L$、$\nabla_W L$ 与 $\nabla_x L$ 的精确数值（保留 4 位小数）。
\end{enumerate}

\exTier{C. AI 核心推演与实用证明}{连接线性代数与后续机器学习的核心公式}
\begin{enumerate}[label=\textbf{5.\arabic*}, resume]
  \item \textbf{岭回归梯度与闭式解推导}：考虑带 $L_2$ 正则项的岭回归损失函数（$\lambda > 0$）：
  \[
    J(w) = \frac{1}{2}\norm{Xw - y}_2^2 + \frac{\lambda}{2}\norm{w}_2^2.
  \]
  用全微分法求出 $\nabla_w J(w)$ 与 Hessian 矩阵 $H_J(w)$，证明对任意矩阵 $X\in\R^{N\times d}$（即使列不满秩甚至 $N < d$），$H_J(w)$ 均严格正定，从而 $J(w)$ 存在唯一全局极小值 $w_{\text{ridge}} = (X^T X + \lambda I)^{-1} X^T y$。
  \exHint{利用第 3 章性质：$X^T X \succeq 0$ 且 $\lambda I \succ 0$，两者之和对任意非零向量 $v$ 满足 $v^T(X^T X + \lambda I)v = \norm{Xv}_2^2 + \lambda\norm{v}_2^2 > 0$。}

  \item \textbf{带约束最小二乘与双线性最大奇异值}：利用拉格朗日乘子法求解优化问题
  \[
    \max_{u\in\R^m,\, v\in\R^n}\; u^T A v \qquad \text{subject to}\qquad \norm{u}_2^2 = 1,\; \norm{v}_2^2 = 1,
  \]
  其中 $A\in\R^{m\times n}$ 为非零矩阵。证明最优解 $(u^*, v^*)$ 必满足 $A v^* = \sigma u^*$ 与 $A^T u^* = \sigma v^*$（即第 4 章的左右奇异向量方程！），且最大值为 $A$ 的最大奇异值 $\sigma_1(A)$。
  \exHint{构造含两个乘子的拉格朗日函数 $\mathcal{L}(u,v,\lambda_1,\lambda_2) = u^T A v - \frac{\lambda_1}{2}(u^T u - 1) - \frac{\lambda_2}{2}(v^T v - 1)$，分别令 $\nabla_u\mathcal{L}=0, \nabla_v\mathcal{L}=0$，再左乘 $u^T$ 与 $v^T$ 证明 $\lambda_1 = \lambda_2 = u^T A v$。}
\end{enumerate}

\exTier{D. 可自检编程实验}{动手写 NumPy 验证公式与计算复杂度}
\begin{enumerate}[label=\textbf{5.\arabic*}, resume]
  \item \textbf{前向模式 vs 反向模式（VJP）计时对比}：构造一条 3 层线性标量链条 $L(x) = c^T W_2 W_1 x$，其中 $x\in\R^d$，$W_1, W_2 \in \R^{d\times d}$，$c\in\R^d$，取 $d = 800$。用 NumPy 分别按两种顺序计算梯度行向量 $g^T = c^T W_2 W_1$：
  \begin{itemize}
    \item 前向结合：\texttt{c @ (W2 @ W1)}；
    \item 反向 VJP 结合：\texttt{(c @ W2) @ W1}。
  \end{itemize}
  \textbf{自检标准}：验证两种结果的相对误差 $\norm{g_{\text{fwd}} - g_{\text{rev}}}_2 / \norm{g_{\text{rev}}}_2 < 10^{-12}$，且反向模式运行时间比前向模式快至少 $10\times$ 以上（实测通常在 $50\times$--$300\times$）。
  \exHint{使用 \texttt{time.perf\_counter()}，在 \texttt{d=800} 时分别重复取平均（前向可跑 5 次，反向可跑 200 次）以避免计时抖动。}

  \item \textbf{手写两层神经网络前向与矩阵反向传播}：给定批量数据 $X\in\R^{5\times 16}$、目标 $Y\in\R^{3\times 16}$、权重 $W_1\in\R^{8\times 5}$ 与 $W_2\in\R^{3\times 8}$，定义前向计算
  \[
    Z_1 = W_1 X,\quad A_1 = \tanh(Z_1),\quad \hat{Y} = W_2 A_1,\quad L = \frac{1}{2}\norm{\hat{Y} - Y}_F^2.
  \]
  请利用本节式~\eqref{eq:ch5-backprop-z}--\eqref{eq:ch5-backprop-wx} 编写纯 NumPy 函数同时返回损失 $L$ 与解析梯度 $(\nabla_{W_1} L, \nabla_{W_2} L)$。
  \textbf{自检标准}：断言 \texttt{grad\_W1.shape == W1.shape} 与 \texttt{grad\_W2.shape == W2.shape}，并用中心有限差分（$\epsilon = 10^{-5}$，详见第 6.4 节）核对两个权重矩阵梯度的相对误差均小于 $10^{-9}$。
  \exHint{先算 $G_{\hat{Y}} = \hat{Y} - Y$，再依次反传得 $\nabla_{W_2} L = G_{\hat{Y}} A_1^T$、$G_{A_1} = W_2^T G_{\hat{Y}}$、$G_{Z_1} = G_{A_1} \had (1 - A_1^2)$、$\nabla_{W_1} L = G_{Z_1} X^T$。}
\end{enumerate}

\chapter[二次型梯度下降与有限差分校验]{优化初体验：二次型梯度下降动力学\\与有限差分校验}
\label{ch:gd-dynamics-gradcheck}

\chapEpigraph{第一原则是你绝不能欺骗自己——而你自己正是最容易被骗的人。}{理查德·费曼（Richard Feynman），《草包族科学》（1974）}

\begin{chapIntro}{本章定位}
至此，我们已经攒齐了第 1 卷的全部核心部件：第 1 章的 NumPy 向量化、第 2 章的线性最小二乘、第 3 章的实对称矩阵谱定理与条件数、第 4 章的奇异值视角，以及第 5 章的多元梯度与矩阵微分。本章将把这五大部件拧成一台完整运转的机器——\textbf{迭代优化的起点：梯度下降（Gradient Descent）}。我们不提前涉足第 4 卷的一般凸优化抽象收敛证明，而是用第 3 章的谱定理\textbf{精确解开二次型梯度下降的动力学全貌}：看清步长门槛 $0 < \eta < 2/\lambda_{\max}$、最优步长 $\eta^*$、病态条件数如何制造锯齿震荡、为什么 Z-score 特征标准化能瞬间把条件数砍掉近 80 万倍，以及如何用\textbf{有限差分梯度校验}为未来所有手写梯度代码装上终极安全网。
\end{chapIntro}

% ====================================================================
\section{为什么解析解不够用？从正规方程走向梯度下降}
\label{sec:ch6-why-gd}

第 2 章里我们用正规方程 $X^T X w = X^T y$（或 \texttt{np.linalg.lstsq}）一步就拿到了线性回归的精确最优权重 $w^*$。读者可能会自然发问：既然能一步算出解析解，为什么现代 AI 几乎清一色采用看似“笨拙”的迭代算法——\textbf{梯度下降}？
\begin{equation}\label{eq:ch6-gd-update}
  w_{t+1} = w_t - \eta\, \nabla f(w_t), \qquad t = 0, 1, 2, \dots
\end{equation}
其中 $\eta > 0$ 称为\textbf{步长（Step Size）}或\textbf{学习率（Learning Rate）}。答案来自计算规模与模型表达力两个无法绕开的现实约束（表~\ref{tab:ch6-closed-vs-gd}）：

\begin{table}[htbp]
\centering
\caption{解析求解（正规方程 / QR / SVD）与迭代一阶方法（梯度下降）的本质对比}
\label{tab:ch6-closed-vs-gd}
\small
\begin{tabularx}{\linewidth}{@{}p{24mm}p{58mm}X@{}}
\toprule
\textbf{对比维度} & \textbf{解析解（正规方程 / \texttt{lstsq}）} & \textbf{梯度下降（Full-batch / Mini-batch）} \\
\midrule
适用目标函数 & \textbf{仅限}线性二次目标（如普通/岭线性回归） & \textbf{任意}可微目标（第 2 卷逻辑回归、第 5 卷深层网络均无闭式解） \\
单步时间复杂度 & 一次性 $O(ND^2 + D^3)$ & 全批量单步 $O(ND)$；小批量（Mini-batch）仅 $O(BD)$ \\
显存 / 内存占用 & 需构造 $D\times D$ 矩阵 $X^T X$（$D=10^5$ 时占 80 GB） & 仅需存参数与梯度 $O(D)$（$D=10^5$ 仅占 0.8 MB） \\
高维可扩展性 & 当 $D > 10^4$ 时 $O(D^3)$ 彻底不可行 & 可扩展至百亿参数大模型（只需一阶梯度） \\
\bottomrule
\end{tabularx}
\end{table}

在给出一般理论之前，我们先按第 1 卷一贯的传统，在一个极简二维椭圆碗上\textbf{亲手推演前 3 步梯度下降}——亲眼目睹步长选得稍大时，算法是如何原地爆炸的。

\begin{example}[二维二次型上 3 步梯度下降的亲手追踪]\label{ex:ch6-hand-trace}
考虑最小化二次型函数
\[
  f(w_1, w_2) = \frac{1}{2}\big(w_1^2 + 10 w_2^2\big) = \frac{1}{2} w^T Q w,
  \qquad Q = \begin{bmatrix} 1 & 0 \\ 0 & 10 \end{bmatrix},
\]
其唯一全局极小点显然为 $w^* = (0, 0)^T$，最小值 $f(w^*) = 0$。梯度为 $\nabla f(w) = Qw = (w_1,\, 10w_2)^T$，因此单步更新公式为
\[
  w_{t+1,1} = (1 - \eta) w_{t,1},
  \qquad
  w_{t+1,2} = (1 - 10\eta) w_{t,2}.
\]
固定初值 $w_0 = (1, 1)^T$（初始损失 $f(w_0) = 5.5$），分别试探两种步长：
\begin{itemize}
  \item \textbf{保守步长 $\eta = 0.05$}：乘子为 $1-\eta = 0.95$，$1-10\eta = 0.50$。两个坐标都单调向 $0$ 收缩；
  \item \textbf{冒进步长 $\eta = 0.21$}：乘子为 $1-\eta = 0.79$，$1-10\eta = \mathbf{-1.10}$！注意第二维乘子绝对值大于 $1$ 且带负号，每一步都跨过谷底并被放大 $1.1$ 倍！
\end{itemize}
表~\ref{tab:ch6-hand-trace} 完整列出了 $t=0,1,2,3$ 的精确数值：
\end{example}

\begin{table}[htbp]
\centering
\caption{例~\ref{ex:ch6-hand-trace} 中 $f(w)=\frac{1}{2}(w_1^2+10w_2^2)$ 从 $w_0=(1,1)^T$ 出发的前 3 步手算对比}
\label{tab:ch6-hand-trace}
\small
\begin{tabular}{@{}crrrrrr@{}}
\toprule
& \multicolumn{3}{c}{$\eta = 0.05$（稳定收敛）} & \multicolumn{3}{c}{$\eta = 0.21$（沿陡轴剧烈发散）} \\
\cmidrule(lr){2-4}\cmidrule(lr){5-7}
步数 $t$ & $w_{t,1}$ & $w_{t,2}$ & $f(w_t)$ & $w_{t,1}$ & $w_{t,2}$ & $f(w_t)$ \\
\midrule
$0$ & $1.0000$ & $1.0000$ & $5.5000$ & $1.0000$ & $1.0000$ & $5.5000$ \\
$1$ & $0.9500$ & $0.5000$ & $1.7013$ & $0.7900$ & $-1.1000$ & $6.3621$ \\
$2$ & $0.9025$ & $0.2500$ & $0.7198$ & $0.6241$ & $+1.2100$ & $7.5153$ \\
$3$ & $0.8574$ & $0.1250$ & $0.4457$ & $0.4930$ & $-1.3310$ & $8.9793$ \\
\bottomrule
\end{tabular}
\end{table}

看看表~\ref{tab:ch6-hand-trace} 中蕴含的两个深刻信号：
\begin{enumerate}
  \item 在 $\eta = 0.21$ 下，尽管平缓轴 $w_{t,1}$ 正在乖乖衰减（$1 \to 0.79 \to 0.624 \to 0.493$），但陡峭轴 $w_{t,2}$ 正以 $(-1.1)^t$ 左右反弹并指数级爆炸，拉着总损失 $f(w_t)$ 一路飙升！临界点恰好卡在 $1 - 10\eta = -1 \iff \eta = 0.20 = 2/\lambda_{\max}$。
  \item 在 $\eta = 0.05$ 下，第 3 步之后 $w_{3,2} = 0.125$ 已经很快贴近谷底，但平缓轴 $w_{3,1} = 0.8574$ 却依然蜗牛般爬行——这就是为什么 $\eta$ 不能太小、也不能太大的两难困境。
\end{enumerate}

% ====================================================================
\section{二次型梯度下降的精确谱动力学}
\label{sec:ch6-spectral-gd}

例~\ref{ex:ch6-hand-trace} 中的矩阵 $Q = \diag(1, 10)$ 恰好是对角阵，坐标轴天然解耦。如果 $Q$ 是一般的实对称正定矩阵（例如旋转过后的倾斜椭圆碗，或线性回归的 Hessian 矩阵 $Q = \frac{1}{N}X^T X$），各分量彼此纠缠，我们还能如此干净地分析每一步迭代吗？

答案是肯定的——这正是第 3 章\textbf{实对称矩阵谱定理}在本卷最华丽的一次应用！

\subsection{特征坐标解耦与收敛/最优步长定理}

考虑一般正定二次型目标函数：
\begin{equation}\label{eq:ch6-quad-obj}
  f(w) = \frac{1}{2} w^T Q w - b^T w + c,
\end{equation}
其中 $Q \in \R^{d\times d}$ 为实对称正定矩阵（$Q = Q^T \succ 0$），$b\in\R^d$，$c\in\R$。由第 5.4 节矩阵求导公式得 $\nabla f(w) = Qw - b$，Hessian 为常矩阵 $H_f(w) = Q \succ 0$。因 $Q$ 正定可逆，函数存在唯一全局极小点 $w^* = Q^{-1} b$（满足 $Qw^* = b$）。

记 $Q$ 的特征值从大到小排列为
\[
  \lambda_{\max} = \lambda_1 \ge \lambda_2 \ge \dots \ge \lambda_d = \lambda_{\min} > 0,
\]
并记第 3 章定义的对称正定矩阵\textbf{条件数}为 $\cond \coloneqq \cond_2(Q) = \dfrac{\lambda_{\max}}{\lambda_{\min}} \ge 1$。

\begin{thmBox}
\begin{theorem}[二次型梯度下降的精确谱收敛与最优步长定理]\label{thm:ch6-gd-spectral}
对式~\eqref{eq:ch6-quad-obj} 中的正定二次型 $f(w)$（$Q=Q^T\succ 0$），从任意初值 $w_0\in\R^d$ 运行固定步长梯度下降 $w_{t+1} = w_t - \eta(Qw_t - b)$。记参数误差向量 $e_t \coloneqq w_t - w^*$，并记谱分解 $Q = V\Lambda V^T$（$V$ 正交，$\Lambda = \diag(\lambda_1,\dots,\lambda_d)$），令特征坐标误差为 $z_t \coloneqq V^T e_t$。则：
\begin{enumerate}
  \item \textbf{精确解耦递推}：对每个特征方向 $i = 1, \dots, d$，误差分量独立按等比数列演化：
  \begin{equation}\label{eq:ch6-eigen-recur}
    (z_t)_i = (1 - \eta \lambda_i)^t\, (z_0)_i, \qquad \norm{e_t}_2 = \norm{z_t}_2 \le \rho(\eta)^t \norm{e_0}_2,
  \end{equation}
  其中单步最坏收缩因子定义为 $\rho(\eta) \coloneqq \max_{1\le i\le d} \abs{1 - \eta \lambda_i}$。
  \item \textbf{收敛的充要步长条件}：梯度下降对\textbf{所有}初值 $w_0\in\R^d$ 均收敛到 $w^*$，当且仅当
  \begin{equation}\label{eq:ch6-eta-bound}
    0 < \eta < \frac{2}{\lambda_{\max}}.
  \end{equation}
  \item \textbf{最优固定步长与最快线性收缩率}：在区间 $(0, 2/\lambda_{\max})$ 内，使最坏收缩因子 $\rho(\eta)$ 达到全局最小的最优固定步长及其收缩因子为
  \begin{equation}\label{eq:ch6-eta-opt}
    \eta^* = \frac{2}{\lambda_{\max} + \lambda_{\min}},
    \qquad
    \rho^* = \rho(\eta^*) = \frac{\lambda_{\max} - \lambda_{\min}}{\lambda_{\max} + \lambda_{\min}}
           = \frac{\cond - 1}{\cond + 1}.
  \end{equation}
\end{enumerate}
\end{theorem}
\end{thmBox}

\begin{proof}
\textbf{(1)} 利用 $b = Q w^*$，将梯度改写为 $\nabla f(w_t) = Qw_t - Qw^* = Q e_t$。在更新式 $w_{t+1} = w_t - \eta \nabla f(w_t)$ 两边同时减去 $w^*$，得误差线性递推：
\[
  e_{t+1} = e_t - \eta Q e_t = (I - \eta Q) e_t.
\]
代入第 3 章的实对称矩阵谱分解 $Q = V\Lambda V^T$（$V^T V = V V^T = I$），在两边同时左乘 $V^T$：
\[
  z_{t+1} = V^T e_{t+1} = V^T (V V^T - \eta V\Lambda V^T) e_t = (I - \eta \Lambda) V^T e_t = (I - \eta \Lambda) z_t.
\]
因为 $I - \eta\Lambda = \diag(1-\eta\lambda_1, \dots, 1-\eta\lambda_d)$ 是对角阵，各个分量彻底解耦，递推 $t$ 步即得 $(z_t)_i = (1-\eta\lambda_i)^t (z_0)_i$。又因正交变换保 $\ell_2$ 范数，故
\[
  \norm{e_t}_2^2 = \norm{z_t}_2^2 = \sum_{i=1}^d (1-\eta\lambda_i)^{2t} (z_0)_i^2
  \le \rho(\eta)^{2t} \sum_{i=1}^d (z_0)_i^2 = \rho(\eta)^{2t} \norm{e_0}_2^2,
\]
开方即得式~\eqref{eq:ch6-eigen-recur}。顺带可得目标函数值误差满足 $f(w_t) - f(w^*) = \frac{1}{2}e_t^T Q e_t = \frac{1}{2}\sum_i \lambda_i (1-\eta\lambda_i)^{2t}(z_0)_i^2 \le \rho(\eta)^{2t}\big(f(w_0)-f(w^*)\big)$。

\textbf{(2)} 序列 $e_t \to 0$ 对任意 $w_0\in\R^d$（即任意 $z_0\in\R^d$）均成立，当且仅当每个特征方向的公比绝对值严格小于 $1$：
\[
  \forall\, i\in\{1,\dots,d\},\quad \abs{1 - \eta\lambda_i} < 1
  \iff -1 < 1 - \eta\lambda_i < 1
  \iff 0 < \eta < \frac{2}{\lambda_i}.
\]
对所有 $i$ 同时成立，取决于右侧最紧的那个上界 $\min_i (2/\lambda_i) = 2/\lambda_{\max}$，故充要条件为 $0 < \eta < 2/\lambda_{\max}$。

\textbf{(3)} 固定 $\eta > 0$，函数 $\lambda \mapsto |1 - \eta\lambda|$ 在区间 $[\lambda_{\min}, \lambda_{\max}]$ 上是一个先减后增的 V 形折线（凸函数），其最大值必定在两个端点 $\lambda_{\min}$ 或 $\lambda_{\max}$ 处取得：
\[
  \rho(\eta) = \max\Big\{\abs{1 - \eta\lambda_{\min}},\; \abs{1 - \eta\lambda_{\max}}\Big\}.
\]
当 $\eta$ 从 $0$ 增大时，$1 - \eta\lambda_{\min}$ 单调递减，而 $\eta\lambda_{\max} - 1$ 单调递增。两条直线的交点恰使两者相等：
\[
  1 - \eta^*\lambda_{\min} = \eta^*\lambda_{\max} - 1
  \implies \eta^* = \frac{2}{\lambda_{\max} + \lambda_{\min}}.
\]
容易验证此时两者的共同取值即为最小值：
\[
  \rho(\eta^*) = 1 - \frac{2\lambda_{\min}}{\lambda_{\max} + \lambda_{\min}}
               = \frac{\lambda_{\max} - \lambda_{\min}}{\lambda_{\max} + \lambda_{\min}}
               = \frac{\cond - 1}{\cond + 1}.
\]
若 $\eta < \eta^*$，则 $1 - \eta\lambda_{\min} > \rho^*$；若 $\eta > \eta^*$，则 $\eta\lambda_{\max} - 1 > \rho^*$。因此 $\eta^*$ 是唯一的最优固定步长。
\end{proof}

\subsection{实测曲线与窄谷锯齿几何图景}

把定理~\ref{thm:ch6-gd-spectral} 代回例~\ref{ex:ch6-hand-trace} 中的 $Q = \diag(1, 10)$（$\lambda_{\max}=10, \lambda_{\min}=1, \cond=10$）：
\begin{itemize}
  \item 收敛上限：$2/\lambda_{\max} = 0.20$；
  \item 最优步长：$\eta^* = \frac{2}{10+1} = \frac{2}{11} \approx 0.1818$，最优单步参数收缩率 $\rho^* = \frac{10-1}{10+1} = \frac{9}{11} \approx 0.8182$，因此每步函数值精确衰减 $(\frac{9}{11})^2 \approx 0.6694$ 倍！
\end{itemize}

图~\ref{fig:ch6-gd-dynamics} 左侧基于本卷脚本 \texttt{code/vol1\_experiments.py} 生成的实测数据（\texttt{data/gd\_rates.csv}）绘制了半对数收敛曲线，右侧展示了在椭圆等高线上的真实迭代轨迹：

\begin{figure}[htbp]
\centering
\begin{minipage}[c]{0.52\linewidth}
\centering
\begin{tikzpicture}
\begin{semilogyaxis}[
  width=0.97\linewidth,
  height=5.25cm,
  xlabel={迭代步数 $t$},
  ylabel={目标函数值 $f(w_t)$（对数刻度）},
  xmin=-0.6, xmax=30.6,
  ymin=1e-5, ymax=2e3,
  grid=both,
  grid style={line width=0.25pt, draw=slate!20},
  major grid style={line width=0.4pt, draw=slate!35},
  tick label style={font=\scriptsize},
  label style={font=\footnotesize},
  legend style={font=\scriptsize, legend columns=2, at={(0.5,1.025)}, anchor=south,
    fill=white, fill opacity=0.9, draw=codeborder, inner sep=2pt},
  legend cell align=left
]
  \addplot[navy, thick, mark=none]
    table[col sep=comma, x=t, y=eta005]{data/gd_rates.csv};
  \addlegendentry{$\eta=0.05$（过保守）}

  \addplot[teal, very thick, mark=*, mark size=1.1pt, mark repeat=3]
    table[col sep=comma, x=t, y=etaopt]{data/gd_rates.csv};
  \addlegendentry{$\eta^*=2/11$（最优）}

  \addplot[plum, thick, dashed]
    table[col sep=comma, x=t, y=eta0195]{data/gd_rates.csv};
  \addlegendentry{$\eta=0.195$（临近震荡）}

  \addplot[ochre, very thick, mark=triangle*, mark size=1.3pt, mark repeat=3]
    table[col sep=comma, x=t, y=eta021]{data/gd_rates.csv};
  \addlegendentry{$\eta=0.21 > 2/\lambda_{\max}$（发散）}
\end{semilogyaxis}
\end{tikzpicture}
\end{minipage}\hfill
\begin{minipage}[c]{0.47\linewidth}
\centering
\begin{tikzpicture}[xscale=2.35, yscale=1.62]
  % Axes
  \draw[-{Stealth[length=2mm]}, slate!65, thin] (-0.28,0) -- (1.42,0) node[right, font=\scriptsize, ink] {$w_1$};
  \draw[-{Stealth[length=2mm]}, slate!65, thin] (0,-1.45) -- (0,1.45) node[above, font=\scriptsize, ink] {$w_2$};

  % Ellipses w1^2 + 10 w2^2 = c
  \begin{scope}
    \clip (-0.25,-1.38) rectangle (1.36,1.38);
    \foreach \c/\col in {0.8/tealLine, 2.5/tealLine, 5.5/tealLine, 11.0/teal} {
      \draw[\col, thin] (0,0) ellipse ({sqrt(\c)} and {sqrt(\c/10)});
    }
  \end{scope}

  % Minimum
  \fill[tealDark] (0,0) circle (1.1pt) node[below left, font=\tiny, tealDark] {$w^*$};

  % Trajectory eta* = 2/11 (rho = 9/11 = 0.818)
  \draw[teal, very thick, -{Stealth[length=1.7mm]}] (1.000, 1.000) -- (0.818, -0.818);
  \draw[teal, very thick, -{Stealth[length=1.7mm]}] (0.818, -0.818) -- (0.669, 0.669);
  \draw[teal, very thick, -{Stealth[length=1.7mm]}] (0.669, 0.669) -- (0.548, -0.548);
  \draw[teal, very thick, -{Stealth[length=1.7mm]}] (0.548, -0.548) -- (0.448, 0.448);
  \draw[teal, very thick, -{Stealth[length=1.7mm]}] (0.448, 0.448) -- (0.367, -0.367);
  \fill[teal] (0.818,-0.818) circle (0.9pt) (0.669,0.669) circle (0.9pt)
              (0.548,-0.548) circle (0.9pt) (0.448,0.448) circle (0.9pt);

  % Trajectory eta = 0.21 (diverging): (1,1) -> (0.79,-1.10) -> (0.624, 1.21) -> (0.493, -1.331)
  \draw[ochre, thick, dashed, -{Stealth[length=1.7mm]}] (1.000, 1.000) -- (0.790, -1.100);
  \draw[ochre, thick, dashed, -{Stealth[length=1.7mm]}] (0.790, -1.100) -- (0.624, 1.210);
  \draw[ochre, thick, dashed, -{Stealth[length=1.7mm]}] (0.624, 1.210) -- (0.493, -1.331);
  \fill[ochre] (0.790,-1.100) circle (0.9pt) (0.624,1.210) circle (0.9pt) (0.493,-1.331) circle (0.9pt);

  % Start point
  \fill[ink] (1.000, 1.000) circle (1.3pt) node[above right, font=\scriptsize, ink] {$w_0=(1,1)$};

  % Clean non-overlapping annotations
  \node[font=\scriptsize, tealDark, align=left, fill=white, fill opacity=0.88, text opacity=1, inner sep=1pt]
    at (1.37, 0.39) {$\eta^*=2/11$\\锯齿收敛};
  \node[font=\scriptsize, ochre, align=left, fill=white, fill opacity=0.88, text opacity=1, inner sep=1pt]
    at (1.08, -1.18) {$\eta=0.21$\\越振越远};
\end{tikzpicture}
\end{minipage}
\caption{二次型 $f(w)=\frac{1}{2}(w_1^2+10w_2^2)$ 上梯度下降的实测动力学与几何轨迹。\textbf{左图}（\texttt{data/gd\_rates.csv} 实测）：最优步长 $\eta^*=2/11$ 在半对数图上呈现最陡直线（$t=30$ 时降至 $3.2\times 10^{-5}$）；$\eta=0.05$ 与 $\eta=0.195$ 均受限于各自较慢一侧的收缩率；而 $\eta=0.21$ 越过 $2/\lambda_{\max}=0.20$ 红线后呈指数发散。\textbf{右图}：椭圆窄谷上的锯齿收敛 vs 越振越远。}
\label{fig:ch6-gd-dynamics}
\end{figure}

\begin{insightBox}{大条件数为什么是梯度下降的头号天敌？}
观察最优收缩率 $\rho^* = \frac{\cond - 1}{\cond + 1} = 1 - \frac{2}{\cond + 1}$：
\begin{itemize}
  \item 当 $\cond = 1$（正圆形等高线）时，$\rho^* = 0$，梯度直接指向圆心，\textbf{一步精准到位}！
  \item 当 $\cond = 10^4$（极扁的狭长山谷）时，$\rho^* \approx 0.9998$：陡峭方向 $\lambda_{\max}$ 死死封死了步长上限（$\eta < 2/\lambda_{\max}$），导致沿平缓谷底 $\lambda_{\min}$ 方向的每步推进幅度被压榨到不足 $0.02\%$！
\end{itemize}
把误差压缩到 $10^{-6}$ 所需的迭代步数大致与条件数 $\cond$ \textbf{成正比}（$t \propto \cond$）。这直接解释了为什么第 5 卷深度神经网络极度依赖 BatchNorm / LayerNorm（动态改善每层特征条件数）与 Adam（按坐标自适应缩放步长）！
\end{insightBox}

% ====================================================================
\section{条件数、特征标准化与 Mini-batch SGD 直觉}
\label{sec:ch6-scaling-sgd}

上一节的理论绝不是纸上谈兵。在真实机器学习数据中，不同特征的量纲常常相差千倍万倍——例如房价预测中“房屋年收入”量级在 $10^5$、“卧室个数”量级在 $1$；或者传感器数据中一个通道标准差为 $100$，另一个通道标准差为 $0.1$。这对线性回归的梯度下降意味着什么？

\subsection{实测：Z-score 特征标准化把条件数砍掉近 80 万倍}

考虑 \texttt{code/vol1\_experiments.py} 中的实测实验：生成 $N=200$ 个二维样本，特征 $x_1 \sim \mathcal{N}(0, 100^2)$、$x_2 \sim \mathcal{N}(0, 0.1^2)$，拼成设计矩阵 $X_{\text{raw}} \in \R^{200\times 2}$。回忆第 5.4 节，线性回归 MSE 的 Hessian 矩阵正是 $H = \frac{1}{N} X^T X$。

现在对每一列特征执行 \textbf{Z-score 标准化（Standardization）}：
\begin{equation}\label{eq:ch6-zscore}
  \tilde{x}_{ij} = \frac{x_{ij} - \mu_j}{\sigma_j},
  \qquad
  \mu_j = \frac{1}{N}\sum_{i=1}^N x_{ij},
  \quad
  \sigma_j = \sqrt{\frac{1}{N}\sum_{i=1}^N (x_{ij} - \mu_j)^2}.
\end{equation}
对比标准化前后 Hessian 矩阵 $H = \frac{1}{N}X^T X$ 的本机实测指标（表~\ref{tab:ch6-scaling}）：

\begin{table}[htbp]
\centering
\caption{特征量纲失衡 vs Z-score 标准化的线性回归 Hessian 条件数与收敛步数实测对比（$N=200$）}
\label{tab:ch6-scaling}
\small
\begin{tabular}{@{}lrrr@{}}
\toprule
\textbf{特征预处理方式} & \textbf{Hessian 条件数 $\cond(H)$} & \textbf{最优收缩率 $\rho^* = \frac{\cond-1}{\cond+1}$} & \textbf{误差降 $10^{-6}$ 最坏步数} \\
\midrule
原始特征 $X_{\text{raw}}$ & $8.22 \times 10^{5}$ & $0.9999976$ & 约 \textbf{568 万步} \\
Z-score 标准化 $X_{\text{std}}$ & $\mathbf{1.036}$ & $\mathbf{0.0174}$ & 仅 \textbf{4 步} \\
\bottomrule
\end{tabular}
\end{table}

仅仅一行两秒钟的 Z-score 预处理，就把 Hessian 条件数从 \textbf{$8.22\times 10^5$} 暴降到 \textbf{$1.036$}（改善了近 \textbf{79.4 万倍}！），把等高线从一根几乎贴成直线的“针”瞬间拉回接近完美的圆碗，最坏收敛步数从\textbf{五百多万步骤降至 4 步}！

\subsubsection{训练完后如何还原原始特征上的权重与偏置？}
在实际部署时，若想直接拿原始未缩放特征 $x$ 做预测，只需把标准化变换代入预测式并展开整理即可。设我们在标准化特征上训练出了模型 $\hat{y} = \tilde{b} + \sum_{j=1}^d \tilde{w}_j \tilde{x}_j$，将 $\tilde{x}_j = (x_j - \mu_j)/\sigma_j$ 代入：
\[
  \hat{y} = \tilde{b} + \sum_{j=1}^d \tilde{w}_j \frac{x_j - \mu_j}{\sigma_j}
          = \underbrace{\left(\tilde{b} - \sum_{j=1}^d \frac{\tilde{w}_j \mu_j}{\sigma_j}\right)}_{b}
            + \sum_{j=1}^d \underbrace{\left(\frac{\tilde{w}_j}{\sigma_j}\right)}_{w_j} x_j.
\]
于是原始坐标系下的等价权重与偏置为
\begin{equation}\label{eq:ch6-recover-w}
  w_j = \frac{\tilde{w}_j}{\sigma_j},
  \qquad
  b = \tilde{b} - \sum_{j=1}^d \frac{\tilde{w}_j \mu_j}{\sigma_j}.
\end{equation}
请牢记一条第 2 卷还将反复强调的工程红线：\textbf{均值 $\mu_j$ 与标准差 $\sigma_j$ 只能从训练集计算}，验证集和测试集必须原封不动套用训练集保存下来的 $(\mu_j, \sigma_j)$，绝不能把测试集统计量混入预处理中（防止数据泄露）。

\subsection{从全批量梯度下降到 Mini-batch SGD 的直觉}

全批量梯度下降（Full-batch GD）每走一步都要对全部 $N$ 条样本求和：
\[
  \nabla f(w) = \frac{1}{N} \sum_{i=1}^N \nabla \ell_i(w),
\]
单步开销为 $O(ND)$。当训练集规模达到 $N = 10^6$ 甚至 $10^9$（如现代图像与语料库）时，全批量算完一次梯度要扫描几十 GB 数据，但参数只挪动了区区一小步——且相邻样本往往高度冗余，前 64 条样本算出的平均梯度方向其实已经与全集梯度八九不离十！

\textbf{小批量随机梯度下降（Mini-batch SGD）}顺势而生：每一步仅随机抽取一小批样本 $\mathcal{B}_t \subset \{1,\dots,N\}$（批量大小 $B = |\mathcal{B}_t|$ 通常取 $32, 64, 128, 256$），用小批量平均梯度代替真实全梯度：
\begin{equation}\label{eq:ch6-sgd}
  g_{\mathcal{B}_t}(w_t) = \frac{1}{B} \sum_{i\in\mathcal{B}_t} \nabla \ell_i(w_t),
  \qquad
  w_{t+1} = w_t - \eta_t\, g_{\mathcal{B}_t}(w_t).
\end{equation}
初读时只需把握 Mini-batch SGD 的三条核心直觉（关于凸目标下随机梯度的精确收敛率证明留待第 4 卷系统展开）：
\begin{enumerate}
  \item \textbf{方向平均无偏}：若每个样本等概率被抽入 $\mathcal{B}_t$，则随机梯度在平均意义上等于真实全梯度（$\E[g_{\mathcal{B}_t}(w)] = \nabla f(w)$），总体趋势依然朝向谷底。
  \item \textbf{以小噪声换取极速更新}：全批量走 1 步的时间，Mini-batch SGD 已经更新了 $N/B$ 步，早期下降速度往往快几个数量级。
  \item \textbf{谷底随机抖动与应对手段}：即使参数已经走到最优解 $w^*$ 附近，单个小批量的梯度 $g_{\mathcal{B}_t}(w^*)$ 依然不为零，因此固定步长 $\eta$ 下参数会在谷底附近持续小幅“嗡嗡震荡”。工程上常用的两大化解手段是\textbf{学习率衰减}（如 $\eta_t = \eta_0 / (1 + \alpha t)$，越靠近谷底步子迈得越小）以及引入\textbf{动量（Momentum）}（对历史梯度做指数滑动平均 $v_{t+1} = \beta v_t + g_{\mathcal{B}_t}$，既平滑随机抖动，又在狭长山谷的平缓轴上累积速度）。
\end{enumerate}

代码~\ref{lst:ch6-gd-sgd} 给出完整可运行的纯 NumPy 实现：对量纲悬殊 1000 倍的数据做 Z-score 标准化，分别用全批量 GD 和 Mini-batch SGD 求解线性回归，并按式~\eqref{eq:ch6-recover-w} 还原原始权重，与 \texttt{np.linalg.lstsq} 精确核对！

\begin{lstlisting}[caption={标准化全批量 GD 与 Mini-batch SGD 线性回归（对比 \texttt{lstsq} 校验）},label={lst:ch6-gd-sgd}]
import numpy as np

rng = np.random.default_rng(7)
N, D = 200, 2
X_raw = np.column_stack([rng.normal(0, 100.0, N), rng.normal(0, 0.1, N)])
w_true, b_true = np.array([0.03, -4.0]), 1.5
y = X_raw @ w_true + b_true + rng.normal(0, 0.2, N)

# 1. 解析解参考 (np.linalg.lstsq)
X_aug = np.column_stack([X_raw, np.ones(N)])
theta_ls, *_ = np.linalg.lstsq(X_aug, y, rcond=None)
w_ls, b_ls = theta_ls[:D], theta_ls[D]

# 2. Z-score 标准化（仅用训练集统计量）
mu, sigma = X_raw.mean(axis=0), X_raw.std(axis=0)
X_std = (X_raw - mu) / sigma

# 3. 标准化后的全批量梯度下降
w_s, b_s, eta = np.zeros(D), 0.0, 0.2
for _ in range(150):
    err = X_std @ w_s + b_s - y
    w_s -= eta * (X_std.T @ err) / N
    b_s -= eta * np.mean(err)

w_rec = w_s / sigma
b_rec = b_s - np.dot(w_s, mu / sigma)
assert np.allclose(w_rec, w_ls, atol=1e-6) and np.isclose(b_rec, b_ls, atol=1e-6)

# 4. 带学习率衰减的 Mini-batch SGD
w_sgd, b_sgd, B = np.zeros(D), 0.0, 32
for epoch in range(100):
    perm = rng.permutation(N)
    lr = 0.15 / (1.0 + 0.08 * epoch)
    for start in range(0, N, B):
        idx = perm[start:start + B]
        err = X_std[idx] @ w_sgd + b_sgd - y[idx]
        w_sgd -= lr * (X_std[idx].T @ err) / len(idx)
        b_sgd -= lr * np.mean(err)

w_sgd_rec = w_sgd / sigma
b_sgd_rec = b_sgd - np.dot(w_sgd, mu / sigma)
rel_err = np.linalg.norm(w_sgd_rec - w_ls) / np.linalg.norm(w_ls)
assert rel_err < 1e-2  # 实测相对误差 ~ 3.0e-3
\end{lstlisting}

% ====================================================================
\section{写梯度代码不出错的终极防线：有限差分梯度校验}
\label{sec:ch6-gradcheck}

从第 2 卷手写逻辑/Softmax 回归，到第 5 卷手写自动微分与反向传播，你会无数次把纸上的矩阵微积分公式敲成 NumPy 代码。梯度代码有一个令所有工程师头皮发麻的特点：\textbf{即使写错了一个转置、漏了一个系数 $2$ 或广播错了轴，损失函数往往仍然会“看起来在慢慢下降”，把 Bug 悄无声息地藏起来！}

怎么在 3 秒钟内 100\% 确信自己写的解析梯度函数 \texttt{grad\_f(w)} 分毫不差？答案是 Stanford CS231n 与工业界公认的终极防线——\textbf{中心有限差分梯度校验（Finite-Difference Gradient Checking）}。

\subsection{为什么中心差分比前向差分精确整整一个阶次？}

固定参数向量 $w\in\R^d$ 和第 $i$ 个标准基向量 $e_i$，考虑沿第 $i$ 坐标轴的一元切片函数 $g(\epsilon) \coloneqq f(w + \epsilon e_i)$（设 $f$ 三阶连续可微）。比较两种数值逼近偏导数 $\frac{\partial f}{\partial w_i}(w) = g'(0)$ 的公式：

\begin{thmBox}
\begin{proposition}[前向差分 $O(\epsilon)$ vs 中心差分 $O(\epsilon^2)$ 截断误差]\label{prop:ch6-fd-order}
对三阶连续可微函数 $f:\R^d\to\R$ 与微小步长 $\epsilon > 0$，记 $D_i^{\text{fwd}}$ 为前向差分、$D_i^{\text{cen}}$ 为中心差分：
\begin{align}
  D_i^{\text{fwd}}(w;\epsilon) &\coloneqq \frac{f(w + \epsilon e_i) - f(w)}{\epsilon}
    = \frac{\partial f}{\partial w_i}(w) + \frac{1}{2}\frac{\partial^2 f}{\partial w_i^2}(w)\,\epsilon + O(\epsilon^2), \label{eq:ch6-fwd-diff}\\
  D_i^{\text{cen}}(w;\epsilon) &\coloneqq \frac{f(w + \epsilon e_i) - f(w - \epsilon e_i)}{2\epsilon}
    = \frac{\partial f}{\partial w_i}(w) + \frac{1}{6}\frac{\partial^3 f}{\partial w_i^3}(w)\,\epsilon^2 + O(\epsilon^4). \label{eq:ch6-cen-diff}
\end{align}
中心差分通过正负对称相减，\textbf{精确消去了二阶截断误差项}，将精度从 $O(\epsilon)$ 提升为 $O(\epsilon^2)$。
\end{proposition}
\end{thmBox}

\begin{proof}
对 $g(\epsilon) = f(w+\epsilon e_i)$ 在 $0$ 处做三阶 Taylor 展开：
\begin{align*}
  g(+\epsilon) &= g(0) + g'(0)\epsilon + \frac{1}{2}g''(0)\epsilon^2 + \frac{1}{6}g'''(0)\epsilon^3 + O(\epsilon^4), \\
  g(-\epsilon) &= g(0) - g'(0)\epsilon + \frac{1}{2}g''(0)\epsilon^2 - \frac{1}{6}g'''(0)\epsilon^3 + O(\epsilon^4).
\end{align*}
第一式减 $g(0)$ 后除以 $\epsilon$ 立即得前向差分式~\eqref{eq:ch6-fwd-diff}（主误差项为 $\frac{1}{2}g''(0)\epsilon$）。将第一式减去第二式，常数项 $g(0)$ 与偶次项 $\frac{1}{2}g''(0)\epsilon^2$ 完美抵消，得
\[
  g(+\epsilon) - g(-\epsilon) = 2 g'(0)\epsilon + \frac{1}{3}g'''(0)\epsilon^3 + O(\epsilon^4),
\]
两边同除以 $2\epsilon$ 即得中心差分式~\eqref{eq:ch6-cen-diff}。
\end{proof}

\subsection{浮点舍入与截断误差的博弈：经典 V 形误差曲线}

既然中心差分截断误差是 $O(\epsilon^2)$，那是不是把 $\epsilon$ 设成 $10^{-14}$ 越小越好？\textbf{绝对不是！}在计算机 IEEE-754 双精度浮点数（\texttt{float64}，机器精度 $u \approx 2.22\times 10^{-16}$）下，每次函数求值都会携带量级约 $u|f(w)|$ 的舍入误差。当 $\epsilon$ 极小时，分子 $f(w+\epsilon e_i) - f(w-\epsilon e_i)$ 发生\textbf{大数相减的灾难性消去}，再除以极小的 $2\epsilon$，舍入误差被放大为 $O(u/\epsilon)$！

因此总数值误差呈现截断误差与舍入误差的拔河：
\[
  \text{总误差}_{\text{cen}}(\epsilon) \approx C_1 \epsilon^2 + C_2 \frac{u}{\epsilon}.
\]
对 $\epsilon$ 求导令其为零，立即算出 \texttt{float64} 下中心差分的理论最优步长：
\[
  \epsilon^*_{\text{cen}} \sim u^{1/3} \approx (2.22\times 10^{-16})^{1/3} \approx 6\times 10^{-6} \;\;(\text{工程常用 } 10^{-5}),
\]
此时最低误差可达 $O(u^{2/3}) \approx 10^{-11}$！同理，前向差分平衡 $O(\epsilon)$ 与 $O(u/\epsilon)$ 的最优步长在 $\epsilon^*_{\text{fwd}} \sim u^{1/2} \approx 10^{-8}$，最低误差仅能到达 $O(u^{1/2}) \approx 10^{-8}$——精度比中心差分足足差了 3 个数量级！

图~\ref{fig:ch6-gradcheck} 展示了 \texttt{code/vol1\_experiments.py} 对光滑非线性函数 $f(w) = \ln(1 + e^{a^T w}) + \frac{1}{4}\norm{w}_2^4$（$w\in\R^{10}$）的本机实测 V 形双对数曲线（读取自 \texttt{data/gradcheck\_eps.csv}）：
\begin{itemize}
  \item 在黄金步长 $\epsilon = 10^{-5}$ 处：前向差分相对误差为 $8.27\times 10^{-6}$，而中心差分相对误差低至 $\mathbf{2.62\times 10^{-11}}$，精度拉开近\textbf{六个数量级}！
  \item 当步长过小滑向 $\epsilon = 10^{-12}$ 时：浮点舍入误差接管全局，两者误差双双恶化到 $\sim 10^{-4}$！
\end{itemize}

\begin{figure}[htbp]
\centering
\begin{tikzpicture}
\begin{loglogaxis}[
  width=0.86\linewidth,
  height=6.5cm,
  xlabel={有限差分扰动步长 $\epsilon$（对数刻度）},
  ylabel={梯度相对误差 $\norm{g_{\text{num}} - g_{\text{ana}}}_2 / \norm{g_{\text{ana}}}_2$},
  xmin=7e-15, xmax=1.5e-1,
  ymin=1e-12, ymax=2e-1,
  grid=both,
  grid style={line width=0.25pt, draw=slate!20},
  major grid style={line width=0.4pt, draw=slate!35},
  tick label style={font=\scriptsize},
  label style={font=\footnotesize},
  legend style={font=\footnotesize, at={(0.03,0.05)}, anchor=south west,
    fill=white, fill opacity=0.92, draw=codeborder, inner sep=3pt},
  legend cell align=left
]
  \addplot[ochre, very thick, mark=square*, mark size=1.8pt]
    table[col sep=comma, x=eps, y=fwd_err]{data/gradcheck_eps.csv};
  \addlegendentry{前向差分 $O(\epsilon)$（右支斜率 $+1$）}

  \addplot[teal, very thick, mark=*, mark size=2.0pt]
    table[col sep=comma, x=eps, y=cen_err]{data/gradcheck_eps.csv};
  \addlegendentry{中心差分 $O(\epsilon^2)$（右支斜率 $+2$）}

  % Vertical reference line at eps = 1e-5
  \draw[navy, dashed, thick] (axis cs:1e-5, 1e-12) -- (axis cs:1e-5, 1e-1);
  \node[font=\scriptsize\sffamily, navy, anchor=south west, align=left,
    fill=white, fill opacity=0.92, text opacity=1, inner sep=1.5pt]
    at (axis cs:3.8e-5, 1.9e-12) {黄金步长 $\epsilon=10^{-5}$\\中心差分 $2.62\times 10^{-11}$};

  % Regime labels
  \node[font=\scriptsize\sffamily, slate, align=center, fill=white, fill opacity=0.88, text opacity=1, inner sep=1pt]
    at (axis cs:4e-11, 2.5e-3) {浮点舍入主导区\\$O(u/\epsilon)$（斜率 $-1$）};
  \node[font=\scriptsize\sffamily, slate, align=center, fill=white, fill opacity=0.88, text opacity=1, inner sep=1pt]
    at (axis cs:8e-6, 2.5e-3) {Taylor 截断主导区\\中心差分快得多};
\end{loglogaxis}
\end{tikzpicture}
\caption{\texttt{data/gradcheck\_eps.csv} 实测有限差分 V 形误差曲线（\texttt{float64}）：左侧为浮点消去舍入误差主导区，右侧为 Taylor 截断误差主导区；$\epsilon=10^{-5}$ 是中心差分的黄金甜蜜点，可将校验精度压低至 $10^{-11}$ 量级。}
\label{fig:ch6-gradcheck}
\end{figure}

\subsection{可直接复用到第 2--5 卷的通用梯度校验工具}

代码~\ref{lst:ch6-gradcheck} 实现了支持任意形状（向量、矩阵、高维张量）参数的 \texttt{numerical\_gradient} 与 \texttt{check\_gradient}，并立即拿它检验第 5.4 节推导的神经网络层梯度 $\nabla_W\big(\frac{1}{2}\norm{\tanh(WX)-Y}_F^2\big) = \big((A-Y)\had(1-A^2)\big)X^T$（本机实测相对误差低至 $5.01\times 10^{-11}$！）。

\begin{lstlisting}[caption={通用中心差分梯度校验工具（支持任意维度 NumPy 张量）},label={lst:ch6-gradcheck}]
import numpy as np

def numerical_gradient(f, w, eps=1e-5):
    """用中心有限差分计算标量函数 f 对任意形状数组 w 的数值梯度。"""
    w = np.asarray(w, dtype=np.float64).copy()
    grad = np.zeros_like(w)
    it = np.nditer(w, flags=["multi_index"])
    while not it.finished:
        idx = it.multi_index
        old_val = w[idx]
        w[idx] = old_val + eps
        fp = f(w)
        w[idx] = old_val - eps
        fm = f(w)
        w[idx] = old_val
        grad[idx] = (fp - fm) / (2.0 * eps)
        it.iternext()
    return grad

def check_gradient(f, grad_f, w, eps=1e-5, tol=1e-7):
    """校验解析梯度 grad_f(w) 的形状与数值正确性，返回 (相对误差, 是否通过)。"""
    g_ana = np.asarray(grad_f(w), dtype=np.float64)
    g_num = numerical_gradient(f, w, eps=eps)
    assert g_ana.shape == w.shape, f"Shape mismatch: {g_ana.shape} vs {w.shape}"
    denom = max(np.linalg.norm(g_ana) + np.linalg.norm(g_num), 1e-12)
    rel_err = float(np.linalg.norm(g_ana - g_num) / denom)
    return rel_err, rel_err < tol

# 实战校验第 5.4 节神经网络层反向传播公式：L(W) = 0.5 * ||tanh(W @ X) - Y||_F^2
rng = np.random.default_rng(123)
W0 = rng.normal(size=(3, 4))
X0 = rng.normal(size=(4, 5))
Y0 = rng.normal(size=(3, 5))

f_layer = lambda W: 0.5 * float(np.sum((np.tanh(W @ X0) - Y0) ** 2))
def grad_layer(W):
    A = np.tanh(W @ X0)
    dZ = (A - Y0) * (1.0 - A ** 2)   # Hadamard 穿过激活层
    return dZ @ X0.T                  # 迹技巧穿过线性层

rel_err, passed = check_gradient(f_layer, grad_layer, W0)
assert passed  # 实测 rel_err = 5.01e-11
\end{lstlisting}

\begin{failBox}{做梯度校验时的三大避坑铁律}
\begin{enumerate}
  \item \textbf{必须使用 \texttt{float64}，绝不用 \texttt{float32}}：\texttt{float32} 的机器精度仅 $u\approx 1.19\times 10^{-7}$，中心差分最低误差只能卡在 $10^{-4}$ 附近，根本分不清代码有没有写错！
  \item \textbf{避开不可微折点（Kinks）}：如果函数包含 $\operatorname{ReLU}(z)=\max(0,z)$ 或 $|z|$，当某个分量恰好落在 $z=0$ 附近 $|z|<\epsilon$ 时，$w+\epsilon e_i$ 与 $w-\epsilon e_i$ 会跨过折点两侧，导致数值差分与解析梯度出现假警报。校验时请用小规模随机高斯权重并避开零点。
  \item \textbf{只在小尺寸张量上做单元测试}：对 $D$ 维参数做中心差分需要调用前向函数 $2D$ 次，千万不要在正式训练循环里跑梯度校验！正确做法是在 $D\sim 10$--$50$ 的迷你测试数据上跑通 \texttt{check\_gradient}（相对误差 $<10^{-7}$），再放心地把同一套解析梯度代码投入大规模训练。
\end{enumerate}
\end{failBox}

% ====================================================================
\section{全卷收束与第 2 卷预告}
\label{sec:ch6-wrapup}

回顾第 1 卷走过的六章阶梯，我们已经把大一数理与 C 语言基础彻底升级为现代人工智能的\textbf{确定性计算语言}：
\begin{itemize}
  \item \textbf{第 1 章（语言与算力）}：打通 Python/NumPy 向量化、广播排错与矩阵乘法四视角（尤其是外积与张量批处理）；
  \item \textbf{第 2 章（几何与回归）}：掌握高维近正交直觉、四大子空间正交投影，并亲手跑出线性最小二乘与岭回归；
  \item \textbf{第 3--4 章（谱分解与降维）}：用实对称谱定理看穿 PSD 矩阵与条件数 $\cond$，用 SVD 与 PCA 揭示数据压缩和最小范数解的几何本质；
  \item \textbf{第 5--6 章（微积分与优化闭环）}：建立 Jacobian 反向 VJP 与矩阵微分迹技巧，通过二次型谱动力学理解步长阈值、特征标准化与 Mini-batch SGD，并装备好有限差分梯度校验武器。
\end{itemize}

\begin{insightBox}{为什么要迈向第 2 卷：《AI 概率论、数理统计、信息论与首个回归/分类闭环》？}
第 1 卷已经能优雅地拟合数据并运行梯度下降，但三个根本问题在确定性线代与微积分框架内\textbf{无法回答}：
\begin{enumerate}
  \item \textbf{为什么回归用平方损失、岭回归加 $L_2$ 惩罚？}第 2 卷将用\textbf{极大似然估计（MLE）与最大后验估计（MAP）}揭示：最小二乘等价于高斯噪声假设下的 MLE，而 $L_2$ 岭回归正是权重服从高斯先验时的 MAP！
  \item \textbf{如何处理“是/否”与多分类问题？}面对离散类别标签，平方损失不再合适；我们需要概率分布、\textbf{信息熵、交叉熵与 KL 散度}，从而自然导出\textbf{逻辑回归与 Softmax 分类}。
  \item \textbf{如何量化模型泛化可靠性？}第 2 卷将把概率统计工具与本卷的梯度下降、有限差分校验合流，用纯 NumPy 打造一套完整的回归/分类训练、交叉验证与精确评测闭环！
\end{enumerate}
\end{insightBox}

% ====================================================================
\section{本章习题}
\label{sec:ch6-exercises}

\exTier{A. 概念辨析与反例查错}{戳破调参与预处理中最常见的直觉误区}
\begin{enumerate}[label=\textbf{6.\arabic*}]
  \item \textbf{判断正误并纠错}：
  (a) “梯度下降在每一步都沿当前负梯度方向移动，因此只要步长 $\eta > 0$，目标函数值 $f(w_{t+1})$ 就一定小于 $f(w_t)$。”
  (b) “在用中心有限差分做梯度校验时，步长 $\epsilon$ 越小结果越准确，因此取 $\epsilon = 10^{-14}$ 远优于 $\epsilon = 10^{-5}$。”
  请判断真假，并分别用本章的具体例子或误差机理说明原因。

  \item \textbf{特征标准化与数据泄露排错}：某同学在做线性回归实验时写了如下代码：先对包含训练集和测试集的整个数据矩阵 \texttt{X\_all} 计算整体均值与标准差 \texttt{mu, sigma} 并标准化，再划分训练集与测试集；训练得到权重 \texttt{w\_std} 后，在部署上线时直接把原始未缩放特征 \texttt{x\_new} 乘以 \texttt{w\_std} 输出预测值。请指出这段流程中的\textbf{两处严重错误}，并写出正确做法。
\end{enumerate}

\exTier{B. 具体小例手算}{亲手算通步长阈值、最优收缩率与权重还原}
\begin{enumerate}[label=\textbf{6.\arabic*}, resume]
  \item \textbf{非对角二次型的步长红线与最优收敛率}：设二次型函数 $f(w) = \frac{1}{2}w^T Q w - b^T w$，其中
  \[
    Q = \begin{bmatrix} 5 & 3 \\ 3 & 5 \end{bmatrix}, \qquad b = \begin{bmatrix} 8 \\ 8 \end{bmatrix}.
  \]
  (a) 求 $Q$ 的特征值 $\lambda_{\max}, \lambda_{\min}$、条件数 $\cond(Q)$ 与唯一极小值点 $w^*$；
  (b) 求梯度下降对任意初值都收敛的最大允许步长上限 $\eta_{\max} = 2/\lambda_{\max}$、最优固定步长 $\eta^*$ 及最优单步误差收缩因子 $\rho^*$；
  (c) 若从 $w_0 = (0, 0)^T$ 出发，分别手算采用最优步长 $\eta^*$ 与步长 $\eta = 1/8$ 时的第一步结果 $w_1$，并从特征坐标角度解释为何 $\eta = 1/8$ 竟能\textbf{一步直接命中} $w^*$（提示：观察 $e_0 = w_0 - w^*$ 落在 $Q$ 的哪一个特征方向上！）。

  \item \textbf{标准化后的权重还原计算}：某房价线性回归模型包含两个原始特征 $x_1$（面积，训练集均值 $\mu_1 = 100$，标准差 $\sigma_1 = 25$）和 $x_2$（房龄，均值 $\mu_2 = 10$，标准差 $\sigma_2 = 5$）。在 Z-score 标准化特征上运行梯度下降得到模型 $\hat{y} = 300 + 50\tilde{x}_1 - 15\tilde{x}_2$。请按式~\eqref{eq:ch6-recover-w} 手算出等价的原始特征线性模型 $\hat{y} = b + w_1 x_1 + w_2 x_2$，并用一条具体样本 $(x_1, x_2) = (125, 5)$ 验证两个表达式预测值完全相等。
\end{enumerate}

\exTier{C. AI 核心推演与实用证明}{夯实数值分析与优化动力学的推导功底}
\begin{enumerate}[label=\textbf{6.\arabic*}, resume]
  \item \textbf{带 $L_2$ 正则化对 Hessian 条件数的改善证明}：设线性回归的原始 Hessian 矩阵 $H_0 = \frac{1}{N}X^T X \succeq 0$ 的最大与最小特征值分别为 $\lambda_{\max} > \lambda_{\min} > 0$（条件数 $\cond_0 = \lambda_{\max}/\lambda_{\min}$）。现加入岭回归正则项 $\frac{\alpha}{2}\norm{w}_2^2$（$\alpha > 0$），新 Hessian 变为 $H_\alpha = H_0 + \alpha I$。
  证明：新条件数 $\cond(H_\alpha) = \dfrac{\lambda_{\max}+\alpha}{\lambda_{\min}+\alpha}$ 严格小于 $\cond_0$，并说明随着 $\alpha$ 增大，梯度下降的最优单步收缩因子 $\rho^*(\alpha)$ 严格单调减小（即正则化不仅抑制过拟合，还同时加速梯度下降收敛！）。
  \exHint{由第 3 章谱映射性质，$H_0 + \alpha I$ 的特征值为 $\lambda_i + \alpha$；直接作差 $\cond_0 - \cond(H_\alpha) = \frac{\alpha(\lambda_{\max}-\lambda_{\min})}{\lambda_{\min}(\lambda_{\min}+\alpha)} > 0$。}

  \item \textbf{中心差分最优步长 $\epsilon^* \sim u^{1/3}$ 的定量推导}：设函数值计算引入的绝对浮点舍入误差上界为 $M u$（其中 $u\approx 2.22\times 10^{-16}$ 为双精度机器精度），三阶导数绝对值上界为 $K = |g'''(0)| > 0$。结合命题~\ref{prop:ch6-fd-order} 写出中心差分总误差上界
  \[
    E(\epsilon) = \frac{K}{6}\epsilon^2 + \frac{M u}{\epsilon}.
  \]
  求使 $E(\epsilon)$ 最小的最优步长 $\epsilon^*$ 及最小误差 $E(\epsilon^*)$ 的显式表达式；当 $K \sim M \sim 1$ 时估算 $\epsilon^*$ 与 $E(\epsilon^*)$ 的数量级，并与图~\ref{fig:ch6-gradcheck} 中的实测底部对比。
  \exHint{令 $E'(\epsilon) = \frac{K}{3}\epsilon - \frac{M u}{\epsilon^2} = 0$ 解出 $\epsilon^* = (3Mu/K)^{1/3}$，再代回 $E(\epsilon^*)$。}
\end{enumerate}

\exTier{D. 可自检编程实验}{动手验证病态加速与完整梯度校验闭环}
\begin{enumerate}[label=\textbf{6.\arabic*}, resume]
  \item \textbf{复现特征标准化对梯度下降的加速奇迹}：使用固定随机种子生成 $N=300$、特征维度 $D=3$ 的合成回归数据，三个特征标准差分别取 $(1000, 1, 0.01)$，真实权重 $w_{\text{true}} = (0.002, 3.0, -50.0)^T$、偏置 $b_{\text{true}} = 1.0$，叠加标准差 $0.1$ 的高斯噪声。
  \begin{enumerate}[label=(\alph*)]
    \item 计算原始 Hessian $\frac{1}{N}X_{\text{raw}}^T X_{\text{raw}}$ 与 Z-score 标准化 Hessian $\frac{1}{N}X_{\text{std}}^T X_{\text{std}}$ 的条件数；
    \item 对标准化特征分别按理论最优步长 $\eta^* = 2/(\lambda_{\max}+\lambda_{\min})$ 与手动步长 $\eta = 0.2$ 运行全批量梯度下降 100 步，并还原原始权重 $(w_{\text{rec}}, b_{\text{rec}})$。
  \end{enumerate}
  \textbf{自检标准}：断言原始条件数 $> 10^8$、标准化条件数 $< 2.0$，且用最优步长 $\eta^*$ 迭代 30 步内还原后的权重满足 $\norm{w_{\text{rec}} - w_{\text{ls}}}_2 / \norm{w_{\text{ls}}}_2 < 10^{-8}$（$w_{\text{ls}}$ 为 \texttt{np.linalg.lstsq} 参考解）。
  \exHint{加入偏置项一起做标准化矩阵时，可以把 $\tilde{X}_{\text{aug}} = [X_{\text{std}},\, \mathbf{1}]$ 作为整体设计矩阵计算 $\frac{1}{N}\tilde{X}_{\text{aug}}^T \tilde{X}_{\text{aug}}$ 的最大最小特征值，直接给出含偏置的精确最优步长 $\eta^*$！}

  \item \textbf{用中心差分校验岭回归与 Softplus 损失梯度}：编写解析梯度并调用代码~\ref{lst:ch6-gradcheck} 中的 \texttt{check\_gradient}，对以下两个第 2 卷高频损失函数完成自动化梯度校验（取 $N=40, D=8$ 随机 \texttt{float64} 数据）：
  \begin{enumerate}[label=(\alph*)]
    \item 岭回归损失：$L_1(w) = \frac{1}{2N}\norm{Xw - y}_2^2 + \frac{\lambda}{2}\norm{w}_2^2$（$\lambda = 0.5$）；
    \item 二分类对数损失（逻辑回归雏形，标签 $y_i\in\{0,1\}$）：
    \[
      L_2(w) = \frac{1}{N}\sum_{i=1}^N \Big(\ln\big(1 + e^{x_i^T w}\big) - y_i (x_i^T w)\Big) + \frac{\lambda}{2}\norm{w}_2^2,
    \]
    提示：其解析梯度为 $\nabla_w L_2(w) = \frac{1}{N}X^T(\sigma(Xw) - y) + \lambda w$（其中 $\sigma(z)=1/(1+e^{-z})$），计算 $\ln(1+e^z)$ 时请使用数值稳定的 \texttt{np.logaddexp(0.0, z)}！
  \end{enumerate}
  \textbf{自检标准}：两个损失函数在 $\epsilon = 10^{-5}$ 下的相对误差均需严格小于 $10^{-9}$；同时故意把梯度中的一个符号改错（例如把 $\lambda w$ 写成 $-\lambda w$），验证你的 \texttt{check\_gradient} 能立刻报警返回 \texttt{passed == False}。
  \exHint{这就是第 2 卷逻辑回归的核心梯度公式！养成“每写完一个新损失函数立刻跑一遍 \texttt{check\_gradient}”的习惯，未来手写神经网络能省下无数小时排错时间。}
\end{enumerate}


% ==================== Back Matter / Appendix ====================
\appendix
\def\inappendix{}
\chapter{全卷 48 道分层习题完整精解}
\label{app:solutions}

\chapEpigraph{学习数学的唯一方法，就是去做数学。}{保罗·哈尔莫斯（Paul R. Halmos）}

\begin{chapIntro}{使用指南}
本附录为第 1--6 章全部 \textbf{48 道四层递进习题}（A 概念反例 / B 手算推演 / C 实用证明 / D 可自检编程）提供完整解答、推导细节与数值自检基准。建议合上附录、独立完成后再逐项对照，尤其注意 A 类反例的边界条件与 D 类断言阈值。
\end{chapIntro}

% ====================================================================
\section{第 1 章习题解答（1.1--1.8）}
% ====================================================================

\textbf{1.1 [A]}
(1) 按广播三步规则（右对齐、维长相等或一侧为 $1$）：
(a) \texttt{(8,1,5)} 与 \texttt{(4,5)} $\to$ 补前导 $1$ 得 \texttt{(1,4,5)}，合法，输出 \texttt{(8,4,5)}；
(b) \texttt{(8,3,5)} 与 \texttt{(3,1)} $\to$ 补成 \texttt{(1,3,1)}，合法，输出 \texttt{(8,3,5)}；
(c) \texttt{(10,4)} 与 \texttt{(10,)} $\to$ 补成 \texttt{(1,10)}，末尾维 $4\neq 10$，\textbf{冲突报错}；
(d) \texttt{(64,1)} 与 \texttt{(64,)} $\to$ 补成 \texttt{(1,64)}，合法，输出 \texttt{(64,64)}。\\
(2) \texttt{y\_pred - y\_true} 静默广播成 $(64,64)$ 矩阵，其中第 $(i,j)$ 项为 $\hat{y}_i - y_j$（把第 $i$ 个预测与全部 $64$ 个样本标签两两相减），导致求出的 MSE 严重偏大且不报错。修复方式：
(i) 统一一维化 \texttt{y\_pred.reshape(-1) - y\_true.reshape(-1)}；
(ii) 显式断言 \texttt{assert y\_pred.shape == y\_true.shape}。

\textbf{1.2 [A]}
输出为 \texttt{A[0, 0] = 1.0}（原数组被污染！），\texttt{C = [[1.0, 0.0], [0.0, 3.0]]}。两处错误：
(1) \texttt{A[:]} 返回原内存的\textbf{视图（view）}，原地除法 \texttt{/=} 同时改写了 \texttt{A}，应改为 \texttt{A\_scaled = A.copy()}；
(2) \texttt{*} 是逐元素 Hadamard 乘积而非矩阵乘法，应改为 \texttt{C = A\_scaled @ B}（正确结果为 $\begin{psmallmatrix}1.0 & 4.0\\ 0.5 & 3.0\end{psmallmatrix}$）。

\textbf{1.3 [B]}
$A = \begin{psmallmatrix}1&2\\0&-1\\3&1\end{psmallmatrix},\; B = \begin{psmallmatrix}2&-1\\1&3\end{psmallmatrix}$。
\begin{enumerate}[label=(\arabic*)]
  \item \textbf{列线性组合}：第 1 列 $2\begin{psmallmatrix}1\\0\\3\end{psmallmatrix} + 1\begin{psmallmatrix}2\\-1\\1\end{psmallmatrix} = \begin{psmallmatrix}4\\-1\\7\end{psmallmatrix}$；第 2 列 $(-1)\begin{psmallmatrix}1\\0\\3\end{psmallmatrix} + 3\begin{psmallmatrix}2\\-1\\1\end{psmallmatrix} = \begin{psmallmatrix}5\\-3\\0\end{psmallmatrix}$。
  \item \textbf{行线性组合}：第 1 行 $1(2,-1)+2(1,3)=(4,5)$；第 2 行 $0(2,-1)+(-1)(1,3)=(-1,-3)$；第 3 行 $3(2,-1)+1(1,3)=(7,0)$。
  \item \textbf{外积秩 1 求和}：
  \[
    AB = \begin{pmatrix}1\\0\\3\end{pmatrix}\begin{pmatrix}2 & -1\end{pmatrix}
       + \begin{pmatrix}2\\-1\\1\end{pmatrix}\begin{pmatrix}1 & 3\end{pmatrix}
       = \begin{pmatrix}2&-1\\0&0\\6&-3\end{pmatrix}
       + \begin{pmatrix}2&6\\-1&-3\\1&3\end{pmatrix}
       = \begin{pmatrix}4&5\\-1&-3\\7&0\end{pmatrix}.
  \]
  每个外积项的每一列都是左向量的标量倍，非零故列空间维数为 $1$，秩必为 $1$。
\end{enumerate}

\textbf{1.4 [B]}
(1) 沿样本轴 \texttt{axis=0} 聚合：\texttt{mean = [[4., 15., 4.]]}，\texttt{std = [[2., 5., 0.]]}（注意第 3 列常值特征方差为 $0$，实际标准化时须防除以零），去中心化矩阵为 $\begin{psmallmatrix}-2 & -5 & 0\\ 2 & 5 & 0\end{psmallmatrix}$。\\
(2) 翻译结果：
(a) $y_i = \sum_j A_{ij} x_j$，等价于 \texttt{A @ x}，形状 \texttt{(m,)}；
(b) $c_n = \sum_i U_{ni} V_{ni}$，等价于 \texttt{np.sum(U * V, axis=1)}，形状 \texttt{(N,)}；
(c) $Y_{btm} = \sum_d H_{btd} W_{dm}$，等价于 \texttt{H @ W}，形状 \texttt{(B, T, M)}。

\textbf{1.5 [C]（证明）}
\begin{enumerate}[label=(\arabic*)]
  \item 矩阵 $M = uv^T$ 的第 $j$ 列为 $v_j u \in \spanop\{u\}$，故 $\col(M) \subseteq \spanop\{u\}$，从而 $\rank(M) \le 1$；又因 $u\neq 0, v\neq 0$，存在 $i,j$ 使 $u_i v_j \neq 0$，即 $M \neq 0$，故 $\rank(M) = 1$。
  \item 当 $m=n$ 时，$\tr(uv^T) = \sum_{i=1}^n (uv^T)_{ii} = \sum_{i=1}^n u_i v_i = u^T v$。
  \item 由外积求和视角 $AB = \sum_{k=1}^p a_{\cdot k} b_{k\cdot}^T$，每项秩至多为 $1$，由秩的次可加性 $\rank(X+Y)\le \rank(X)+\rank(Y)$ 立得 $\rank(AB) \le \sum_{k=1}^p \rank(a_{\cdot k} b_{k\cdot}^T) \le p$。$\square$
\end{enumerate}

\textbf{1.6 [C]}
方案一先算 $BA \in \R^{d\times d}$ 需 $d^2 r$ 次乘法，再乘 $x$ 需 $d^2$ 次，合计 $d^2(r+1)$；方案二先算 $Ax \in \R^r$ 需 $dr$ 次，再左乘 $B$ 需 $dr$ 次，合计 $2dr$。加速比为 $\frac{d^2(r+1)}{2dr} = \frac{d(r+1)}{2r}$。代入 $d=4096, r=16$ 得 $\frac{4096\times 17}{32} = 2176$ 倍！这正是大模型 LoRA 推理与训练中绝不显式展开 $BA$ 的原因。

\textbf{1.7--1.8 [D]（编程自检要点）}
\begin{itemize}
  \item \textbf{1.7}：展开法核心公式为 \texttt{D2 = np.sum(X**2, axis=1, keepdims=True)} + \texttt{np.sum(Y**2, axis=1)} $-$ \texttt{2.0 * (X @ Y.T)}，再取 \texttt{np.sqrt(np.maximum(D2, 0.0))} 防浮点负零导致 \texttt{NaN}。$N=M=1000, D=64$ 时广播法分配 $1000\times 1000\times 64$ 约 $512\,\text{MB}$ 临时数组，而展开法走 BLAS 仅占 $8\,\text{MB}$，快数十倍以上。
  \item \textbf{1.8}：数值稳定行 Softmax 分三步：\texttt{S = S - S.max(axis=-1, keepdims=True)}，随后 \texttt{W = np.exp(S)}，再除以 \texttt{W.sum(axis=-1, keepdims=True)}。确认 \texttt{@} 版本（\texttt{Q @ K.swapaxes(-1,-2) / np.sqrt(D\_h)}）与 \texttt{einsum} 版本差 $<10^{-12}$，且加 $1000$ 后输出逐位一致。
\end{itemize}

% ====================================================================
\section{第 2 章习题解答（2.1--2.8）}
% ====================================================================

\textbf{2.1 [A]}
(1) 不是范数。取 $u=(1,0)^T, v=(0,1)^T$，有 $f(u)=1, f(v)=1$，但 $f(u+v) = (\sqrt{1}+\sqrt{1})^2 = 4 > f(u)+f(v)=2$，违背三角不等式（对应非凸单位球）。\\
(2) 取 $P = \begin{psmallmatrix}1&1\\0&0\end{psmallmatrix}$，显然 $P^2=P$ 但 $P^T\neq P$。取 $b=(0,1)^T$，得 $Pb=(1,0)^T$、$b-Pb=(-1,1)^T$，内积 $\inner{b-Pb}{Pb}=-1\neq 0$（斜投影）。\\
(3) 错误。$\cos\theta(u,v)=1$ 仅蕴含同向共线，即存在正标量 $\alpha>0$ 使 $u = \alpha v$（如 $u=(2,0)^T, v=(1,0)^T$）。

\textbf{2.2 [A]}
(1) 设 $A$ 第 1、2 列相同，则 $z=(1,-1,0,\dots,0)^T \in \nul(A)$。若 $\hat{w}$ 是最小二乘解，则 $\hat{w}+tz$ 对任意 $t\in\R$ 有相同拟合向量 $A(\hat{w}+tz)=A\hat{w}$，故\textbf{$\hat{w}$ 不唯一}；但由正交分解唯一性，$\hat{b}=A\hat{w}$ 是 $b$ 在 $\col(A)$ 上的唯一正交投影，\textbf{$\hat{b}$ 必唯一}。\\
(2) 疏漏一：用训练 MSE 而非未见测试集 MSE 选模型；疏漏二：忽视了权重范数 $\|w_{\text{OLS}}\|_2 \approx 9.94\times 10^6$ 的剧烈爆炸——高次多项式靠相邻高次项巨额正负抵消强行穿过 15 个含噪点，导致区间内剧烈震荡。

\textbf{2.3 [B]}
对 $A = \begin{psmallmatrix}1&0&1\\0&1&1\\1&1&2\end{psmallmatrix}$ 做行化简得 $\begin{psmallmatrix}1&0&1\\0&1&1\\0&0&0\end{psmallmatrix}$，$\rank(A)=2$。
\begin{itemize}
  \item $\row(A) = \spanop\{(1,0,1)^T,\; (0,1,1)^T\}$（维数 $2$）；$\nul(A) = \spanop\{(1,1,-1)^T\}$（维数 $1$），验证两两内积为 $0$，$2+1=3$。
  \item $\col(A) = \spanop\{(1,0,1)^T,\; (0,1,1)^T\}$（维数 $2$）；$\nul(A^T) = \spanop\{(1,1,-1)^T\}$（维数 $1$），验证两两内积为 $0$，$2+1=3$。
\end{itemize}

\textbf{2.4 [B]}
$A = \begin{psmallmatrix}1&-1\\1&0\\1&1\\1&2\end{psmallmatrix},\; b=\begin{psmallmatrix}0\\1\\1\\3\end{psmallmatrix}$。
(1) 正规方程：$A^TA = \begin{psmallmatrix}4&2\\2&6\end{psmallmatrix},\; A^Tb = \begin{psmallmatrix}5\\7\end{psmallmatrix}$，解得 $\hat{w} = (\frac{4}{5},\; \frac{9}{10})^T = (0.8, 0.9)^T$。拟合值 $(-0.1, 0.8, 1.7, 2.6)^T$，残差 $e = (0.1, 0.2, -0.7, 0.4)^T$，验算 $A^T e = (0.1+0.2-0.7+0.4,\; -0.1+0-0.7+0.8)^T = (0,0)^T$。\\
(2) Gram-Schmidt：$q_1 = \frac{1}{2}(1,1,1,1)^T$，$r_{11}=2$；$r_{12} = q_1^T a_2 = 1$，$u_2 = a_2 - q_1 = (-\frac{3}{2},-\frac{1}{2},\frac{1}{2},\frac{3}{2})^T$，$\|u_2\|_2 = \sqrt{5} = r_{22}$，$q_2 = \frac{1}{2\sqrt{5}}(-3,-1,1,3)^T$。得 $R = \begin{psmallmatrix}2&1\\0&\sqrt{5}\end{psmallmatrix}$，$Q^T b = (\frac{5}{2},\; \frac{9}{2\sqrt{5}})^T$。回代 $R\hat{w} = Q^T b$ 得 $w_1 = \frac{9}{10},\; w_0 = \frac{4}{5}$，完全一致。

\textbf{2.5 [C]（证明）}
(1) 设 $k$ 为使 $|x_k| = \norm{x}_\infty$ 的下标，则 $\norm{x}_\infty^2 = x_k^2 \le \sum_{i=1}^n x_i^2 = \norm{x}_2^2 \implies \norm{x}_\infty \le \norm{x}_2$（取 $e_1$ 取等）；再因 $\norm{x}_2^2 = \sum x_i^2 \le (\sum |x_i|)^2 = \norm{x}_1^2 \implies \norm{x}_2 \le \norm{x}_1$（取 $e_1$ 取等）；对全 $1$ 向量 $\mathbf{1}$ 与 $|x|$ 应用 Cauchy-Schwarz：$\norm{x}_1 = \inner{\mathbf{1}}{|x|} \le \norm{\mathbf{1}}_2 \norm{x}_2 = \sqrt{n}\norm{x}_2$（取 $\mathbf{1}$ 取等）。\\
(2) $(I-P)^T = I-P^T = I-P$，$(I-P)^2 = I-2P+P^2 = I-P$，故 $I-P$ 为正交投影矩阵。又 $\inner{Pb}{(I-P)b} = b^T P^T(I-P)b = b^T(P-P^2)b = 0$，由勾股定理立得 $\norm{b}_2^2 = \norm{Pb}_2^2 + \norm{(I-P)b}_2^2$。$\square$

\textbf{2.6 [C]（证明）}
由最优性定义，分别比较 $\hat{w}_{\lambda_1}$ 与 $\hat{w}_{\lambda_2}$ 在两个目标上的取值：
\begin{align*}
  \norm{A\hat{w}_{\lambda_1}-y}_2^2 + \lambda_1 \norm{\hat{w}_{\lambda_1}}_2^2
  &\le \norm{A\hat{w}_{\lambda_2}-y}_2^2 + \lambda_1 \norm{\hat{w}_{\lambda_2}}_2^2, \\
  \norm{A\hat{w}_{\lambda_2}-y}_2^2 + \lambda_2 \norm{\hat{w}_{\lambda_2}}_2^2
  &\le \norm{A\hat{w}_{\lambda_1}-y}_2^2 + \lambda_2 \norm{\hat{w}_{\lambda_1}}_2^2.
\end{align*}
两式相加消去拟合残差项得 $(\lambda_2 - \lambda_1)\norm{\hat{w}_{\lambda_2}}_2^2 \le (\lambda_2 - \lambda_1)\norm{\hat{w}_{\lambda_1}}_2^2$。因 $\lambda_2 > \lambda_1$，立即得到 $\norm{\hat{w}_{\lambda_2}}_2 \le \norm{\hat{w}_{\lambda_1}}_2$。代回第一式移项即得 $\norm{A\hat{w}_{\lambda_2}-y}_2^2 - \norm{A\hat{w}_{\lambda_1}-y}_2^2 \ge \lambda_1(\norm{\hat{w}_{\lambda_1}}_2^2 - \norm{\hat{w}_{\lambda_2}}_2^2) \ge 0$。$\square$

\textbf{2.7--2.8 [D]（编程自检要点）}
\begin{itemize}
  \item \textbf{2.7}：$\varepsilon=10^{-8}$ 时 $A^TA = \begin{psmallmatrix}1+10^{-16} & 1\\ 1 & 1+10^{-16}\end{psmallmatrix}$ 在 float64 中退化为全 $1$ 奇异矩阵，\texttt{solve} 抛 \texttt{LinAlgError}；而 \texttt{np.linalg.qr} 与 \texttt{lstsq} 不平方条件数，全程保持误差 $<10^{-14}$。
  \item \textbf{2.8}：(1) 实测 $d=1000$ 时标准差 $0.0317$（理论 $0.0316$），$P(|\cos|<0.1)=99.9\%$；(2) 增广矩阵 $A_{\text{aug}} = \begin{psmallmatrix}A\\ \sqrt{\lambda} P\end{psmallmatrix}$（$P$ 去掉首行或首对角置 $0$）用 \texttt{lstsq} 求解，$\lambda=10^{-3}$ 时测试 MSE $= 0.02161$，$\norm{\hat{w}}_2 = 14.73$。
\end{itemize}

% ====================================================================
\section{第 3 章习题解答（3.1--3.8）}
% ====================================================================

\textbf{3.1 [A]}
取非对称阵 $A = \begin{psmallmatrix}1&4\\0&1\end{psmallmatrix}$，特征值均为 $1 > 0$，但对称化 $S = \frac{1}{2}(A+A^T) = \begin{psmallmatrix}1&2\\2&1\end{psmallmatrix}$ 特征值为 $3$ 与 $-1$，在 $x=(1,-1)^T$ 处 $x^T A x = x^T S x = -2 < 0$（根源：二次型只感知对称部分，非对称阵特征值不能判定二次型符号）。\\
半正定反例：取 $Q = \begin{psmallmatrix}0&0\\0&-1\end{psmallmatrix}$，顺序主子式 $\Delta_1=0, \Delta_2=0$ 均 $\ge 0$，但 $e_2^T Q e_2 = -1 < 0$（负半定！半正定须检查\textbf{全部主子式} $\ge 0$）。

\textbf{3.2 [A]}
(a) \textbf{假}。取 $A = \begin{psmallmatrix}1&1\\0&0\end{psmallmatrix}$（特征值 $1,0$，可对角化）与 $B = \begin{psmallmatrix}0&0\\0&1\end{psmallmatrix}$（对角阵），则 $A+B = \begin{psmallmatrix}1&1\\0&1\end{psmallmatrix}$ 不可对角化。\\
(b) \textbf{真}（实对称阵谱定理引理二已证）。\\
(c) \textbf{假}。反例 $J_2(0.9)=\begin{psmallmatrix}0.9&1\\0&0.9\end{psmallmatrix}$，特征值全为 $0.9<1$，但 $\|J_2(0.9)^k\|_F$ 先瞬态放大至 $k=9$ 的 $3.9128$ 才衰减。\\
(d) \textbf{真}。$P_1, P_2$ 元素非负故 $P_1P_2$ 非负，且 $(P_1 P_2)\mathbf{1} = P_1(P_2\mathbf{1}) = P_1\mathbf{1} = \mathbf{1}$。

\textbf{3.3 [B]}
(1) 特征多项式 $(3-\lambda)^2 - 1 = 0 \implies \lambda_1=4,\; \lambda_2=2$。单位特征向量 $q_1 = \frac{1}{\sqrt{2}}\binom{1}{-1},\; q_2 = \frac{1}{\sqrt{2}}\binom{1}{1}$，谱分解 $S = 4 q_1 q_1^T + 2 q_2 q_2^T$。\\
(2) 条件数 $\cond(S) = 4/2 = 2$。在特征坐标 $y = Q^T x$ 下 $4y_1^2 + 2y_2^2 = 8 \iff \frac{y_1^2}{2} + \frac{y_2^2}{4} = 1$：短半轴沿 $q_1$ 长 $\sqrt{2}$，长半轴沿 $q_2$ 长 $2$。\\
(3) $S^{10}$ 最大特征值为 $4^{10} = 1{,}048{,}576$；$S^{-1} = \frac{1}{4}q_1 q_1^T + \frac{1}{2}q_2 q_2^T = \frac{1}{8}\begin{psmallmatrix}3&1\\1&3\end{psmallmatrix}$。

\textbf{3.4 [B]}
(1) $W = \begin{psmallmatrix}0&1&0\\1&0&1\\0&1&0\end{psmallmatrix}$，$D = \diag(1,2,1)$，$L = \begin{psmallmatrix}1&-1&0\\-1&2&-1\\0&-1&1\end{psmallmatrix}$。\\
(2) 特征多项式 $-\lambda(\lambda-1)(\lambda-3)=0 \implies \lambda_1=0, \lambda_2=1, \lambda_3=3$（均 $\ge 0$ 且 $0$ 为单根，对应图连通）。单位特征基：$q_1 = \frac{1}{\sqrt{3}}(1,1,1)^T$，$q_2 = \frac{1}{\sqrt{2}}(1,0,-1)^T$，$q_3 = \frac{1}{\sqrt{6}}(1,-2,1)^T$。\\
(3) $P = D^{-1}W = \begin{psmallmatrix}0&1&0\\0.5&0&0.5\\0&1&0\end{psmallmatrix}$，每行非负和为 $1$。解 $\pi^T P = \pi^T$ 且 $\sum \pi_i = 1$（或按度数比例 $\pi_i \propto d_i$）立得平稳分布 $\pi^* = (\frac{1}{4}, \frac{1}{2}, \frac{1}{4})$。

\textbf{3.5 [C]（证明）}
(1) 对任意 $x\in\R^n$，$x^T(\alpha A+\beta B)x = \alpha(x^TAx) + \beta(x^TBx) \ge 0$，且 $\alpha A+\beta B$ 对称，故半正定。\\
(2) $H$ 对称且对 $x\neq 0$ 有 $x^T H x = \norm{Xx}_2^2 + \lambda\norm{x}_2^2 \ge \lambda\norm{x}_2^2 > 0$，故 $H\succ 0$。若 $X^TX$ 特征值为 $\mu_i \ge 0$，则 $H$ 特征值为 $\mu_i+\lambda$，故 $\cond(H) = \frac{\mu_{\max}+\lambda}{\mu_{\min}+\lambda}$。\\
(3) 谱分解 $A = Q\Lambda Q^T$（$\lambda_i\ge 0$），令 $A^{1/2} := Q\Lambda^{1/2}Q^T$（$\Lambda^{1/2}=\diag(\sqrt{\lambda_i})$），显然对称半正定且 $(A^{1/2})^2 = Q\Lambda Q^T = A$。唯一性：若 $B\succeq 0$ 满足 $B^2 = A$，则对每个 $A$ 的特征子空间 $E_\mu(A)$，因 $BA=B^3=AB$，$E_\mu(A)$ 是 $B$-不变子空间；在其上 $B$ 对称半正定且特征值平方均为 $\mu$，故特征值必全为 $\sqrt{\mu}$，即 $B|_{E_\mu} = \sqrt{\mu} I = A^{1/2}|_{E_\mu}$，从而 $B = A^{1/2}$ 唯一。$\square$

\textbf{3.6 [C]（证明）}
任取非零 $x\perp q_1$，在标准正交基 $\{q_1,\dots,q_n\}$ 下展开为 $x = \sum_{i=2}^n c_i q_i$（因 $c_1 = q_1^T x = 0$）。代入 Rayleigh 商：
\[
  \frac{x^T S x}{x^T x} = \frac{\sum_{i=2}^n \lambda_i c_i^2}{\sum_{i=2}^n c_i^2} \ge \frac{\lambda_2 \sum_{i=2}^n c_i^2}{\sum_{i=2}^n c_i^2} = \lambda_2,
\]
且取 $x = q_2$ 时等号成立。对连通图拉普拉斯 $L$，$\lambda_1=0$ 对应 $q_1 = \frac{1}{\sqrt{n}}\mathbf{1}$，约束 $\sum x_i = 0$ 恰为 $x \perp q_1$，因此极小化图平滑能量 $x^T L x = \sum_{\{i,j\}\in E} w_{ij}(x_i-x_j)^2$ 的单位信号必为 Fiedler 向量 $q_2$。$\square$

\textbf{3.7--3.8 [D]（编程自检要点）}
\begin{itemize}
  \item \textbf{3.7}：实测 $m=2$ 在 $k=9$ 达峰 $3.9128$；$m=3$ 在 $k\approx 19$ 达峰 $\approx 28.5$；$m=4$ 在 $k\approx 29$ 达峰 $>160$。\texttt{eigh} 利用对称性比 \texttt{eig} 快 3--5 倍、零虚部且正交误差 $<10^{-13}$。
  \item \textbf{3.8}：用 \texttt{P @ L @ P.T} 置换节点后，Fiedler 向量 \texttt{q2} 的正负号严格分离两个社群（允许整体取反号），准确率恒为 $100\%$。
\end{itemize}

% ====================================================================
\section{第 4 章习题解答（4.1--4.8）}
% ====================================================================

\textbf{4.1 [A]}
(a) \textbf{假}。取错切阵 $A=\begin{psmallmatrix}1&1\\0&1\end{psmallmatrix}$，特征值全为 $1$，但 $A^TA=\begin{psmallmatrix}1&1\\1&2\end{psmallmatrix}$ 奇异值为 $\frac{\sqrt{5}\pm 1}{2}\approx 1.618, 0.618 \neq 1$（仅正规阵满足 $\sigma_i = |\lambda_i|$）。\\
(b) \textbf{真}（未去中心化时原点到质心的位移能量主导总平方和）。\\
(c) \textbf{真}（显式形成 $\widetilde{X}^T\widetilde{X}$ 会使条件数平方，淹没小奇异值）。\\
(d) \textbf{真}（定理~\ref{thm:ch4-min-norm} 直接推论）。

\textbf{4.2 [A]}
取 $A = \begin{psmallmatrix}1&1\\0&0\end{psmallmatrix},\; B = \begin{psmallmatrix}1&0\\0&0\end{psmallmatrix}$。算得 $AB = \begin{psmallmatrix}1&0\\0&0\end{psmallmatrix} \implies (AB)^+ = \begin{psmallmatrix}1&0\\0&0\end{psmallmatrix}$；但 $A^+ = \frac{1}{2}\begin{psmallmatrix}1&0\\1&0\end{psmallmatrix},\; B^+ = \begin{psmallmatrix}1&0\\0&0\end{psmallmatrix} \implies B^+ A^+ = \frac{1}{2}\begin{psmallmatrix}1&0\\0&0\end{psmallmatrix} \neq (AB)^+$。同例中 $AA^+ = \begin{psmallmatrix}1&0\\0&0\end{psmallmatrix} \neq \frac{1}{2}\begin{psmallmatrix}1&1\\1&1\end{psmallmatrix} = A^+A$（因 $\col(A)\neq \row(A)$）。

\textbf{4.3 [B]}
(1) $AA^T = \begin{psmallmatrix}2&0\\0&4\end{psmallmatrix}$，按降序得 $\sigma_1 = 2, u_1 = \binom{0}{1}$；$\sigma_2 = \sqrt{2}, u_2 = \binom{1}{0}$。由 $v_i = \frac{1}{\sigma_i}A^T u_i$ 得 $v_1 = (0,0,1)^T,\; v_2 = \frac{1}{\sqrt{2}}(1,1,0)^T$，零空间补基 $v_3 = \frac{1}{\sqrt{2}}(1,-1,0)^T$。秩 1 近似 $A_1 = 2 u_1 v_1^T = \begin{psmallmatrix}0&0&0\\0&0&2\end{psmallmatrix}$，相对 Frobenius 误差 $\frac{\sigma_2}{\sqrt{\sigma_1^2+\sigma_2^2}} = \frac{\sqrt{2}}{\sqrt{6}} = \frac{1}{\sqrt{3}}\approx 0.5774$。\\
(2) $A^+ = V_2\Sigma_2^{-1}U^T = \begin{psmallmatrix}1/2 & 0\\ 1/2 & 0\\ 0 & 1/2\end{psmallmatrix}$，$x^+ = A^+\binom{2}{4} = (1,1,2)^T$。\\
(3) $\nul(A) = \spanop\{v_3\}$，$v_3^T x^+ = \frac{1}{\sqrt{2}}(1 - 1 + 0) = 0$，精确正交。

\textbf{4.4 [B]}
(1) $\mu = (4,5)$，$\widetilde{X} = \begin{psmallmatrix}-3&-1\\-1&-3\\1&3\\3&1\end{psmallmatrix}$。\\
(2) $\Sigma_{\text{cov}} = \frac{1}{4}\widetilde{X}^T\widetilde{X} = \begin{psmallmatrix}5&3\\3&5\end{psmallmatrix}$，特征值 $\lambda_1=8, \lambda_2=2$，主方向 $v_1 = \frac{1}{\sqrt{2}}\binom{1}{1},\; v_2 = \frac{1}{\sqrt{2}}\binom{-1}{1}$。\\
(3) 投影坐标 $Z_1 = \widetilde{X}v_1 = (-2\sqrt{2}, -2\sqrt{2}, 2\sqrt{2}, 2\sqrt{2})^T$；重构 $\widehat{X}_1 = \mathbf{1}\mu^T + Z_1 v_1^T$ 得四点 $(2,3),\allowbreak(2,3),\allowbreak(6,7),\allowbreak(6,7)$；PC1 解释方差比例为 $\frac{8}{8+2} = 80\%$。

\textbf{4.5 [C]（证明）}
(1) $Q_1 A Q_2 = (Q_1 U)\Sigma (Q_2^T V)^T$。因正交阵乘积仍正交（$(Q_1U)^T(Q_1U)=I$，$(Q_2^TV)^T(Q_2^TV)=I$），此即 $Q_1 A Q_2$ 的合法 SVD，故奇异值集合 $\Sigma$ 完全不变，从而 2-范数 $\sigma_1$ 与 Frobenius 范数 $\sqrt{\sum\sigma_i^2}$ 均不变。\\
(2) 令 $Q = UV^T,\; S = V\Sigma V^T$。显然 $Q^TQ = VU^TUV^T = I_n$（正交），$S$ 对称且特征值 $\sigma_i\ge 0$（半正定），且 $QS = UV^TV\Sigma V^T = U\Sigma V^T = A$。$\square$

\textbf{4.6 [C]（证明）}
将 $A = U\Sigma V^T$ 代入得 $A^TA + \lambda I_n = V(\Sigma^T\Sigma + \lambda I_n)V^T$，求逆后右乘 $A^T b = V\Sigma^T U^T b$：
\[
  x_\lambda = V(\Sigma^T\Sigma + \lambda I_n)^{-1}\Sigma^T U^T b
  = \sum_{i=1}^r \frac{\sigma_i}{\sigma_i^2 + \lambda}(u_i^T b) v_i
  = \sum_{i=1}^r \underbrace{\frac{\sigma_i^2}{\sigma_i^2+\lambda}}_{f_\lambda(\sigma_i)} \frac{u_i^T b}{\sigma_i} v_i.
\]
令 $\lambda \to 0^+$，每一项 $f_\lambda(\sigma_i) \to 1$，故极限恰为 $\sum_{i=1}^r \frac{u_i^T b}{\sigma_i}v_i = A^+ b$。滤波解释：大奇异值（$\sigma_i \gg \sqrt{\lambda}$）处 $f_\lambda \approx 1$ 几乎无损保留信号；微小奇异值（$\sigma_i \ll \sqrt{\lambda}$）处 $f_\lambda \approx \sigma_i^2/\lambda \to 0$，强力压制噪声被 $1/\sigma_i$ 放大的灾难。$\square$

\textbf{4.7--4.8 [D]（编程自检要点）}
\begin{itemize}
  \item \textbf{4.7}：由于每步梯度 $A^T(Ax-b) \in \row(A)$，从 $x^{(0)}=0$ 出发的轨迹永远留在 $\row(A)$ 内，绝不产生 $\nul(A)$ 分量，故必收敛到唯一最小范数解 $x^+$；随机初值则保留了初值的零空间分量 $(I-A^+A)x^{(0)}$，使范数更大但行空间投影与 $x^+$ 完全重合。
  \item \textbf{4.8}：能量守恒误差 $<10^{-14}$；$\varepsilon=10^{-8}$ 时 \texttt{svd} 精准给出 $\sigma_2 \approx 10^{-8}$，而 \texttt{eigh(X.T @ X)} 因 $1 + 10^{-16} == 1$ 丢掉全部精度。
\end{itemize}

% ====================================================================
\section{第 5 章习题解答（5.1--5.8）}
% ====================================================================

\textbf{5.1 [A]}
(a) \textbf{假}。梯度只决定\textbf{局部}瞬时变化率最快的方向，对非凸多峰函数或偏心椭圆等高线，负梯度既不一定指向全局最优点，甚至不指向局部极小点（如图 5.1 所示，负梯度垂直接触椭圆切线而非直指圆心）。\\
(b) \textbf{假}。1D 反例 $f(x) = x^3$ 在 $x^*=0$ 处 $f'(0)=0$ 且 $f''(0)=0 \ge 0$（半正定），但 $0$ 是拐点而非局部极小值（2D 反例 $f(x_1,x_2)=x_1^2 + x_2^3$ 在 $(0,0)$ 处 $H=\diag(2,0)\succeq 0$ 但非极小值）。

\textbf{5.2 [A]}
致命错误在于把 $\tr(AWB(dW)^T)$ 错误化简成了 $\tr(BW^TA\,dW)$！对非对称 $A,B$，正确转置应为：
\[
  \tr(AWB(dW)^T) = \tr\big((AWB(dW)^T)^T\big) = \tr\big((dW) B^T W^T A^T\big) = \tr(B^T W^T A^T\,dW).
\]
因此 $df = \tr\big((BW^TA + B^TW^TA^T)\,dW\big) = \tr\big((A^TWB^T + AWB)^T dW\big)$，正确梯度为 $\nabla_W f = A^T W B^T + A W B$（仅当 $A,B$ 均对称时才简化为 $2AWB$）。

\textbf{5.3 [B]}
(a) $\nabla f(x_1, x_2) = (3x_1^2 - 3x_2,\; -3x_1 + 3x_2^2)^T$。在 $(1,0)^T$ 处 $\nabla f(1,0) = (3, -3)^T$，故 $D_u f(1,0) = 3\cdot\frac{3}{5} + (-3)\cdot\frac{4}{5} = -\frac{3}{5} = -0.6$。\\
(b) 令 $\nabla f = 0 \implies x_2 = x_1^2$ 且 $x_1 = x_2^2 = x_1^4 \implies x_1(x_1^3-1)=0$，得两实驻点 $(0,0)^T$ 与 $(1,1)^T$。Hessian $H_f(x) = \begin{psmallmatrix}6x_1 & -3\\ -3 & 6x_2\end{psmallmatrix}$：
\begin{itemize}
  \item 在 $(0,0)^T$：$H = \begin{psmallmatrix}0&-3\\-3&0\end{psmallmatrix}$，特征值 $\pm 3$（不定），为\textbf{鞍点}；
  \item 在 $(1,1)^T$：$H = \begin{psmallmatrix}6&-3\\-3&6\end{psmallmatrix}$，特征值 $9, 3 > 0$（正定），为\textbf{严格局部极小值点}。
\end{itemize}

\textbf{5.4 [B]}
(a) 前向：$z = Wx = (1, 0)^T$；$a = \tanh(z) \approx (0.7616,\; 0.0000)^T$；$L = \frac{1}{2}\big((0.7616-1)^2 + (0-1)^2\big) \approx 0.5284$。\\
(b) 反向 VJP：$\nabla_a L = a - y \approx (-0.2384,\; -1.0000)^T$；因 $\tanh'(z)=1-a^2 \approx (0.4200,\; 1.0000)^T$，得 $\nabla_z L = (\nabla_a L)\had(1-a^2) \approx (-0.1001,\; -1.0000)^T$；于是
\[
  \nabla_W L = (\nabla_z L) x^T \approx \begin{pmatrix}-0.1001 & 0\\ -1.0000 & 0\end{pmatrix},\qquad
  \nabla_x L = W^T \nabla_z L \approx \begin{pmatrix}-0.1001\\ -1.8999\end{pmatrix}.
\]

\textbf{5.5 [C]（证明）}
取全微分：$dJ = (Xw-y)^T X\,dw + \lambda w^T dw = \big(X^T(Xw-y) + \lambda w\big)^T dw$，故 $\nabla_w J(w) = (X^TX + \lambda I)w - X^Ty$。再微分得 $d(\nabla_w J) = (X^TX + \lambda I)\,dw$，故 $H_J(w) = X^TX + \lambda I$。对一切非零 $v$，$v^T H_J v = \norm{Xv}_2^2 + \lambda\norm{v}_2^2 \ge \lambda\norm{v}_2^2 > 0$，故无论 $X$ 秩如何均有 $H_J \succ 0$（可逆）。令 $\nabla_w J = 0$ 得唯一驻点 $w_{\text{ridge}} = (X^TX+\lambda I)^{-1}X^Ty$，由二阶展开 $J(w) = J(w_{\text{ridge}}) + \frac{1}{2}(w-w_{\text{ridge}})^T H_J (w-w_{\text{ridge}}) \ge J(w_{\text{ridge}})$ 知其为唯一全局极小点。$\square$

\textbf{5.6 [C]（证明）}
构造拉格朗日函数 $\mathcal{L}(u, v, \alpha, \beta) = u^T A v - \frac{\alpha}{2}(u^T u - 1) - \frac{\beta}{2}(v^T v - 1)$。分别对 $u, v$ 求梯度置零：
\[
  \nabla_u \mathcal{L} = Av - \alpha u = 0 \implies Av = \alpha u,\qquad
  \nabla_v \mathcal{L} = A^T u - \beta v = 0 \implies A^T u = \beta v.
\]
左乘 $u^T$ 于第一式、左乘 $v^T$ 于第二式，结合单位长约束得 $\alpha = u^T A v = \beta$。记 $\sigma := \alpha = \beta$，代回得 $A^T A v = \sigma A^T u = \sigma^2 v$，即 $v$ 为 $A^T A$ 的单位特征向量，$\sigma^2$ 为对应特征值；目标值 $u^T A v = \sigma$ 最大化时，必取最大奇异值 $\sigma = \sigma_1(A)$。$\square$

\textbf{5.7--5.8 [D]（编程自检要点）}
\begin{itemize}
  \item \textbf{5.7}：\texttt{c @ (W2 @ W1)} 需 $O(d^3)$（约 $5\times 10^8$ 次乘加），而反向 \texttt{(c @ W2) @ W1} 仅做两次向量-矩阵乘 $O(d^2)$（约 $1.3\times 10^6$ 次），在 $d=800$ 时通常加速 $50\text{--}200\times$。
  \item \textbf{5.8}：解析反传核心四式：\texttt{dY = Y\_hat - Y}、\texttt{dW2 = dY @ A1.T}、\texttt{dZ1 = (W2.T @ dY) * (1.0 - A1**2)}、\texttt{dW1 = dZ1 @ X.T}。用 \texttt{check\_gradient} 校验两矩阵相对误差均 $<10^{-10}$。
\end{itemize}

% ====================================================================
\section{第 6 章习题解答（6.1--6.8）}
% ====================================================================

\textbf{6.1 [A]}
(a) \textbf{假}。负梯度仅保证\textbf{充分小步长}下函数下降；一旦步长跨过 $\eta_{\max}=2/\lambda_{\max}$（如正文 6.1 节 $f(w)=\frac{1}{2}(w_1^2+10w_2^2)$ 取 $\eta=0.21 > 0.2$），就会越过谷底飞到更高处，$f(w_{t+1}) > f(w_t)$ 且指数发散。\\
(b) \textbf{假}。当 $\epsilon=10^{-14}$ 时，$f(w+\epsilon e_i)-f(w-\epsilon e_i)$ 发生严重灾难性抵消，浮点舍入误差 $O(u/\epsilon)\sim 10^{-2}$ 远大于 $\epsilon=10^{-5}$ 处的最佳平衡点（$\sim 10^{-11}$）。

\textbf{6.2 [A]}
错误一：\textbf{数据泄露}——用含测试集的 \texttt{X\_all} 算 \texttt{mu, sigma} 会将测试分布信息提前泄露给训练过程，应先切分训练集/测试集，仅在训练集上 \texttt{fit} 得到 \texttt{mu\_tr, sigma\_tr} 再变换两边。\\
错误二：\textbf{推理坐标系错配}——\texttt{w\_std} 是在标准化空间学出的权重，直接乘未缩放的 \texttt{x\_new} 会导致量纲彻底错乱；正确做法是推理时先用同一组 \texttt{(mu\_tr, sigma\_tr)} 标准化 \texttt{x\_new}，或按式~\eqref{eq:ch6-recover-w} 解析还原出原始空间权重 $(w, b)$。

\textbf{6.3 [B]}
(a) $Q=\begin{psmallmatrix}5&3\\3&5\end{psmallmatrix}$ 特征值 $\lambda_{\max}=8,\; \lambda_{\min}=2$，条件数 $\cond(Q)=4$；解 $Qw^*=b$ 得 $w^*=(1,1)^T$。\\
(b) $\eta_{\max} = 2/8 = 0.25$；最优步长 $\eta^* = 2/(8+2) = 0.2$；最优单步收缩因子 $\rho^* = \frac{4-1}{4+1} = 0.6$。\\
(c) 在 $w_0=(0,0)^T$ 处 $\nabla f(w_0) = -b = (-8,-8)^T$。取 $\eta^*=0.2$ 得 $w_1 = (1.6, 1.6)^T$（误差 $e_1=(0.6,0.6)^T = -0.6 e_0$）；取 $\eta=1/8=0.125$ 得 $w_1 = (1,1)^T = w^*$。\textbf{原因}：初始误差 $e_0 = (-1,-1)^T$ 恰好完全落在 $\lambda_{\max}=8$ 的特征方向 $v_1=\frac{1}{\sqrt{2}}(1,1)^T$ 上（在 $\lambda_{\min}=2$ 方向分量为 $0$），故取 $\eta=1/\lambda_{\max}=1/8$ 使因子 $1-\eta\lambda_{\max} = 0$，一步直接归零！

\textbf{6.4 [B]}
由 $\tilde{x}_1 = \frac{x_1-100}{25},\; \tilde{x}_2 = \frac{x_2-10}{5}$，还原得：
\[
  w_1 = \frac{50}{25} = 2,\qquad w_2 = \frac{-15}{5} = -3,\qquad
  b = 300 - 2\times 100 - (-3)\times 10 = 130,
\]
即 $\hat{y} = 130 + 2x_1 - 3x_2$。代入 $(125, 5)$：标准化坐标 $\tilde{x}=(1, -1)$，标准化式得 $300+50-(-15)=365$；原始式得 $130+250-15=365$，完全相等。

\textbf{6.5 [C]（证明）}
由第 3 章谱映射性质，$H_\alpha = H_0 + \alpha I$ 的最大、最小特征值为 $\lambda_{\max}+\alpha$ 与 $\lambda_{\min}+\alpha$，故 $\cond(H_\alpha) = \frac{\lambda_{\max}+\alpha}{\lambda_{\min}+\alpha}$。作差得
\[
  \cond_0 - \cond(H_\alpha) = \frac{\lambda_{\max}}{\lambda_{\min}} - \frac{\lambda_{\max}+\alpha}{\lambda_{\min}+\alpha}
  = \frac{\alpha(\lambda_{\max} - \lambda_{\min})}{\lambda_{\min}(\lambda_{\min}+\alpha)} > 0,
\]
故 $\cond(H_\alpha) < \cond_0$。进一步，最优收缩因子 $\rho^*(\alpha) = \frac{\cond(H_\alpha)-1}{\cond(H_\alpha)+1} = \frac{\lambda_{\max}-\lambda_{\min}}{\lambda_{\max}+\lambda_{\min}+2\alpha}$，分子为正常数、分母随 $\alpha>0$ 严格增大，故 $\rho^*(\alpha)$ 严格单调减小。$\square$

\textbf{6.6 [C]（推导）}
对 $E(\epsilon) = \frac{K}{6}\epsilon^2 + \frac{Mu}{\epsilon}$ 求导并置零：
\[
  E'(\epsilon) = \frac{K}{3}\epsilon - \frac{Mu}{\epsilon^2} = 0 \implies \epsilon^* = \left(\frac{3Mu}{K}\right)^{1/3}.
\]
因 $E''(\epsilon) = \frac{K}{3} + \frac{2Mu}{\epsilon^3} > 0$，此为唯一极小值点。代回得最小总误差：
\[
  E(\epsilon^*) = \frac{K}{6}\left(\frac{3Mu}{K}\right)^{2/3} + Mu\left(\frac{K}{3Mu}\right)^{1/3}
  = \frac{3}{2}\left(\frac{K}{3}\right)^{1/3}(Mu)^{2/3}.
\]
代入 $K\sim M\sim 1$ 与 $u\approx 2.22\times 10^{-16}$：$\epsilon^* \approx (6.66\times 10^{-16})^{1/3} \approx 8.7\times 10^{-6} \sim 10^{-5}$，$E(\epsilon^*) \approx 3.8\times 10^{-11}$，与图~\ref{fig:ch6-gradcheck} 实测底部（$\epsilon=10^{-5}$ 处误差 $2.62\times 10^{-11}$）数量级完全吻合。$\square$

\textbf{6.7--6.8 [D]（编程自检要点）}
\begin{itemize}
  \item \textbf{6.7}：原始特征尺度差 $10^5$ 倍，Hessian 条件数约 $10^{10}$；标准化后条件数降至 $\sim 1.1$--$1.3$，最优步长 $\eta^* \approx 1.0$，不到 $20$ 步相对误差即可跌破 $10^{-10}$。
  \item \textbf{6.8}：注意 $L_2(w)$ 前向写成 \texttt{np.mean(np.logaddexp(0.0, X @ w) - y * (X @ w)) + 0.5 * lam * np.dot(w, w)}，梯度为 \texttt{(X.T @ (1.0 / (1.0 + np.exp(-X @ w)) - y)) / N + lam * w}；两损失在 $\epsilon=10^{-5}$ 下相对误差均 $\sim 10^{-11} < 10^{-9}$，改错符号后相对误差跳升至 $O(1)$ 并立即触发 \texttt{passed == False}。
\end{itemize}


\end{document}
